\chapter{积分理论}
上一章我们已经熟悉了测度论的基础，现在转向积分理论。本章第一部分研究一般测度空间上的积分，后半部分则利用勒贝格测度的特殊性质。

关于任意测度的积分不仅在众多应用中很重要，而且在最后一章也将不可或缺，那时底层的集合不再是“平坦”的，而是一个流形。正因如此，像本书这样的入门教材也必须处理这一课题。

在第 1 节中，我们引入 $ \mu $-可测函数并研究它们的基本性质。在分析中，那些使每个连续函数都可测的自然测度备受关注。一个例子是拉东测度类，我们在本节引入它，并且在第 XII 章还会再次遇到。

在分析中——而且不只是在分析中——能够处理向量值函数，即以巴拿赫空间为值的映射，将变得越来越重要。我们在前两卷中已经沿着这条路线做过工作，你也会注意到，由此得到的叙述不仅在优雅性上有所收获，而且在许多情形下还更加简单。积分理论的情况也是如此。因此我们一开始就决定按向量值函数来发展这套理论，从而处理 Bochner–Lebesgue 积分。这样做不需要付出多少额外的努力。少数例外之一，是证明向量值函数为 $ \mu $-可测当且仅当它按通常意义可测并且是 $ \mu $-几乎可分值的。当然，你完全可以忽略这个结果，只考虑标量值函数。但这并不值得推荐，因为那样你会错过工具箱中一项重要而高效的补充。

除向量值映射之外，我们还将较为详细地研究取值于扩充实数轴 $ [0,\infty] $ 的函数。这主要是出于技术上的原因；在后面的各节中，它会使我们不必总是把特殊情形单独挑出来处理。

在第 2 节中，我们引入一般的 Bochner–Lebesgue 积分，做法是取简单函数空间的 $ \mathcal{L}_1 $-完备化。这一途径不仅基本上原封不动地推广到向量值函数，而且为勒贝格收敛定理的证明奠定了基础。该定理以及其他重要的收敛定理，我们在第 3 节中处理。

第 4 节讨论勒贝格空间的初等理论。我们证明它们的完备性，并说明：如果把几乎处处相等的函数等同起来，它们就成为巴拿赫空间。由于根据我们的经验，这种等同是初学者困难的来源，我们在全章中都细致地区分函数的等价类与它们各自的代表元。

尽管到此为止我们考虑的都是关于任意测度的积分，第 5 节将处理 $ \mathbb{R}^n $ 中勒贝格测度这一特殊情形。我们证明一维勒贝格积分是柯西–黎曼积分在绝对可积函数上的推广。这使我们能够把第 II 卷中学到的关于积分的内容纳入一般理论的框架。这在富比尼定理的语境中尤为重要，该定理给出了计算高维积分的一种归约程序。

第 6 节讨论富比尼定理。我们决定不对任意的乘积测度空间证明它，而只对勒贝格测度空间证明。这大大简化了叙述，并且实际上足以满足分析的一切需要——只要再通过推广到乘积流形而得到加强，而这一推广将在第 XII 章处理。

向量值情形下富比尼定理的证明需要细致的可测性论证。为此我们先研究标量情形。我们在第 6 节末尾证明向量值版本，并给出一些重要应用。初读时可以跳过这一部分，因为它的结果在之后并没有以任何本质的方式被用到，而且也因为读者大概会在日后的某个时候熟悉泛函分析中的哈恩–巴拿赫定理：借助它，向量值函数的富比尼定理很容易由标量版本推出。

第 7 节研究卷积。这一运算使我们能够以非凡的效率证明一些基本的逼近定理，例如关于光滑单位分解的定理，它在最后一章中起着重要作用。在本节的后半部分，我们讨论卷积与逼近定理在分析和数学物理中的意义，并第一次瞥见经典微分学的一个非常重要的推广，即所谓的分布理论。

除勒贝格与富比尼的收敛定理之外，变换定理构成整个积分学的第三根支柱。它将在第 8 节证明，在那里我们还讨论它较为基本的应用。

在最后一节中，我们通过证明关于傅里叶变换的几个基本事实来说明刚刚发展起来的这套理论的力量。与第 7 节的后半部分一样，这一部分让人得以一窥分析中一个相关的领域，你日后在更深入的学习中可能会遇到它。

\section{可测函数}

设 $ (X,\mathcal{A},\mu) $ 是测度空间且 $ A \in \mathcal{A} $。与初等几何构造的类比使人把特征函数 $ \chi_A $ 关于测度 $ \mu $ 在 $ X $ 上的积分定义为 $ \int_X \chi_A \, d\mu := \mu(A) $。显然，这只有在 $ A $ 属于 $ \mathcal{A} $ 时才有意义。因此函数 $ f = \chi_A $ 必须在这个意义下与底层的测度空间 $ (\mathcal{A},\mu) $“相容”。对于更复杂的函数，适当的逼近论证使人们可以推广函数与测度之间的这种相容性概念，从而引出函数的可测性这一概念。

本节中总设

- $ (X,\mathcal{A},\mu) $ 是完备的 $ \sigma $-有限测度空间；$ E = (E,|\cdot|) $ 是 Banach 空间。

\needspace{3\baselineskip}\subsection*{简单函数与可测函数}

设 $ \mathrm{E} $ 是 $ X $ 的每个点上或真或假的性质。若存在 $ \mu $-零集 $ N $ 使得 $ \mathrm{E}(x) $ 对每个 $ x \in N^c $ 为真，我们就说 $ \mathrm{E} $ \textbf{$ \mu $-几乎处处}成立，或者说\textbf{对 $ \mu $-几乎每个 $ x \in X $} 成立。“几乎每个”与“几乎处处”都缩写为“a.e.”。

\begin{example}

\begin{items}
\item

对 $ f, g \in \mathbb{R}^X $，若存在 $ \mu $-零集 $ N $ 使得 $ f(x) \ge g(x) $ 对每个 $ x \in N^c $ 成立，我们就记“$ f \ge g $ $ \mu $-a.e.”。

\item

设 $ f_j, f \in E^X $，$ j \in \mathbb{N} $。则 $ (f_j) $ $ \mu $-a.e. 收敛于 $ f $ 当且仅当存在 $ \mu $-零集 $ N $ 使得 $ f_j(x) \to f(x) $ 对 $ x \in N^c $ 成立。

\item

函数 $ f \in E^X $ 是 $ \mu $-a.e. 有界的，当且仅当存在 $ \mu $-零集 $ N $ 与 $ M \ge 0 $ 使得 $ |f(x)| \le M $ 对每个 $ x \in N^c $ 成立。

\item

若 $ \mathrm{E} $ 在 $ \mu $-a.e. 成立，则集合 $ \{ x \in X ;\ \mathrm{E}(x) \text{ 不成立} \} $ 是 $ \mu $-零集。

\proofhead
这由 $ (X,\mathcal{A},\mu) $ 的完备性推出。\nobreak\hspace{0.35em}■
\item

设 $ (X,\mathcal{B},\nu) $ 是不完备的测度空间。则 $ X $ 上存在性质 $ \mathrm{E} $，它 $ \nu $-几乎处处成立，而集合 $ \{ x \in X ;\ \mathrm{E}(x) \text{ 不成立} \} $ 却不是 $ \nu $-零集。

\proofhead
由于 $ (X,\mathcal{B},\nu) $ 不完备，存在 $ \nu $-零集 $ N $ 与 $ M \subset N $ 使得 $ M \notin \mathcal{B} $。若令 $ f := \chi_M $，则 $ f = 0 $ 在 $ \nu $-几乎处处成立，而 $ \{ x \in X ;\ f(x) \neq 0 \} = M $ 不是 $ \nu $-零集。\nobreak\hspace{0.35em}■
\end{items}

我们称 $ f \in E^X $ 是 \textbf{$ \mu $-简单的}\footnote{若测度空间的标识是清楚的，我们就把函数称为\textbf{简单}的而不说 $ \mu $-简单的；对即将引入的 $ \mu $-可测函数也类似。}，若 $ f(X) $ 有限、$ f^{-1}(e) \in \mathcal{A} $ 对每个 $ e \in E $ 成立且 $ \mu\bigl(f^{-1}(E \setminus \{0\})\bigr) < \infty $。我们用 $ \mathcal{S}(X,\mu,E) $ 表示从 $ X $ 到 $ E $ 的所有 $ \mu $-简单函数之集\footnote{我们曾把跳跃连续函数空间记作 $ \mathcal{S}(I,E) $，但这不会引起混淆。}。

称函数 $ f \in E^X $ 为 \textbf{$ \mu $-可测的}，若存在 $ \mathcal{S}(X,\mu,E) $ 中的序列 $ (f_j) $ 使得 $ f_j \to f $ $ \mu $-几乎处处（$ j \to \infty $）。我们设\footnote{显然，函数的可测性定义即使在不完备的测度空间上也有意义。}

\[
\mathcal{L}_0(X,\mu,E) := \{ f \in E^X ;\ f \text{ 是 } \mu\text{-可测的} \} .
\]
\end{example}

\begin{remark}

\begin{items}
\item

我们有向量子空间的包含关系

\[
\mathcal{S}(X,\mu,E) \subset \mathcal{L}_0(X,\mu,E) \subset E^X .
\]
\item

设 $ j = 0,\ldots,m $（其中 $ m \in \mathbb{N} $），考虑满足 $ \mu(A_j) < \infty $ 的 $ e_j \in E $ 与 $ A_j \in \mathcal{A} $。则 $ f := \sum_{j=0}^{m} e_j \chi_{A_j} $ 属于 $ \mathcal{S}(X,\mu,E) $。若

\[
e_j \neq 0 \quad \text{对 } j = 0,\ldots,m ,
\]
\[
e_j \neq e_k \quad \text{对 } j \neq k ,
\]
\[
A_j \cap A_k = \varnothing \quad \text{对 } j \neq k ,
\]

我们就称这是 $ f $ 的 \textbf{正规形}。

\item

每个简单函数都有唯一的正规形，且\footnote{比较习题 VI.6.8 的脚注。}

\[
\mathcal{S}(X,\mu,E) = \left\{ \sum_{j=1}^{m} e_j \chi_{A_j} \ ;\ m \in \mathbb{N},\ e_j \in E \setminus \{0\},\ A_j \in \mathcal{A},\ \mu(A_j) < \infty,\ A_j \cap A_k = \varnothing \text{ 对 } j \neq k \right\} .
\]

\proofhead
设 $ f \in \mathcal{S}(X,\mu,E) $。则存在 $ m \in \mathbb{N} $ 与 $ E $ 中两两不同的元素 $ e_0,\ldots,e_m $ 使得 $ f(X) \setminus \{0\} = \{e_0,\ldots,e_m\} $。令 $ A_j := f^{-1}(e_j) $，则 $ A_j \in \mathcal{A} $、$ \mu(A_j) < \infty $ 且 $ A_j \cap A_k = \varnothing $（$ j \neq k $）。容易验证
\[
\sum_{j=0}^{m} e_j \chi_{A_j}
\]

是 $ f $ 的唯一正规形。第二部分现在由 (b) 推出。\nobreak\hspace{0.35em}■

\item

设 $ f \in E^X $ 与 $ g \in \mathbb{K}^X $ 是 $ \mu $-简单的［或 $ \mu $-可测的］。则 $ |f| \in \mathbb{R}^X $ 与 $ gf \in E^X $ 也是 $ \mu $-简单的［或 $ \mu $-可测的］。特别地，$ \mathcal{S}(X,\mu,\mathbb{K}) $ 与 $ \mathcal{L}_0(X,\mu,\mathbb{K}) $ 是 $ \mathbb{K}^X $ 的子代数。

\item

设 $ A \in \mathcal{A} $ 与 $ f \in E^X $，作限制 $ \nu := \mu \mid (\mathcal{A} \mid A) $（见习题 IX.1.7）。则

\[
f \mid A \in \mathcal{S}(A,\nu,E) \iff \chi_A f \in \mathcal{S}(X,\mu,E) ,
\]
\[
f \mid A \in \mathcal{L}_0(A,\nu,E) \iff \chi_A f \in \mathcal{L}_0(X,\mu,E) .
\]

\proofhead
验证很简单，留给读者。\nobreak\hspace{0.35em}■
\item

设 $ f \in \mathcal{L}_0(X,\mu,\mathbb{K}) $ 且 $ A := [f \neq 0] $。再通过

\begin{equation*}
g(x) := \begin{cases} 1/f(x) & \text{若 } x \in A , \\ 0 & \text{若 } x \notin A \end{cases}
\end{equation*}

定义 $ g \in \mathbb{K}^X $。
则 $ g $ 是 $ \mu $-可测的。

\proofhead
$ f $ 的可测性蕴含存在 $ \mu $-零集 $ N $ 与 $ \mathcal{S}(X,\mu,\mathbb{K}) $ 中的序列 $ (\varphi_j) $ 使得 $ \varphi_j(x) \to f(x) $ 对 $ x \in N^c $ 成立。我们设
\begin{equation*}
\psi_j(x) := \begin{cases} 1/\varphi_j(x) & \text{若 } \varphi_j(x) \neq 0 , \\ 0 & \text{若 } \varphi_j(x) = 0 , \end{cases}
\end{equation*}

对 $ x \in X $、$ j \in \mathbb{N} $。由 (c) 与 (d)，$ (\chi_A \psi_j) $ 是 $ \mathcal{S}(X,\mu,\mathbb{K}) $ 中的序列，且容易验证 $ (\chi_A \psi_j)(x) \to g(x) $ 对每个 $ x \in N^c $ 成立（见命题 II.2.6）。\nobreak\hspace{0.35em}■

\item

设 $ e \in E \setminus \{0\} $，并设 $ \mu(X) = \infty $。则 $ e\chi_X $ 属于 $ \mathcal{L}_0(X,\mu,E) $ 但不属于 $ \mathcal{S}(X,\mu,E) $。

\proofhead
显然 $ e\chi_X $ 不是 $ \mu $-简单的。由于 $ X $ 是 $ \sigma $-有限的，存在 $ \mathcal{A} $ 中的序列 $ (A_j) $ 使得 $ \bigcup_j A_j = X $ 且 $ \mu(A_j) < \infty $ 对 $ j \in \mathbb{N} $ 成立。对 $ j \in \mathbb{N} $，令 $ X_j := \bigcup_{k=0}^{j} A_k $ 与 $ \varphi_j := e\chi_{X_j} $。则 $ (\varphi_j) $ 是 $ \mathcal{S}(X,\mu,E) $ 中逐点收敛于 $ e\chi_X $ 的序列。\nobreak\hspace{0.35em}■
\end{items}

\end{remark}

\needspace{3\baselineskip}\subsection*{一个可测性判别法}

称函数 $ f \in E^X $ 为 \textbf{$ \mathcal{A} $-可测的}，若 $ E $ 的开集在 $ f $ 下的原像可测，即若 $ f^{-1}(\mathcal{T}_E) \subset \mathcal{A} $，其中 $ \mathcal{T}_E $ 是 $ E $ 上的范数拓扑。若存在 $ \mu $-零集 $ N $ 使得 $ f(N^c) $ 可分，我们就说 $ f $ 是 \textbf{$ \mu $-几乎可分值的}。

\begin{remark}

\begin{items}
\item

习题 IX.1.6 表明 $ \mathcal{A} $-可测函数之集与 $ \mathcal{A} $-$ \mathcal{B}(E) $-可测函数之集一致。

\item

可分度量空间的每个子空间都是可分的。

\proofhead
由命题 IX.1.8，可分性等价于具有可数基。但对任意给定的子空间，拓扑空间的基经限制就给出该子空间的一个基（基数不更大）；见命题 III.2.26。\nobreak\hspace{0.35em}■
\item

设 $ E $ 可分且 $ f \in E^X $。则 $ f $ 是 $ \mu $-几乎可分值的。

\proofhead
这由 (b) 推出。\nobreak\hspace{0.35em}■
\item

每个有限维赋范向量空间都是可分的。\footnote{比较例 V.4.3(e)。} \nobreak\hspace{0.35em}■

\end{items}

下一个结果给出 $ \mu $-可测函数的一个刻画，它除了具有理论意义之外，在判定可测性时也非常实用。

\end{remark}

\begin{theorem}

$ E^X $ 中的函数是 $ \mu $-可测的当且仅当它是 $ \mathcal{A} $-可测的且 $ \mu $-几乎可分值的。

\proofhead
“$ \Rightarrow $” 设 $ f \in \mathcal{L}_0(X,\mu,E) $。

\begin{subitems}
\item

存在 $ \mu $-零集 $ N $ 与 $ \mathcal{S}(X,\mu,E) $ 中的序列 $ (\varphi_j) $ 使得

\begin{equation}
\varphi_j(x) \to f(x) \quad (j \to \infty) \quad \text{对 } x \in N^c .\tag{1.1}
\end{equation}

由命题 I.6.8，$ F := \bigcup_{j=0}^{\infty} \varphi_j(X) $ 可数，因此其闭包 $ \overline{F} $ 可分。由 (1.1) 有 $ f(N^c) \subset \overline{F} $。注 1.3(b) 现在表明 $ f $ 是 $ \mu $-几乎可分值的。

\item

设 $ O $ 是 $ E $ 中的开集，并定义 $ O_n := \{ y \in O ;\ \operatorname{dist}(y,O^c) > 1/n \} $，$ n \in \mathbb{N}^\times $。则 $ O_n $ 开且 $ \overline{O}_n \subset O $。又设 $ x \in N^c $。由 (1.1)，$ f(x) $ 属于 $ O $ 当且仅当存在 $ n \in \mathbb{N}^\times $ 与 $ m = m(n) \in \mathbb{N}^\times $ 使得 $ \varphi_j(x) \in O_n $ 对 $ j \ge m $ 成立。因此

\begin{equation}
f^{-1}(O) \cap N^c = \bigcup_{m,n \in \mathbb{N}^\times} \bigcap_{j \ge m} \varphi_j^{-1}(O_n) \cap N^c .\tag{1.2}
\end{equation}

但 $ \varphi_j^{-1}(O_n) \in \mathcal{A} $（$ n \in \mathbb{N}^\times $，$ j \in \mathbb{N} $），因为 $ \varphi_j $ 是 $ \mu $-简单的。因此 (1.2) 说明 $ f^{-1}(O) \cap N^c \in \mathcal{A} $。
此外，$ \mu $ 的完备性表明 $ f^{-1}(O) \cap N $ 是 $ \mu $-零集，合起来我们得到

\[
f^{-1}(O) = (f^{-1}(O) \cap N) \cup (f^{-1}(O) \cap N^c) \in \mathcal{A} .
\]

“$ \Leftarrow $” 设 $ f $ 是 $ \mu $-几乎可分值的且 $ \mathcal{A} $-可测的。

\item

先考虑 $ \mu(X) < \infty $ 的情形。取 $ n \in \mathbb{N} $。按假设，存在 $ \mu $-零集 $ N $ 使得 $ f(N^c) $ 可分。若 $ \{ e_j ;\ j \in \mathbb{N} \} $ 是 $ f(N^c) $ 的可数稠密子集，则族 $ \{ \mathbb{B}(e_j,1/(n+1)) ;\ j \in \mathbb{N} \} $ 覆盖集合 $ f(N^c) $，因此

\[
X = N \cup \bigcup_{j \in \mathbb{N}} f^{-1}\bigl(\mathbb{B}(e_j,1/(n+1))\bigr) .
\]

由于 $ f $ 是 $ \mathcal{A} $-可测的，$ X_{j,n} := f^{-1}\bigl(\mathbb{B}(e_j,1/(n+1))\bigr) $ 对每个 $ (j,n) \in \mathbb{N}^2 $ 属于 $ \mathcal{A} $。$ \mu $ 的下连续性以及假设 $ \mu(X) < \infty $ 于是蕴含存在 $ m_n \in \mathbb{N}^\times $ 与 $ Y_n \in \mathcal{A} $ 使得

\[
\bigcup_{j=0}^{m_n} X_{j,n} = Y_n^c \quad \text{且} \quad \mu(Y_n) < \frac{1}{2^{n+1}} .
\]

现在通过

\begin{equation*}
\varphi_n(x) := \begin{cases} e_0 & \text{若 } x \in X_{0,n} , \\ e_j & \text{若 } x \in X_{j,n} \setminus \bigcup_{k=0}^{j-1} X_{k,n} \text{ 且 } 1 \le j \le m_n , \\ 0 & \text{其他情形} \end{cases}
\end{equation*}

定义 $ \varphi_n \in E^X $。
显然 $ \varphi_n \in \mathcal{S}(X,\mu,E) $（$ n \in \mathbb{N} $），且

\[
|\varphi_n(x) - f(x)| < 1/(n+1) \quad \text{对 } x \in Y_n^c .
\]

递减序列 $ Z_n := \bigcup_{k=0}^{\infty} Y_{n+k} $ 满足

\[
\mu(Z_n) \le \sum_{k=0}^{\infty} \mu(Y_{n+k}) \le \frac{1}{2^n} \quad \text{对 } n \in \mathbb{N} .
\]

因此由 $ \mu $ 的上连续性推出 $ Z := \bigcap_{n \in \mathbb{N}} Z_n $ 是 $ \mu $-零集。现在我们设

\begin{equation*}
\psi_n(x) := \begin{cases} \varphi_n(x) & \text{若 } x \in Z_n^c , \\ 0 & \text{若 } x \in Z_n . \end{cases}
\end{equation*}

则 $ (\psi_n) $ 是 $ \mathcal{S}(X,\mu,E) $ 中的序列。又对每个 $ x \in Z^c = \bigcup_n Z_n^c $ 存在 $ m \in \mathbb{N} $ 使得 $ x \in Z_m^c $。由于 $ Z_m^c \subset Z_n^c $（$ n \ge m $）推出

\[
|\psi_n(x) - f(x)| = |\varphi_n(x) - f(x)| < 1/(n+1) .
\]

总之，$ \lim \psi_n(x) = f(x) $ 对每个 $ x \in Z^c $ 成立。因此 $ f $ 是 $ \mu $-可测的。

\item

最后我们考虑 $ \mu(X) = \infty $ 的情形。注 IX.2.4(c) 表明存在 $ \mathcal{A} $ 中的不交序列 $ (X_j) $ 使得 $ \bigcup_j X_j = X $ 且 $ \mu(X_j) < \infty $。由 (iii)，对每个 $ j \in \mathbb{N} $ 存在 $ \mathcal{S}(X,\mu,E) $ 中的序列 $ (\varphi_{j,k})_{k \in \mathbb{N}} $ 与 $ \mu $-零集 $ N_j $ 使得 $ \lim_k \varphi_{j,k}(x) = f(x) $ 对每个 $ x \in X_j \cap N_j^c $ 成立。令 $ N := \bigcup_j N_j $ 及

\begin{equation*}
\varphi_k(x) := \begin{cases} \varphi_{j,k}(x) & \text{若 } x \in X_j ,\ j \in \{0,\ldots,k\} , \\ 0 & \text{若 } x \notin \bigcup_{j=0}^{k} X_j \end{cases}
\end{equation*}

对 $ k \in \mathbb{N} $，我们有 $ \varphi_k \in \mathcal{S}(X,\mu,E) $ 且 $ \lim_k \varphi_k(x) = f(x) $ 对 $ x \in N^c $ 成立。结果成立，因为 $ N $ 是 $ \mu $-零集。\nobreak\hspace{0.35em}■

\end{subitems}

\end{theorem}

\begin{corollary}

设 $ E $ 可分且 $ f \in E^X $。则下列陈述等价：

\begin{subitems}
\item

$ f $ 是 $ \mu $-可测的。

\item

$ f $ 是 $ \mathcal{A} $-可测的。

\item

$ f^{-1}(\mathcal{S}) \subset \mathcal{A} $ 对某个满足 $ \mathcal{A}_\sigma(\mathcal{S}) = \mathcal{B}(E) $ 的 $ \mathcal{S} \subset \mathfrak{P}(E) $ 成立。

\item

$ f^{-1}(\mathcal{S}) \subset \mathcal{A} $ 对任何满足 $ \mathcal{A}_\sigma(\mathcal{S}) = \mathcal{B}(E) $ 的 $ \mathcal{S} \subset \mathfrak{P}(E) $ 成立。

\proofhead
这由定理 1.4、注 1.3(c) 与习题 IX.1.6 推出。\nobreak\hspace{0.35em}■
\end{subitems}

\end{corollary}

\begin{remark}

定理 1.4 与注 1.3(c) 的证明表明推论 1.5 对不完备的测度空间仍然成立。\nobreak\hspace{0.35em}■

毫不费力地，我们从推论 1.5 得到 $ \mu $-可测函数的下列性质。

\end{remark}

\begin{theorem}

\begin{subitems}
\item

若 $ E $ 与 $ F $ 是可分 Banach 空间，且有映射 $ f \in \mathcal{L}_0(X,\mu,E) $ 与 $ g \in C(f(X),F) $，则 $ g \circ f $ 属于 $ \mathcal{L}_0(X,\mu,F) $。特别地，$ |f| \in \mathcal{L}_0(X,\mu,\mathbb{R}) $。

\item

映射 $ f = (f_1,\ldots,f_n): X \to \mathbb{K}^n $ 是 $ \mu $-可测的当且仅当它的每个分量 $ f_j $ 都是。

\item

设 $ g, h \in \mathbb{R}^X $。则 $ f = g + ih $ 是 $ \mu $-可测的当且仅当 $ g $ 与 $ h $ 都是。

\item

若 $ f \in \mathcal{L}_0(X,\mu,E) $ 且 $ g \in \mathcal{L}_0(X,\mu,F) $，则 $ (f,g) \in \mathcal{L}_0(X,\mu,E \times F) $。

\proofhead
(i) 设 $ O $ 是 $ F $ 中的开集。由于 $ g $ 连续，$ g^{-1}(O) $ 在 $ f(X) $ 中开。由命题 III.2.26，存在 $ E $ 的开子集 $ U $ 使得 $ g^{-1}(O) = f(X) \cap U $。由于 $ f $ 是 Lebesgue 可测的，由推论 1.5，$ f^{-1}(U) $ 属于 $ \mathcal{A} $。因为
\[
(g \circ f)^{-1}(O) = f^{-1}\bigl(g^{-1}(O)\bigr) = f^{-1}(f(X) \cap U) = f^{-1}(U) ,
\]

再次应用推论 1.5 即得结论。

\end{subitems}

\begin{subitems}
\item[（ii）]

蕴含 “$ \Rightarrow $” 由 (i) 推出，因为 $ f_j = \mathrm{pr}_j \circ f $（$ 1 \le j \le n $）。
“$ \Leftarrow $” 先考虑 $ \mathbb{K} = \mathbb{R} $ 的情形。取 $ I \in \mathbb{J}(n) $，并把它写成 $ I = \prod_{j=1}^{n} I_j $，其中 $ I_j \in \mathbb{J}(1) $（$ 1 \le j \le n $）。因为每个 $ f_j^{-1}(I_j) $ 属于 $ \mathcal{A} $，所以 $ f^{-1}(I) = \bigcap_{j=1}^{n} f_j^{-1}(I_j) $ 也属于 $ \mathcal{A} $，即 $ f^{-1}(\mathbb{J}(n)) \subset \mathcal{A} $。又由定理 IX.1.11 知 $ \mathcal{A}_\sigma(\mathbb{J}(n)) = \mathcal{B}^n $。因此推论 1.5 蕴含 $ f $ 是 $ \mu $-可测的。
利用等同 $ \mathbb{C}^n = \mathbb{R}^{2n} $，$ \mathbb{K} = \mathbb{C} $ 的情形立即由刚刚证明的结果推出。

\item[（iii）]

是 (ii) 的特殊情形，而 (iv) 我们留作练习。\nobreak\hspace{0.35em}■

\end{subitems}

\end{theorem}

\needspace{3\baselineskip}\subsection*{可测的扩充实值函数}

在积分理论中，不仅考虑实值函数，而且考虑映入扩充实数轴 $ \overline{\mathbb{R}} $ 的映射是有用的。这样的映射称为 \textbf{$ \overline{\mathbb{R}} $ 值函数}。称 $ \overline{\mathbb{R}} $ 值函数 $ f: X \to \overline{\mathbb{R}} $ 为 \textbf{$ \mu $-可测的}，若 $ \mathcal{A} $ 包含 $ f^{-1}(-\infty) $、$ f^{-1}(\infty) $ 以及每个开子集 $ O \subset \overline{\mathbb{R}} $ 的 $ f^{-1}(O) $。我们用 $ \mathcal{L}_0(X,\mu,\overline{\mathbb{R}}) $ 表示 $ X $ 上所有 $ \mu $-可测 $ \overline{\mathbb{R}} $ 值函数之集。

\begin{remark}

\begin{items}
\item

任何实值函数 $ f: X \to \mathbb{R} $ 都可以看作 $ \overline{\mathbb{R}} $ 值函数。因此原则上适用于 $ f $ 的可测性概念有两个。但由于 $ f^{-1}(\{-\infty,\infty\}) = \varnothing $，推论 1.5 蕴含 $ f $ 作为实值函数是 $ \mu $-可测的当且仅当它作为 $ \overline{\mathbb{R}} $ 值函数是 $ \mu $-可测的。

\item

注意 $ \mathcal{L}_0(X,\mu,\overline{\mathbb{R}}) $ 不是向量空间。\nobreak\hspace{0.35em}■

\end{items}

在下一个结果中，我们列出 $ \overline{\mathbb{R}} $ 值函数的简单可测性判别法。

\end{remark}

\begin{proposition}

对 $ \overline{\mathbb{R}} $ 值函数 $ f: X \to \overline{\mathbb{R}} $，下列陈述等价：

\begin{subitems}
\item

$ f \in \mathcal{L}_0(X,\mu,\overline{\mathbb{R}}) $。

\item

$ [f < \alpha] \in \mathcal{A} $ 对每个 $ \alpha \in \mathbb{Q} $［或 $ \alpha \in \mathbb{R} $］。

\item

$ [f \le \alpha] \in \mathcal{A} $ 对每个 $ \alpha \in \mathbb{Q} $［或 $ \alpha \in \mathbb{R} $］。

\item

$ [f > \alpha] \in \mathcal{A} $ 对每个 $ \alpha \in \mathbb{Q} $［或 $ \alpha \in \mathbb{R} $］。

\item

$ [f \ge \alpha] \in \mathcal{A} $ 对每个 $ \alpha \in \mathbb{Q} $［或 $ \alpha \in \mathbb{R} $］。

\proofhead
“(i)$ \Rightarrow $(ii)” 集合 $ f^{-1}(-\infty) $ 与 $ f^{-1}\bigl((-\infty,\alpha)\bigr) $（$ \alpha \in \mathbb{Q} $［或 $ \alpha \in \mathbb{R} $］）属于 $ \mathcal{A} $。因为
\[
[f < \alpha] = f^{-1}\bigl([-\infty,\alpha)\bigr) = f^{-1}(-\infty) \cup f^{-1}\bigl((-\infty,\alpha)\bigr) ,
\]

所以 $ [f < \alpha] $ 也如此。
蕴涵 “(ii)$ \Rightarrow $(iii)$ \Rightarrow $(iv)$ \Rightarrow $(v)” 由下列恒等式推出

\[
[f \le \alpha] = \bigcap_{j=1}^{\infty} [f < \alpha + 1/j] , \quad [f > \alpha] = [f \le \alpha]^c , \quad [f \ge \alpha] = \bigcap_{j=1}^{\infty} [f > \alpha - 1/j] .
\]

“(v)$ \Rightarrow $(i)” 设 $ O $ 在 $ \overline{\mathbb{R}} $ 中开。由命题 IX.5.6，存在 $ (\alpha_j), (\beta_j) \in \mathbb{Q}^{\mathbb{N}} $ 使得 $ O = \bigcup_j [\alpha_j,\beta_j) $。因此

\[
f^{-1}(O) = \bigcup_{j \in \mathbb{N}} f^{-1}\bigl([\alpha_j,\beta_j)\bigr) = \bigcup_{j \in \mathbb{N}} \bigl([f \ge \alpha_j] \cap [f < \beta_j]\bigr) ,
\]

又因为 $ [f < \alpha] = [f \ge \alpha]^c $，我们断定 $ f^{-1}(O) $ 属于 $ \mathcal{A} $。此外

\[
f^{-1}(-\infty) = \bigcap_{j \in \mathbb{N}} [f < -j] \quad \text{且} \quad f^{-1}(\infty) = \bigcap_{j \in \mathbb{N}} [f > j] .
\]

于是 $ f^{-1}(\pm\infty) $ 也落在 $ \mathcal{A} $ 中。\nobreak\hspace{0.35em}■

\end{subitems}

\end{proposition}

\needspace{3\baselineskip}\subsection*{可测扩充实值函数构成的格}

称有序集 $ V = (V,\le) $ 为 \textbf{格}，若对每对 $ (a,b) \in V \times V $，下确界 $ a \wedge b $ 与上确界 $ a \vee b $ 都在 $ V $ 中存在。子集 $ U \subset V $ 称为 $ V $ 的 \textbf{子格}，若 $ U $ 赋以 $ V $ 诱导的序后是格。一个同时是格的有序向量空间称为 \textbf{向量格}。若向量格的向量子空间是子格，我们称它为 \textbf{向量子格}。

\begin{example}

\begin{items}
\item

设 $ V $ 是格［或向量格］。则 $ V^X $ 关于逐点序是格［或向量格］。

\item

$ \overline{\mathbb{R}} $ 是格，而 $ \mathbb{R} $ 是向量格。

\item

向量格 $ \mathbb{R}^X $ 满足

\[
f \vee g = (f + g + |f - g|)/2 , \quad f \wedge g = (f + g - |f - g|)/2 .
\]
\item

$ B(X,\mathbb{R}) $ 是 $ \mathbb{R}^X $ 的向量子格。

\item

设 $ X $ 是拓扑空间。则 $ C(X,\mathbb{R}) $ 是 $ \mathbb{R}^X $ 的向量子格。

\proofhead
这由 (c) 以及“$ f $ 连续则 $ |f| $ 连续”这一事实推出。\nobreak\hspace{0.35em}■
\item

$ \mathcal{S}(X,\mu,\mathbb{R}) $ 与 $ \mathcal{L}_0(X,\mu,\mathbb{R}) $ 是 $ \mathbb{R}^X $ 的向量子格。

\proofhead
第一个陈述是显然的。第二个由 (c) 与定理 1.7 或注 1.2(d) 推出。\nobreak\hspace{0.35em}■
\item

设 $ V $ 是向量格，$ x, y, z \in V $。则

\[
(x \vee y) + z = (x + z) \vee (y + z) ,
\]
\[
(-x) \vee (-y) = -(x \wedge y) ,
\]
\[
x + y = (x \vee y) + (x \wedge y) .
\]

\proofhead
若 $ u \in V $ 满足 $ u \ge x $ 与 $ u \ge y $，则显然 $ u + z \ge (x + z) \vee (y + z) $。因此
\[
(x \vee y) + z \ge (x + z) \vee (y + z) .
\]

设 $ v \ge (x + z) \vee (y + z) $。则 $ v - z \ge x $ 且 $ v - z \ge y $，从而 $ v \ge (x \vee y) + z $。由于这对 $ \{x + z, y + z\} $ 的每个上界 $ v $ 都成立，推出

\[
(x + z) \vee (y + z) \ge (x \vee y) + z .
\]

这证明了第一个等式。第二个不过是平凡关系

\[
\sup\{-x,-y\} = \sup(-\{x,y\}) = -\inf\{x,y\} .
\]

利用它，我们现在得到

\[
x \vee y = \bigl(-y + (x + y)\bigr) \vee \bigl(-x + (x + y)\bigr) = \bigl((-y) \vee (-x)\bigr) + (x + y) = -(x \wedge y) + (x + y) ,
\]

这证明了最后一个断言。\nobreak\hspace{0.35em}■

\item

设 $ V $ 是向量格。对 $ x \in V $，我们设

\[
x^{+} := x \vee 0 , \quad x^{-} := (-x) \vee 0 , \quad |x| := x \vee (-x) .
\]

于是我们有\footnote{见 II.8 节脚注 8。}

\[
x = x^{+} - x^{-} , \quad |x| = x^{+} + x^{-} , \quad x^{+} \wedge x^{-} = 0 .
\]

\proofhead
第一个断言容易由 (g) 推出。由此及 (g)，我们得到
\[
x^{+} + x^{-} = x + 2x^{-} = x + \bigl((-2x) \vee 0\bigr) = (-x) \vee x = |x| .
\]

类似地，我们有

\[
(x^{+} \wedge x^{-}) - x^{-} = (x^{+} - x^{-}) \wedge (x^{-} - x^{-}) = x \wedge 0 = -x^{-} ,
\]

因此 $ x^{+} \wedge x^{-} = 0 $。\nobreak\hspace{0.35em}■

\end{items}

若 $ V $ 是向量格且 $ x \in V $，我们称 $ x^{+} $ 为 $ x $ 的 \textbf{正部}、$ x^{-} $ 为 $ x $ 的 \textbf{负部}，而 $ |x| $ 是 $ x $ 的 \textbf{模}\footnote{若 $ V $ 是赋范向量空间，这不要与向量 $ x $ 的范数相混淆。$ x \in V $ 的模总是 $ V $ 中的向量，而范数是一个非负数。}。显然 $ x^{+} \ge 0 $、$ x^{-} \ge 0 $、$ |x| \ge 0 $。

下图说明向量格 $ \mathbb{R}^X $ 的元素 $ f $ 的正部与负部。

\bookfigure[0.95]{figures/43333abe4cbdb589.png}

设 $ f \in \overline{\mathbb{R}}^X $。则 $ f^{+} := f \vee 0 $ 称为 $ f $ 的 \textbf{正部}，而 $ f^{-} := 0 \vee (-f) $ 称为 $ f $ 的 \textbf{负部}。这些术语显然是比照向量格 $ \mathbb{R}^X $ 的情形选定的。\footnote{记住 $ \overline{\mathbb{R}}^X $ 是格但不是向量格。} 这里同样有

\[
f^{+} \ge 0 , \quad f^{-} \ge 0 , \quad f = f^{+} - f^{-} , \quad |f| = f^{+} + f^{-} .
\]

下一个结果表明 $ \mathcal{L}_0(X,\mu,\overline{\mathbb{R}}) $ 是 $ \overline{\mathbb{R}}^X $ 的子格，并且对可数多个格运算封闭。

\end{example}

\begin{proposition}

设 $ f \in \mathcal{L}_0(X,\mu,\overline{\mathbb{R}}) $，$ (f_j) $ 是 $ \mathcal{L}_0(X,\mu,\overline{\mathbb{R}}) $ 中的序列，$ k \in \mathbb{N} $。则下列每个 $ \overline{\mathbb{R}} $ 值函数

\[
f^{+} , \quad f^{-} , \quad |f| , \quad \max_{0 \le j \le k} f_j , \quad \min_{0 \le j \le k} f_j , \quad \sup_j f_j , \quad \inf_j f_j , \quad \overline{\lim}_j f_j , \quad \underline{\lim}_j f_j
\]

都属于 $ \mathcal{L}_0(X,\mu,\overline{\mathbb{R}}) $。

证明

\begin{subitems}
\item

设 $ \alpha \in \mathbb{R} $。由命题 1.9 知 $ [f_j > \alpha] $ 对 $ j \in \mathbb{N} $ 属于 $ \mathcal{A} $。因此

\[
[\sup_j f_j > \alpha] = \bigcup_j [f_j > \alpha]
\]

也如此，而命题 1.9 蕴含 $ \sup_j f_j $ 是 $ \mu $-可测的。

\item

由于 $ f_j $ 属于 $ \mathcal{L}_0(X,\mu,\overline{\mathbb{R}}) $，$ -f_j $ 也属于。于是由 (i) 推出函数 $ \inf_j f_j = -\sup_j(-f_j) $ 是 $ \mu $-可测的。

\item

对 $ j \in \mathbb{N} $，设

\begin{equation*}
g_j := \begin{cases} f_j & \text{若 } 0 \le j \le k , \\ f_k & \text{若 } j > k . \end{cases}
\end{equation*}

由 (i)，$ \sup_j g_j = \max_{0 \le j \le k} f_j $ 属于 $ \mathcal{L}_0(X,\mu,\overline{\mathbb{R}}) $。类似地可以证明 $ \min_{0 \le j \le k} f_j $ 是 $ \mu $-可测的。

\item

由 (iii) 推出 $ f^{+} $、$ f^{-} $ 与 $ |f| $ 属于 $ \mathcal{L}_0(X,\mu,\overline{\mathbb{R}}) $。

\item

我们有

\[
\overline{\lim}_j f_j = \inf_j \sup_{k \ge j} f_k \quad \text{且} \quad \underline{\lim}_j f_j = \sup_j \inf_{k \ge j} f_k .
\]

因此由 (i) 与 (ii)，$ \overline{\lim}_j f_j $ 与 $ \underline{\lim}_j f_j $ 也属于 $ \mathcal{L}_0(X,\mu,\overline{\mathbb{R}}) $。\nobreak\hspace{0.35em}■

\end{subitems}

\end{proposition}

$ \mathcal{S}(X,\mu,\mathbb{R}) $ 的正锥 $ \mathcal{S}(X,\mu,\mathbb{R}^{+}) $ 是所有满足 $ f(X) \subset \mathbb{R}^{+} $ 的 $ f \in \mathcal{S}(X,\mu,\mathbb{R}) $ 之集；见注 VI.4.7(b) 与 (d)。因此把它记作 $ \mathcal{S}(X,\mu,\mathbb{R}^{+}) $ 是自然的。类似地，若 $ \overline{\mathbb{R}}^{+} := [0,\infty] $ 是扩充实数轴的非负部分，我们用 $ \mathcal{L}_0(X,\mu,\overline{\mathbb{R}}^{+}) $ 表示 $ X $ 上所有处处非负的 $ \mu $-可测 $ \overline{\mathbb{R}} $ 值函数之集。

集合 $ \mathcal{L}_0(X,\mu,\overline{\mathbb{R}}^{+}) $ 有一个有趣的刻画：

\begin{theorem}

对 $ f: X \to \overline{\mathbb{R}}^{+} $，下列陈述等价：

\begin{subitems}
\item

$ f \in \mathcal{L}_0(X,\mu,\overline{\mathbb{R}}^{+}) $。

\item

存在 $ \mathcal{S}(X,\mu,\overline{\mathbb{R}}^{+}) $ 中的递增序列 $ (f_j) $ 使得 $ f_j \to f $（$ j \to \infty $）。

\proofhead
“(i)$ \Rightarrow $(ii)” 由 $ (\mathcal{A},\mu) $ 的 $ \sigma $-有限性，只需考虑 $ \mu(X) < \infty $ 的情形（比较定理 1.4 证明中的 (iv) 部分）。于是对 $ j, k \in \mathbb{N} $，设
\begin{equation*}
A_{j,k} := \begin{cases} [k2^{-j} \le f < (k+1)2^{-j}] & \text{若 } k = 0,\ldots,j2^{j} - 1 , \\ [f \ge j] & \text{若 } k = j2^{j} . \end{cases}
\end{equation*}

集合 $ A_{j,k} $ 对 $ k = 0,\ldots,j2^{j} $ 显然两两不交，且由命题 1.9 它们落在 $ \mathcal{A} $ 中。由于 $ \mu(X) < \infty $，每个 $ A_{j,k} $ 有有限测度。于是由注 1.2(b)，

\[
f_j := \sum_{k=0}^{j2^{j}} k2^{-j} \chi_{A_{j,k}} \quad \text{对 } j \in \mathbb{N}
\]

属于 $ \mathcal{S}(X,\mu,\mathbb{R}) $。此外可以验证 $ 0 \le f_j \le f_{j+1} $（$ j \in \mathbb{N} $）。
现在设 $ x \in X $。若 $ f(x) = \infty $，我们有 $ f_j(x) = j $，故 $ \lim_j f_j(x) = f(x) $。另一方面，若 $ f(x) < \infty $，则当 $ j > f(x) $ 时 $ f_j(x) \le f(x) < f_j(x) + 2^{-j} $，故此时同样有 $ \lim_j f_j(x) = f(x) $。这表明 $ (f_j) $ 逐点收敛于 $ f $。
“(ii)$ \Rightarrow $(i)” 这由命题 1.11 推出。\nobreak\hspace{0.35em}■

\end{subitems}

\end{theorem}

下面说明定理 1.12 的证明中 $ A_{j,k} $ 的构造：

\bookfigure[0.89]{figures/283135c1b9ab82fe.png}

\begin{corollary}

\begin{subitems}
\item

对每个 $ f \in \mathcal{L}_0(X,\mu,\overline{\mathbb{R}}) $，存在 $ \mathcal{S}(X,\mu,\mathbb{R}) $ 中的序列 $ (f_j) $ 使得 $ f_j \to f $。

\item

设 $ f \in \mathcal{L}_0(X,\mu,\overline{\mathbb{R}}^{+}) $ 有界。则存在 $ \mathcal{S}(X,\mu,\overline{\mathbb{R}}^{+}) $ 中的递增序列 $ (f_j) $ 一致收敛于 $ f $。

\item

设 $ (f_j) $ 是 $ \mathcal{L}_0(X,\mu,\overline{\mathbb{R}}^{+}) $ 中的序列。则 $ \sum_j f_j \in \mathcal{L}_0(X,\mu,\overline{\mathbb{R}}^{+}) $。

\proofhead
(i) 鉴于分解 $ f = f^{+} - f^{-} $，这由定理 1.12 与注 1.2(a) 推出。
\end{subitems}

\begin{subitems}
\item[（ii）]

设 $ f \in \mathcal{L}_0(X,\mu,\overline{\mathbb{R}}^{+}) $ 有界。对定理 1.12 的证明中构造的序列 $ (f_j) $，我们有

\[
f_j(x) \le f(x) < f_j(x) + 2^{-j} \quad \text{对 } j > \|f\|_\infty .
\]

因此 $ (f_j) $ 一致收敛于 $ f $。

\item[（iii）]

由定理 1.12，对每个 $ j \in \mathbb{N} $ 存在 $ \mathcal{S}(X,\mu,\overline{\mathbb{R}}^{+}) $ 中的递增序列 $ (\varphi_{j,k})_{k \in \mathbb{N}} $ 使得 $ \varphi_{j,k} \uparrow f_j $（$ k \to \infty $）。对 $ k, n \in \mathbb{N} $ 令 $ s_{k,n} := \sum_{j=0}^{k} \varphi_{j,n} $。则 $ (s_{k,n})_{n \in \mathbb{N}} $ 是 $ \mathcal{S}(X,\mu,\overline{\mathbb{R}}^{+}) $ 中的递增序列，当 $ n \to \infty $ 时收敛于 $ s_k := \sum_{j=0}^{k} f_j $。于是由定理 1.12，$ (s_k) $ 是 $ \mathcal{L}_0(X,\mu,\overline{\mathbb{R}}^{+}) $ 中的序列且 $ \lim_k s_k = \sup_k s_k = \sum_{j=0}^{\infty} f_j $。结论现在由命题 1.11 推出。\nobreak\hspace{0.35em}■

\end{subitems}

\end{corollary}

\needspace{3\baselineskip}\subsection*{可测函数的逐点极限}

设 $ (f_j) $ 是 $ \mathcal{L}_0(X,\mu,\mathbb{R}) $ 中的逐点收敛序列。由命题 1.11，$ f := \lim_j f_j $ 也属于 $ \mathcal{L}_0(X,\mu,\mathbb{R}) $。下面我们推导向量值函数序列的类似结论。

\begin{theorem}

设 $ (f_j) $ 是 $ \mathcal{L}_0(X,\mu,E) $ 中的序列，且 $ f \in E^X $。若 $ (f_j) $ 关于 $ \mu $ 几乎处处收敛于 $ f $，则 $ f $ 是 $ \mu $-可测的。

证明

\begin{subitems}
\item

先证 $ f $ 是 $ \mu $-几乎可分值的。按假设，存在 $ \mu $-零集 $ M $，使得对任意 $ x \in M^c $，当 $ j \to \infty $ 时 $ f_j(x) \to f(x) $。对每个 $ j \in \mathbb{N} $，由定理 1.4，存在 $ \mu $-零集 $ N_j $，使得 $ f_j(N_j^c) $ 可分，从而还存在在 $ f_j(N_j^c) $ 中稠密的可数集 $ B_j $：

\[
B_j \subset f_j(N_j^c) \subset \overline{B_j} \qquad \text{对 } j \in \mathbb{N} .
\]

令 $ B := \bigcup_j B_j $，由推论 III.2.13(i) 可知 $ \bigcup_j \overline{B_j} \subset \overline{B} $，并且得到

\[
\bigcup_{j\in\mathbb{N}} f_j(N_j^c) \subset \bigcup_{j\in\mathbb{N}} \overline{B_j} \subset \overline{B} .
\]

最后令 $ N := M \cup \bigcup_j N_j $。则 $ N $ 是 $ \mu $-零集，且对任意 $ k \in \mathbb{N} $ 有 $ N^c = M^c \cap \bigcap_j N_j^c \subset N_k^c $。因为对 $ x \in M^c $ 有 $ \lim_j f_j(x) = f(x) $，故

\[
f(N^c) \subset \overline{\bigcup_{j\in\mathbb{N}} f_j(N_j^c)} \subset \overline{B} = \overline{B} .
\]

因为 $ B $ 可数，注 1.3(b) 表明 $ f(N^c) $ 可分。

\item

现在证 $ f $ 是 $ \mathcal{A} $-可测的。设 $ O $ 是 $ E $ 中的开集，并对 $ n \in \mathbb{N}^\times $ 定义 $ O_n := \{ x \in O ; \operatorname{dist}(x,O^c) > 1/n \} $。如同 (1.2) 那样，可得

\[
f^{-1}(O) \cap M^c = \bigcup_{m,n\in\mathbb{N}^\times} \bigcap_{j\ge m} f_j^{-1}(O_n) \cap M^c .
\]

由定理 1.4，对每个 $ j,n \in \mathbb{N}^\times $，$ f_j^{-1}(O_n) $ 属于 $ \mathcal{A} $。因此这对 $ f^{-1}(O) \cap M^c $ 也成立。此外，$ \mu $ 的完备性蕴含 $ f^{-1}(O) \cap M $ 是 $ \mu $-零集，合起来便得

\[
f^{-1}(O) = (f^{-1}(O)\cap M^c) \cup (f^{-1}(O)\cap M) \in \mathcal{A} .
\]

现在论断由定理 1.4 得出。\nobreak\hspace{0.35em}■

\end{subitems}

\end{theorem}

\begin{remark}

定理 1.14 对不完备的测度空间一般不成立。

\proofhead
设 $ C $ 是康托尔集。在推论 IX.5.29 的证明中已说明，$ C $ 含有博雷尔不可测子集 $ N \subset C $。取 $ f_j := \chi_C $（$ j \in \mathbb{N} $）与 $ f := \chi_N $。注 1.2(b) 与 $ C $ 的紧性蕴含 $ \chi_C \in \mathcal{S}(\mathbb{R},\beta_1,\mathbb{R}) $。又对 $ x \in C^c \subset N^c $ 与 $ j \in \mathbb{N} $ 有 $ f_j(x) = f(x) $。因为 $ C $ 的测度为零，$ (f_j) $ 关于 $ \beta_1 $ 几乎处处收敛于 $ f $。然而因为 $ [f > 0] = N \notin \mathcal{B}^1 $，命题 1.9 表明 $ f $ 不属于 $ \mathcal{L}_0(\mathbb{R},\beta_1,\mathbb{R}) $。\nobreak\hspace{0.35em}■
\end{remark}

\needspace{3\baselineskip}\subsection*{拉东测度}

本节最后我们探讨向量值函数中可测性与连续性如何相互关联。除了证明一个简单的可测性判别法，我们还证明卢津定理，它揭示了连续函数与博雷尔可测函数之间出人意料的紧密联系。

称度量空间 $ X = (X,d) $ 是\textbf{σ-紧的}，若 $ X $ 局部紧，且存在由 $ X $ 的紧子集构成的序列 $ (X_j)_{j\in\mathbb{N}} $，使得 $ X = \bigcup_j X_j $。

设 $ X $ 是 σ-紧度量空间。$ X $ 上的\textbf{拉东测度}是 $ X $ 上某个 σ-代数 $ \mathcal{A} $ 上的正则、局部有限测度，其中 $ \mathcal{A} \supset \mathcal{B}(X) $。称拉东测度 $ \mu $ 是 \textbf{massive} 的，若 $ \mu $ 完备且 $ X $ 的每个非空开子集 $ O $ 都满足 $ \mu(O) > 0 $。

\begin{remark}

\begin{items}
\item

每个 σ-紧度量空间都是 IX.5 节定义意义下的 σ-紧集；然而，度量空间中紧子集的可数并却未必是 σ-紧度量空间。

\proofhead
第一个论断显然。对第二个论断，考虑 $ \mathbb{Q} \subset \mathbb{R} $。\nobreak\hspace{0.35em}■
\item

每个拉东测度都是 σ-有限的。

\proofhead
这由注 IX.5.3(b) 得出。\nobreak\hspace{0.35em}■
\item

设 $ X $ 是局部紧度量空间。则对 $ X $ 的每个紧子集 $ K $，都存在 $ K $ 的一个相对紧\footnote{若拓扑空间的子集 $ A $ 的闭包 $ \overline{A} $ 是紧的，则称 $ A $ 为\textbf{相对紧}的。}开超集。

\proofhead
对每个 $ x \in X $，我们找到 $ x $ 的一个相对紧开邻域 $ O(x) $。因为 $ K $ 紧，存在 $ x_0,\ldots,x_m \in K $，使得 $ O := \bigcup_{j=0}^m O(x_j) $ 是 $ K $ 的开超集。推论 III.2.13(iii) 蕴含 $ \overline{O} = \bigcup_{j=0}^m \overline{O(x_j)} $。因此 $ \overline{O} $ 紧。\nobreak\hspace{0.35em}■
\item

$ \mathbb{R}^n $ 的每个开子集都是 σ-紧度量空间。

\proofhead
设 $ X \subset \mathbb{R}^n $ 开且非空。对每个 $ x \in X $，存在 $ r > 0 $ 使得 $ \mathbb{B}(x,r) \subset X $。因为 $ \overline{\mathbb{B}}(x,r) $ 紧，$ X $ 便是局部紧度量空间。对 $ j \in \mathbb{N}^\times $，定义\footnote{这里约定 $ \operatorname{dist}(x,\varnothing) := \infty $。}
\begin{equation}
U_j := \{ x \in X ; \operatorname{dist}(x,U^c) > 1/j \} \cap \mathbb{B}(0,j) .\tag{1.3}
\end{equation}

由例 III.1.3(l) 与 III.2.22(c)，集合 $ U_j $ 开。又 $ U_j \subset \overline{U_j} \subset U_{j+1} $，且 $ \bigcup_j \overline{U_j} \subset \bigcup_j U_j = X $。特别地，存在 $ j_0 \in \mathbb{N}^\times $，使得对 $ j \ge j_0 $ 有 $ U_j \ne \varnothing $。因为由海涅–博雷尔定理 $ \overline{U_j} $ 紧，论断得证。\nobreak\hspace{0.35em}■

\item

对局部紧度量空间 $ X $，下列陈述等价：

\begin{subitems}
\item

$ X $ 是 σ-紧的。

\item

$ X $ 是可数个相对紧开子集构成的序列 $ (U_j)_{j\in\mathbb{N}} $ 的并，且对 $ j \in \mathbb{N} $ 有 $ \overline{U_j} \subset U_{j+1} $。

\item

$ X $ 是林德勒夫空间。

\item

$ X $ 满足第二可数公理。

\item

$ X $ 可分。

\proofhead
“（i）$ \Rightarrow $（ii）” 设 $ (X_j)_{j\in\mathbb{N}} $ 是 $ X $ 中紧集构成的序列，使得 $ X = \bigcup_j X_j $。由 (c)，$ X_0 $ 有一个相对紧开超集 $ U_0 $。归纳地选取相对紧开子集 $ U_j $，使得对 $ j \ge 1 $ 有 $ U_j \supset \overline{U_{j-1}} \cup X_j $。显然 $ X = \bigcup_j U_j $。
“（ii）$ \Rightarrow $（iii）” 设 $ \mathcal{O} := \{ O_\alpha ; \alpha \in A \} $ 是 $ X $ 的开覆盖。对每个 $ j \in \mathbb{N} $，归纳地选取 $ m(j) \in \mathbb{N} $ 与 $ \alpha_0,\ldots,\alpha_{m(j)} \in A $，使得 $ \overline{U_j} \subset \bigcup_{k=0}^{m(j)} O_{\alpha_k} $。则 $ \{ O_{\alpha_k} ; k = 0,\ldots,m(j),\ j \in \mathbb{N} \} $ 是 $ \mathcal{O} $ 对 $ X $ 的可数子覆盖。
“（iii）$ \Rightarrow $（i）” 按假设，存在 $ X $ 中的序列 $ (x_j) $ 与 $ x_j $ 的相对紧开邻域 $ O(x_j) $（$ j \in \mathbb{N} $），使得 $ X = \bigcup_{j\in\mathbb{N}} O(x_j) $。于是 $ X = \bigcup_{j\in\mathbb{N}} \overline{O(x_j)} $，这表明 $ X $ 是 σ-紧的。
其余等价关系由命题 IX.1.8 得出。\nobreak\hspace{0.35em}■
\end{subitems}

\item

σ-紧度量空间上的每个局部有限博雷尔测度都是正则的，因而也是拉东测度。

\proofhead
这由 (e) 与 [Els99] 中的推论 VIII.1.12 得出。\nobreak\hspace{0.35em}■
\item

在（非可度量化）紧拓扑空间上，有限博雷尔测度未必正则。

\proofhead
见 [Flo81, Example A4.5, S. 350]。\nobreak\hspace{0.35em}■
\item

勒贝格 n-测度 $ \lambda_n $ 是 $ \mathbb{R}^n $ 上的 massive 拉东测度。

\proofhead
这由定理 IX.5.1 与 IX.5.4 得出。\nobreak\hspace{0.35em}■
\item

s 维豪斯多夫测度 $ \mathcal{H}^s $ 仅在 $ s \ge n $ 时才是 $ \mathbb{R}^n $ 上的拉东测度。它是 massive 的当且仅当 $ s = n $。

\proofhead
由例 IX.4.4(c) 与定理 IX.4.3，每个博雷尔集都是 $ \mathcal{H}^s $-可测的。$ s > 0 $ 时 $ \mathcal{H}^s $ 的正则性由推论 IX.5.22 与定理 IX.5.4 得出。
设 $ O $ 是 $ \mathbb{R}^n $ 中的非空开集。因为 $ O $ 的豪斯多夫维数为 $ n $（习题 IX.3.6），可得
\begin{equation*}
\mathcal{H}^s(O) = \begin{cases} 0 & \text{若 } s > n , \\ \infty & \text{若 } s < n . \end{cases}
\end{equation*}

因此若 $ s < n $，$ \mathcal{H}^s $ 不可能是 $ \mathbb{R}^n $ 上的拉东测度。另一方面，若 $ s > n $，则 $ \mathcal{H}^s $ 是非 massive 的拉东测度。
引理 IX.5.21 表明 $ \mathcal{H}^n $ 局部有限，因而是 $ \mathbb{R}^n $ 上的拉东测度。最后，推论 IX.5.22 蕴含 $ \mathcal{H}^n(O) > 0 $，证毕。\nobreak\hspace{0.35em}■

\item

设 $ F: \mathbb{R} \to \mathbb{R} $ 是测度生成函数，记 $ \mu_F $ 为 $ F $ 诱导的 $ \mathbb{R} $ 上的勒贝格–斯蒂尔杰斯测度。则 $ \mu_F $ 是 $ \mathbb{R} $ 上的拉东测度，且它是 massive 的当仅当 $ F $ 严格递增。

\proofhead
这由例 IX.4.4(b)、定理 IX.4.3、习题 IX.5.19 与命题 IX.3.5 得出。\nobreak\hspace{0.35em}■
\end{items}

\end{remark}

\begin{theorem}

设 $ \mu $ 是 $ X $ 上的完备拉东测度。则 $ C(X,E) $ 是 $ \mathcal{L}_0(X,\mu,E) $ 的向量子空间。

\proofhead
取 $ f \in C(X,E) $，设 $ (X_j) $ 是 $ X $ 中紧集构成的序列，使得 $ X = \bigcup_j X_j $。据习题 IX.1.6(b)，$ f $ 是博雷尔可测的，从而是 $ \mathcal{A} $-可测的，其中 $ \mathcal{A} $ 是 $ \mu $ 的定义域。由注 1.16(e)，$ f(X_j) $ 作为 $ E $ 的紧子集可分。因此 $ f(X) = \bigcup_j f(X_j) $ 也可分，论断由定理 1.4 得出。\nobreak\hspace{0.35em}■
\end{theorem}

\begin{theorem}

（卢津） 设 $ X $ 是 σ-紧度量空间，$ \mu $ 是 $ X $ 上的完备拉东测度，$ f \in \mathcal{L}_0(X,\mu,E) $。则对每个有限测度的 $ \mu $-可测集 $ A $ 与每个 $ \varepsilon > 0 $，存在 $ X $ 的紧子集 $ K $，使得 $ \mu(A\setminus K) < \varepsilon $ 且 $ f \mid K \in C(K,E) $。

证明

\begin{subitems}
\item

因为 $ X $ 是 σ-紧的，我们能找到紧集 $ \widetilde{X} $，使得 $ \mu(A\setminus \widetilde{X}) < \varepsilon/2 $。令 $ \widetilde{f} := f \mid \widetilde{X} $，$ \widetilde{A} := A \cap \widetilde{X} $。则 $ \mu(\widetilde{X}) < \infty $。

\item

由定理 1.4，存在 $ \widetilde{X} $ 的 $ \mu $-零集 $ N $，使得 $ \widetilde{f}(N^c) $ 可分。于是由命题 IX.1.8，存在 $ \widetilde{f}(N^c) $ 的可数基 $ \{ \widetilde{V_j} ; j \in \mathbb{N} \} $，并且由命题 III.2.26，存在 $ E $ 中的开子集 $ V_j $，使得 $ \widetilde{V_j} = V_j \cap \widetilde{f}(N^c) $。

\item

据定理 1.4，对每个 $ j \in \mathbb{N} $，$ \widetilde{f}^{-1}(V_j) $ 是 $ \mu $-可测的。因此由 $ \mu $ 的正则性与 $ \mu(\widetilde{X}) $ 的有限性可知，对每个 $ j \in \mathbb{N} $，存在紧集 $ K_j $ 与开集 $ U_j $，使得 $ K_j \subset \widetilde{f}^{-1}(V_j) \subset U_j $ 且 $ \mu(U_j \setminus K_j) < \varepsilon 2^{-(j+3)} $。令 $ U := \bigcup_j (U_j \setminus K_j) $，则 $ \mu(U) < \varepsilon/4 $。

\item

令 $ Y := (U \cup N)^c $，并证明 $ \widetilde{f} \mid Y $ 连续。为此，设 $ V $ 是 $ E $ 中的开集。则存在 $ \{ V_j ; j \in \mathbb{N} \} $ 的子集 $ \{ V_{j_k} ; k \in \mathbb{N} \} $，使得 $ V \cap \widetilde{f}(N^c) = \bigcup_k V_{j_k} \cap \widetilde{f}(N^c) $。这意味着

\[
\widetilde{f}^{-1}(V) \cap N^c = \bigcup_k \widetilde{f}^{-1}(V_{j_k}) \cap N^c .
\]

显然对 $ \ell \in \mathbb{N} $ 有 $ \widetilde{f}^{-1}(V_\ell) \cap Y \subset U_\ell \cap Y $。因为

\[
Y = U^c \cap N^c = \bigcap_j (U_j^c \cup K_j) \cap N^c \subset \bigcap_j (U_j^c \cup \widetilde{f}^{-1}(V_j)) \subset U_\ell^c \cup \widetilde{f}^{-1}(V_\ell) ,
\]

故 $ \widetilde{f}^{-1}(V_\ell) \cap Y = U_\ell \cap Y $，并且得到

\[
(\widetilde{f} \mid Y)^{-1}(V) = \widetilde{f}^{-1}(V) \cap N^c \cap U^c = \bigcup_k U_{j_k} \cap Y .
\]

因为 $ \bigcup_k U_{j_k} $ 在 $ X $ 中开，从而 $ \bigcup_k U_{j_k} \cap Y $ 在 $ Y $ 中开，所以 $ \widetilde{f} \mid Y $ 连续。

\item

我们再次应用 $ \mu $ 的正则性，推得 $ \mu $-可测集 $ Y $ 存在紧子集 $ K $，使得 $ \mu(Y \setminus K) < \varepsilon/4 $。然后
$ \widetilde{f} \mid K $ 属于 $ C(K,E) $，并且

\[
\mu(\widetilde{A}\setminus K) \le \mu(Y\setminus K) + \mu(Y^{c}\setminus K) \le \mu(Y\setminus K) + \mu(U) < \varepsilon/2 .
\]

因为 $ \mu(A\setminus K) \le \mu(\widetilde{A}\setminus K) + \mu(A\setminus\widetilde{X}) < \varepsilon $，于是完成证明。\nobreak\hspace{0.35em}■

\end{subitems}

\end{theorem}

\exerciseshead

\begin{exercises}
\item

设 $ H $ 是可分的希尔伯特空间。若对每个 $ e \in H $，$ (f \mid e) $ 都属于 $ \mathcal{L}_0(X,\mu,\mathbb{K}) $，则称 $ f \in H^{X} $ 是\textbf{弱 μ-可测的}。证明：

\begin{items}
\item

若 $ f $ 是弱 μ-可测的，则 $ |f| $ 是 μ-可测的。

\item

$ f $ 是 μ-可测的当且仅当 $ f $ 是弱 μ-可测的。
（提示：(a) 设 $ \{e_j \ ;\ j \in \mathbb{N}\} $ 是 $ \overline{\mathbb{B}}_H $ 的稠密子集。则

\[
[\,|f| \le \alpha\,] = \bigcap_{j}\big[\,|(f \mid e_j)| \le \alpha\,\big] \qquad \text{对 } \alpha \in \mathbb{R} .
\]
\end{items}

（b）“$ \Leftarrow $” 利用 (a)，我们可以像定理 1.4 的证明那样构造一列 μ-简单函数，使其 μ-几乎处处收敛到 $ f $。）

\item

记 $ \mathcal{S}(\mathbb{R},E) $ 为 $ \mathbb{R} $ 上所有 $ E $ 值容许函数构成的向量空间（见第 VI.8 节）。证明或否定：

\begin{items}
\item

$ \mathcal{S}(\mathbb{R},E) \subset \mathcal{L}_0(\mathbb{R},\beta_1,E) $；

\item

$ \mathcal{S}(\mathbb{R},E) \supset \mathcal{L}_0(\mathbb{R},\beta_1,E) $。

\end{items}

\item

证明注 1.2(e) 的陈述。

\item

证明每个单调的扩充实值函数都是博雷尔可测的。

\item

设 $ f, g \in \mathcal{L}_0(X,\mu,\overline{\mathbb{R}}) $。证明集合 $ [f<g] $、$ [f\le g] $、$ [f=g] $ 与 $ [f\ne g] $ 都属于 $ \mathcal{A} $。

\item

设 $ (f_j) $ 是 $ \mathcal{L}_0(X,\mu,\overline{\mathbb{R}}) $ 中的序列。证明

\[
K := \{\, x \in X \ ;\ \lim_{j} f_{j}(x) \text{ 在 } \overline{\mathbb{R}} \text{ 中存在} \,\}
\]

是 μ-可测的。

\item

设 $ f: X \to \overline{\mathbb{R}} $。证明或否定：

\begin{items}
\item

$ f \in \mathcal{L}_0(X,\mu,\overline{\mathbb{R}}) \iff f^{+}, f^{-} \in \mathcal{L}_0(X,\mu,\overline{\mathbb{R}}^{+}) $。

\item

$ f \in \mathcal{L}_0(X,\mu,\overline{\mathbb{R}}) \iff |f| \in \mathcal{L}_0(X,\mu,\overline{\mathbb{R}}^{+}) $。

\end{items}

\item

设 $ B $ 是 $ \overline{\mathbb{R}}^{X} $ 的非空子集，若下列陈述成立，则称 $ B $ 为\textbf{贝尔函数空间}：

\begin{subitems}
\item

$ \alpha \in \mathbb{R} $ 且 $ f \in B $ 蕴含 $ \alpha f \in B $。

\item

若对 $ f, g \in B $，$ f+g $ 在 $ \overline{\mathbb{R}}^{X} $ 中存在，则 $ f+g \in B $。

\item

对 $ B $ 中的每个序列 $ (f_j) $，$ \sup_j f_j $ 都属于 $ B $。
证明：

\end{subitems}

\begin{items}
\item

$ \overline{\mathbb{R}}^{X} $ 与 $ \mathcal{L}_0(X,\mu,\overline{\mathbb{R}}) $ 都是贝尔函数空间。

\item

若 $ \{ B_\alpha \subset \overline{\mathbb{R}}^{X} \ ;\ \alpha \in \mathsf{A} \} $ 是一族贝尔函数空间，则 $ \bigcap_{\alpha \in \mathsf{A}} B_\alpha $ 也是贝尔函数空间。

\end{items}

\item

对 $ C \subset \overline{\mathbb{R}}^{X} $，称

\[
\sigma(C) := \bigcap \{\, B \subset \overline{\mathbb{R}}^{X} \ ;\ B \supset C,\ B \text{ 是贝尔函数空间} \,\}
\]

为由 $ C $ 生成的\textbf{贝尔函数空间}。由习题 8(b)，$ \sigma(C) $ 是良定义的贝尔函数空间。证明

\[
\sigma\big(\mathcal{S}(X,\mu,\mathbb{R})\big) = \mathcal{L}_0(X,\mu,\overline{\mathbb{R}}) .
\]
\item

证明 $ \sigma\big(C(\mathbb{R}^{n},\mathbb{R})\big) = \mathcal{L}_0(\mathbb{R}^{n},\beta_n,\mathbb{R}) $。

\item

证明不可数族可测实值函数的上确界一般不是可测的。

\item

若对每个 $ \delta>0 $ 都存在 $ A \in \mathcal{A} $ 使 $ \mu(A^{c})<\delta $ 且序列 $ (f_j \mid A) $ 一致收敛，则称 $ E^{X} $ 中的序列 $ (f_j) $ 是 \textbf{μ-几乎一致收敛的}。

\begin{items}
\item

设 $ (f_j) $ 是 $ \mathcal{L}_0(X,\mu,E) $ 中 μ-几乎一致收敛的序列。证明存在 $ f \in \mathcal{L}_0(X,\mu,E) $ 使 $ f_j \to f $ μ-几乎处处成立。

\item

定义 $ f_j(x) := x^{j} $（$ j \in \mathbb{N} $，$ x \in [0,1] $）。验证 $ (f_j) $ 几乎一致收敛（关于勒贝格测度），尽管不存在测度为零的集合 $ N \subset [0,1] $ 使 $ (f_j \mid N^{c}) $ 一致收敛。

\end{items}

\item

设 $ (X,\mathcal{A},\mu) $ 是有限测度空间，$ f_j, f \in \mathcal{L}_0(X,\mu,E) $，且 $ f_j \to f $ μ-几乎处处成立。证明：

\begin{items}
\item

对 $ \varepsilon>0 $ 与 $ \delta>0 $，存在 $ k \in \mathbb{N} $ 与 $ A \in \mathcal{A} $，使得 $ \mu(A^{c})<\delta $，且对 $ x \in A $ 与 $ j \ge k $ 有 $ |f_j(x)-f(x)|<\varepsilon $。

\item

序列 $ (f_j) $ μ-几乎一致收敛到 $ f $（\textbf{叶戈罗夫定理}）。

\item

当 $ \mu(X)=\infty $ 时，(b) 一般不再成立。
（提示：(a) 考虑 $ K := [f_j \to f] $ 与 $ K_k := [\,|f_j-f|<\varepsilon \ ;\ j \ge k\,] $，并利用测度的上连续性。(b) 为得到 $ A_j $，在 (a) 中取 $ \varepsilon := 1/j $ 与 $ \delta := \delta 2^{-j} $，并令 $ A := \bigcup_j A_j $。(c) 考虑测度空间 $ (X,\mathcal{A},\mu) = (\mathbb{R},\lambda_1,\mathcal{L}(1)) $，并令 $ f_j := \chi_{[j,j+1]} $。）

\end{items}

\item

设 $ (X,\mathcal{A},\mu) $ 是测度空间，$ f_j, f \in \mathcal{L}_0(X,\mu,E) $。若对每个 $ \varepsilon>0 $ 都有 $ \lim_{j\to\infty} \mu\big(\big[\,|f_j-f|\ge\varepsilon\,\big]\big)=0 $，则称 $ (f_j) $ \textbf{依测度收敛}于 $ f $。证明：

\begin{items}
\item

$ f_j \to f $ μ-几乎一致收敛 $ \Rightarrow $ $ f_j \to f $ 依测度收敛。

\item

若 $ (f_j) $ 依测度收敛到 $ f $，也收敛到 $ g $，则 $ f=g $ μ-几乎处处成立。

\item

存在 $ [0,1] $ 上的勒贝格可测函数序列，它依测度收敛，但在任何点处都不逐点收敛。

\item

存在 $ \mathbb{R} $ 上的勒贝格可测函数序列，它逐点收敛，但不依测度收敛。
（提示：(c) 令 $ f_j := \chi_{I_j} $，其中区间 $ I_j \subset [0,1] $ 的选取使得 $ \lambda_1(I_j) \to 0 $，且对每个 $ x \in [0,1] $，序列 $ (f_j(x)) $ 都有两个聚点。(d) 考虑 $ f_j := \chi_{[j,j+1]} $。）

\end{items}

\item

设 $ (f_j) $ 是 $ \mathcal{L}_0(X,\mu,E) $ 中依测度收敛到 $ f \in \mathcal{L}_0(X,\mu,E) $ 的序列。证明 $ (f_j) $ 有子列 μ-几乎处处收敛到 $ f $。

（提示：存在递增序列 $ (j_k)_{k\in\mathbb{N}} $，使得

\[
\mu\big(\big[\,|f_m-f_n|\ge 2^{-k}\,\big]\big) \le 2^{-k} \qquad \text{对 } m, n \ge j_k .
\]

借助 $ B_\ell := \bigcup_{k=\ell}^{\infty}\big[\,|f_{n_{k+1}}-f_{n_k}|\ge 2^{-k}\,\big] $，可推出 $ (f_{j_k})_{k\in\mathbb{N}} $ μ-几乎一致收敛。另见习题 12、14(a) 与 14(b)。）

\item

对 $ x = (x_j) \in \mathbb{K}^{\mathbb{N}} $ 与 $ p \in [1,\infty] $，定义\footnote{另见命题 IV.2.17。}

\begin{equation*}
\|x\|_p := \begin{cases} \Big(\sum_{j=0}^{\infty} |x_j|^{p}\Big)^{1/p} & \text{若 } p \in [1,\infty), \\ \sup_j |x_j| & \text{若 } p = \infty, \end{cases}
\end{equation*}

以及

\[
\ell_p := \ell_p(\mathbb{K}) := \big( \{\, x \in \mathbb{K}^{\mathbb{N}} \ ;\ \|x\|_p < \infty \,\},\ \|\cdot\|_p \big) .
\]

证明：

\begin{items}
\item

若 $ p \in [1,\infty) $，则 $ \ell_p $ 是可分的赋范向量空间。

\item

$ \ell_\infty $ 不可分。

\end{items}

\item

设 $ x \in \overline{\mathbb{R}} $。若 $ x \in \mathbb{R} $，则称 $ U \subset \overline{\mathbb{R}} $ 是 $ x $ 在 $ \overline{\mathbb{R}} $ 中的\textbf{邻域}，是指 $ U $ 包含 $ x $ 在 $ \mathbb{R} $ 中的一个邻域。对 $ x \in \overline{\mathbb{R}}\setminus\mathbb{R} $，邻域在第 II.5 节中定义为那些包含相应类型的半无限区间的集合。设 $ O \subset \overline{\mathbb{R}} $。若每个 $ x \in O $ 都有 $ \overline{\mathbb{R}} $ 中的邻域 $ U $ 使 $ U \subset O $，则称 $ O $ 在 $ \overline{\mathbb{R}} $ 中是\textbf{开的}。现定义 $ \overline{\mathcal{T}} := \{ O \subset \overline{\mathbb{R}} \ ;\ O \text{ 在 } \overline{\mathbb{R}} \text{ 中是开的} \} $。证明：

\begin{items}
\item

$ O $ 在 $ \overline{\mathbb{R}} $ 中是开的，当且仅当 $ O \cap \mathbb{R} $ 在 $ \mathbb{R} $ 中是开的，并且当 $ \infty \in O $（或 $ -\infty \in O $）时，存在 $ a \in \mathbb{R} $ 使 $ (a,\infty] \subset O $（或 $ [-\infty,a) \subset O $）。

\item

$ (\overline{\mathbb{R}},\overline{\mathcal{T}}) $ 是紧拓扑空间。

\item

$ \mathcal{B}(\overline{\mathbb{R}}) = \{\, B \cup F \ ;\ B \in \mathcal{B}^{1},\ F \subset \{-\infty,\infty\} \,\} $。

\item

$ \mathcal{B}(\overline{\mathbb{R}}) \mid \mathbb{R} = \mathcal{B}^{1} $。

\item

$ \overline{\mathbb{R}}^{X} $ 中的元素属于 $ \mathcal{L}_0(X,\mu,\overline{\mathbb{R}}) $ 当且仅当它是 $ \mathcal{A} \text{-} \mathcal{B}(\overline{\mathbb{R}}) $-可测的。

\end{items}

\item

验证：若 $ S $ 是 $ E $ 的可分子集，则其张成的闭包是可分的巴拿赫空间。

\item

对 $ f \in \mathbb{K}^{X} $，定义

\begin{equation*}
(\operatorname{sign} f)(x) := \begin{cases} f(x)/|f(x)| & \text{若 } f(x) \ne 0, \\ 0 & \text{若 } f(x) = 0. \end{cases}
\end{equation*}

证明 $ f \in \mathcal{L}_0(X,\mu,\mathbb{K}) $ 蕴含 $ \operatorname{sign} f \in \mathcal{L}_0(X,\mu,\mathbb{K}) $。

\end{exercises}

\section{可积函数}

本节我们定义一般的 Bochner–Lebesgue 积分，并叙述它的基本性质。我们还证明，可积函数构成的向量空间关于由该积分诱导的半范数是完备的。

与上一节一样，假设

- $ (X, \mathcal{A}, \mu) $ 是完备的 $ \sigma $-有限测度空间；$ E = (E, \left|\cdot\right|) $ 是巴拿赫空间。

\needspace{3\baselineskip}\subsection*{简单函数的积分}

在注 1.2(c) 中我们已经知道，每个简单函数都有唯一的规范形式。这个形式在下文中很有用，我们优先使用它。

约定 我们总是用规范形式表示 $ \mu $-简单函数，除非另有说明。此外，我们令\footnote{约定 (2.1) 在积分论中很常见，它使例如在简单函数的整个定义域上对其积分成为可能。在 $ E = \mathbb{R} $ 的情形，这\textbf{不}应理解为 $ \overline{\mathbb{R}} $ 中的又一条计算规则，而应理解为 $ \infty, -\infty \in \overline{\mathbb{R}} $ 与 $ \mathbb{R} $ 的零向量的“外部”乘法。}

\begin{equation}
\infty \cdot 0_E := -\infty \cdot 0_E := 0_E ,\tag{2.1}
\end{equation}

其中 $ 0_E $ 是 $ E $ 的零向量。

对 $ \varphi = \sum_{j=0}^{m} e_j \chi_{A_j} \in \mathcal{S}(X,\mu,E) $，我们把\textbf{$ \varphi $ 在 $ X $ 上关于测度 $ \mu $ 的积分}定义为和

\[
\int_X \varphi \, d\mu := \int \varphi \, d\mu := \sum_{j=0}^{m} e_j \mu(A_j) .
\]

若 $ A $ 是 $ \mu $-可测集，我们把\textbf{$ \varphi $ 在 $ A $ 上关于测度 $ \mu $ 的积分}定义为

\[
\int_A \varphi \, d\mu := \int_X \chi_A \varphi \, d\mu .
\]
\bookfigure[0.74]{figures/76a0c233b03a4dde.png}

\begin{remark}

\begin{items}
\item

对 $ \varphi \in \mathcal{S}(X,\mu,E) $ 和 $ A \in \mathcal{A} $，积分 $ \int_A \varphi \, d\mu $ 是良定义的。

\proofhead
这由注 1.2(c) 与 (d) 推出。\nobreak\hspace{0.35em}■
\end{items}

\begin{items}
\item[（b）]

设 $ \varphi = \sum_{k=0}^{n} f_k \chi_{B_k} $，其中 $ f_0, \ldots, f_n \in E $ 非零，$ B_0, \ldots, B_n \in \mathcal{A} $ 两两不交，是一个不必处于规范形式的 $ \mu $-简单函数。则

\[
\int_X \varphi \, d\mu = \sum_{k=0}^{n} f_k \mu(B_k) .
\]

\proofhead
设 $ \sum_{j=0}^{m} e_j \chi_{A_j} $ 是 $ \varphi $ 的规范形式。令
\begin{equation}
A_{m+1} := \bigcap_{j=0}^{m} A_j^{c} , \quad B_{n+1} := \bigcap_{k=0}^{n} B_k^{c} , \quad e_{m+1} := 0 , \quad f_{n+1} := 0 .\tag{2.2}
\end{equation}

于是 $ X = \bigcup_{j=0}^{m+1} A_j = \bigcup_{k=0}^{n+1} B_k $，故

\[
A_j = \bigcup_{k=0}^{n+1} (A_j \cap B_k) \quad \text{且} \quad B_k = \bigcup_{j=0}^{m+1} (A_j \cap B_k) \quad \text{对 } j = 0, \ldots, m+1 \text{ 和 } k = 0, \ldots, n+1 .
\]

因为集合 $ A_j \cap B_k $ 两两不交，我们有

\[
\mu(A_j) = \sum_{k=0}^{n+1} \mu(A_j \cap B_k) \quad \text{且} \quad \mu(B_k) = \sum_{j=0}^{m+1} \mu(A_j \cap B_k) .
\]

若 $ A_j \cap B_k \ne \varnothing $，则 $ e_j = f_k $，并且我们得到

\[
\int_X \varphi \, d\mu = \sum_{j=0}^{m} e_j \mu(A_j) = \sum_{j=0}^{m+1} e_j \sum_{k=0}^{n+1} \mu(A_j \cap B_k) = \sum_{k=0}^{n+1} f_k \sum_{j=0}^{m+1} \mu(A_j \cap B_k)
\]
\[
= \sum_{k=0}^{n} f_k \mu(B_k) .\quad\blacksquare
\]
\item[（c）]

积分 $ \int_X \cdot \, d\mu : \mathcal{S}(X,\mu,E) \to E $ 是线性的。

\proofhead
设 $ \varphi = \sum_{j=0}^{m} e_j \chi_{A_j} $ 与 $ \psi = \sum_{k=0}^{n} f_k \chi_{B_k} $ 是 $ \mu $-简单函数，$ \alpha \in \mathbb{K} $。容易验证 $ \int_X \alpha \varphi \, d\mu = \alpha \int_X \varphi \, d\mu $。利用关系式 (2.2)，像 (b) 中那样论证可得
\[
\chi_{A_j} = \sum_{k=0}^{n+1} \chi_{A_j \cap B_k} \quad \text{且} \quad \chi_{B_k} = \sum_{j=0}^{m+1} \chi_{A_j \cap B_k} ,
\]

于是我们有\footnote{一般而言，(2.3) \textbf{不}给出 $ \varphi + \psi $ 的规范形式。}

\begin{equation}
\varphi + \psi = \sum_{j=0}^{m+1} \sum_{k=0}^{n+1} (e_j + f_k) \chi_{A_j \cap B_k} .\tag{2.3}
\end{equation}

结论现在由 (b) 推出。\nobreak\hspace{0.35em}■

\item[（d）]

对 $ A, B \in \mathcal{A} $ 且 $ A \cap B = \varnothing $，我们有

\[
\int_{A \cup B} \varphi \, d\mu = \int_A \varphi \, d\mu + \int_B \varphi \, d\mu \quad \text{对 } \varphi \in \mathcal{S}(X,\mu,E) .
\]

\proofhead
这由 (c) 以及等式 $ \chi_{A \cup B} \varphi = \chi_A \varphi + \chi_B \varphi $ 推出。\nobreak\hspace{0.35em}■
\item[（e）]

对 $ \varphi \in \mathcal{S}(X,\mu,E) $ 和 $ A \in \mathcal{A} $，我们有

\[
\left| \int_A \varphi \, d\mu \right| \le \int_A |\varphi| \, d\mu \le \|\varphi\|_\infty \, \mu(A) .
\]

\proofhead
这由注 1.2(d) 与三角不等式推出。\nobreak\hspace{0.35em}■
\item[（f）]

若 $ \varphi, \psi \in \mathcal{S}(X,\mu,\mathbb{R}) $ 满足 $ \varphi \le \psi $，则 $ \int_A \varphi \, d\mu \le \int_A \psi \, d\mu $。

\proofhead
显然 $ \int_A \eta \, d\mu \ge 0 $ 对 $ \eta \in \mathcal{S}(X,\mu,\mathbb{R}^{+}) $。结论现在由 (c) 推出。\nobreak\hspace{0.35em}■
\end{items}

\end{remark}

\needspace{3\baselineskip}\subsection*{$ \mathcal{L}_1 $-半范数}

设 $ V $ 是 $ \mathbb{K} $ 上的向量空间。若映射 $ p : V \to \mathbb{R} $ 满足以下性质，则称它是 $ V $ 上的\textbf{半范数}：

\begin{subitems}
\item

$ p(v) \ge 0 $ 对 $ v \in V $；

\item

$ p(\lambda v) = |\lambda| \, p(v) $ 对 $ v \in V $ 和 $ \lambda \in \mathbb{K} $；

\item

$ p(v + w) \le p(v) + p(w) $ 对 $ v, w \in V $。
对 $ v \in V $ 和 $ r > 0 $，我们用

\[
\mathbb{B}_p(v,r) := \{ w \in V ; \; p(v - w) < r \}
\]

记 $ (V,p) $ 中以 $ v $ 为心、半径为 $ r $ 的\textbf{开球}。称 $ V $ 的子集 $ O $ 是 \textbf{p-开}的，若对每个 $ v \in O $ 都存在 $ r > 0 $ 使得 $ \mathbb{B}_p(v,r) \subset O $。

\end{subitems}

\begin{remark}

设 $ V $ 是向量空间，$ p $ 是 $ V $ 上的半范数。

\begin{items}
\item

半范数 $ p $ 是范数当且仅当 $ p^{-1}(0) = \{0\} $。

\item

设 $ K \subset \mathbb{R}^n $ 紧，$ k \in \mathbb{N} \cup \{\infty\} $，且

\[
p_K(f) := \max_{x \in K} |f(x)| \quad \text{对 } f \in C^{k}(\mathbb{R}^n, E) .
\]

则 $ p_K $ 是 $ C^{k}(\mathbb{R}^n, E) $ 上的半范数，但不是范数。

\proofhead
容易验证 $ p_K $ 是 $ C^{k}(\mathbb{R}^n, E) $ 上的半范数。设 $ U $ 是 $ K $ 的开邻域。则习题 VII.6.7 表明存在 $ f \in C^{\infty}(\mathbb{R}^n, \mathbb{R}) $，使得 $ f(x) = 1 $ 对 $ x \in K $ 且 $ f(x) = 0 $ 对 $ x \in U^{c} $。对 $ e \in E \setminus \{0\} $，令 $ g := (\chi_{\mathbb{R}^n} - f)e $。则 $ g $ 属于 $ C^{\infty}(\mathbb{R}^n, E) $，并且 $ p_K(g) = 0 $，但 $ g \ne 0 $。因此 $ p_K $ 不是 $ C^{k}(\mathbb{R}^n, E) $ 上的范数。\nobreak\hspace{0.35em}■
\item

令

\[
\|\varphi\|_1 := \int_X |\varphi| \, d\mu \quad \text{对 } \varphi \in \mathcal{S}(X,\mu,E) .
\]

则 $ \|\cdot\|_1 $ 是 $ \mathcal{S}(X,\mu,E) $ 上的半范数。若 $ \mathcal{A} $ 中存在非空的 $ \mu $-零集，则 $ \|\cdot\|_1 $ 不是 $ \mathcal{S}(X,\mu,E) $ 上的范数。

\proofhead
显然 $ \|\cdot\|_1 $ 是 $ \mathcal{S}(X,\mu,E) $ 上的半范数。取非空 $ \mu $-零集 $ N $，我们有 $ \|\chi_N\|_1 = 0 $，但 $ \chi_N \ne 0 $。\nobreak\hspace{0.35em}■
\item

$ \mathcal{T}_p := \{ O \subset V ; \; O \text{ 是 } p\text{-开} \} $ 是 $ V $ 上的拓扑，称为\textbf{由 $ p $ 生成的拓扑}。

\proofhead
容易验证，命题 III.2.4 的证明中所用的论证可以移植到此情形。\nobreak\hspace{0.35em}■
\item

拓扑 $ \mathcal{T}_p $ 不必是豪斯多夫的。若不是，则 $ V $ 上不存在生成 $ \mathcal{T}_p $ 的度量。

\proofhead
我们使用 (b) 的记号并令 $ K := \{0\} $。再设 $ f \in C^{k}(\mathbb{R}^n, E) $ 满足 $ f(0) = 0 $ 且 $ f \ne 0 $。则 $ \mathbb{B}_{p_K}(f,\varepsilon) = \mathbb{B}_{p_K}(0,\varepsilon) $ 对每个 $ \varepsilon > 0 $。因此 $ \mathcal{T}_{p_K} $ 不是豪斯多夫的。第二个结论由命题 III.2.17 推出。\nobreak\hspace{0.35em}■
\item

称线性映射 $ A : V \to E $ 是 \textbf{(p)-有界}的，若存在 $ M \ge 0 $ 使得 $ |Av| \le M p(v) $ 对 $ v \in V $。对线性映射 $ A : V \to E $，以下结论等价：

\begin{subitems}
\item

$ A $ 连续。

\item

$ A $ 在 0 处连续。

\item

$ A $ 有界。

\proofhead
这由定理 VI.2.5 的证明推出，那里只用到半范数的性质。\nobreak\hspace{0.35em}■
\end{subitems}

\item

$ \int \cdot \, d\mu : \mathcal{S}(X,\mu,E) \to E $ 连续。

\proofhead
这由 (c)、(f) 以及注 2.1(c) 推出。\nobreak\hspace{0.35em}■
\end{items}

设 $ p $ 是 $ V $ 上的半范数。由注 2.2(e) 可知，$ (V,p) $ 的拓扑未必由 $ V $ 上的度量生成，此时柯西序列与完备性这些度量概念都不可用。因此我们需要新的定义：若对每个 $ \varepsilon > 0 $ 都存在 $ N \in \mathbb{N} $ 使得 $ p(v_j - v_k) < \varepsilon $ 对 $ j, k \ge N $，则称序列 $ (v_j) \in V^{\mathbb{N}} $ 是 $ (V,p) $ 中的\textbf{柯西序列}。若 $ (V,p) $ 中每个柯西序列都收敛，则称 $ (V,p) $ \textbf{完备}。

\end{remark}

\begin{remark}

\begin{items}
\item

若 $ (V,p) $ 是赋范向量空间，这些概念与第 II.6 节的一致。

\item

设 $ (v_j) \in V^{\mathbb{N}} $ 且 $ v \in V $。则 $ v_j \to v $ 当且仅当 $ p(v - v_j) \to 0 $。然而收敛序列的极限一般并不唯一确定：若 $ p $ 不是范数，则 $ v_j \to v $ 蕴含 $ v_j \to w $ 对每个满足 $ p(v - w) = 0 $ 的 $ w \in V $。

\item

$ (V,p) $ 中所有柯西序列的集合构成 $ V^{\mathbb{N}} $ 的向量子空间。\nobreak\hspace{0.35em}■

\end{items}

下文我们总给空间 $ \mathcal{S}(X,\mu,E) $ 装配由 $ \|\cdot\|_1 $ 诱导的拓扑。这时我们也可以把 $ \mathcal{S}(X,\mu,E) $ 中的柯西序列称为 \textbf{$ \mathcal{L}_1 $-柯西序列}。

称函数 $ f \in E^{X} $ 是 \textbf{μ-可积}的，若 $ f $ 是 $ \mathcal{S}(X,\mu,E) $ 中某个 $ \mathcal{L}_1 $-柯西序列 $ (\varphi_j) $ 的 $ \mu $-几乎处处极限。我们把 $ X $ 的 $ E $ 值 $ \mu $-可积函数之集记作 $ \mathcal{L}_1(X,\mu,E) $。

\end{remark}

\begin{proposition}

作为向量子空间，我们有包含关系

\[
\mathcal{S}(X,\mu,E) \subset \mathcal{L}_1(X,\mu,E) \subset \mathcal{L}_0(X,\mu,E) .
\]

\proofhead
显然每个 $ \mu $-简单函数都是 $ \mu $-可积的。每个 $ \mu $-可积函数都是 $ \mu $-可测的，这由注 1.2(a) 与定理 1.14 推出。还需证明 $ \mathcal{L}_1(X,\mu,E) $ 是 $ \mathcal{L}_0(X,\mu,E) $ 的向量子空间。取 $ f, g \in \mathcal{L}_1(X,\mu,E) $ 与 $ \alpha \in \mathbb{K} $。存在 $ \mathcal{S}(X,\mu,E) $ 中的 $ \mathcal{L}_1 $-柯西序列 $ (\varphi_j) $ 与 $ (\psi_j) $，使得当 $ j \to \infty $ 时 $ \varphi_j \to f $ 且 $ \psi_j \to g $ 都是 $ \mu $-几乎处处的。由三角不等式可知 $ (\alpha \varphi_j + \psi_j)_{j \in \mathbb{N}} $ 是 $ \mathcal{S}(X,\mu,E) $ 中的 $ \mathcal{L}_1 $-柯西序列，且 $ \mu $-几乎处处收敛到 $ \alpha f + g $。因此 $ \alpha f + g $ 是 $ \mu $-可积的，如所要求。\nobreak\hspace{0.35em}■
\end{proposition}

\needspace{3\baselineskip}\subsection*{Bochner–Lebesgue 积分}

设 $ f \in \mathcal{L}_1(X,\mu,E) $。则存在 $ \mathcal{S}(X,\mu,E) $ 中的 $ \mathcal{L}_1 $-柯西序列 $ (\varphi_j) $ 使得 $ \varphi_j \to f $ 是 $ \mu $-几乎处处的。我们将看到序列 $ \left( \int_X \varphi_j \, d\mu \right)_{j \in \mathbb{N}} $ 在 $ E $ 中收敛。于是自然地，把 $ f $ 关于 $ \mu $ 的积分定义为这个积分序列的极限。当然我们必须验证该极限与逼近的简单函数序列无关；也就是说，必须证明：若 $ (\psi_j) $ 是 $ \mathcal{S}(X,\mu,E) $ 中另一个满足 $ \psi_j \to f $（$ \mu $-几乎处处）的柯西序列，则 $ \lim_j \int \varphi_j \, d\mu = \lim_j \int \psi_j \, d\mu $。

\begin{lemma}

设 $ (\varphi_j) $ 是 $ \mathcal{S}(X,\mu,E) $ 中的柯西序列。则存在 $ (\varphi_j) $ 的子列 $ (\varphi_{j_k})_{k \in \mathbb{N}} $ 与 $ f \in \mathcal{L}_1(X,\mu,E) $，使得

\begin{subitems}
\item

当 $ k \to \infty $ 时 $ \varphi_{j_k} \to f $ 是 $ \mu $-几乎处处的；

\item

对每个 $ \varepsilon > 0 $ 都存在 $ A_\varepsilon \in \mathcal{A} $，使得 $ \mu(A_\varepsilon) < \varepsilon $ 且 $ (\varphi_{j_k})_{k \in \mathbb{N}} $ 在 $ A_\varepsilon^{c} $ 上一致收敛到 $ f $。

\proofhead
（α）按假设，对任意 $ k \in \mathbb{N} $ 都存在 $ j_k \in \mathbb{N} $，使得界 $ \|\varphi_\ell - \varphi_m\|_1 < 2^{-2k} $ 对 $ \ell, m \ge j_k $ 成立。不妨设序列 $ (j_k)_{k \in \mathbb{N}} $ 递增。令 $ \psi_k := \varphi_{j_k} $，我们便得到
\[
\|\psi_\ell - \psi_m\|_1 < 2^{-2\ell} \quad \text{对 } m \ge \ell \ge 0 .
\]

（β）令 $ B_\ell := \left[ |\psi_{\ell+1} - \psi_\ell| \ge 2^{-\ell} \right] $ 对 $ \ell \in \mathbb{N} $。则 $ B_\ell $ 属于 $ \mathcal{A} $，并且 $ \mu(B_\ell) < \infty $ 对 $ \ell \in \mathbb{N} $，因为每个 $ \psi_m $ 都是 $ \mu $-简单的。因此 $ \chi_{B_\ell} $ 也是 $ \mu $-简单的，且注 2.1(f) 蕴含

\[
2^{-\ell} \mu(B_\ell) = 2^{-\ell} \int_X \chi_{B_\ell} \, d\mu \le \int_X |\psi_{\ell+1} - \psi_\ell| \, d\mu = \|\psi_{\ell+1} - \psi_\ell\|_1 < 2^{-2\ell} .
\]

这导致 $ \mu(B_\ell) < 2^{-\ell} $ 对 $ \ell \in \mathbb{N} $。
令 $ A_n := \bigcup_{k=0}^{\infty} B_{n+k} $，我们有 $ \mu(A_n) \le 2^{-n+1} $ 对 $ n \in \mathbb{N} $，并且我们看到 $ A := \bigcap_{n=0}^{\infty} A_n $ 是 $ \mu $-零集。
（γ）若 $ x $ 位于 $ A_n^{c} = \bigcap_{k=0}^{\infty} B_{n+k}^{c} $ 中，则

\[
|\psi_{\ell+1}(x) - \psi_\ell(x)| < 2^{-\ell} \quad \text{对 } \ell \ge n .
\]

由魏尔斯特拉斯优级数判别法（定理 V.1.6），级数

\[
\psi_0 + \sum (\psi_{\ell+1} - \psi_\ell)
\]

在 $ A_n^{c} $ 上一致收敛于 $ E $ 中。现在我们令

\begin{equation*}
f(x) := \begin{cases} \lim_k \psi_k(x) & \text{若 } x \in A^{c} , \\ 0 & \text{若 } x \in A . \end{cases}
\end{equation*}

则当 $ k \to \infty $ 时 $ \varphi_{j_k} \to f $ 是 $ \mu $-几乎处处的。此外，对每个 $ \varepsilon > 0 $ 都存在 $ n \in \mathbb{N} $ 使得 $ \mu(A_n) \le 2^{-n+1} < \varepsilon $，且当 $ k \to \infty $ 时 $ (\varphi_{j_k})_{k \in \mathbb{N}} $ 在 $ A_n^{c} $ 上一致收敛到 $ f $。\nobreak\hspace{0.35em}■

\end{subitems}

\end{lemma}

\begin{lemma}

设 $ (\varphi_j) $ 与 $ (\psi_j) $ 是 $ \mathcal{S}(X,\mu,E) $ 中 $ \mu $-几乎处处收敛到同一函数的 $ \mathcal{L}_1 $-柯西序列。则 $ \lim \|\varphi_j - \psi_j\|_1 = 0 $。

证明

\begin{subitems}
\item

取 $ \varepsilon > 0 $ 并令 $ \eta_j := \varphi_j - \psi_j $ 对 $ j \in \mathbb{N} $。由注 2.3(c)，$ (\eta_j) $ 是 $ \mathcal{S}(X,\mu,E) $ 中的 $ \mathcal{L}_1 $-柯西序列。因此存在自然数 $ N $，使得 $ \|\eta_j - \eta_k\| < \varepsilon/8 $ 对 $ j, k \ge N $。

\item

因为 $ \eta_N $ 是 $ \mu $-简单的，$ A := [\eta_N \ne 0] $ 属于 $ \mathcal{A} $ 且 $ \mu(A) < \infty $。此外 $ (\eta_j) $ 是 $ \mu $-几乎处处收敛到零的。于是引理 2.5 说明存在 $ B \in \mathcal{A} $ 使得 $ \mu(B) < \varepsilon/8(1 + \|\eta_N\|_\infty) $，并且存在 $ (\eta_j) $ 的子列 $ (\eta_{j_k})_{k \in \mathbb{N}} $，它在 $ B^{c} $ 上一致收敛到 0。因此存在 $ K \ge N $ 使得

\[
|\eta_{j_K}(x)| \le \varepsilon/8(1 + \mu(A)) \quad \text{对 } x \in A \setminus B .
\]

这表明 $ \int_{A \setminus B} |\eta_{j_K}| \, d\mu \le \varepsilon/8 $。

\item

$ B $ 与 $ K $ 的性质蕴含

\[
\int_B |\eta_{j_K}| \, d\mu \le \int_B |\eta_{j_K} - \eta_N| \, d\mu + \int_B |\eta_N| \, d\mu
\]
\[
\le \|\eta_{j_K} - \eta_N\|_1 + \|\eta_N\|_\infty \, \mu(B) < \varepsilon/4 .
\]

由 $ A $ 的定义，我们有

\[
\int_{A^{c}} |\eta_{j_K}| \, d\mu = \int_{A^{c}} |\eta_{j_K} - \eta_N| \, d\mu \le \|\eta_{j_K} - \eta_N\|_1 < \varepsilon/8 .
\]

综合起来，利用注 2.1(d) 得到

\[
\|\eta_{j_K}\|_1 \le \int_{A^{c} \cup (A \setminus B) \cup B} |\eta_{j_K}| \, d\mu < \varepsilon/2 ,
\]

因此 $ \|\eta_j\|_1 \le \|\eta_{j_K}\|_1 + \|\eta_j - \eta_{j_K}\|_1 < \varepsilon $ 对 $ j \ge N $。由于 $ \varepsilon > 0 $ 任意，一切得证。\nobreak\hspace{0.35em}■

\end{subitems}

\end{lemma}

\begin{corollary}

设 $ (\varphi_j) $ 与 $ (\psi_j) $ 是 $ \mathcal{S}(X,\mu,E) $ 中 $ \mu $-几乎处处收敛到同一函数的柯西序列。则序列 $ \left( \int_X \varphi_j \, d\mu \right) $ 与 $ \left( \int_X \psi_j \, d\mu \right) $ 在 $ E $ 中收敛，且

\[
\lim_j \int_X \varphi_j \, d\mu = \lim_j \int_X \psi_j \, d\mu .
\]

\proofhead
因为
\[
\left| \int_X \varphi_j \, d\mu - \int_X \varphi_k \, d\mu \right| \le \|\varphi_j - \varphi_k\|_1 \quad \text{对 } j, k \in \mathbb{N} ,
\]

$ \left( \int \varphi_j \, d\mu \right)_{j \in \mathbb{N}} $ 是 $ E $ 中的柯西序列；因此当 $ j \to \infty $ 时 $ \int \varphi_j \, d\mu \to e $，其中 $ e \in E $。同样存在 $ e' \in E $ 使得当 $ j \to \infty $ 时 $ \int \psi_j \, d\mu \to e' $。应用引理 2.6 与 $ E $ 的范数的连续性，我们看到

\[
|e - e'| = \lim_j \left| \int_X \varphi_j \, d\mu - \int_X \psi_j \, d\mu \right| \le \lim_j \int_X |\varphi_j - \psi_j| \, d\mu = \lim_j \|\varphi_j - \psi_j\|_1 = 0 ,
\]

于是证毕。\nobreak\hspace{0.35em}■

\end{corollary}

在这些准备之后，我们以自然的方式定义可积函数的积分，它是简单函数积分的推广。设 $ f \in \mathcal{L}_1(X,\mu,E) $。则存在 $ \mathcal{S}(X,\mu,E) $ 中的 $ \mathcal{L}_1 $-柯西序列 $ (\varphi_j) $ 使得 $ \varphi_j \to f $ 是 $ \mu $-几乎处处的。根据推论 2.7，量

\[
\int_X f \, d\mu := \lim_j \int_X \varphi_j \, d\mu
\]

在 $ E $ 中存在，且与序列 $ (\varphi_j) $ 无关。这称为 $ f $ 在 $ X $ 上关于测度 $ \mu $ 的 \textbf{Bochner–Lebesgue 积分}。除 $ \int_X f \, d\mu $ 外还常用其他记号，即

\[
\int f \, d\mu , \quad \int_X f(x) \, d\mu(x) , \quad \int_X f(x) \, \mu(dx) .
\]

显然，对简单函数而言，Bochner–Lebesgue 积分与本节开头（第 80 页）定义的积分一致。

\needspace{3\baselineskip}\subsection*{$ \mathcal{L}_1 $ 的完备性}

借助积分，我们现在在 $ \mathcal{L}_1(X,\mu,E) $ 上定义一个半范数，并证明 $ \mathcal{L}_1(X,\mu,E) $ 关于此半范数是完备的。

\begin{lemma}

若 $ f \in \mathcal{L}_1(X,\mu,E) $，则 $ |f| $ 属于 $ \mathcal{L}_1(X,\mu,\mathbb{R}) $。若 $ (\varphi_j) $ 是 $ \mathcal{S}(X,\mu,E) $ 中的 $ \mathcal{L}_1 $-柯西序列且 $ \varphi_j \to f $ 是 $ \mu $-几乎处处的，则 $ \int |f| \, d\mu = \lim_j \int |\varphi_j| \, d\mu $。

\proofhead
反向三角不等式（它显然对半范数也成立）给出
\[
\left\| |\varphi_j| - |\varphi_k| \right\|_1 \le \|\varphi_j - \varphi_k\|_1 \quad \text{且} \quad |\,|\varphi_j| - |\varphi_k|\,| \le |\varphi_j - \varphi_k| \quad \text{对 } j, k \in \mathbb{N} .
\]

因此 $ (|\varphi_j|)_{j \in \mathbb{N}} $ 是 $ \mathcal{S}(X,\mu,\mathbb{R}) $ 中的 $ \mathcal{L}_1 $-柯西序列，且 $ \mu $-几乎处处收敛到 $ |f| $。所以 $ |f| $ 属于 $ \mathcal{L}_1(X,\mu,\mathbb{R}) $，并且 $ \int |f| \, d\mu = \lim_j \int |\varphi_j| \, d\mu $。\nobreak\hspace{0.35em}■

\end{lemma}

\begin{corollary}

对 $ f \in \mathcal{L}_1(X,\mu,E) $，令 $ \|f\|_1 := \int_X |f| \, d\mu $。则 $ \|\cdot\|_1 $ 是 $ \mathcal{L}_1(X,\mu,E) $ 上的半范数，称为 \textbf{$ \mathcal{L}_1 $-半范数}。

\proofhead
取 $ f, g \in \mathcal{L}_1(X,\mu,E) $，并设 $ (\varphi_j) $ 与 $ (\psi_j) $ 是 $ \mathcal{S}(X,\mu,E) $ 中的 $ \mathcal{L}_1 $-柯西序列，使得 $ \varphi_j \to f $ 且 $ \psi_j \to g $ 都是 $ \mu $-几乎处处的。由引理 2.8 与注 2.2(c)，我们有
\[
\|f\|_1 = \int_X |f| \, d\mu = \lim_j \int_X |\varphi_j| \, d\mu = \lim_j \|\varphi_j\|_1 \ge 0 ,
\]

以及

\[
\|f + g\|_1 = \lim_j \|\varphi_j + \psi_j\|_1 \le \lim_j (\|\varphi_j\|_1 + \|\psi_j\|_1) = \|f\|_1 + \|g\|_1 ,
\]

和

\[
\|\alpha f\|_1 = \lim_j \|\alpha \varphi_j\| = |\alpha| \lim_j \|\varphi_j\| = |\alpha| \, \|f\|_1
\]

对任意 $ \alpha \in \mathbb{K} $。\nobreak\hspace{0.35em}■

\end{corollary}

我们将总给 $ \mathcal{L}_1(X,\mu,E) $ 装配由半范数 $ \|\cdot\|_1 $ 诱导的拓扑。

\begin{theorem}

\begin{subitems}
\item

$ \mathcal{S}(X,\mu,E) $ 在 $ \mathcal{L}_1(X,\mu,E) $ 中稠密。

\item

空间 $ \mathcal{L}_1(X,\mu,E) $ 是完备的。

\proofhead
(i) 设 $ f \in \mathcal{L}_1(X,\mu,E) $，并令 $ (\varphi_j) $ 表示简单函数构成的 $ \mathcal{L}_1 $-柯西序列，使得当 $ j \to \infty $ 时 $ \varphi_j \to f $ 是 $ \mu $-几乎处处的。再设 $ k \in \mathbb{N} $。则 $ (\varphi_j - \varphi_k)_{j \in \mathbb{N}} $ 是 $ \mathcal{S}(X,\mu,E) $ 中的 $ \mathcal{L}_1 $-柯西序列，使得 $ (\varphi_j - \varphi_k) \to (f - \varphi_k) $
当 $ j \to \infty $ 时 $ \mu $-几乎处处。由引理 2.8，
\[
\|f - \varphi_k\|_1 = \lim_{j} \|\varphi_j - \varphi_k\|_1 \qquad \text{对 } k \in \mathbb{N} .
\]

设 $ \varepsilon > 0 $。则存在 $ N \in \mathbb{N} $，使得对 $ j, k \ge N $ 有 $ \|\varphi_j - \varphi_k\|_1 < \varepsilon $；令 $ j \to \infty $ 取极限，得到 $ \|f - \varphi_N\|_1 \le \varepsilon $。这说明 $ \mathcal{S}(X,\mu,E) $ 在 $ \mathcal{L}_1(X,\mu,E) $ 中稠密。

\end{subitems}

(ii) 设 $ (f_j) $ 是 $ \mathcal{L}_1(X,\mu,E) $ 中的柯西序列，并取 $ \varepsilon > 0 $。选取 $ M \in \mathbb{N} $，使得对 $ j, k \ge M $ 有 $ \|f_j - f_k\|_1 < \varepsilon/2 $。由 (i) 我们知道，对任意 $ j \in \mathbb{N} $ 都可找到 $ \varphi_j \in \mathcal{S}(X,\mu,E) $，使得 $ \|f_j - \varphi_j\|_1 < 2^{-j} $。于是

\[
\|\varphi_j - \varphi_k\|_1 \le \|\varphi_j - f_j\|_1 + \|f_j - f_k\|_1 + \|f_k - \varphi_k\|_1 < 2^{-j} + 2^{-k} + \varepsilon/2
\]

对 $ j, k \ge M $。这说明 $ (\varphi_j) $ 是 $ \mathcal{S}(X,\mu,E) $ 中的 $ \mathcal{L}_1 $-柯西序列。因此由引理 2.5，存在 $ (\varphi_j) $ 的子列 $ (\varphi_{j_k})_{k \in \mathbb{N}} $ 以及 $ f \in \mathcal{L}_1(X,\mu,E) $，使得当 $ k \to \infty $ 时 $ \varphi_{j_k} \to f $ $ \mu $-几乎处处。(i) 的证明表明，存在 $ N \ge M $ 使得 $ \|f - \varphi_{j_N}\|_1 < \varepsilon/4 $，于是得到

\[
\|f - f_j\|_1 \le \|f - \varphi_{j_N}\|_1 + \|\varphi_{j_N} - f_{j_N}\|_1 + \|f_{j_N} - f_j\|_1 < \varepsilon \qquad \text{对 } j \ge N ,
\]

即 $ (f_j) $ 在 $ \mathcal{L}_1(X,\mu,E) $ 中收敛到 $ f $。\nobreak\hspace{0.35em}■

\end{theorem}

\needspace{3\baselineskip}\subsection*{积分的基本性质}

我们已经看到，简单函数空间上的积分是连续的、线性的，并且当 $ E = \mathbb{R} $ 时还是单调的——参见注 2.2(g)、2.1(c) 与 2.1(f)。现在我们要说明，这些性质在把积分从简单函数空间延拓到可积函数空间之后仍然保持。

\begin{theorem}

\begin{subitems}
\item

$ \int_X \cdot \, d\mu: \mathcal{L}_1(X,\mu,E) \to E $ 是线性且连续的，且

\[
\left| \int_X f \, d\mu \right| \le \int_X |f| \, d\mu = \|f\|_1 .
\]
\item

$ \int_X \cdot \, d\mu: \mathcal{L}_1(X,\mu,\mathbb{R}) \to \mathbb{R} $ 是连续的正线性泛函。

\item

设 $ F $ 是巴拿赫空间，$ T \in \mathcal{L}(E,F) $。则

\[
Tf \in \mathcal{L}_1(X,\mu,F) \qquad \text{且} \qquad T \int_X f \, d\mu = \int_X Tf \, d\mu
\]

对 $ f \in \mathcal{L}_1(X,\mu,E) $。

\proofhead
(i) 命题 2.4 表明，$ \mu $-可积函数构成一个向量空间。取 $ f, g \in \mathcal{L}_1(X,\mu,E) $ 与 $ \alpha \in \mathbb{K} $。则存在 $ \mathcal{S}(X,\mu,E) $ 中的 $ \mathcal{L}_1 $-柯西序列 $ (\varphi_j) $ 与 $ (\psi_j) $，使得 $ \varphi_j \to f $ 且 $ \psi_j \to g $ $ \mu $-几乎处处。由注 2.1(c)，
\[
\int_X (\alpha\varphi_j + \psi_j) \, d\mu = \alpha \int_X \varphi_j \, d\mu + \int_X \psi_j \, d\mu \qquad \text{对 } j \in \mathbb{N} .
\]

在 $ \mathcal{L}_1(X,\mu,E) $ 上积分的线性由取极限 $ j \to \infty $ 得到。由推论 2.9，$ \|\cdot\|_1 $ 是 $ \mathcal{L}_1(X,\mu,E) $ 上的半范数，而注 2.1(e) 给出

\[
\left| \int_X \varphi_j \, d\mu \right| \le \int_X |\varphi_j| \, d\mu = \|\varphi_j\|_1 \qquad \text{对 } j \in \mathbb{N} .
\]

由引理 2.8，我们可以取极限 $ j \to \infty $，并得到

\[
\left| \int_X f \, d\mu \right| \le \int_X |f| \, d\mu = \|f\|_1 .
\]

连续性现在由注 2.2(f) 推出。
刚才所用的方法可以毫无困难地改用来证明 (ii) 与 (iii)。我们把它留给读者作为练习。\nobreak\hspace{0.35em}■

\end{subitems}

\end{theorem}

\begin{corollary}

\begin{subitems}
\item

映射 $ f = (f_1,\ldots,f_n): X \to \mathbb{K}^n $ 是 $ \mu $-可积的，当且仅当其每个坐标 $ f_j $ 都是 $ \mu $-可积的。此时，

\[
\int_X f \, d\mu = \left( \int_X f_1 \, d\mu, \ldots, \int_X f_n \, d\mu \right) .
\]
\item

设 $ g, h \in \mathbb{R}^X $ 并定义 $ f := g + ih $。则 $ f $ 属于 $ \mathcal{L}_1(X,\mu,\mathbb{C}) $ 当且仅当 $ g $ 与 $ h $ 属于 $ \mathcal{L}_1(X,\mu,\mathbb{R}) $。此时，

\[
\int_X f \, d\mu = \int_X g \, d\mu + i \int_X h \, d\mu .
\]
\item

函数 $ f \in \mathbb{R}^X $ 是 $ \mu $-可积的，当且仅当 $ f^+ $ 与 $ f^- $ 都是 $ \mu $-可积的。此时，

\[
\int_X f \, d\mu = \int_X f^+ \, d\mu - \int_X f^- \, d\mu , \qquad \int_X |f| \, d\mu = \int_X f^+ \, d\mu + \int_X f^- \, d\mu .
\]

\proofhead
(i) “$ \Rightarrow $” 取 $ f \in \mathcal{L}_1(X,\mu,\mathbb{K}^n) $。由于对 $ j = 1,\ldots,n $ 有 $ \mathrm{pr}_j \in \mathcal{L}(\mathbb{K}^n,\mathbb{K}) $，由定理 2.11(iii) 可得 $ f_j = \mathrm{pr}_j \circ f $ 属于 $ \mathcal{L}_1(X,\mu,\mathbb{K}) $。此外 $ \int f_j \, d\mu = \mathrm{pr}_j \int f \, d\mu $，故
\[
\int_X f \, d\mu = \left( \int_X f_1 \, d\mu, \ldots, \int_X f_n \, d\mu \right) .
\]

“$ \Leftarrow $” 对 $ j = 1,\ldots,n $，我们考虑映射

\[
b_j: \mathbb{K} \to \mathbb{K}^n , \qquad y \mapsto (0,\ldots,0,y,0,\ldots,0) ,
\]

其中 $ y $ 位于右侧第 $ j $ 个位置。于是

\[
b_j \in \mathcal{L}(\mathbb{K},\mathbb{K}^n) \qquad \text{且} \qquad f := \sum_{j=1}^{n} b_j \circ f_j .
\]

该断言现在由定理 2.11(i) 与 (iii) 推出。

\end{subitems}

\begin{subitems}
\item[（ii）]

这由 (i) 以及把 $ \mathbb{C} $ 与 $ \mathbb{R}^2 $ 等同得到。

\item[（iii）]

对 $ f \in \mathbb{R}^X $，我们有

\[
f^+ = (f + |f|)/2 , \qquad f^- = (|f| - f)/2 , \qquad f = f^+ - f^- , \qquad |f| = f^+ + f^- .
\]

于是定理 2.11(i) 与引理 2.8 蕴含结论。\nobreak\hspace{0.35em}■

\end{subitems}

\end{corollary}

\begin{lemma}

对 $ f \in \mathcal{L}_1(X,\mu,E) $ 与 $ A \in \mathcal{A} $，有 $ \chi_A f \in \mathcal{L}_1(X,\mu,E) $。

\proofhead
设 $ (\varphi_j) $ 是 $ \mathcal{S}(X,\mu,E) $ 中 $ \mu $-几乎处处收敛到 $ f $ 的 $ \mathcal{L}_1 $-柯西序列。则由注 1.2(d)，$ \chi_A \varphi_j $ 是 $ \mu $-简单的，且 $ (\chi_A \varphi_j)_{j \in \mathbb{N}} $ 显然 $ \mu $-几乎处处收敛到 $ \chi_A f $。由注 2.1(f)，我们有
\[
\int_X |\chi_A \varphi_j - \chi_A \varphi_k| \, d\mu = \int_X \chi_A |\varphi_j - \varphi_k| \, d\mu \le \int_X |\varphi_j - \varphi_k| \, d\mu \qquad \text{对 } j, k \in \mathbb{N} .
\]

因此 $ (\chi_A \varphi_j)_{j \in \mathbb{N}} $ 是 $ \mathcal{S}(X,\mu,E) $ 中的 $ \mathcal{L}_1 $-柯西序列。这说明 $ \chi_A f $ 是 $ \mu $-可积的。\nobreak\hspace{0.35em}■

\end{lemma}

对 $ f \in \mathcal{L}_1(X,\mu,E) $ 与 $ A \in \mathcal{A} $，我们如下定义 $ f $ 在 $ A $ 上的 \textbf{$ \mu $-积分}：

\[
\int_A f \, d\mu := \int_X \chi_A f \, d\mu .
\]

这由引理 2.13 是良定义的。

\begin{remark}

设 $ f \in \mathcal{L}_1(X,\mu,E) $ 与 $ A \in \mathcal{A} $。

\begin{items}
\item

$ \int_A \cdot \, d\mu: \mathcal{L}_1(X,\mu,E) \to E $ 是线性且连续的，且

\[
\left| \int_A f \, d\mu \right| \le \int_A |f| \, d\mu = \|\chi_A f\|_1 .
\]
\item

设 $ \mathcal{B} := \mathcal{A}|A $ 且 $ \nu := \mu|_{\mathcal{B}} $。则 $ \int_A f \, d\mu = \int_A f|A \, d\nu $。

\proofhead
证明很简单，留给读者（习题 1）。\nobreak\hspace{0.35em}■
\item

若 $ E = \mathbb{R} $ 且 $ f \ge 0 $，则映射

\[
\mathcal{A} \to [0,\infty) , \qquad A \mapsto \int_A f \, d\mu
\]

是有限测度（见习题 11）。\nobreak\hspace{0.35em}■

\end{items}

\end{remark}

\begin{lemma}

设 $ f \in \mathcal{L}_1(X,\mu,E) $ 与 $ g \in E^X $ 满足 $ f = g $ $ \mu $-几乎处处。则 $ g $ 也属于 $ \mathcal{L}_1(X,\mu,E) $，且 $ \int_X f \, d\mu = \int_X g \, d\mu $。

\proofhead
设 $ (\varphi_j) $ 是 $ \mathcal{S}(X,\mu,E) $ 中使 $ \varphi_j \to f $ $ \mu $-几乎处处的 $ \mathcal{L}_1 $-柯西序列。又设 $ M $ 与 $ N $ 是 $ \mu $-零集，使得在 $ M^c $ 上 $ \varphi_j \to f $，在 $ N^c $ 上 $ f = g $。则 $ (\varphi_j) $ $ \mu $-几乎处处收敛到 $ g $，因为对 $ x \in (M \cup N)^c $ 有 $ \varphi_j(x) \to g(x) $。因此 $ g $ 属于 $ \mathcal{L}_1(X,\mu,E) $，且 $ \int g \, d\mu = \lim_{j} \int \varphi_j \, d\mu = \int f \, d\mu $。\nobreak\hspace{0.35em}■
\end{lemma}

\begin{corollary}

\begin{subitems}
\item

设 $ f \in E^X $ $ \mu $-几乎处处为零。则 $ f $ 是 $ \mu $-可积的，且 $ \int_X f \, d\mu = 0 $。

\item

设 $ f, g \in \mathcal{L}_1(X,\mu,\mathbb{R}) $ $ \mu $-几乎处处满足 $ f \le g $。则 $ \int_X f \, d\mu \le \int_X g \, d\mu $。

\proofhead
(i) 这立刻由引理 2.15 推出。
\end{subitems}

(ii) 定理 2.11(ii) 与引理 2.15 蕴含 $ 0 \le \int_X (g - f) \, d\mu $，因此 $ \int_X f \, d\mu \le \int_X g \, d\mu $。\nobreak\hspace{0.35em}■

\end{corollary}

\begin{proposition}

对 $ f \in \mathcal{L}_1(X,\mu,E) $ 与 $ \alpha > 0 $，有 $ \mu([|f| \ge \alpha]) < \infty $。

\proofhead
引理 2.5 表明，存在 $ \mathcal{S}(X,\mu,E) $ 中的 $ \mathcal{L}_1 $-柯西序列 $ (\varphi_j) $ 以及 $ \mu $-可测集 $ A $，使得 $ \mu(A) \le 1 $ 且 $ (\varphi_j) $ 在 $ A^c $ 上一致收敛到 $ f $。由于 $ |f| $ 是 $ \mu $-可测的，$ B := A^c \cap [|f| \ge \alpha] $ 属于 $ \mathcal{A} $。此外存在 $ N \in \mathbb{N} $，使得对 $ x \in A^c $ 有 $ |\varphi_N(x) - f(x)| \le \alpha/2 $。因此
\[
|\varphi_N(x)| \ge |f(x)| - |\varphi_N(x) - f(x)| \ge \alpha/2 \qquad \text{对 } x \in B .
\]

特别地，$ B $ 包含于 $ [\varphi_N \ne 0] $。于是 $ \mu(B) \le \mu([\varphi_N \ne 0]) < \infty $，因为 $ \varphi_N $ 是 $ \mu $-简单的。由于

\[
[|f| \ge \alpha] = B \cup (A \cap [|f| \ge \alpha]) \subset B \cup A ,
\]

可得 $ \mu([|f| \ge \alpha]) \le \mu(B) + 1 < \infty $。\nobreak\hspace{0.35em}■

\end{proposition}

\needspace{3\baselineskip}\subsection*{$ \mathcal{L}_1 $ 中的收敛}

对每个可积函数 $ f $，都存在由简单函数构成的 $ \mathcal{L}_1 $-柯西序列，它几乎处处收敛到 $ f $。下面我们说明，$ \mathcal{L}_1(X,\mu,E) $ 中的每个柯西序列实际上都有一个子列，几乎处处收敛到该序列的 $ \mathcal{L}_1 $ 极限。

\begin{theorem}

设 $ (f_j) $ 是 $ \mathcal{L}_1(X,\mu,E) $ 中收敛到 $ f \in \mathcal{L}_1(X,\mu,E) $ 的序列。

\begin{subitems}
\item

存在 $ (f_j) $ 的子列 $ (f_{j_k})_{k \in \mathbb{N}} $，具有以下性质：
(α) 当 $ k \to \infty $ 时 $ f_{j_k} \to f $ $ \mu $-几乎处处。
(β) 对每个 $ \varepsilon > 0 $，存在 $ A_\varepsilon \in \mathcal{A} $ 使 $ \mu(A_\varepsilon) < \varepsilon $，且 $ (f_{j_k})_{k \in \mathbb{N}} $ 在 $ A_\varepsilon^c $ 上一致收敛到 $ f $。

\item

积分 $ \int_X f_j \, d\mu $ 当 $ j \to \infty $ 时收敛到 $ \int_X f \, d\mu $。

\proofhead
(i) 只需处理 $ f = 0 $ 的情形，因为若 $ f \ne 0 $，我们可以考虑序列 $ (f_j - f)_{j \in \mathbb{N}} $。
如同引理 2.5 的证明，存在 $ (f_j) $ 的子列 $ (g_k) $，使得对 $ m \ge \ell \ge 0 $ 有 $ \|g_\ell - g_m\|_1 < 2^{-2\ell} $。极限 $ m \to \infty $ 给出 $ \|g_\ell\|_1 \le 2^{-2\ell} $，对 $ \ell \in \mathbb{N} $。我们令 $ B_\ell := [|g_\ell| \ge 2^{-\ell}] $。由引理 2.8、命题 2.4 与命题 1.9，$ B_\ell $ 属于 $ \mathcal{A} $，而且我们得到
\[
2^{-\ell} \mu(B_\ell) \le \int_{B_\ell} |g_\ell| \, d\mu \le \int_X |g_\ell| \, d\mu = \|g_\ell\|_1 \le 2^{-2\ell} \qquad \text{对 } \ell \in \mathbb{N}
\]

（比较定理 2.11(ii)）。因此对 $ \ell \in \mathbb{N} $ 有 $ \mu(B_\ell) \le 2^{-\ell} $。取 $ A_n := \bigcup_{k=0}^{\infty} B_{n+k} $，我们有 $ \mu(A_n) \le 2^{-n+1} $，并得出 $ A := \bigcap_{n=0}^{\infty} A_n $ 是 $ \mu $-零集。容易验证 $ (g_k) $ 在 $ A_n^c $ 上一致收敛到 0，并在 $ A^c $ 上逐点收敛到 0（这方面参见引理 2.5 的证明）。

\end{subitems}

(ii) 由定理 2.11(i) 可得

\[
\left| \int_X f_j \, d\mu - \int_X f \, d\mu \right| \le \int_X |f_j - f| \, d\mu = \|f_j - f\|_1 \qquad \text{对 } j \in \mathbb{N} ,
\]

故左端当 $ j \to \infty $ 时的极限为 0。\nobreak\hspace{0.35em}■

\end{theorem}

\begin{corollary}

对 $ f \in \mathcal{L}_1(X,\mu,E) $，

\[
\|f\|_1 = 0 \iff f = 0 \quad \mu\text{-几乎处处} .
\]

\proofhead
“$ \Rightarrow $” 因为 $ \|f\|_1 = 0 $，由 $ f_j := 0 $（$ j \in \mathbb{N} $）定义的序列 $ (f_j) $ 在 $ \mathcal{L}_1(X,\mu,E) $ 中收敛到 $ f $。由定理 2.18，于是存在 $ (f_j) $ 的子列 $ (f_{j_k})_{k \in \mathbb{N}} $，使得当 $ k \to \infty $ 时 $ f_{j_k} \to f $ $ \mu $-几乎处处。因此 $ f = 0 $ $ \mu $-几乎处处。

“$ \Leftarrow $” 按假设，$ |f| = 0 $ $ \mu $-几乎处处；断言由推论 2.16(i) 推出。\nobreak\hspace{0.35em}■
\end{corollary}

下面我们通过一个特别简单的情形来说明本节的概念与结果，以此结束本节。

\begin{example}

（可和级数空间） 设 $ X $ 表示 $ \mathbb{N} $ 或 $ \mathbb{Z} $，并设 $ \mathcal{H}^0 $ 为其 0 维豪斯多夫测度，或计数测度。由 $ \mathbb{R} $ 诱导的拓扑显然使 $ X $ 成为 $ \sigma $-紧度量空间，其中每个单点集都是开集。因此 $ X $ 的拓扑与 $ \mathfrak{P}(X) $ 一致：$ X $ 的每个子集都是开集。从而 $ X $ 的每个映射都是连续的 $ E $ 值映射，即 $ C(X,E) = E^X $。

进一步可得 $ \mathcal{B}(X) = \mathfrak{P}(X) $，且 $ \mathcal{H}^0 $ 是 $ X $ 上的 massive 拉东测度。于是由定理 1.17，

\[
\mathcal{L}_0(X,\mathcal{H}^0,E) = C(X,E) = E^X .
\]

此外，$ \mathcal{H}^0 $ 没有非空零集。因此 $ \mathcal{H}^0 $-几乎处处收敛与逐点收敛相同。

对 $ \varphi \in E^X $，我们定义 $ \varphi $ 的\textbf{支集}为集合

\[
\operatorname{supp}(\varphi) := \{ x \in X ; \varphi(x) \ne 0 \} ,
\]

并记 $ C_c(X,E) $ 为 $ X $ 上具有紧支集的连续 $ E $ 值函数空间：

\[
C_c(X,E) := \{ \varphi \in C(X,E) ; \operatorname{supp}(\varphi) \text{ 是紧的} \} .
\]

显然 $ \varphi \in C(X,E) $ 属于 $ C_c(X,E) $ 当且仅当 $ \operatorname{supp}(\varphi) $ 是有限集。又 $ C_c(X,E) $ 是 $ C(X,E) $ 的向量子空间，且容易验证 $ C_c(X,E) = \mathcal{S}(X,\mathcal{H}^0,E) $。

对 $ \varphi \in C_c(X) $，由注 2.1(b) 可得

\begin{equation}
\int_X \varphi \, d\mathcal{H}^0 = \sum_{x \in \operatorname{supp}(\varphi)} \varphi(x) .\tag{2.4}
\end{equation}

现在我们令

\[
\ell_1(X,E) := \left\{ f \in E^X ; \sum_{x \in X} |f(x)| < \infty \right\} .
\]

对 $ f \in \ell_1(X,E) $ 与 $ n \in \mathbb{N} $，令

\begin{equation*}
\varphi_n(x) := \begin{cases} f(x) & \text{若 } |x| \le n , \\ 0 & \text{若 } |x| > n . \end{cases}
\end{equation*}

则 $ \varphi_n $ 属于 $ C_c(X,E) $，且当 $ n \to \infty $ 时 $ \varphi_n \to f $。对 $ m > n $，由 (2.4) 得到

\[
\|\varphi_n - \varphi_m\|_1 = \sum_{n < |x| \le m} |f(x)| .
\]

因此 $ (\varphi_n) $ 是 $ \mathcal{S}(X,\mathcal{H}^0,E) $ 中的 $ \mathcal{L}_1 $-柯西序列，这说明 $ f $ 属于 $ \mathcal{L}_1(X,\mathcal{H}^0,E) $。因此 $ \ell_1(X,E) \subset \mathcal{L}_1(X,\mathcal{H}^0,E) $，且

\begin{equation}
\int_X f \, d\mathcal{H}^0 = \sum_{x \in X} f(x) \qquad \text{对 } f \in \ell_1(X,E) .\tag{2.5}
\end{equation}

现在设 $ f \in \mathcal{L}_1(X,\mathcal{H}^0,E) $。则存在 $ \mathcal{S}(X,\mu,E) $ 中的 $ \mathcal{L}_1 $-柯西序列 $ (\psi_j) $——从而也是 $ C_c(X,E) $ 中的——它逐点收敛到 $ f $。由引理 2.8，$ |f| $ 属于 $ \mathcal{L}_1(X,\mathcal{H}^0,\mathbb{R}) $，且

\[
\|f\|_1 = \int_X |f| \, d\mathcal{H}^0 = \lim_{j \to \infty} \int_X |\psi_j| \, d\mathcal{H}^0 = \lim_{j \to \infty} \sum_{x \in X} |\psi_j(x)| .
\]

于是存在 $ k \in \mathbb{N} $ 使得

\[
\left| \int_X |f| \, d\mathcal{H}^0 - \sum_{x \in X} |\psi_j(x)| \right| \le 1 \qquad \text{对 } j \ge k .
\]

这说明

\[
\sum_{x \in X} |\psi_j(x)| \le 1 + \int_X |f| \, d\mathcal{H}^0 =: K < \infty \qquad \text{对 } j \ge k .
\]

因此对每个 $ m \in \mathbb{N} $ 都有

\[
\sum_{|x| \le m} |\psi_j(x)| \le K \qquad \text{对 } j \ge k ,
\]

由此，因为当 $ j \to \infty $ 时 $ \psi_j \to f $，我们得到

\[
\sum_{|x| \le m} |f(x)| \le K \qquad \text{对 } m \in \mathbb{N} .
\]

现在定理 II.7.7 蕴含 $ f $ 属于 $ \ell_1(X,E) $（且满足 $ \|f\|_1 \le K $）。因此我们已证明

\[
\mathcal{L}_1(X,\mathcal{H}^0,E) = \ell_1(X,E) ,
\]

由此我们从 (2.5) 得到关系\footnote{见定理 II.8.9。}

\[
\|f\|_1 = \sum_{x \in X} |f(x)| .
\]

最后，由定理 2.10 与注 2.2(a) 可得

\[
\ell_1(X,E) := (\ell_1(X,E), \|\cdot\|_1)
\]

是巴拿赫空间，即\textbf{可和}（$ E $ 值）级数的空间。

若 $ E = \mathbb{K} $，通常用 $ \ell_1(\mathbb{Z}) $ 与 $ \ell_1(\mathbb{N}) $ 表示 $ \ell_1(X,\mathbb{K}) $，缩写 $ \ell_1 := \ell_1(\mathbb{N}) $ 也很常见。\footnote{比较习题 II.8.6。} \nobreak\hspace{0.35em}■

\end{example}

\exerciseshead

\begin{exercises}
\item

设 $ A \in \mathcal{A} $，$ \mathcal{B} := \mathcal{A} \mid A $，且 $ \nu := \mu \mid \mathcal{B} $。验证：对 $ f \in E^X $ 有

\[
\chi_A f \in \mathcal{L}_1(X,\mu,E) \iff f \mid A \in \mathcal{L}_1(A,\nu,E) .
\]

对这样的 $ f $，证明

\[
\int_X \chi_A f \, d\mu = \int_A f \mid A \, d\nu .
\]
\item

设 $ (f_j) $ 是 $ \mathcal{L}_1(X,\mu,E) $ 中的序列，且一致收敛于 $ f \in E^X $。又设 $ \mu(X) < \infty $。证明：$ f $ 属于 $ \mathcal{L}_1(X,\mu,E) $，$ f_j \to f $ 于 $ \mathcal{L}_1(X,\mu,E) $ 中，且 $ \lim_j \int_X f_j \, d\mu = \int_X f \, d\mu $。

\item

验证：对 $ f \in \mathcal{L}_1(X,\mu,\mathbb{R}^+) $，

\[
\int_X f \, d\mu = \sup \left\{ \int_X \varphi \, d\mu ;\ \varphi \in \mathcal{S}(X,\mu,\mathbb{R}^+),\ \varphi \le f\ \mu\text{-几乎处处} \right\} .
\]
\item

设 $ X $ 是任意非空集，$ a \in X $，$ \delta_a $ 是支集在 $ a $ 处的狄拉克测度。证明 $ \mathcal{L}_1(X,\delta_a,\mathbb{R}) = \mathbb{R}^X $，并对 $ \mathbb{R}^X $ 计算 $ \int f \, d\delta_a $。

\item

设 $ \mu_F $ 是习题 IX.4.10 中的勒贝格–斯蒂尔杰斯测度。确定 $ \mathcal{L}_1(\mathbb{R},\mu_F,\mathbb{K}) $，并对 $ f \in \mathcal{L}_1(\mathbb{R},\mu_F,\mathbb{K}) $ 计算 $ \int f \, d\mu_F $。

\item

证明定理 2.11 的陈述 (ii) 与 (iii)。

\item

设 $ f \in \mathcal{L}_0(X,\mu,E) $ 是 $ \mu $-几乎处处有界的，且 $ \mu(X) < \infty $。证明或否证 $ f $ 是 $ \mu $-可积的。

\item

设 $ (f_j) $ 是 $ \mathcal{L}_1(X,\mu,\mathbb{R}) $ 中的递增序列，且 $ f_j \ge 0 $，并设它 $ \mu $-几乎处处收敛于 $ f \in \mathcal{L}_1(X,\mu,\mathbb{R}) $。则 $ \int_X f_j \, d\mu \uparrow \int_X f \, d\mu $。（这就是 $ \mathcal{L}_1 $ 中的\textbf{单调收敛定理}。）

（提示：证明 $ (f_j) $ 是 $ \mathcal{L}_1(X,\mu,\mathbb{R}) $ 中的柯西序列，并指出其极限。）

\item

设 $ (f_j) $ 是 $ \mathcal{L}_1(X,\mu,\mathbb{R}) $ 中的序列，$ \mu $-几乎处处 $ f_j \ge 0 $，且 $ \sum_{j=0}^{\infty} f_j \in \mathcal{L}_1(X,\mu,\mathbb{R}) $。证明 $ \sum_{j=0}^{\infty} \int f_j \, d\mu = \int \left( \sum_{j=0}^{\infty} f_j \right) d\mu $。（提示：习题 8。）

\item

设 $ f \in \mathcal{L}_1(X,\mu,\mathbb{R}) $ 满足 $ \mu $-几乎处处 $ f > 0 $。证明：对每个满足 $ \mu(A) > 0 $ 的 $ A \in \mathcal{A} $，都有 $ \int_A f \, d\mu > 0 $。

\item

给定 $ f \in \mathcal{L}_1(X,\mu,\mathbb{R}) $ 且 $ f \ge 0 $，定义 $ \varphi_f(A) := \int_A f \, d\mu $，其中 $ A \in \mathcal{A} $。证明：

\begin{items}
\item

$ (X,\varphi_f,\mathcal{A}) $ 是有限测度空间。

\item

$ \mathcal{N}_\mu \subset \mathcal{N}_{\varphi_f} $。

\item

若 $ \mu $-几乎处处 $ f > 0 $，则 $ \mathcal{N}_\mu = \mathcal{N}_{\varphi_f} $。
特别地，证明：若 $ \mu $-几乎处处 $ f > 0 $，则 $ (X,\mathcal{A},\varphi_f) $ 是完备的有限测度空间。
（提示：(a) 习题 9。(b) 习题 10。）

\end{items}

\item

设 $ f \in \mathcal{L}_1(X,\mu,\mathbb{R}) $ 满足 $ \mu $-几乎处处 $ f > 0 $，并取 $ g \in \mathcal{L}_0(X,\mu,\mathbb{R}) $。证明：$ g $ 是 $ \varphi_f $-可积的当且仅当 $ gf $ 是 $ \mu $-可积的。此时，证明

\[
\int_X g \, d\varphi_f = \int_X fg \, d\mu .
\]
\item

对 $ f \in \mathcal{L}_1(X,\mu,\overline{\mathbb{R}}^+) $，证明\textbf{切比雪夫不等式}：

\[
\mu\left( [f \ge \alpha] \right) \le \frac{1}{\alpha} \int_X f \, d\mu \quad \text{对 } \alpha > 0 .
\]
\item

设 $ \mu(X) < \infty $，$ I $ 是 $ \mathbb{R} $ 中的完全区间。又设 $ \varphi \in C^1(I,\mathbb{R}) $ 是凸的。证明\textbf{詹森不等式}：若 $ f \in \mathcal{L}_1(X,\mu,\mathbb{R}) $ 满足 $ f(X) \subset I $ 且 $ \varphi \circ f \in \mathcal{L}_1(X,\mu,\mathbb{R}) $，则

\[
\varphi\left( \fint_X f \, d\mu \right) \le \fint_X \varphi \circ f \, d\mu , \quad \text{其中 } \fint_X f \, d\mu := \frac{1}{\mu(X)} \int_X f \, d\mu .
\]

（提示：证明 $ \alpha := \fint f \, d\mu $ 落在 $ I $ 中，从而界 $ \varphi(y) \ge \varphi(\alpha) + \varphi'(\alpha)(y-\alpha) $ 适用。）

\item

设 $ f \in \mathcal{L}_1(X,\mu,E) $。证明：对每个 $ \varepsilon > 0 $ 都存在 $ \delta > 0 $，使得对一切满足 $ \mu(A) < \delta $ 的 $ A \in \mathcal{A} $ 都有 $ \left| \int_A f \, d\mu \right| < \varepsilon $。（提示：考虑定理 2.10。）

\end{exercises}

\section{收敛定理}

勒贝格积分理论与第 VI 章的黎曼理论形成鲜明对照，因为它包含非常一般且多用途的判别准则，用于判定取极限与积分是否可交换。因此，对于分析的需要而言，Bochner–Lebesgue 积分比（更简单的）黎曼积分更为适用。

与通常一样，我们在本节中总假定 $ (X, \mathcal{A}, \mu) $ 是完备的 σ-有限测度空间；$ E = (E, |\cdot|) $ 是巴拿赫空间。

\needspace{3\baselineskip}\subsection*{非负扩充实值函数的积分}

在数学、自然科学以及其他领域中积分的许多应用里，实值函数都扮演着重要角色。通常人们感兴趣的是这类情形下的可积函数，也就是有限的积分。然而我们已经提到，如果使理论涵盖 $ \overline{\mathbb{R}} $ 值函数的积分，既不排除函数取无穷值、也不排除积分取无穷值，那么理论在简洁性与优雅性上都会显著提升。作为从这种更具包容性的处理中获益的结果，我们举出单调收敛定理（习题 2.8 与定理 3.4）以及关于积分可交换性的富比尼–托内利定理（定理 6.11）。

由于实值情形的重要性，也由于它能提供依赖于 $ \mathbb{R} $ 与 $ \overline{\mathbb{R}} $ 的全序的有用补充结果，我们现在就来发展一套针对 $ \overline{\mathbb{R}} $ 值函数——特别是实值函数——的积分理论，作为 Bochner–Lebesgue 积分的补充。\footnote{针对 $ \overline{\mathbb{R}} $ 值函数的理论几乎是所有积分论教材的核心内容。它在某些方面比更一般的 Bochner–Lebesgue 理论更简单，而且如果只关心实值与复值函数，它已经够用；但它不能满足现代高等分析的需要，这正是我们选择更一般的处理方式的原因。}

根据定理 1.12，对每个 $ f \in \mathcal{L}_0(X,\mu,\overline{\mathbb{R}}^+) $，都存在 $ \mathcal{S}(X,\mu,\mathbb{R}^+) $ 中的递增序列 $ (f_j) $，它逐点收敛于 $ f $。把 $ f $ 的积分定义为递增序列 $ (\int_X f_j\,d\mu)_{j\in\mathbb{N}} $ 在 $ \overline{\mathbb{R}}^+ $ 中的极限，是自然的。只要我们能够保证该极限不依赖于 $ (f_j) $ 的选取，这一做法就是有意义的。

\begin{lemma}

设对 $ j \in \mathbb{N} $ 有 $ \varphi_j, \psi \in \mathcal{S}(X,\mu,\mathbb{R}^+) $。又设 $ (\varphi_j) $ 递增，且 $ \psi \le \lim_j \varphi_j $。则

\[
\int_X \psi\,d\mu \le \lim_j \int_X \varphi_j\,d\mu .
\]

\proofhead
设 $ \sum_{j=0}^m \alpha_j \chi_{A_j} $ 是 $ \psi $ 的标准形式，并固定 $ \lambda > 1 $。对 $ k \in \mathbb{N} $ 定义 $ B_k := [\lambda\varphi_k \ge \psi] $。由于 $ (\varphi_k) $ 递增且 $ \lambda > 1 $，我们有 $ B_k \subset B_{k+1} $ 对 $ k \in \mathbb{N} $ 成立，并且 $ \bigcup_{k\in\mathbb{N}} B_k = X $。测度的下连续性于是给出
\[
\int_X \psi\,d\mu = \sum_{j=0}^m \alpha_j \mu(A_j) = \lim_k \sum_{j=0}^m \alpha_j \mu(A_j \cap B_k) = \lim_k \int_X \psi\chi_{B_k}\,d\mu .
\]

由 $ B_k $ 的定义有 $ \lambda\varphi_k \ge \psi\chi_{B_k} $，于是得到

\[
\int_X \psi\,d\mu = \lim_k \int_X \psi\chi_{B_k}\,d\mu \le \lambda \lim_k \int_X \varphi_k\,d\mu .
\]

令 $ \lambda \downarrow 1 $ 取极限即完成证明。 \nobreak\hspace{0.35em}■

\end{lemma}

\begin{corollary}

设 $ (\varphi_j) $ 与 $ (\psi_j) $ 是 $ \mathcal{S}(X,\mu,\mathbb{R}^+) $ 中的递增序列，且 $ \lim_j \varphi_j = \lim_j \psi_j $。则

\[
\lim_j \int_X \varphi_j\,d\mu = \lim_j \int_X \psi_j\,d\mu \quad \text{在 } \overline{\mathbb{R}}^+ \text{ 中} .
\]

\proofhead
按假设，对 $ k \in \mathbb{N} $ 有 $ \psi_k \le \lim_j \psi_j = \lim_j \varphi_j $。由引理 3.1 得
\[
\int_X \psi_k\,d\mu \le \lim_j \int_X \varphi_j\,d\mu \quad \text{对 } k \in \mathbb{N} ,
\]

而令 $ k \to \infty $ 便得

\[
\lim_k \int_X \psi_k\,d\mu \le \lim_j \int_X \varphi_j\,d\mu .
\]

交换 $ (\varphi_j) $ 与 $ (\psi_j) $，即得反向的不等式，从而得到所要求的等式。 \nobreak\hspace{0.35em}■

\end{corollary}

设 $ (\varphi_j) $ 是 $ \mathcal{S}(X,\mu,\mathbb{R}^+) $ 中的递增序列，它逐点收敛于 $ f \in \mathcal{L}_0(X,\mu,\overline{\mathbb{R}}^+) $。我们称

\[
\int_X f\,d\mu := \lim_j \int_X \varphi_j\,d\mu
\]

为 $ f $ 关于测度 $ \mu $ 在 $ X $ 上的（勒贝格）积分。对 $ A \in \mathcal{A} $，

\[
\int_A f\,d\mu := \int_X \chi_A f\,d\mu
\]

为 $ f $ 在可测集 $ A $ 上的（勒贝格）积分。

\begin{remark}

\begin{items}
\item

对一切 $ f \in \mathcal{L}_0(X,\mu,\overline{\mathbb{R}}^+) $ 与 $ A \in \mathcal{A} $，$ \int_A f\,d\mu $ 都是良好定义的。

\proofhead
这由定理 1.12 与推论 3.2 推出。 \nobreak\hspace{0.35em}■
\item

若 $ f, g \in \mathcal{L}_0(X,\mu,\overline{\mathbb{R}}^+) $ 满足 $ f \le g $ μ-几乎处处，则 $ \int_X f\,d\mu \le \int_X g\,d\mu $。

\item

对 $ f \in \mathcal{L}_0(X,\mu,\overline{\mathbb{R}}^+) $，下列陈述等价：

\begin{subitems}
\item

$ \int_X f\,d\mu = 0 $。

\item

$ [f > 0] $ 是 μ-零集。

\item

$ f = 0 $ μ-几乎处处。

\proofhead
“（i）⇒（ii）” 令 $ A := [f > 0] $、$ A_j := [f > 1/j] $，$ j \in \mathbb{N}^\times $。则 $ (A_j) $ 是 $ \mathcal{A} $ 中的递增序列，且 $ A = \bigcup_j A_j $。又有 $ \chi_{A_j} \le jf $。由此推出
\[
0 \le \mu(A_j) = \int_X \chi_{A_j}\,d\mu \le j \int_X f\,d\mu = 0 \quad \text{对 } j \in \mathbb{N}^\times ,
\]

而由下连续性得 $ \mu(A) = \lim_j \mu(A_j) = 0 $。
“（ii）⇒（iii）” 显然。
“（iii）⇒（i）” 设 $ N $ 是 μ-零集，且对 $ x \in N^c $ 有 $ f(x) = 0 $。则\footnote{回忆约定 (2.1)。} $ f\chi_{N^c} = 0 $ 且 $ f\chi_N \le \infty\chi_N $。结合积分的定义（另见下面的 (d)）便得到

\[
0 \le \int_X f\,d\mu = \int_X f\chi_N\,d\mu + \int_X f\chi_{N^c}\,d\mu \le \infty\mu(N) = 0 .\quad\blacksquare
\]
\end{subitems}

\item

设 $ f, g \in \mathcal{L}_0(X,\mu,\overline{\mathbb{R}}^+) $ 且 $ \alpha \in [0,\infty] $。则

\[
\int_X (\alpha f + g)\,d\mu = \alpha \int_X f\,d\mu + \int_X g\,d\mu .
\]

\proofhead
我们考虑 $ \alpha = \infty $、$ g = 0 $ 的情形。令 $ \varphi_j := j\chi_{[f>0]} $，$ j \in \mathbb{N} $，则 $ f_j \uparrow \infty f $，从而
\begin{equation*}
\int_X (\infty f)\,d\mu = \begin{cases} 0 & \text{若 } \mu([f > 0]) = 0 , \\ \infty & \text{若 } \mu([f > 0]) > 0 . \end{cases}
\end{equation*}

由 (c) 现在可得 $ \int_X (\infty f)\,d\mu = \infty \int_X f\,d\mu $。其余结论容易由积分的定义推出，留作练习。 \nobreak\hspace{0.35em}■

\item

(i) 设 $ f \in \mathcal{L}_0(X,\mu,\mathbb{R}^+) $ 的勒贝格积分 $ \int_X f\,d\mu $ 有限。则 $ f $ 属于 $ \mathcal{L}_1(X,\mu,\mathbb{R}^+) $，且 $ f $ 在 $ X $ 上的勒贝格积分与 Bochner–Lebesgue 积分一致。

\begin{subitems}
\item

对 $ f \in \mathcal{L}_1(X,\mu,\mathbb{R}^+) $，勒贝格积分 $ \int_X f\,d\mu $ 有限，且与 Bochner–Lebesgue 积分一致。

\proofhead
(i) 定理 1.12 保证存在 $ \mathcal{S}(X,\mu,\mathbb{R}^+) $ 中的序列 $ (\varphi_j) $，使得 $ \varphi_j \uparrow f $。按假设，对每个 $ \varepsilon > 0 $ 都存在 $ N \in \mathbb{N} $，使得当 $ j \ge N $ 时有 $ \int_X f\,d\mu - \int_X \varphi_j\,d\mu < \varepsilon $。对 $ k \ge j \ge N $，由 $ \int_X f\,d\mu $ 的有限性得
\[
\int_X |\varphi_k - \varphi_j|\,d\mu = \int_X (\varphi_k - \varphi_j)\,d\mu \le \int_X (f - \varphi_j)\,d\mu = \int_X f\,d\mu - \int_X \varphi_j\,d\mu < \varepsilon .
\]

因此 $ (\varphi_j) $ 是 $ \mathcal{S}(X,\mu,\mathbb{R}^+) $ 中的 $ \mathcal{L}_1 $-柯西序列。这说明 $ f $ 属于 $ \mathcal{L}_1(X,\mu,\mathbb{R}^+) $。第二个结论是习题 2.8 的推论。

\end{subitems}

\begin{subitems}
\item

这由定理 1.12 与习题 2.8 推出。 \nobreak\hspace{0.35em}■

\end{subitems}

\item

对每个 $ f \in \mathcal{L}_0(X,\mu,\overline{\mathbb{R}}^+) $，有

\[
\int_X f\,d\mu = \sup \left\{ \int_X \varphi\,d\mu ; \varphi \in \mathcal{S}(X,\mu,\mathbb{R}^+) \text{ 且 } \varphi \le f \ \mu\text{-几乎处处} \right\} .\quad\blacksquare
\]
\end{items}

\end{remark}

\needspace{3\baselineskip}\subsection*{单调收敛定理}

我们现在证明习题 2.8 中对 $ \mathcal{L}_1(X,\mu,\mathbb{R}) $ 中函数所述的单调收敛定理的一个实质性加强。我们将看到，对于 $ \mathcal{L}_0(X,\mu,\overline{\mathbb{R}}^+) $ 中的递增序列，勒贝格积分与取极限可交换。

\begin{theorem}

（单调收敛） 设 $ (f_j) $ 是 $ \mathcal{L}_0(X,\mu,\overline{\mathbb{R}}^+) $ 中的递增序列。则

\[
\int_X \lim_j f_j\,d\mu = \lim_j \int_X f_j\,d\mu \quad \text{在 } \overline{\mathbb{R}}^+ \text{ 中} .
\]

证明

\begin{subitems}
\item

令 $ f := \lim_j f_j $。由命题 1.11，$ f $ 属于 $ \mathcal{L}_0(X,\mu,\overline{\mathbb{R}}^+) $，且对 $ j \in \mathbb{N} $ 有 $ f_j \le f $。于是由注 3.3(b)，对 $ j \in \mathbb{N} $ 有 $ \int f_j\,d\mu \le \int f\,d\mu $，从而 $ \lim_j \int f_j\,d\mu \le \int f\,d\mu $。

\item

设 $ \varphi \in \mathcal{S}(X,\mu,\mathbb{R}^+) $ 且 $ \varphi \le f $。取 $ \lambda > 1 $ 并令 $ A_j := [\lambda f_j \ge \varphi] $，$ j \in \mathbb{N} $。则 $ (A_j) $ 是 $ \mathcal{A} $ 中的递增序列，且 $ \bigcup_j A_j = X $、$ \lambda f_j \ge \varphi\chi_{A_j} $。此外 $ \varphi\chi_{A_j} \uparrow \varphi $，故

\[
\int_X \varphi\,d\mu = \lim_j \int_X \varphi\chi_{A_j}\,d\mu \le \lambda \lim_j \int_X f_j\,d\mu .
\]

令 $ \lambda \downarrow 1 $ 取极限，则对所有满足 $ \varphi \le f $ 的 μ-简单函数 $ \varphi $ 都有 $ \int_X \varphi\,d\mu \le \lim_j \int_X f_j\,d\mu $。由注 3.3(f) 可得 $ \int_X f\,d\mu \le \lim_j \int_X f_j\,d\mu $，于是证毕。 \nobreak\hspace{0.35em}■

\end{subitems}

\end{theorem}

\begin{corollary}

设 $ (f_j) $ 是 $ \mathcal{L}_0(X,\mu,\overline{\mathbb{R}}^+) $ 中的序列。则

\[
\sum_{j=0}^{\infty} \int_X f_j\,d\mu = \int_X \left( \sum_{j=0}^{\infty} f_j \right) d\mu \quad \text{在 } \overline{\mathbb{R}}^+ \text{ 中} .
\]

\proofhead
这由推论 1.13(iii) 与定理 3.4 推出。 \nobreak\hspace{0.35em}■
\end{corollary}

\begin{remark}

\begin{items}
\item

若序列不递增，定理 3.4 的结论可能不成立。

\proofhead
取 $ f_j := (1/j)\chi_{[0,j]} $，$ j \in \mathbb{N}^\times $。则 $ (f_j) $ 是 $ \mathcal{S}(\mathbb{R},\lambda_1,\mathbb{R}^+) $ 中一个（不增的）序列，它一致收敛于 0。但是序列 $ \int f_j\,d\lambda_1 $ 并不收敛于 0，因为对 $ j \in \mathbb{N}^\times $ 有 $ \int f_j\,d\lambda_1 = 1 $。 \nobreak\hspace{0.35em}■
\item

设 $ a_{j,k} \in \mathbb{R}^+ $，$ j, k \in \mathbb{N} $。则

\[
\sum_{j=0}^{\infty} \sum_{k=0}^{\infty} a_{jk} = \sum_{k=0}^{\infty} \sum_{j=0}^{\infty} a_{jk} .
\]

\proofhead
令 $ (X,\mu) := (\mathbb{N},\mathcal{H}^0) $，并定义 $ f_j : X \to \mathbb{R}^+ $ 为 $ f_j(k) := a_{jk} $，$ j, k \in \mathbb{N} $。则 $ (f_j) $ 是 $ \mathcal{L}_0(X,\mathcal{H}^0,\overline{\mathbb{R}}^+) $ 中的序列（见例 2.20），结论由推论 3.5 推出。 \nobreak\hspace{0.35em}■
\end{items}

对非负二重级数，这个结果比定理 II.8.10 更强，因为不再假设 $ \sum_{jk} a_{jk} $ 可和。

\end{remark}

\needspace{3\baselineskip}\subsection*{法图引理}

我们现在对 $ \mathcal{L}_0(X,\mu,\overline{\mathbb{R}}^+) $ 中的任意（不必递增的）序列证明单调收敛定理的一个推广。

\begin{theorem}

（法图引理） 对 $ \mathcal{L}_0(X,\mu,\overline{\mathbb{R}}^+) $ 中的每个序列 $ (f_j) $，有

\[
\int_X \left( \varliminf_j f_j \right) d\mu \le \varliminf_j \int_X f_j\,d\mu \quad \text{在 } \overline{\mathbb{R}}^+ \text{ 中} .
\]

\proofhead
令 $ g_j := \inf_{k\ge j} f_k $。由命题 1.11，$ g_j $ 属于 $ \mathcal{L}_0(X,\mu,\overline{\mathbb{R}}^+) $，且递增序列 $ (g_j) $ 收敛于 $ \varliminf_j f_j $。由定理 3.4 便得 $ \lim_j \int g_j\,d\mu = \int (\varliminf_j f_j)\,d\mu $。又 $ g_j \le f_k $，因此对 $ k \ge j $ 有 $ \int g_j\,d\mu \le \int f_k\,d\mu $。由此可得 $ \int g_j\,d\mu \le \inf_{k\ge j} \int f_k\,d\mu $，令 $ j \to \infty $ 取极限即完成证明。 \nobreak\hspace{0.35em}■
\end{theorem}

\qyxffntextdef{1}{假设 $ \int_X g\,d\mu < \infty $ 不能放宽（见习题 1）。}
\begin{corollary}

设 $ (f_j) $ 是 $ \mathcal{L}_0(X,\mu,\overline{\mathbb{R}}^+) $ 中的序列，$ g \in \mathcal{L}_0(X,\mu,\overline{\mathbb{R}}^+) $ 满足 $ \int_X g\,d\mu < \infty $，且对 $ j \in \mathbb{N} $ 有 $ f_j \le g $ μ-几乎处处。则\qyxffnmark{1}

\[
\varlimsup_j \int_X f_j\,d\mu \le \int_X \left( \varlimsup_j f_j \right) d\mu \quad \text{在 } \overline{\mathbb{R}}^+ \text{ 中} .
\]

\proofhead
设 $ N $ 是 μ-零集，使得对 $ x \in N^c $ 与 $ j \in \mathbb{N} $ 有 $ f_j(x) \le g(x) $。则 $ f_j \le g + \infty\chi_N $ 在 $ X $ 上成立，且 $ \int_X (g + \infty\chi_N)\,d\mu = \int_X g\,d\mu $（见注 3.3(c) 与 (d)）。因此我们可以不失一般性地假设对 $ j \in \mathbb{N} $ 有 $ f_j \le g $。令 $ g_j := g - f_j $，由法图引理得到
\[
\int_X \left( \varliminf_j g_j \right) d\mu = \int_X g\,d\mu - \int_X \left( \varlimsup_j f_j \right) d\mu \le \varliminf_j \int_X g_j\,d\mu = \int_X g\,d\mu - \varlimsup_j \int_X f_j\,d\mu .
\]

由于 $ \int_X g\,d\mu < \infty $，结论随之成立。 \nobreak\hspace{0.35em}■

\end{corollary}

作为第一个应用，我们证明可积函数的一个基本刻画。

\begin{theorem}

对 $ f \in \mathcal{L}_0(X,\mu,E) $，下列条件等价：

\begin{subitems}
\item

$ f \in \mathcal{L}_1(X,\mu,E) $；

\item

$ |f| \in \mathcal{L}_1(X,\mu,\mathbb{R}) $；

\item

$ \int_X |f|\,d\mu < \infty $。
若这些条件得到满足，则 $ \left| \int_X f\,d\mu \right| \le \|f\|_1 < \infty $。

\proofhead
“（i）⇒（ii）” 由引理 2.8 推出，“（ii）⇒（iii）” 显然。“（iii）⇒（ii）” 已在注 3.3(e) 中证明。
“（ii）⇒（i）” 设 $ (\varphi_j) $ 是 $ \mathcal{S}(X,\mu,E) $ 中 μ-几乎处处收敛于 $ f $ 的序列。令 $ A_j := [\,|\varphi_j| \le 2|f|\,] $、$ f_j := \varphi_j\chi_{A_j} $，$ j \in \mathbb{N} $。定理 1.7 与命题 1.9 表明 $ A_j $ 属于 $ \mathcal{A} $。于是 $ (f_j) $ 是 $ \mathcal{S}(X,\mu,E) $ 中的序列。
取 $ N \in \mathcal{A} $ 使 $ \mu(N) = 0 $，且对 $ x \in N^c $ 有 $ \varphi_j(x) \to f(x) $。若对某个 $ x \in N^c $ 有 $ f(x) \ne 0 $，则存在 $ k := k(x) \in \mathbb{N} $，使得对 $ j \ge k $ 有 $ |\varphi_j(x) - f(x)| \le 3|f(x)| $。因此当 $ j \ge k(x) $ 时 $ x \in N^c \cap [\,|f| > 0\,] $ 属于 $ A_j $。这意味着当 $ j \ge k(x) $ 时 $ f_j(x) = \varphi_j(x) $，从而对 $ x \in N^c \cap [\,|f| > 0\,] $ 有 $ f_j(x) \to f(x) $。若对某个 $ x \in N^c $ 有 $ f(x) = 0 $，则同样有 $ j \to \infty $ 时 $ f_j(x) \to f(x) $。由于 $ x $ 属于某个 $ A_k $（$ k \in \mathbb{N} $），我们得到 $ f_k(x) = \varphi_k(x) = 0 $，因为 $ |\varphi_k(x)| \le 2|f(x)| = 0 $。对 $ x \notin A_k $，同样有 $ f_k(x) = \chi_{A_k}(x)\varphi_k(x) = 0 $。这意味着 $ |f - f_j| \to 0 $ μ-几乎处处。现在显然对 $ j \in \mathbb{N} $ 有 $ |f - f_j| \le 3|f| $，故推论 3.8 给出
\[
\varlimsup_j \int_X |f - f_j|\,d\mu \le \int_X \varlimsup_j |f - f_j|\,d\mu = 0 .
\]

因此对每个 $ \varepsilon > 0 $ 都能找到 $ m \in \mathbb{N} $，使得当 $ j \ge m $ 时 $ \int |f - f_j|\,d\mu < \varepsilon/2 $。由此可知，对满足 $ j, k \ge m $ 的 $ j, k \in \mathbb{N} $，

\[
\|f_j - f_k\|_1 = \int_X |f_j - f_k|\,d\mu \le \int_X |f_j - f|\,d\mu + \int_X |f - f_k|\,d\mu < \varepsilon .
\]

因此 $ (f_j) $ 是 $ \mathcal{S}(X,\mu,E) $ 中的 $ \mathcal{L}_1 $-柯西序列，$ f $ 是 μ-可积的。
最后一条结论由定理 2.11(i) 推出。 \nobreak\hspace{0.35em}■

\end{subitems}

\end{theorem}

\begin{conclusion}

\begin{items}
\item

设 $ f \in \mathcal{L}_0(X,\mu,E) $，并设存在 $ \mathcal{L}_1(X,\mu,E) $ 中的序列 $ (f_j) $，使得 $ f_j \to f $ μ-几乎处处，且 $ \varliminf_j \|f_j\|_1 < \infty $。则 $ f $ 属于 $ \mathcal{L}_1(X,\mu,E) $，且 $ \|f\|_1 \le \varliminf_j \|f_j\|_1 $。

\proofhead
由引理 2.15，我们可以假设 $ (f_j) $ 在整个 $ X $ 上收敛于 $ f $。利用法图引理得
\[
\int_X |f|\,d\mu = \int_X \varliminf_j |f_j|\,d\mu \le \varliminf_j \int_X |f_j|\,d\mu < \infty ,
\]

再由定理 3.9 即得结论。 \nobreak\hspace{0.35em}■

\item

设 $ (f_j) $ 是 $ \mathcal{L}_1(X,\mu,\mathbb{R}^+) $ 中的序列。设存在 $ f \in \mathcal{L}_1(X,\mu,\mathbb{R}) $ 使得

\[
f_j \to f \ \mu\text{-几乎处处} \quad \text{且} \quad \int_X f_j\,d\mu \to \int_X f\,d\mu \quad (j \to \infty) .
\]

于是可推出\footnote{比较定理 2.18 的陈述。} $ (f_j) $ 在 $ \mathcal{L}_1(X,\mu,\mathbb{R}) $ 中收敛于 $ f $。

\proofhead
这里同样可以假设 $ (f_j) $ 在整个 $ X $ 上收敛于 $ f $。则 $ f \ge 0 $ 且 $ |f_j - f| \le f_j + f $。由定理 3.7 得
\[
2\int_X f\,d\mu = \int_X \varliminf_j (f_j + f - |f_j - f|)\,d\mu \le \varliminf_j \int_X (f_j + f - |f_j - f|)\,d\mu = 2\int_X f\,d\mu - \varlimsup_j \int_X |f_j - f|\,d\mu .
\]

按定理 3.9，$ \int_X f\,d\mu $ 有限，于是得到 $ \lim_j \int_X |f_j - f|\,d\mu = 0 $。 \nobreak\hspace{0.35em}■

\end{items}

\end{conclusion}

\needspace{3\baselineskip}\subsection*{扩充实值函数的积分}

将 $ \overline{\mathbb{R}} $ 值函数分解为正部与负部，也使我们能把勒贝格积分推广到可取负值的可测 $ \overline{\mathbb{R}} $ 值函数。若 $ \int_X f^+\,d\mu < \infty $ 且 $ \int_X f^-\,d\mu < \infty $，则称 $ f \in \mathcal{L}_0(X,\mu,\overline{\mathbb{R}}) $ 关于 $ \mu $ 是\textbf{勒贝格可积}的。此时称

\[
\int_X f\,d\mu := \int_X f^+\,d\mu - \int_X f^-\,d\mu
\]

为 $ f $ 关于测度 $ \mu $ 在 $ X $ 上的（勒贝格）积分。

\begin{remark}

证明 该假设蕴含 $ A := [|f| = \infty] $ 是 $ \mu $-可测的，并且也蕴含 $ \int_X |f| \, d\mu < \infty $。进一步有 $ \infty\chi_A \le |f| $，而由注 3.3(b) 与

\begin{items}
\item

对 $ f \in \mathcal{L}_0(X,\mu,\overline{\mathbb{R}}) $，下列三个陈述等价：

\begin{subitems}
\item

$ f $ 关于 $ \mu $ 勒贝格可积。

\item

$ \int_X |f|\,d\mu < \infty $；

\item

存在 $ g \in \mathcal{L}_1(X,\mu,\mathbb{R}) $，使得 $ |f| \le g $ μ-几乎处处。

\proofhead
“（i）⇒（ii）” 这是 $ |f| = f^+ + f^- $ 的推论。
“（ii）⇒（iii）” 定理 3.9 说明 $ |f| \in \mathcal{L}_1(X,\mu,\mathbb{R}) $。因此 (iii) 在 $ g = |f| $ 时成立。
“（iii）⇒（i）” 这由 $ f^+ \vee f^- \le |f| \le g $ 与注 3.3(b) 推出。 \nobreak\hspace{0.35em}■
\end{subitems}

\item

设 $ f \in \mathcal{L}_0(X,\mu,\overline{\mathbb{R}}) $。则 $ f $ 关于 $ \mu $ 勒贝格可积当且仅当 $ f $ 是 μ-可积的。此时 $ f $ 在 $ X $ 上的勒贝格积分等于 Bochner–Lebesgue 积分。换句话说，若考虑实值映射，则 $ \overline{\mathbb{R}} $ 值函数的勒贝格可积性定义与第 2 节中的定义一致。\footnote{另见推论 2.12(iii)。}

\proofhead
这由 (a)、定理 3.9 与注 3.3(e) 推出。 \nobreak\hspace{0.35em}■
\item

若 $ f \in \mathcal{L}_0(X,\mu,\overline{\mathbb{R}}) $ 关于 $ \mu $ 勒贝格可积，则 $ [|f| = \infty] $ 是 $ \mu $-零集。

\item

我们得出

\[
\infty\mu(A) = \int_X (\infty\chi_A) \, d\mu \le \int_X |f| \, d\mu < \infty.
\]

因此 $ \mu(A) = 0 $。\nobreak\hspace{0.35em}■

\end{items}

\end{remark}

\needspace{3\baselineskip}\subsection*{勒贝格控制收敛定理}

我们现在证明一个关于交换极限与积分的极其通用而实用的定理，它由 Henri Lebesgue 证明。它是勒贝格积分理论的基石之一，有着不计其数的应用。

\qyxffntextdef{1}{也称为“勒贝格定理”。}
\begin{theorem}

（控制收敛\qyxffnmark{1}） 设 $ (f_j) $ 是 $ \mathcal{L}_1(X,\mu,E) $ 中的序列，并假设存在 $ g \in \mathcal{L}_1(X,\mu,\mathbb{R}) $ 使得

\[
\text{(a)} \qquad |f_j| \le g \quad \mu\text{-几乎处处，对 } j \in \mathbb{N}.
\]

还假设对某个 $ f \in E^X $，

\[
\text{(b)} \qquad f_j \to f \quad \mu\text{-几乎处处，对 } j \to \infty.
\]

则 $ f $ 是 $ \mu $-可积的，$ f_j \to f $ 在 $ \mathcal{L}_1(X,\mu,E) $ 中，且 $ \int_X f_j \, d\mu \to \int_X f \, d\mu $ 在 $ E $ 中。

\proofhead
定义
\[
g_j := \sup_{k,\ell \ge j} |f_k - f_\ell|
\]

对 $ j \in \mathbb{N} $。由命题 1.11，$ (g_j) $ 是 $ \mathcal{L}_0(X,\mu,\overline{\mathbb{R}}^+) $ 中的序列，且 $ \mu $-几乎处处收敛到 0。又 $ |f_k - f_\ell| \le 2g $ $ \mu $-几乎处处对 $ k,\ell \in \mathbb{N} $，从而 $ |g_j| \le 2g $ $ \mu $-几乎处处对 $ j \in \mathbb{N} $。由推论 3.8 可得

\[
0 \le \varliminf_{j} \int_X g_j \, d\mu \le \int_X \varlimsup_{j} g_j \, d\mu = 0.
\]

因此 $ \left( \int_X g_j \, d\mu \right)_{j \in \mathbb{N}} $ 是一个（递减的）零序列。这意味着对每个 $ \varepsilon > 0 $ 存在 $ N \in \mathbb{N} $ 使得

\[
\int_X |f_k - f_\ell| \, d\mu \le \int_X \sup_{k,\ell \ge j} |f_k - f_\ell| \, d\mu < \varepsilon \quad \text{对 } k,\ell \ge j \ge N.
\]

因此 $ (f_j) $ 是 $ \mathcal{L}_1(X,\mu,E) $ 中的柯西序列，而结论由 $ \mathcal{L}_1(X,\mu,E) $ 的完备性与定理 2.18 得出。\nobreak\hspace{0.35em}■

\end{theorem}

\begin{remark}

注 3.6(a) 中的例子表明，可积控制函数 $ g $ 的存在性对定理 3.12 是本质的。\nobreak\hspace{0.35em}■

作为控制收敛定理的第一个应用，我们证明可测函数可积性的一个简单判别法。

\end{remark}

\begin{theorem}

（可积性判别法） 设 $ f \in \mathcal{L}_0(X,\mu,E) $ 与 $ g \in \mathcal{L}_1(X,\mu,\mathbb{R}) $ 满足 $ |f| \le g $ $ \mu $-几乎处处。则 $ f $ 属于 $ \mathcal{L}_1(X,\mu,E) $。

\proofhead
设 $ (\varphi_j) $ 是 $ \mathcal{S}(X,\mu,E) $ 中的序列，使得当 $ j \to \infty $ 时 $ \varphi_j \to f $ $ \mu $-几乎处处。令 $ A_j := [|\varphi_j| \le 2g] $ 与 $ f_j := \chi_{A_j}\varphi_j $，$ j \in \mathbb{N} $。则 $ (f_j) $ 是 $ \mathcal{S}(X,\mu,E) $ 中 $ \mu $-几乎处处收敛到 $ f $ 的序列（见定理 3.9 的证明）。因为 $ |f_j| \le 2g $ 对 $ j \in \mathbb{N} $，结论由控制收敛定理得出。\nobreak\hspace{0.35em}■
\end{theorem}

\begin{corollary}

\begin{subitems}
\item

取 $ f \in \mathcal{L}_1(X,\mu,E) $、$ g \in \mathcal{L}_0(X,\mu,\mathbb{K}) $ 与 $ \alpha \in [0,\infty) $，满足 $ |g| \le \alpha $ $ \mu $-几乎处处。则 $ gf $ 是 $ \mu $-可积的，且

\[
\left| \int_X gf \, d\mu \right| \le \alpha \|f\|_1.
\]
\item

取 $ f \in \mathcal{L}_0(X,\mu,E) $ 与 $ \alpha \in [0,\infty) $。若 $ |f| \le \alpha $ $ \mu $-几乎处处且 $ \mu(X) < \infty $，则 $ f $ 是 $ \mu $-可积的，且

\[
\left| \int_X f \, d\mu \right| \le \|f\|_1 \le \alpha\mu(X).
\]
\item

设 $ X $ 是 $ \sigma $-紧度量空间，$ \mu $ 是 $ X $ 上完备的拉东测度。假设 $ f \in C(X,E) $ 且 $ K \subset X $ 紧。则 $ \chi_K f $ 属于 $ \mathcal{L}_1(X,\mu,E) $，且

\[
\left| \int_K f \, d\mu \right| \le \|\chi_K f\|_\infty \mu(K).
\]

\proofhead
(i) 由注 1.2(d)，$ gf $ 是 $ \mu $-可测的。又 $ |gf| \le \alpha|f| $ $ \mu $-几乎处处，且 $ \alpha|f| $ 是 $ \mu $-可积的。因此定理 3.14 表明 $ gf $ 是 $ \mu $-可积的；定理 2.11(i) 与推论 2.16(ii) 蕴含
\[
\left| \int_X gf \, d\mu \right| \le \int_X |gf| \, d\mu \le \int_X \alpha|f| \, d\mu = \alpha\|f\|_1.
\]
\end{subitems}

\begin{subitems}
\item[（ii）]

因为 $ \mu(X) $ 有限，$ \chi_X $ 属于 $ \mathcal{L}_1(X,\mu,\mathbb{R}) $。由定理 1.7(i)，$ |f| $ 是 $ \mu $-可测的。因此 (i) 表明（取 $ g := |f| $ 与 $ f := \chi_X $）$ |f| $ 是 $ \mu $-可积的，且

\[
\int_X |f| \, d\mu \le \alpha\|\chi_X\|_1 = \alpha\mu(X) < \infty.
\]

结论现在由定理 3.9 得出。

\item[（iii）]

由定理 1.17，$ f $ 是 $ \mu $-可测的。此外 $ \chi_K $ 是 $ \mu $-简单的，因为由注 IX.5.3(a) $ \mu(K) $ 有限。因此 $ \chi_K f $ 是 $ \mu $-可测的，而结论由 (ii) 取 $ \alpha := \max_{x \in K}|f(x)| $ 得出。\nobreak\hspace{0.35em}■

\end{subitems}

\end{corollary}

在处理并非在整个 $ X $ 上有定义的函数时，积分理论中偶尔有用的是：在尚未定义之处把它定义为 0 来延拓其定义。可测性与可积性问题于是可以相对于测度空间 $ (X,\mathcal{A},\mu) $ 来讨论。为此我们提出以下约定。

对 $ f: \operatorname{dom}(f) \subset X \to E $，定义 $ f $ 到 $ X $ 的\textbf{平凡延拓} $ \tilde{f} \in E^X $ 为

\begin{equation*}
\tilde{f}(x) := \begin{cases} f(x), & \text{若 } x \in \operatorname{dom}(f), \\ 0, & \text{若 } x \notin \operatorname{dom}(f). \end{cases}
\end{equation*}

我们说 $ f $ 是 \textbf{$ \mu $-可测的}或 \textbf{$ \mu $-可积的}，如果 $ \tilde{f} $ 分别属于 $ \mathcal{L}_0(X,\mu,E) $ 或 $ \mathcal{L}_1(X,\mu,E) $。若 $ f $ 是 $ \mu $-可积的，我们令 $ \int_X f \, d\mu := \int_X \tilde{f} \, d\mu $。

\begin{theorem}

（级数的逐项积分） 设 $ (f_j) $ 是 $ \mathcal{L}_1(X,\mu,E) $ 中的序列，使得 $ \sum_{j=0}^{\infty} \int_X |f_j| \, d\mu < \infty $。则 $ \sum_j f_j $ $ \mu $-几乎处处绝对收敛，$ \sum_j f_j $ 是 $ \mu $-可积的，且

\[
\int_X \left( \sum_{j=0}^{\infty} f_j \right) d\mu = \sum_{j=0}^{\infty} \int_X f_j \, d\mu.
\]

证明

\begin{subitems}
\item

由定理 1.7(i) 与推论 1.13(iii)，$ \overline{\mathbb{R}} $ 值函数 $ g := \sum_{j=0}^{\infty} |f_j| $ 是 $ \mu $-可测的。推论 3.5 蕴含

\[
\int_X g \, d\mu = \sum_{j=0}^{\infty} \int_X |f_j| \, d\mu < \infty.
\]

\begin{items}
\item

因此由注 3.11(a) 与 (c) 可得，$ [g = \infty] $ 是 $ \mu $-零集，这证明了 $ \sum_j f_j $ 对几乎每个 $ x \in X $ 绝对收敛。

\end{items}

\item

令 $ g_k := \sum_{j=0}^{k} f_j $，并对 $ x \in [g < \infty] $ 令 $ f(x) := \sum_{j=0}^{\infty} f_j(x) $。序列 $ (g_k) $ $ \mu $-几乎处处收敛到 $ f $，且我们有估计 $ |g_k| \le \sum_{j=0}^{k} |f_j| \le g $。由控制收敛定理，$ f $ 属于 $ \mathcal{L}_1(X,\mu,E) $，且

\begin{equation*}
\begin{aligned}
\sum_{j=0}^{\infty} \int_X f_j \, d\mu = \lim_{k \to \infty} \int_X g_k \, d\mu = \int_X \lim_{k \to \infty} g_k \, d\mu = \int_X \left( \sum_{j=0}^{\infty} f_j \right) d\mu. \\quad\blacksquare
\end{aligned}
\end{equation*}
\end{subitems}

\end{theorem}

\needspace{3\baselineskip}\subsection*{含参积分}

作为控制收敛定理的另一个应用，我们研究含参积分的连续性与可微性。

\begin{theorem}

（含参积分的连续性） 设 $ M $ 是度量空间，$ f: X \times M \to E $ 满足

\begin{items}
\item

对每个 $ m \in M $，$ f(\cdot,m) \in \mathcal{L}_1(X,\mu,E) $；

\item

对 $ \mu $-几乎每个 $ x \in X $，$ f(x,\cdot) \in C(M,E) $；

\item

存在 $ g \in \mathcal{L}_1(X,\mu,E) $ 使得 $ |f(x,m)| \le g(x) $ 对 $ (x,m) \in X \times M $。
则

\[
F: M \to E, \quad m \mapsto \int_X f(x,m) \, \mu(dx)
\]

是良定义且连续的。

\proofhead
第一个断言立即由 (a) 得出。设 $ m \in M $，并设 $ (m_j) $ 是 $ M $ 中收敛到 $ m $ 的序列。我们令 $ f_j := f(\cdot,m_j) $，$ j \in \mathbb{N} $。由 (b) 可得 $ f_j \to f $ $ \mu $-几乎处处。因此由 (a) 与 (c)，我们可以对序列 $ (f_j) $ 应用控制收敛定理，并得到
\[
\lim_{j \to \infty} F(m_j) = \lim_{j \to \infty} \int_X f_j \, d\mu = \int_X \lim_{j \to \infty} f_j \, d\mu = \int_X f(x,m) \, \mu(dx) = F(m).
\]

结论现在由定理 III.1.4 得出。\nobreak\hspace{0.35em}■

\end{items}

\end{theorem}

\begin{theorem}

（含参积分的可微性） 设 $ U $ 在 $ \mathbb{R}^n $ 中开，或 $ U \subset \mathbb{K} $ 是完全且凸的，并设 $ f: X \times U \to E $ 满足

\begin{items}
\item

对每个 $ y \in U $，$ f(\cdot,y) \in \mathcal{L}_1(X,\mu,E) $；

\item

对 $ \mu $-几乎每个 $ x \in X $，$ f(x,\cdot) \in C^1(U,E) $；

\item

存在 $ g \in \mathcal{L}_1(X,\mu,\mathbb{R}) $ 使得

\[
\left| \frac{\partial}{\partial y^j} f(x,y) \right| \le g(x) \quad \text{对 } (x,y) \in X \times U \text{ 与 } 1 \le j \le n.
\]

则

\[
F: U \to E, \quad y \mapsto \int_X f(x,y) \, \mu(dx)
\]

是连续可微的，且

\[
\partial_j F(y) = \int_X \frac{\partial}{\partial y^j} f(x,y) \, \mu(dx) \quad \text{对 } y \in U \text{ 与 } 1 \le j \le n.
\]

\proofhead
取 $ y \in U $ 与 $ j \in \{1,\ldots,n\} $。设 $ (h_k) $ 是 $ \mathbb{K} $ 中的零序列，使得 $ h_k \ne 0 $ 且 $ y + h_k e_j \in U $ 对 $ k \in \mathbb{N} $。最后令
\[
f_k(x) := \frac{f(x, y + h_k e_j) - f(x, y)}{h_k} \quad \text{对 } x \in X \text{ 与 } k \in \mathbb{N},
\]

中值定理（定理 VII.3.9）于是给出

\[
|f_k(x)| \le \sup_{z \in U} \left| \frac{\partial}{\partial y^j} f(x,z) \right| \le g(x) \quad \mu\text{-几乎处处}.
\]

因为 $ (f_k) $ $ \mu $-几乎处处收敛到 $ \partial f(\cdot,y)/\partial y^j $，由定理 3.12 可得

\[
\lim_{k \to \infty} \frac{F(y + h_k e_j) - F(y)}{h_k} = \lim_{k \to \infty} \int_X f_k \, d\mu = \int_X \frac{\partial}{\partial y^j} f(x,y) \, \mu(dx).
\]

因此 $ F $ 是偏可微的，且 $ \partial_j F(y) = \int_X (\partial/\partial y^j) f(x,y) \, \mu(dx) $。结果现在由定理 3.17 与 VII.2.10 得出。\nobreak\hspace{0.35em}■

\end{items}

\end{theorem}

\begin{corollary}

设 $ U $ 在 $ \mathbb{C} $ 中开，$ f: X \times U \to \mathbb{C} $ 满足

\begin{items}
\item

对每个 $ z \in U $，$ f(\cdot,z) \in \mathcal{L}_1(X,\mu,\mathbb{C}) $；

\item

对 $ \mu $-几乎每个 $ x \in X $，$ f(x,\cdot) \in C^{\omega}(U,\mathbb{C}) $；

\item

存在 $ g \in \mathcal{L}_1(X,\mu,\mathbb{R}) $ 使得 $ |f(x,z)| \le g(x) $ 对 $ (x,z) \in X \times U $。
则

\[
F: U \to \mathbb{C}, \quad z \mapsto \int_X f(x,z) \, \mu(dx)
\]

是全纯的，且

\begin{equation}
F^{(n)}(z) = \int_X \frac{\partial^n}{\partial z^n} f(x,z) \, \mu(dx)\tag{3.1}
\end{equation}

对每个 $ n \in \mathbb{N} $。

\proofhead
取 $ z_0 \in U $ 与 $ r > 0 $ 使得 $ \overline{\mathbb{D}}(z_0,r) \subset U $。柯西导数公式（推论 VIII.5.12）给出
\[
\frac{\partial}{\partial z} f(x,z) = \frac{1}{2\pi i} \int_{\partial\mathbb{D}(z,r)} \frac{f(x,\zeta)}{(\zeta - z)^2} \, d\zeta \quad \text{对 } \mu\text{-几乎每个 } x \in X \text{ 与 } z \in \mathbb{D}(z_0,r),
\]

并由 (c) 与命题 VIII.4.3(iv) 得出

\[
\left| \frac{\partial}{\partial z} f(x,z) \right| \le \frac{g(x)}{r} \quad \text{对 } \mu\text{-几乎每个 } x \in X \text{ 与 } z \in \mathbb{D}(z_0,r).
\]

定理 3.18 现在表明 $ F|\mathbb{D}(z_0,r) $ 属于 $ C^1(\mathbb{D}(z_0,r),\mathbb{C}) $ 且满足

\[
F'(z) = \int_X \frac{\partial}{\partial z} f(x,z) \, \mu(dx) \quad \text{对 } z \in \mathbb{D}(z_0,r).
\]

全纯性是一个局部性质，因此定理 VIII.5.11 蕴含 $ F $ 属于 $ C^{\omega}(U,\mathbb{C}) $。(3.1) 的有效性现在由一个简单的归纳论证得出。\nobreak\hspace{0.35em}■

\end{items}

\end{corollary}

\exerciseshead

\begin{exercises}
\item

求一个测度空间 $ (X,\mathcal{A},\mu) $、$ \mathcal{L}_0(X,\mu,\mathbb{R}^+) $ 中的一个序列 $ (f_j) $ 以及 $ \mathcal{L}_0(X,\mu,\overline{\mathbb{R}}^+) $ 中的一个函数 $ g $，使得

\[
f_j \le g \text{ 对 } j \in \mathbb{N}, \quad \text{且} \quad \overline{\lim}_j \int_X f_j\,d\mu > \int_X \left(\overline{\lim}_j f_j\right) d\mu .
\]
\item

设 $ f \in \mathcal{L}_1(X,\mu,E) $ 且 $ \varepsilon > 0 $。证明存在 $ A \in \mathcal{A} $，使得

\[
\mu(A) < \infty \quad \text{且} \quad \left| \int_X f\,d\mu - \int_B f\,d\mu \right| < \varepsilon
\]

对每个满足 $ B \supset A $ 的 $ B \in \mathcal{A} $ 都成立。

\item

设 $ (f_j) $ 是 $ \mathcal{L}_1(X,\mu,E) $ 中的序列，且依测度收敛于 $ f \in \mathcal{L}_0(X,\mu,E) $。又设存在 $ g \in \mathcal{L}_1(X,\mu,\mathbb{R}) $，使得 $ |f_j| \le g $ 对一切 $ j \in \mathbb{N} $ 均 $ \mu $-几乎处处成立。则 $ f $ 属于 $ \mathcal{L}_1(X,\mu,E) $，

\[
f_j \to f \text{ 在 } \mathcal{L}_1(X,\mu,E) \text{ 中}, \quad \text{且} \quad \int_X f_j\,d\mu \to \int_X f\,d\mu \text{ 在 } E \text{ 中} .
\]

（提示：若 $ \left(\int_X f_j\,d\mu\right) $ 不收敛于 $ \int_X f\,d\mu $，则存在子列 $ (f_{j_k})_{k \in \mathbb{N}} $ 及 $ \delta > 0 $，使得

\begin{equation}
\|f_{j_k} - f\|_1 \ge \delta \quad \text{对 } k \in \mathbb{N} .\tag{3.2}
\end{equation}

利用习题 1.15 与定理 3.12 由 (3.2) 导出矛盾。）

\item

设 $ f, g \in \mathcal{L}_0(X,\mu,\overline{\mathbb{R}}) $ 是勒贝格可积函数。证明：

\begin{subitems}
\item

若 $ f \le g $ $ \mu $-几乎处处成立，则 $ \int_X f\,d\mu \le \int_X g\,d\mu $。

\item

\[ \left| \int_X f\,d\mu \right| \le \int_X |f|\,d\mu . \]

\item

$ f \wedge g $ 与 $ f \vee g $ 均勒贝格可积，且

\[
- \int_X (|f| + |g|)\,d\mu \le \int_X (f \wedge g)\,d\mu \le \int_X (f \vee g)\,d\mu \le \int_X (|f| + |g|)\,d\mu .
\]
\end{subitems}

\item

设 $ \mathcal{L}_0(X,\mu,\overline{\mathbb{R}}^+) $ 中的序列 $ (f_j) $ 依测度收敛于 $ f \in \mathcal{L}_0(X,\mu,\overline{\mathbb{R}}^+) $。证明

\[
\int_X f\,d\mu \le \underline{\lim}_j \int_X f_j\,d\mu .
\]
\item

对 $ x \in \mathbb{R}^n \setminus \{0\} $，定义

\begin{equation*}
k_n(x) := \begin{cases} x^+ & \text{若 } n = 1, \\ \log|x| & \text{若 } n = 2, \\ |x|^{2-n} & \text{若 } n \ge 3. \end{cases}
\end{equation*}

再设 $ U \subset \mathbb{R}^n $ 是非空开集，$ A \in \mathcal{L}(n) $ 满足 $ A \subset U^c $，且 $ f \in C_c(\mathbb{R}^n) $。

\begin{items}
\item

映射 $ A \to \mathbb{R} $，$ x \mapsto f(x)k_n(|y-x|) $ 对每个 $ y \in U $ 都是 $ \lambda_n $-可积的。

\item

映射 $ U \to \mathbb{R} $，$ y \mapsto \int_A f(x)k_n(|y-x|)\,\lambda_n(dx) $ 是光滑且调和的。

\end{items}

\item

验证：

\begin{subitems}
\item

$ \mathcal{L}_1(\mathbb{R}^n,\lambda_n,E) \cap BC(\mathbb{R}^n,E) \subsetneq C_0(\mathbb{R}^n,E) $；

\item

$ \mathcal{L}_1(\mathbb{R}^n,\lambda_n,E) \cap BUC(\mathbb{R}^n,E) \subseteq C_0(\mathbb{R}^n,E) $。

\end{subitems}

\end{exercises}

\section{勒贝格空间}

我们在推论 VI.7.4 中已经看到，紧区间 $ I $ 上连续的 $ \mathbb{K} $ 值函数构成的空间关于 $ L_2 $ 范数不完备。勒贝格积分理论的框架现在为我们提供了完备化内积空间 $ (C(I,\mathbb{K}),(\cdot\mid\cdot)_2) $ 的手段：我们将构造一个向量空间 $ L_2 $ 以及 $ (\cdot\mid\cdot)_2 $ 到 $ L_2 \times L_2 $ 上的延拓——仍记作 $ (\cdot\mid\cdot)_2 $——使得 $ (L_2,(\cdot\mid\cdot)_2) $ 是包含 $ C(I,\mathbb{K}) $ 作为稠密子空间的希尔伯特空间。

这一构造可以自然地推广，由此得到一族新的巴拿赫空间，即勒贝格 $ L_p $ 空间。这些空间在数学的许多领域中极为重要。

以下我们假设：

- $ (X,\mathcal{A},\mu) $ 是完备的 σ-有限测度空间；$ E = (E,|\cdot|) $ 是巴拿赫空间。

\needspace{3\baselineskip}\subsection*{本质有界函数}

称函数 $ f \in \mathcal{L}_0(X,\mu,E) $ 是 \textbf{μ-本质有界}的，若存在 $ \alpha \ge 0 $ 使得 $ \mu([|f| > \alpha]) = 0 $。此时 $ f $ 的 \textbf{μ-本质上确界}为\footnote{注意，现在 $ \|\cdot\|_\infty $ 有两种含义，即可测函数的本质上确界与有界函数的上确界范数。这两个值**未必**相同；见注 4.1 中的 (d) 与 (e)。必要时我们用 $ \|\cdot\|_{B(X,E)} $ 表示上确界范数。}

\[
\|f\|_\infty := \text{ess-sup}_{x \in X}|f(x)| := \inf\{\alpha \ge 0\ ;\ \mu([|f| > \alpha]) = 0\} .
\]

\begin{remark}

\begin{items}
\item

设 $ f \in \mathcal{L}_0(X,\mu,E) $。以下各项等价：

\begin{subitems}
\item

$ f $ 是 μ-本质有界的；

\item

$ \|f\|_\infty < \infty $；

\item

$ f $ μ-几乎处处有界。

\proofhead
“$ \text{(i)} \Rightarrow \text{(ii)} \Rightarrow \text{(iii)} $” 是显然的。
“$ \text{(iii)} \Rightarrow \text{(i)} $” 设 $ N $ 是 μ-零集，并取 $ \alpha \ge 0 $ 使得对 $ x \in N^c $ 有 $ |f(x)| \le \alpha $。则 $ [|f| > \alpha] \subset N $，而 $ \mu $ 的完备性蕴含 $ \mu([|f| > \alpha]) = 0 $。\nobreak\hspace{0.35em}■
\end{subitems}

\item

设 $ f \in \mathcal{L}_0(X,\mu,E) $。则 $ |f| \le \|f\|_\infty $ μ-几乎处处。

\proofhead
$ \|f\|_\infty = \infty $ 的情形是显然的。若 $ \|f\|_\infty < \infty $，则对每个 $ j \in \mathbb{N} $，$ [|f| > \|f\|_\infty + 2^{-j}] $ 是 μ-零集，从而集合 $ [|f| > \|f\|_\infty] = \bigcup_{j \in \mathbb{N}}[|f| > \|f\|_\infty + 2^{-j}] $ 也是 μ-零集。\nobreak\hspace{0.35em}■
\item

设 $ f $ 与 $ g $ 是 μ-本质有界的，且 $ \alpha \in \mathbb{K} $。则 $ \alpha f + g $ 也是 μ-本质有界的，且

\[
\|\alpha f + g\|_\infty \le |\alpha|\,\|f\|_\infty + \|g\|_\infty .
\]

\proofhead
由 (a) 与 (b)，存在 μ-零集 $ M $ 与 $ N $，使得对 $ x \in M^c $ 有 $ |f(x)| \le \|f\|_\infty $，对 $ x \in N^c $ 有 $ |g(x)| \le \|g\|_\infty $。因此
\[
|\alpha f(x) + g(x)| \le |\alpha|\,\|f\|_\infty + \|g\|_\infty \quad \text{对 } x \in (M \cup N)^c = M^c \cap N^c .
\]

于是 $ \alpha f + g $ 是 μ-本质有界的，且 $ \|\alpha f + g\|_\infty \le |\alpha|\,\|f\|_\infty + \|g\|_\infty $。\nobreak\hspace{0.35em}■

\item

设 $ f \in \mathcal{L}_0(X,\mu,E) $ 是有界的。则 $ \|f\|_\infty \le \|f\|_{B(X,E)} $（上确界范数）。若 $ N $ 是非空的 μ-零集，则 $ \|\chi_N\|_\infty = 0 $ 而 $ \|\chi_N\|_{B(X,E)} = 1 $。

\item

设 $ X $ 是 σ-紧度量空间，$ \mu $ 是 $ X $ 上的 massive 拉东测度。则

\[
\|f\|_\infty = \|f\|_{B(X,E)} \quad \text{对 } f \in BC(X,E) .
\]

\proofhead
由定理 1.17，任何 $ f \in BC(X,E) $ 都是 μ-可测的，而由 (d) 我们只需证明 $ \|f\|_{B(X,E)} \le \|f\|_\infty $。假定不然。则存在 $ x \in X $ 使得
\[
\|f\|_\infty < |f(x)| \le \|f\|_{B(X,E)} ,
\]

并且由 $ f $ 的连续性，存在 $ x $ 在 $ X $ 中的开邻域 $ O $ 使得对 $ y \in O $ 有 $ \|f\|_\infty < |f(y)| $。由 (b) 推出 $ \mu(O) = 0 $，这与 $ \mu $ 是 massive 的假设矛盾。\nobreak\hspace{0.35em}■

\end{items}

\end{remark}

\needspace{3\baselineskip}\subsection*{赫尔德不等式与闵可夫斯基不等式}

设 $ f \in \mathcal{L}_0(X,\mu,E) $。对 $ p \in (0,\infty) $，令

\[
\|f\|_p := \left(\int_X |f|^p\,d\mu\right)^{1/p}
\]

并约定 $ \infty^{1/p} := \infty $。我们定义 \textbf{$ X $ 上关于测度 $ \mu $ 的勒贝格空间}为\footnote{定理 3.9 表明，当 $ p = 1 $ 时记号 $ \mathcal{L}_p(X,\mu,E) $ 与第 2 节中的记号一致。以下我们集中讨论 $ p \in [1,\infty] $ 的勒贝格空间 $ \mathcal{L}_p $。$ p \in (0,1) $ 的情形将在习题 13 中处理。}

\[
\mathcal{L}_p(X,\mu,E) := \{f \in \mathcal{L}_0(X,\mu,E) ;\ \|f\|_p < \infty\} \quad \text{对 } p \in (0,\infty] .
\]

对 $ p \in [1,\infty] $，我们定义 \textbf{$ p $ 的对偶指数}为

\begin{equation*}
p' := \begin{cases} \infty, & \text{若 } p = 1, \\ p/(p-1), & \text{若 } p \in (1,\infty), \\ 1, & \text{若 } p = \infty. \end{cases}
\end{equation*}

按此约定，显然有

\[
\frac{1}{p} + \frac{1}{p'} = 1 \quad \text{对 } p \in [1,\infty] .
\]

现在我们可以陈述并证明两个重要的不等式。

\qyxffnnumdef{1}
\afterpage{\qyxffntextof{1}{当 $ p = 2 $ 时，这就是柯西–施瓦茨不等式。}}
\begin{theorem}

设 $ p \in [1,\infty] $。

\begin{subitems}
\item

对 $ f \in \mathcal{L}_p(X,\mu,\mathbb{K}) $ 与 $ g \in \mathcal{L}_{p'}(X,\mu,\mathbb{K}) $，有 $ fg \in \mathcal{L}_1(X,\mu,\mathbb{K}) $，且\qyxffnmark{1}

\[
\left|\int_X fg\,d\mu\right| \le \int_X |fg|\,d\mu \le \|f\|_p\|g\|_{p'} \qquad \text{（赫尔德不等式）} .
\]
\item

设 $ f, g \in \mathcal{L}_p(X,\mu,E) $。则 $ f + g \in \mathcal{L}_p(X,\mu,E) $，且

\[
\|f + g\|_p \le \|f\|_p + \|g\|_p \qquad \text{（闵可夫斯基不等式）} .
\]

\proofhead
(i) 先考虑 $ p = 1 $ 的情形。由注 4.1(b)，存在 μ-零集 $ N $ 使得对 $ x \in N^c $ 有 $ |g(x)| \le \|g\|_\infty $。于是由注 1.2(d)、注 3.3(b) 与引理 2.15 得到
\[
\int_{N^c}|fg|\,d\mu \le \|g\|_\infty \int_{N^c}|f|\,d\mu = \|f\|_1\|g\|_\infty < \infty .
\]

因此注 3.11(a)、定理 3.9 与引理 2.15 推出 $ fg $ 可积，而定理 2.11(i) 蕴含

\[
\left|\int_X fg\,d\mu\right| \le \int_X |fg|\,d\mu = \int_{N^c}|fg|\,d\mu \le \|f\|_1\|g\|_\infty .
\]

现在设 $ p \in (1,\infty) $。若

\begin{equation}
f = 0 \ \mu\text{-a.e.} \quad \text{或} \quad g = 0 \ \mu\text{-a.e.},\tag{4.1}
\end{equation}

则 $ fg $ 也 μ-几乎处处为零，结论由推论 2.16 推出。另一方面，若 (4.1) 不成立，推论 2.19 给出 $ \|f\|_p > 0 $ 与 $ \|g\|_{p'} > 0 $。此时令 $ \xi := |f|/\|f\|_p $、$ \eta := |g|/\|g\|_{p'} $，并由杨不等式（定理 IV.2.15）得到

\[
\frac{|fg|}{\|f\|_p\|g\|_{p'}} \le \frac{1}{p}\frac{|f|^p}{\|f\|_p^p} + \frac{1}{p'}\frac{|g|^{p'}}{\|g\|_{p'}^{p'}} .
\]

由此推出

\begin{equation*}
\begin{aligned}
\int_X |fg|\,d\mu &\le \frac{1}{p}\|f\|_p^{1-p}\|g\|_{p'}\int_X |f|^p\,d\mu + \frac{1}{p'}\|f\|_p\|g\|_{p'}^{1-p'}\int_X |g|^{p'}\,d\mu \\ &= \|f\|_p\|g\|_{p'} ,
\end{aligned}
\end{equation*}

于是由定理 3.9 可知 $ fg $ 属于 $ \mathcal{L}_1(X,\mu,E) $。因此

\[
\left|\int_X fg\,d\mu\right| \le \|fg\|_1 \le \|f\|_p\|g\|_{p'} .
\]

$ p = \infty $ 的情形与 $ p = 1 $ 的情形类似处理。

\end{subitems}

(ii) 由推论 2.9 与注 4.1(c)，只需考虑 $ p \in (1,\infty) $ 的情形。此外，我们可以不妨设 $ \|f + g\|_p > 0 $。我们先证明 $ f + g $ 属于 $ \mathcal{L}_p(X,\mu,E) $。注意不等式

\begin{equation}
|a + b|^p \le (2(|a| \vee |b|))^p \le 2^p(|a|^p + |b|^p) \quad \text{对 } a, b \in E,\tag{4.2}
\end{equation}

我们得到

\[
\int_X |f + g|^p\,d\mu \le 2^p\left(\int_X |f|^p\,d\mu + \int_X |g|^p\,d\mu\right) < \infty ,
\]

因为 $ f, g \in \mathcal{L}_p(X,\mu,E) $。因此 $ \|f + g\|_p < \infty $。由于等价关系

\[
|f + g|^{p-1} \in \mathcal{L}_{p'}(X,\mu,\mathbb{R}) \iff |f + g| \in \mathcal{L}_p(X,\mu,\mathbb{R}) ,
\]

由赫尔德不等式推出，对 $ h \in \mathcal{L}_p(X,\mu,E) $ 有

\[
\int_X |h|\,|f + g|^{p-1}\,d\mu \le \|h\|_p\left\||f + g|^{p-1}\right\|_{p'} = \|h\|_p\|f + g\|_p^{p/p'} ,
\]

并且我们得到

\begin{equation}
\begin{aligned}
\int_X |f + g|^p\,d\mu &\le \int_X |f|\,|f + g|^{p-1}\,d\mu + \int_X |g|\,|f + g|^{p-1}\,d\mu \\ &\le (\|f\|_p + \|g\|_p)\|f + g\|_p^{p/p'} .
\end{aligned}\tag{4.3}
\end{equation}

结论随之成立，因为 $ \|f + g\|_p < \infty $ 且 $ p/p' = p - 1 $。\nobreak\hspace{0.35em}■

\end{theorem}

\begin{corollary}

设 $ p \in [1,\infty] $。则 $ \mathcal{L}_p(X,\mu,E) $ 是 $ \mathcal{L}_0(X,\mu,E) $ 的向量子空间，且 $ \|\cdot\|_p $ 是 $ \mathcal{L}_p(X,\mu,E) $ 上的半范数。

\end{corollary}

\begin{remark}

\begin{items}
\item

令 $ \mathcal{N} := \{f \in \mathcal{L}_0(X,\mu,E) ;\ f = 0 \ \mu\text{-a.e.}\} $。对 $ f \in \mathcal{L}_0(X,\mu,E) $，以下命题等价：

\begin{subitems}
\item

对一切 $ p \in [1,\infty] $ 有 $ \|f\|_p = 0 $。

\item

对某个 $ p \in [1,\infty] $ 有 $ \|f\|_p = 0 $。

\item

$ f \in \mathcal{N} $。

\proofhead
“$ \text{(i)} \Rightarrow \text{(ii)} $” 是平凡的。“$ \text{(ii)} \Rightarrow \text{(iii)} $” 由推论 2.19 与注 4.1(b) 推出。
“$ \text{(iii)} \Rightarrow \text{(i)} $” 对 $ p \in [1,\infty) $，用引理 2.15。$ p = \infty $ 的情形是显然的。\nobreak\hspace{0.35em}■
\end{subitems}

\item

对每个 $ p \in [1,\infty] \cup \{0\} $，$ \mathcal{N} $ 是 $ \mathcal{L}_p(X,\mu,E) $ 的向量子空间。

\proofhead
$ p = 0 $ 的情形是显然的；特别地，$ \mathcal{N} $ 是向量空间。对 $ p \in [1,\infty] $，结论随后由 (a) 的 “$ \text{(iii)} \Rightarrow \text{(i)} $” 推出。\nobreak\hspace{0.35em}■
\item

对 $ p \in [1,\infty] $，我们有下列向量子空间的包含关系：

\[
\mathcal{S}(X,\mu,E) \subset \mathcal{L}_p(X,\mu,E) \subset \mathcal{L}_0(X,\mu,E) .
\]

\proofhead
显然每个 μ-简单函数都是 μ-本质有界的。取 $ p \in [1,\infty) $，并设 $ \varphi \in \mathcal{S}(X,\mu,E) $ 的标准形为 $ \sum_{j=0}^{m} e_j\chi_{A_j} $。则 $ |\varphi|^p \le \sum_{j=0}^{m}|e_j|^p\chi_{A_j} $，故 $ \|\varphi\|_p < \infty $。结论由注 1.2(a) 与推论 4.3 推出。\nobreak\hspace{0.35em}■
\end{items}

\end{remark}

\needspace{3\baselineskip}\subsection*{勒贝格空间的完备性}

我们现在推广定理 2.10(ii)，证明一切满足 $ p \in [1,\infty] $ 的勒贝格空间 $ \mathcal{L}_p(X,\mu,E) $ 都是完备的。对 $ p \in (1,\infty) $，这依赖于下面的引理。

\begin{lemma}

设 $ V $ 是向量空间，$ q $ 是 $ V $ 上的半范数。以下命题等价：

\begin{subitems}
\item

$ (V,q) $ 是完备的。

\item

对每个满足 $ \sum_{j=0}^{\infty} q(v_j) < \infty $ 的序列 $ (v_j) \in V^{\mathbb{N}} $，级数 $ \sum_j v_j $ 在 $ V $ 中收敛。

\proofhead
“$ \text{(i)} \Rightarrow \text{(ii)} $” 设 $ (v_j) \in V^{\mathbb{N}} $ 且 $ \sum_{j=0}^{\infty} q(v_j) < \infty $。对每个 $ \varepsilon > 0 $，存在 $ K \in \mathbb{N} $，使得对 $ \ell \ge K $ 有 $ \sum_{j=\ell+1}^{\infty} q(v_j) < \varepsilon $（见习题 II.7.4）。令 $ w_k := \sum_{j=0}^{k} v_j $（$ k \in \mathbb{N} $），得到
\[
q(w_m - w_\ell) = q\left(\sum_{j=\ell+1}^{m} v_j\right) \le \sum_{j=\ell+1}^{m} q(v_j) \le \sum_{j=\ell+1}^{\infty} q(v_j) < \varepsilon \quad \text{对 } m > \ell \ge K .
\]

因此 $ (w_k) $ 是 $ V $ 中的柯西序列，由 $ V $ 的完备性它收敛于某个 $ v \in V $。于是级数 $ \sum_j v_j $ 收敛。
“$ \text{(ii)} \Rightarrow \text{(i)} $” 设 $ (v_j) $ 是 $ V $ 中的柯西序列。对 $ k \in \mathbb{N} $，取 $ j_k \in \mathbb{N} $ 使得 $ q(v_{j_{k+1}} - v_{j_k}) < 2^{-(k+1)} $。令 $ w_k := v_{j_{k+1}} - v_{j_k} $，则 $ \sum_{k=0}^{\infty} q(w_k) \le 1 $，且由假设可以找到 $ v \in V $，使得当 $ \ell \to \infty $ 时 $ q(v - \sum_{k=0}^{\ell} w_k) \to 0 $。取 $ \varepsilon > 0 $ 与 $ L \in \mathbb{N} $，使得对 $ \ell \ge L $ 有 $ q(v - \sum_{k=0}^{\ell} w_k) < \varepsilon/2 $。因为 $ (v_j) $ 是 $ V $ 中的柯西序列，存在 $ K \ge L $，使得对 $ k, \ell \ge K $ 有 $ q(v_{j_{\ell+1}} - v_k) < \varepsilon/2 $。最后令 $ \tilde{v} := v + v_{j_0} $，则对 $ k \ge K $ 有

\begin{equation*}
\begin{aligned}
q(\tilde{v} - v_k) &= q(v + v_{j_0} - v_{j_{K+1}} + v_{j_{K+1}} - v_k) \\ &\le q\left(v - \sum_{k=0}^{K} w_k\right) + q(v_{j_{K+1}} - v_k) < \varepsilon .
\end{aligned}
\end{equation*}

这表明 $ (v_k) $ 收敛于 $ \tilde{v} $。\nobreak\hspace{0.35em}■

\end{subitems}

\end{lemma}

\begin{theorem}

对 $ p \in [1,\infty] $，$ \mathcal{L}_p(X,\mu,E) $ 是完备的。

证明

\begin{subitems}
\item

先考虑 $ p \in (1,\infty) $ 的情形。设 $ (f_j) $ 是 $ \mathcal{L}_p(X,\mu,E) $ 中满足 $ \sum_{j=0}^{\infty}\|f_j\|_p < \infty $ 的序列。令 $ g_k := \sum_{j=0}^{k}|f_j| $（$ k \in \mathbb{N} $）与 $ g := \sum_{k=0}^{\infty}|f_j| $。由推论 1.13(iii)，$ g $ 属于 $ \mathcal{L}_0(X,\mu,\overline{\mathbb{R}}^+) $，且 $ |g_k|^p \to |g|^p $。由于

\[
\|g_k\|_p \le \sum_{j=0}^{k}\|f_j\|_p \le \sum_{j=0}^{\infty}\|f_j\|_p < \infty ,
\]

结论 3.10(a) 告诉我们 $ g \in \mathcal{L}_p(X,\mu,\mathbb{R}) $。于是由注 3.11(c)，存在 μ-零集 $ N $，使得对 $ x \in N^c $ 有 $ g(x) < \infty $。因此，对每个 $ x \in N^c $，$ f(x) := \sum_{j=0}^{\infty} f_j(x) $ 由魏尔斯特拉斯判别法（定理 V.1.6）有定义。又因为 $ |f|^p \le g^p $ μ-几乎处处且 $ g \in \mathcal{L}_p(X,\mu,\mathbb{R}) $，定理 3.14 蕴含 $ f $ 属于 $ \mathcal{L}_p(X,\mu,E) $。最后，法图引理表明

\[
\left\|\tilde{f} - \sum_{j=0}^{k} f_j\right\|_p^p = \int_X \left|\varliminf_{\ell \to \infty}\sum_{j=k+1}^{\ell} f_j\right|^p d\mu \le \varliminf_{\ell \to \infty}\int_X \left|\sum_{j=k+1}^{\ell} f_j\right|^p d\mu = \varliminf_{\ell \to \infty}\left\|\sum_{j=k+1}^{\ell} f_j\right\|_p^p ,
\]

并且我们得到

\[
\left\|\tilde{f} - \sum_{j=0}^{k} f_j\right\|_p \le \varliminf_{\ell \to \infty}\sum_{j=k+1}^{\ell}\|f_j\|_p = \sum_{j=k+1}^{\infty}\|f_j\|_p \quad \text{对 } k \in \mathbb{N} .
\]

因为 $ \sum_{j=0}^{\infty}\|f_j\|_p < \infty $，序列 $ \left(\sum_{j=k+1}^{\infty}\|f_j\|_p\right)_{k \in \mathbb{N}} $ 收敛于零。因此 $ \left(\left\|\tilde{f} - \sum_{j=0}^{k} f_j\right\|_p\right)_{k \in \mathbb{N}} $ 也收敛于零。现在由引理 4.5 推出 $ \mathcal{L}_p(X,\mu,E) $ 是完备的。

\item

现在设 $ (f_j) $ 是 $ \mathcal{L}_\infty(X,\mu,E) $ 中的柯西序列。令

\[
A_j := [|f_j| > \|f_j\|_\infty] , \quad B_{k,\ell} := [|f_k - f_\ell| > \|f_k - f_\ell\|_\infty] \quad \text{对 } j, k, \ell \in \mathbb{N}
\]

以及 $ N := \bigcup_j A_j \cup \bigcup_{k,\ell} B_{k,\ell} $。由注 4.1(b) 与注 IX.2.5(b)，$ N $ 是零集，且

\[
|f_j(x)| \le \|f_j\|_\infty , \quad |f_k(x) - f_\ell(x)| \le \|f_k - f_\ell\|_\infty \quad \text{对 } j, k, \ell \in \mathbb{N},\ x \in N^c .
\]

因此 $ (f_j\,|\,N^c) $ 是巴拿赫空间 $ B(N^c,E) $ 中的柯西序列，我们可以找到 $ f \in B(N^c,E) $ 使得 $ (f_j\,|\,N^c) $ 一致收敛于 $ f $。于是 $ (f_j) $ μ-几乎处处收敛于 $ \tilde{f} $。我们知道函数 $ \tilde{f} $ 是 μ-本质有界的，因为 $ [|wtf| > \|f\|_{B(N^c,E)}] = \varnothing $，并且有

\[
|\tilde{f}(x) - f_j(x)| \le \|f - f_j\,|\,N^c\|_{B(N^c,E)} \quad \text{对 } x \in N^c \text{ 与 } j \in \mathbb{N} .
\]

于是 $ (f_j) $ 在 $ \mathcal{L}_\infty(X,\mu,E) $ 中收敛于 $ \tilde{f} $。

\item

$ p = 1 $ 的情形已在定理 2.10(ii) 中处理。\nobreak\hspace{0.35em}■

\end{subitems}

\end{theorem}

\begin{corollary}

设 $ p \in [1,\infty] $，并设 $ f_j, f \in \mathcal{L}_p(X,\mu,E) $ 满足在 $ \mathcal{L}_p(X,\mu,E) $ 中 $ f_j \to f $。

\begin{subitems}
\item

若 $ p = \infty $，则 $ (f_j) $ μ-几乎处处收敛于 $ f $。

\item

若 $ p \in [1,\infty) $，则存在 $ (f_j) $ 的子列 $ (f_{j_k})_{k \in \mathbb{N}} $，它 μ-几乎处处收敛于 $ f $。

\proofhead
因为 $ (f_j) $ 在 $ \mathcal{L}_p(X,\mu,E) $ 中收敛于 $ f $，我们知道 $ (f_j) $ 是 $ \mathcal{L}_p(X,\mu,E) $ 中的柯西序列。命题 (i) 现在可由定理 4.6 的证明直接推出。
若 $ p \in (1,\infty) $，选取 $ (f_j) $ 的子列 $ (f_{j_k})_{k \in \mathbb{N}} $ 使得 $ \|f_{j_{k+1}} - f_{j_k}\|_p < 2^{-(k+1)} $。则定理 4.6 的证明表明，存在 $ g \in \mathcal{L}_p(X,\mu,E) $，使得在 $ \mathcal{L}_p(X,\mu,E) $ 中 $ (f_{j_k} - f_{j_0}) \to g $，且当 $ k \to \infty $ 时 $ (f_{j_k} - f_{j_0}) \to g $ μ-几乎处处。因为 $ (f_j) $ 在 $ \mathcal{L}_p(X,\mu,E) $ 中收敛于 $ f $，我们有 $ \|f - (g + f_{j_0})\|_p = 0 $。注 4.4(a) 蕴含 $ f = g + f_{j_0} $ μ-几乎处处，由此即得结论。
$ p = 1 $ 的情形已在定理 2.18 中处理。\nobreak\hspace{0.35em}■
\end{subitems}

\end{corollary}

\qyxffntextdef{1}{当 $ p = \infty $ 时该命题可能不成立；见习题 8（另见习题 9）。}
\begin{proposition}

对 $ p \in [1,\infty) $，$ \mathcal{S}(X,\mu,E) $ 在 $ \mathcal{L}_p(X,\mu,E) $ 中稠密。\qyxffnmark{1}

\proofhead
设 $ f \in \mathcal{L}_p(X,\mu,E) $。则由注 4.4(c)，$ f $ 是 μ-可测的。因此存在 $ \mathcal{S}(X,\mu,E) $ 中的序列 $ (\varphi_j) $，使得当 $ j \to \infty $ 时 $ \varphi_j \to f $ μ-几乎处处。令 $ A_j := [|\varphi_j| \le 2|f|] $ 与 $ \psi_j := \chi_{A_j}\varphi_j $。则 $ (\psi_j) $ 是 $ \mathcal{S}(X,\mu,E) $ 中的序列，且 μ-几乎处处收敛于 $ f $。此外，
\[
|\psi_j - f|^p \le (|\psi_j| + |f|)^p \le 3^p|f|^p \quad \text{对 } j \in \mathbb{N} .
\]

因为 $ 3^p|f|^p $ 属于 $ \mathcal{L}_1(X,\mu,\mathbb{R}) $，我们可以应用控制收敛定理，得到

\[
\|\psi_j - f\|_p^p = \int_X |\psi_j - f|^p\,d\mu \to 0 \quad \text{当 } j \to \infty ,
\]

由此即得结论。\nobreak\hspace{0.35em}■

\end{proposition}

\needspace{3\baselineskip}\subsection*{$ L_p $ 空间}

我们在注 4.4(b) 中证明了

\[
\mathcal{N} := \{f \in \mathcal{L}_0(X,\mu,E) ;\ f = 0 \ \mu\text{-a.e.}\}
\]

对 $ p \in \{0\} \cup [1,\infty] $ 是 $ \mathcal{L}_p(X,\mu,E) $ 的向量子空间。因此商空间

\[
L_p(X,\mu,E) := \mathcal{L}_p(X,\mu,E)/\mathcal{N} \quad \text{对 } p \in \{0\} \cup [1,\infty]
\]

由例 I.12.3(i) 是 $ \mathbb{K} $ 上定义良好的向量空间。由注 4.4(c)，我们还有

\[
L_p(X,\mu,E) \subset L_0(X,\mu,E) \quad \text{对 } p \in [1,\infty] ,
\]

在向量子空间的意义下。设 $ [f] \in L_0(X,\mu,E) $，并设 $ g $ 是 $ [f] $ 的一个代表元。则 $ f - g \in \mathcal{N} $，即 $ f $ 与 $ g $ μ-几乎处处一致。由注 4.4(a)，映射

\[
|||\cdot||| : L_0(X,\mu,E) \to \overline{\mathbb{R}}^+ ,\quad [f] \mapsto \|f\|_p
\]

对每个 $ p \in [1,\infty] $ 都是定义良好的，且对 $ [f] \in L_p(X,\mu,E) $ 我们有

\begin{equation}
|||[f]|||_p = \|f\|_p = 0 \iff f = 0 \ \mu\text{-a.e.} \iff [f] = 0 .\tag{4.4}
\end{equation}

由于 $ |||\cdot|||_p $ 显然继承了半范数 $ \|\cdot\|_p $ 的性质，(4.4) 表明 $ |||\cdot|||_p $ 是 $ L_p(X,\mu,E) $ 上的范数。因此 $ L_p(X,\mu,E) $ 是赋范向量空间，而我们所由构造它的空间 $ \mathcal{L}_p(X,\mu,E) $ 只是半范空间。所以 $ \mathcal{L}_p(X,\mu,E) $ 中的极限一般并不唯一，而 $ L_p(X,\mu,E) $ 中的极限是唯一的。\footnote{见注 2.3(b)。} 我们为 $ L_p(X,\mu,E) $ 更好的拓扑结构所付出的代价是：它的元素不是 $ X $ 上的函数，而是 $ \mathcal{L}_p(X,\mu,E) $ 的向量子空间 $ \mathcal{N} $ 的陪集。换言之，我们把 μ-几乎处处相等的函数等同起来。经验表明，以下的简化记号不会引起误解。

约定 设 $ p \in \{0\} \cup [1, \infty] $。则我们把 $ L_p(X, \mu, E) $ 中的陪集 $ [f] = f + \mathcal{N} $ 写作 $ f $，并把彼此 $ \mu $-几乎处处相等的函数等同起来。此外，若 $ p \in [1, \infty] $，我们用 $ \|\cdot\|_p $ 表示 $ L_p(X, \mu, E) $ 中的范数，并令

\[
L_p(X, \mu, E) := (L_p(X, \mu, E), \|\cdot\|_p) \qquad \text{对 } p \in [1, \infty] .
\]

\begin{remark}

\begin{items}
\item

对 $ f \in L_0(X, \mu, E) $ 与 $ x \in X $，若 $ \mu $ 有非空零集，则 $ f(x) $ 无定义。也就是说，$ L_0(X, \mu, E) $ 的元素不能“逐点求值”。（当然，若选取 $ f $ 的一个代表元 $ f^{*} $，则 $ f^{*}(x) $ 有定义。）

\item

对 $ p \in [1, \infty] $，

\[
L_p(X, \mu, E) = \{ f \in L_0(X, \mu, E) ; \|f\|_p < \infty \} .
\]

\proofhead
“$ \subseteq $” 设 $ f \in L_p(X, \mu, E) $。$ f $ 的任一代表元 $ f^{*} $ 属于 $ \mathcal{L}_p(X, \mu, E) $，即它 $ \mu $-可测且满足 $ \|f^{*}\|_p < \infty $。因此 $ f $ 属于 $ L_0(X, \mu, E) $，且 $ \|f\|_p < \infty $。
“$ \supseteq $” 设 $ f \in L_0(X, \mu, E) $ 满足 $ \|f\|_p < \infty $。$ f $ 的每个代表元 $ f^{*} $ 都 $ \mu $-可测，且 $ \|f\|_p = \|f^{*}\|_p < \infty $。因此 $ f^{*} $ 属于 $ \mathcal{L}_p(X, \mu, E) $，从而 $ f $ 属于 $ L_p(X, \mu, E) $。\nobreak\hspace{0.35em}■
\item

设 $ f, g \in L_0(X, \mu, \mathbb{R}) $，并令 $ f^{*}, g^{*} $ 为 $ f, g $ 的代表元。若我们写

\[
f \le g :\iff f^{*} \le g^{*} \quad \mu\text{-几乎处处} ,
\]

就得到 $ L_0(X, \mu, \mathbb{R}) $ 上一个良定义的序 $ \le $，它使该空间成为向量格。

\proofhead
我们把这一简单证明留作习题。\nobreak\hspace{0.35em}■
\item

设 $ (F, \le) $ 是向量格，$ (F, \|\cdot\|) $ 是巴拿赫空间。若 $ |x| \le |y| $ 蕴含 $ \|x\| \le \|y\| $，则称 $ (F, \le, \|\cdot\|) $ 为\textbf{巴拿赫格}。

\item

对每个 $ p \in [1, \infty] $，$ (L_p(X, \mu, \mathbb{R}), \le, \|\cdot\|_p) $ 都是巴拿赫格。

\proofhead
显然 $ L_p(X, \mu, \mathbb{R}) $ 是 $ L_0(X, \mu, \mathbb{R}) $ 的向量子格。此外，由积分的单调性以及映射 $ t \mapsto t^p $ 的单调性立即推出，在 $ p \in [1, \infty) $ 的情形 $ L_p(X, \mu, \mathbb{R}) $ 是巴拿赫格。
设 $ f, g \in L_\infty(X, \mu, \mathbb{R}) $ 满足 $ |f| \le |g| $，并令 $ f^{*}, g^{*} $ 为它们的代表元。则 $ |f^{*}| \le |g^{*}| $ $ \mu $-几乎处处。此外，注 4.1(b) 表明 $ |g^{*}| \le \|g\|_\infty $ $ \mu $-几乎处处。因此 $ |f^{*}| \le \|g\|_\infty $ $ \mu $-几乎处处，从而 $ \|f\|_\infty \le \|g\|_\infty $。\nobreak\hspace{0.35em}■
\end{items}

\end{remark}

\begin{theorem}

\begin{subitems}
\item

对每个 $ p \in [1, \infty] $，$ L_p(X, \mu, E) $ 是巴拿赫空间。

\item

若 $ H $ 是希尔伯特空间，则 $ L_2(X, \mu, H) $ 关于内积

\[
(\cdot|\cdot)_2 : L_2(X, \mu, H) \times L_2(X, \mu, H) \to \mathbb{K} , \quad (f, g) \mapsto \int_X (f|g)_H \, d\mu .
\]

也是希尔伯特空间。

\proofhead
(i) 设 $ p \in [1, \infty] $。我们已经知道 $ L_p(X, \mu, E) $ 是赋范向量空间。设 $ (f_j) $ 是 $ L_p(X, \mu, E) $ 中的柯西序列，$ (f_j^{*}) $ 是相应的代表元序列。则 $ (f_j^{*}) $ 是 $ \mathcal{L}_p(X, \mu, E) $ 中的柯西序列。由定理 4.6，存在 $ f^{*} \in \mathcal{L}_p(X, \mu, E) $，使得当 $ j \to \infty $ 时 $ \|f_j^{*} - f^{*}\|_p \to 0 $。令 $ f := f^{*} + \mathcal{N} $，则 $ f \in L_p(X, \mu, E) $ 且 $ \|f_j - f\|_p = \|f_j^{*} - f^{*}\|_p \to 0 $。因此 $ L_p(X, \mu, E) $ 是完备的。
\end{subitems}

(ii) 利用定理 1.7 的 (i) 与 (iv) 以及赫尔德不等式，我们容易证明 $ (\cdot|\cdot)_2 $ 是 $ L_2(X, \mu, H) $ 上的内积，且对 $ f \in L_2(X, \mu, H) $ 有 $ |(f|f)_2| = \|f\|_2^2 $。于是由 (i) 即得断言。\nobreak\hspace{0.35em}■

\end{theorem}

\begin{corollary}

$ L_2(X, \mu, \mathbb{K}) $ 关于内积

\[
(f|g)_2 = \int_X f \bar{g} \, d\mu \qquad \text{对 } f, g \in L_2(X, \mu, \mathbb{K}) .
\]

是希尔伯特空间。

\end{corollary}

\needspace{3\baselineskip}\subsection*{具有紧支集的连续函数}

设 $ Y $ 是拓扑空间。对 $ f \in E^Y $，我们称

\[
\operatorname{supp}(f) := \overline{\{ x \in Y ; f(x) \ne 0 \}}
\]

为 $ f $ 的\textbf{支集}。这里照例用横线表示（在 $ Y $ 中的）闭包。具有紧支集的连续函数特别重要。因此我们定义

\[
C_c(Y, E) := \{ f \in C(Y, E) ; \operatorname{supp}(f) \text{ 是紧的} \} .
\]

\begin{example}

\begin{items}
\item

对例 III.1.3(c) 中的狄利克雷函数 $ \chi_{\mathbb{Q}} \in \mathbb{R}^{\mathbb{R}} $，我们有

\[
\operatorname{supp}(\chi_{\mathbb{Q}}) = \operatorname{supp}(\chi_{\mathbb{R}-\mathbb{Q}}) = \mathbb{R} .
\]

\proofhead
这由命题 I.10.8 与 I.10.11 推出。\nobreak\hspace{0.35em}■
\item

设 $ X = \mathbb{Z} $ 或 $ X = \mathbb{N} $，并为 $ X $ 配备由 $ \mathbb{R} $ 诱导的度量。设 $ \mathcal{H}^0 $ 是 $ \mathfrak{P}(X) $ 上的计数测度。则\footnote{比较例 2.20。}

\[
C_c(X, E) = \mathcal{S}(X, \mathcal{H}^0, E) = \{ \varphi \in E^X ; \operatorname{Num}[\varphi \ne 0] < \infty \} .
\]
\item

设 $ X $ 是度量空间。则 $ C_c(X, E) $ 是 $ BC(X, E) $ 的向量子空间。若 $ X $ 是紧的，则 $ C_c(X, E) = C(X, E) = BC(X, E) $。

\proofhead
第一个断言由推论 III.3.7 推出。第二个由习题 III.3.2 与推论 III.3.7 推出。\nobreak\hspace{0.35em}■
\end{items}

\end{example}

\begin{proposition}

设 $ X $ 是度量空间，$ A $ 与 $ B $ 是 $ X $ 的闭的、不交的非空子集。则存在 $ \varphi \in C(X) $ 使得 $ 0 \le \varphi \le 1 $，$ \varphi|A = 1 $，且 $ \varphi|B = 0 $。这样的函数称为\textbf{乌雷松函数}。

\proofhead
若 $ D \subset X $ 非空，例 III.1.3(l) 表明距离函数 $ d(\cdot, D) $ 属于 $ C(X) $。若 $ D $ 还是闭的，则 $ d(x, D) = 0 $ 当且仅当 $ x \in D $。利用这些性质，我们容易证明由
\[
\varphi(x) := \frac{d(x, B)}{d(x, A) + d(x, B)} \qquad \text{对 } x \in X ,
\]

定义的函数具有所述性质。\nobreak\hspace{0.35em}■

\end{proposition}

借助乌雷松函数，我们现在可以证明一个重要的逼近定理。

\begin{theorem}

设 $ X $ 是 $ \sigma $-紧度量空间，$ \mu $ 是 $ X $ 上的拉东测度。则对 $ p \in [1, \infty) $，$ C_c(X, E) $ 是 $ \mathcal{L}_p(X, \mu, E) $ 的稠密向量子空间。

\proofhead
设 $ \varepsilon > 0 $。根据命题 4.8，$ \mathcal{S}(X, \mu, E) $ 在 $ \mathcal{L}_p(X, \mu, E) $ 中稠密。因此，由于定理 1.17 与闵可夫斯基不等式（即三角不等式），只需验证：对每个具有有限测度的 $ \mu $-可测集 $ A $ 与每个 $ e \in E \setminus \{0\} $，存在 $ f \in C_c(X, E) $ 使得 $ \|f - \chi_A e\|_p < \varepsilon $。

于是设 $ A \in \mathcal{A} $ 且 $ \mu(A) < \infty $。因为 $ \mu $ 是正则的，我们可以找到 $ X $ 的紧子集 $ K $ 与开子集 $ U $，使得 $ K \subset A \subset U $ 且
\[
\mu(U \setminus K) = \mu(U) - \mu(K) < (\varepsilon/|e|)^p .
\]

命题 4.13 保证存在 $ X $ 上的乌雷松函数 $ \varphi $，满足 $ \varphi|K = 1 $ 且 $ \varphi|U^c = 0 $。令 $ f := \varphi e $，我们得到所需的

\[
\|\chi_A e - f\|_p^p \le |e|^p \int_X \chi_{U \setminus K} \, d\mu \le |e|^p \mu(U \setminus K) < \varepsilon^p .\quad\blacksquare
\]
\end{theorem}

\needspace{3\baselineskip}\subsection*{嵌入}

设 $ X $ 与 $ Y $ 是拓扑空间，且 $ X $ 是 $ Y $ 的子集。用 $ j : X \to Y $，$ x \mapsto x $ 表示 $ X $ 到 $ Y $ 的包含映射\footnote{见例 I.3.2(b)。}，我们说 $ X $ \textbf{连续嵌入} $ Y $，若 $ j $ 是连续的\footnote{当 $ X $ 未配备 $ Y $ 诱导的拓扑时，这些概念变得重要；见注 4.15(a)。}。在这种情形，我们写 $ X \hookrightarrow Y $。若 $ X $ 还是 $ Y $ 的稠密子集，我们写 $ X \overset{d}{\hookrightarrow} Y $。若 $ X $ 与 $ Y $ 是向量空间，则记号 $ X \hookrightarrow Y $（以及术语“连续嵌入”）将总是指：除此之外，$ X $ 还是 $ Y $ 的向量子空间，而不仅仅是某个任意子集。

\begin{remark}

\begin{items}
\item

设 $ V $ 与 $ W $ 是赋范向量空间。$ V $ 连续嵌入 $ W $ 当且仅当 $ V $ 是 $ W $ 的向量子空间，且存在 $ \alpha > 0 $ 使得对 $ v \in V $ 有 $ \|v\|_W \le \alpha \|v\|_V $，也就是说，$ V $ 的范数比 $ W $ 在 $ V $ 上诱导的范数更强。
若 $ V $ 携带 $ W $ 诱导的范数，则总有 $ V \hookrightarrow W $。

\item

设 $ X $ 在 $ \mathbb{R}^n $ 中是开的。则

\[
BUC^k(X, E) \hookrightarrow BUC^\ell(X, E) \qquad \text{对 } k \ge \ell .
\]

若 $ X $ 还是有界的，则

\[
BUC^k(X, \mathbb{K}) \overset{d}{\hookrightarrow} BUC(X, \mathbb{K}) \qquad \text{对 } k \in \mathbb{N} .
\]

\proofhead
第一个断言是显然的。第二个由斯通-魏尔斯特拉斯逼近定理（推论 V.4.8）以及应用 VI.2.2 推出。\nobreak\hspace{0.35em}■
\end{items}

简单的例子（见习题 5.1）表明，勒贝格空间一般并不互相包含。在对测度空间 $ (X, \mathcal{A}, \mu) $ 作适当的附加假设后，勒贝格空间之间存在连续嵌入。例如，若 $ \mathcal{H}^0 $ 是 $ \mathfrak{P}(\mathbb{N}) $ 上的计数测度，则习题 1.16 中引入的空间 $ \ell_p $ 与 $ \mathcal{L}_p(\mathbb{N}, \mathcal{H}^0, \mathbb{K}) $ 一致（$ 1 \le p \le \infty $），并且我们有嵌入

\[
\ell_1 \hookrightarrow \ell_p \hookrightarrow \ell_q \hookrightarrow \ell_\infty \qquad \text{对 } 1 \le p \le q \le \infty ,
\]

（见习题 11）。

有限测度空间呈现出完全不同的情形：

\end{remark}

\begin{theorem}

设 $ (X, \mathcal{A}, \mu) $ 是有限完备测度空间。则

\[
L_q(X, \mu, E) \overset{d}{\hookrightarrow} L_p(X, \mu, E) \qquad \text{对 } 1 \le p < q \le \infty
\]

且

\begin{equation}
\|f\|_p \le \mu(X)^{1/p-1/q} \|f\|_q \qquad \text{对 } f \in L_q(X, \mu, E) .\tag{4.5}
\end{equation}

证明

\begin{subitems}
\item

取 $ f \in L_q(X, \mu, E) $ 并令 $ r := q/p $。设 $ g \in \mathcal{L}_q(X, \mu, E) $ 是 $ f $ 的代表元。则 $ |g|^p $ 属于 $ \mathcal{L}_r(X, \mu, \mathbb{R}) $，且 $ 1/r' = (q-p)/q $。此外 $ \chi_X $ 属于 $ \mathcal{L}_{r'}(X, \mu, \mathbb{R}) $，因为 $ \mu $ 是有限测度。于是在 $ q < \infty $ 的情形，赫尔德不等式给出

\[
\|g\|_p^p = \int_X \chi_X |g|^p \, d\mu \le \left( \int_X \chi_X^{r'} \, d\mu \right)^{1/r'} \left( \int_X |g|^{pr} \, d\mu \right)^{1/r} = \mu(X)^{(q-p)/q} \|g\|_q^p ,
\]

并且我们得到 $ \|g\|_p \le \mu(X)^{1/p-1/q} \|g\|_q $；这显然在 $ q = \infty $ 的情形也成立。由于 $ g $ 是 $ f $ 的任意代表元，我们看到 $ f $ 属于 $ L_p(X, \mu, E) $ 且 (4.5) 成立。由注 4.15(a) 推出 $ L_q(X, \mu, E) \hookrightarrow L_p(X, \mu, E) $。

\item

$ M := \{ [\varphi] \in L_0(X, \mu, E) ; \varphi \in \mathcal{S}(X, \mu, E) \} $ 满足 $ M \subset L_q(X, \mu, E) $，且由于 $ p < \infty $，命题 4.8 蕴含 $ M $ 在 $ L_p(X, \mu, E) $ 中稠密。因此 $ L_q(X, \mu, E) $ 在 $ L_p(X, \mu, E) $ 中也稠密。\nobreak\hspace{0.35em}■

\end{subitems}

\end{theorem}

下一个定理表明，在 massive 拉东测度 $ \mu $ 的情形，$ L_0(X, \mu, E) $ 的元素至多有一个连续代表元。于是在这种情形，我们可以把 $ C(X, E) $ 中的每个函数与它在 $ L_0(X, \mu, E) $ 中生成的等价类等同起来，并把 $ C(X, E) $ 看作 $ L_0(X, \mu, E) $ 的向量子空间。

\begin{proposition}

设 $ \mu $ 是 $ \sigma $-紧空间 $ X $ 上的 massive 拉东测度。则映射

\begin{equation}
C(X, E) \to L_0(X, \mu, E) , \quad f \mapsto [f]\tag{4.6}
\end{equation}

是线性的且是单射。

\proofhead
定理 1.17 表明映射 (4.6) 是良定义的且线性的。

取 $ f, g \in C(X, E) $ 使得 $ [f] = [g] $。存在 $ h \in \mathcal{N} $ 使得 $ f - g = h $，即 $ f - g = 0 $ $ \mu $-几乎处处。反设对某个 $ x \in X $ 有 $ f(x) \ne g(x) $。由连续性，$ (f - g)(y) \ne 0 $ 对 $ x $ 的某个开邻域 $ U $ 中的所有 $ y $ 成立。但 $ \mu(U) > 0 $，与 $ f - g = 0 $ $ \mu $-几乎处处的假设矛盾。因此 $ f = g $，这证明了所说的单射性。\nobreak\hspace{0.35em}■
\end{proposition}

约定 设 $ \mu $ 是 $ \sigma $-紧空间 $ X $ 上的 massive 拉东测度。我们把 $ C(X, E) $ 与它在单射 (4.6) 下的像等同起来，从而把 $ C(X, E) $ 看作 $ L_0(X, \mu, E) $ 的向量子空间。则

\[
\|f\|_{B(X,E)} = \|f\|_\infty \qquad \text{对 } f \in BC(X, E) .
\]

下述结果是这一约定的一个简单推论。

\begin{theorem}

设 $ \mu $ 是 $ \sigma $-紧度量空间 $ X $ 上的 massive 拉东测度。

\begin{subitems}
\item

对每个 $ p \in [1, \infty) $，$ C_c(X, E) $ 是 $ L_p(X, \mu, E) $ 的稠密向量子空间。

\item

$ BC(X, E) $ 是 $ L_\infty(X, \mu, E) $ 的闭向量子空间。

\proofhead
第一个断言由定理 4.14 推出。第二个是显然的。\nobreak\hspace{0.35em}■
\end{subitems}

\end{theorem}

\needspace{3\baselineskip}\subsection*{$ L_p $ 上的连续线性泛函}

本节余下部分，我们对 $ p \in [1, \infty] $ 使用缩写

\[
L_p(X) := L_p(X, \mu, \mathbb{K}) \quad \text{与} \quad L'_p(X) := (L_p(X))' ,
\]

右边的撇号表示对偶空间（注 VII.2.13(a)）。由赫尔德不等式推出，对每个 $ f \in L_{p'}(X) $，映射

\[
T_f : L_p(X) \to \mathbb{K} , \quad g \mapsto \int_X fg \, d\mu
\]

是 $ L_p(X) $ 上的连续线性泛函，即 $ L'_p(X) $ 的元素；它满足

\begin{equation}
\|T_f\|_{L'_p(X)} \le \|f\|_{p'} .\tag{4.7}
\end{equation}

事实上 (4.7) 中等号成立：

\begin{proposition}

映射

\[
T : L_{p'}(X) \to L'_p(X) , \quad f \mapsto T_f
\]

对每个 $ p \in [1, \infty] $ 都是线性等距。

证明

\begin{subitems}
\item

显然 $ T $ 是线性的。此外，鉴于 (4.7)，我们只需证明：对每个满足 $ f \ne 0 $ 的 $ f \in L_{p'}(X) $ 与每个 $ \varepsilon > 0 $，存在 $ g \in L_p(X) $ 使得

\[
\|g\|_p = 1 \quad \text{且} \quad \|f\|_{p'} < \left| \int_X fg \, d\mu \right| + \varepsilon .
\]
\item

先设 $ p \in (1, \infty) $，从而 $ p' \in (1, \infty) $。因此

\[
g := \overline{\operatorname{sign} f} \, \|f\|_{p'}^{1-p'} |f|^{p'-1}
\]

是良定义的且 $ \mu $-可测（见习题 1.19 与定理 1.7(i)）。此外

\[
\int_X |g|^p \, d\mu = \|f\|_{p'}^{p(1-p')} \int_X |f|^{p(p'-1)} \, d\mu = \|f\|_{p'}^{-p'} \|f\|_{p'}^{p'} = 1
\]

且 $ fg = \|f\|_{p'}^{1-p'} |f|^{p'} $。因此 $ \|f\|_{p'} = \int_X fg \, d\mu $。
对 $ p = \infty $，我们令 $ g := \overline{\operatorname{sign} f} $。则

\[
\|g\|_\infty = 1 \quad \text{且} \quad \|f\|_1 = \int_X fg \, d\mu .
\]
\item

现在设 $ p = 1 $。设 $ 0 < \varepsilon < \|f\|_\infty $ 并令 $ \alpha := \|f\|_\infty - \varepsilon $。因为 $ [\,|f| > \alpha\,] $ 有正测度且 $ \mu $ 是 $ \sigma $-有限的，我们可以找到 $ A \in \mathcal{A} $ 使得 $ A \subset [\,|f| > \alpha\,] $ 且 $ \mu(A) \in (0, \infty) $。因此 $ g := \overline{\operatorname{sign} f} \, (1/\mu(A)) \chi_A $ 是良定义的且 $ \mu $-可测。显然 $ \|g\|_1 = 1 $ 且

\[
\int_X fg \, d\mu = \frac{1}{\mu(A)} \int_A |f| \, d\mu \ge \alpha = \|f\|_\infty - \varepsilon .
\]

这就完成了证明。\nobreak\hspace{0.35em}■

\end{subitems}

\end{proposition}

\begin{remark}

\begin{items}
\item

可以证明，命题 4.19 中的映射 $ T $ 对每个 $ p \in [1,\infty) $ 都是满射；也就是说，$ L_p(X) $ 上每个连续线性泛函都可表示为适当选取的 $ f \in L_{p'}(X) $ 的 $ T_f $ 的形式，例如参见 [Rud83, Theorem 6.1.6]。从而对每个 $ p \in [1,\infty) $，$ T: L_{p'}(X) \to L'_p(X) $ 是等距同构。这个同构允许我们对 $ p \in [1,\infty) $ 把 $ L_{p'}(X) $ 与 $ L'_p(X) $ 等同起来。对偶配对 $ \langle\cdot,\cdot\rangle_{L_p}: L'_p(X) \times L_p(X) \to \mathbb{K} $ 满足

\[
\langle g,f\rangle_{L_p} = \int_X f g \, d\mu \quad \text{对 } (g,f) \in L_{p'}(X) \times L_p(X) .
\]
\item

当 $ p = \infty $ 时，映射 $ T: L_1(X) \to L'_\infty(X) $ 一般不是满射；参见 [Fol99, S. 191]。

\item

记 $ \langle\cdot,\cdot\rangle_E: E' \times E \to \mathbb{K} $ 为 $ E $ 与 $ E' $ 之间的对偶配对。则映射

\[
\kappa: E \to [E']', \quad e \mapsto \langle\cdot,e\rangle_E
\]

是线性且有界的，其范数不超过 1。

\proofhead
显然 $ \kappa $ 是线性的。设 $ e \in E $ 满足 $ \|e\|_E \le 1 $。则
\[
|\langle\kappa(e),e'\rangle_{E'}| = |\langle e',e\rangle_E| \le \|e'\|_{E'} \quad \text{对 } e' \in E' ,
\]

并且得到 $ \|\kappa(e)\|_{(E')'} \le 1 $，由此即得结论。\nobreak\hspace{0.35em}■

\item

借助泛函分析中的工具可以证明，$ \kappa $ 是等距，从而是单射。我们称 $ \kappa $ 为 $ E $ 到 $ E $ 的\textbf{双对偶空间} $ E'' := (E')' $ 中的\textbf{典范单射}。若 $ \kappa $ 同时也是满射，从而是等距同构，则称 $ E $ 是\textbf{自反的}。此时，典范同构 $ \kappa $ 允许我们把 $ E $ 与其双对偶 $ E'' $ 等同起来。

\item

对 $ p \in (1,\infty) $，$ L_p(X) $ 自反。

\proofhead
这由 (a) 推出。\nobreak\hspace{0.35em}■
\item

空间 $ L_1(X) $ 与 $ L_\infty(X) $ 一般不是自反的；例如参见 [Ada75, Theorem 2.35]。\nobreak\hspace{0.35em}■

\end{items}

\end{remark}

\exerciseshead

\begin{exercises}
\item

设 $ S(X,\mu,E) := \{ [f] \in L_0(X,\mu,E) ; [f] \cap \mathcal{S}(X,\mu,E) \ne \varnothing \} $。证明 $ S(X,\mu,E) $ 是 $ L_p(X,\mu,E) $ 的稠密向量子空间（$ 1 \le p < \infty $）。

\item

对 $ a \in \mathbb{R}^n $，我们定义 $ \tau_a: E^{(\mathbb{R}^n)} \to E^{(\mathbb{R}^n)} $，即沿 $ a $ 的\textbf{右平移}，如下：

\[
(\tau_a\varphi)(x) := \varphi(x-a) \quad \text{对 } x \in \mathbb{R}^n, \ \varphi \in E^{(\mathbb{R}^n)} .
\]

令 $ \tau_a[f] := [\tau_a f] $，$ [f] \in L_p $。证明：

\begin{subitems}
\item

$ (\mathbb{R}^n,+) \to (\mathcal{L}\mathrm{aut}(L_p(\mathbb{R}^n,\lambda_n,E)),\circ) $，$ a \mapsto \tau_a $ 是群同态，且对每个 $ p \in [1,\infty] $ 都有 $ \|\tau_a\|_{\mathcal{L}(L_p)} = 1 $。

\item

对 $ p \in [1,\infty) $ 与 $ f \in L_p(\mathbb{R}^n,\lambda_n,E) $，有 $ \lim_{a\to0}\|\tau_a f - f\|_p = 0 $。

\item

若 $ \lim_{a\to0}\|\tau_a f - f\|_\infty = 0 $，则存在 $ g \in BUC(\mathbb{R}^n,E) $ 使得 $ f = g $ $ \mu $-几乎处处。

\end{subitems}

\item

设 $ \mu $ 是 $ \sigma $ 紧空间 $ X $ 上的完备拉东测度，$ (X_j)_{j\in\mathbb{N}} $ 是 $ X $ 的一列覆盖 $ X $ 的相对紧开子集。对 $ p \in [1,\infty] $，置

\[
q_{j,p}(f) := \|\chi_{X_j} f\|_p \quad \text{对 } j \in \mathbb{N}, \ f \in L_0(X,\mu,E) ,
\]
\[
L_{p,\mathrm{loc}}(X,\mu,E) := \{ f \in L_0(X,\mu,E) ; q_{j,p}(f) < \infty, \ j \in \mathbb{N} \} .
\]

最后，定义

\[
d_p(f,g) := \sum_{j=0}^{\infty} \frac{2^{-j} q_{j,p}(f-g)}{1 + q_{j,p}(f-g)} \quad \text{对 } f,g \in L_{p,\mathrm{loc}}(X,\mu,E) .
\]

\begin{subitems}
\item

$ L_{p,\mathrm{loc}}(X,\mu,E) $ 是良定义的，即与具体的序列 $ (X_j) $ 无关。

\item

$ (L_{p,\mathrm{loc}}(X,\mu,E),d_p) $ 是完备度量空间。

\item

$ L_p(X,\mu,E) \xrightarrow{d} L_{p,\mathrm{loc}}(X,\mu,E) \xrightarrow{d} L_{1,\mathrm{loc}}(X,\mu,E) $。

\item

由 $ d_p $ 生成的拓扑与序列 $ (X_j) $ 无关。

\end{subitems}

\item

设 $ p,q \in [1,\infty] $，定义

\[
L_p \cap L_q := (L_p \cap L_q)(X,\mu,E) := L_p(X,\mu,E) \cap L_q(X,\mu,E) ,
\]
\[
L_p + L_q := (L_p + L_q)(X,\mu,E) := L_p(X,\mu,E) + L_q(X,\mu,E) .
\]

又对 $ f \in L_p \cap L_q $ 置 $ \|f\|_{L_p\cap L_q} := \|f\|_p + \|f\|_q $，并令

\[
\|f\|_{L_p+L_q} := \inf\{ \|g\|_p + \|h\|_q ; g \in L_p(X,\mu,E), \ h \in L_q(X,\mu,E), \ f = g+h \} \quad \text{对 } f \in L_p + L_q .
\]

\begin{subitems}
\item

验证\textbf{插值不等式}

\[
\|f\|_r \le \|f\|_p^{1-\theta}\|f\|_q^{\theta} , \qquad \text{其中 } \frac{1}{r} := \frac{1-\theta}{p} + \frac{\theta}{q} ,
\]

对 $ f \in L_p \cap L_q $ 与 $ \theta \in [0,1] $ 成立。

\item

$ (L_p\cap L_q, \|\cdot\|_{L_p\cap L_q}) $ 与 $ (L_p+L_q, \|\cdot\|_{L_p+L_q}) $ 是巴拿赫空间，且

\[
(L_p\cap L_q)(X,\mu,E) \hookrightarrow L_r(X,\mu,E) \hookrightarrow (L_p+L_q)(X,\mu,E) \hookrightarrow L_{1,\mathrm{loc}}(X,\mu,E)
\]

其中 $ 1 \le p \le r \le q \le \infty $。
（提示：(i) 赫尔德不等式。(ii) 取 $ f \in L_p+L_q $ 使 $ \|f\|_{L_p+L_q} = 0 $。为证明它为零，注意对 $ r \in [1,\infty] $ 有 $ L_r \hookrightarrow L_{1,\mathrm{loc}} $（见习题 3）。为证明 $ L_p+L_q $ 的完备性，应用引理 4.5。嵌入 $ L_p\cap L_q \hookrightarrow L_r $ 由 (a) 推出。）

\end{subitems}

\item

设 $ p \in [1,\infty) $ 且 $ f \in (L_p \cap L_\infty)(X,\mu,E) $。证明 $ \lim_{q\to\infty}\|f\|_q = \|f\|_\infty $。

\item

证明映射

\[
L_\infty(X,\mu,\mathbb{K}) \times L_p(X,\mu,E) \to L_p(X,\mu,E) , \quad ([\varphi],[f]) \mapsto [\varphi f]
\]

是双线性且连续的，其范数不超过 1。

\item

设 $ \mu(X) < \infty $，并对 $ f,g \in L_0(X,\mu,E) $ 置

\[
d_0(f,g) := \int_X \frac{|f-g|}{1+|f-g|} \, d\mu .
\]

\begin{subitems}
\item

$ (L_0(X,\mu,E),d_0) $ 是度量空间。

\item

$ (f_j) $ 在 $ (L_0(X,\mu,E),d_0) $ 中收敛到 0，当且仅当它依测度收敛到 0。

\end{subitems}

\item

设 $ \mu $ 是 $ \sigma $ 紧空间 $ X $ 上的拉东测度，且 $ E $ 可分。证明：

\begin{subitems}
\item

$ C_c(X,\mathbb{K}) $ 可分。

\item

$ C_c(X,E) $ 可分。

\item

对 $ p \in [1,\infty) $，$ L_p(X,\mu,E) $ 可分。

\item

$ L_\infty(X,\mu,E) $ 一般不可分。

\item

$ \mathcal{S}(X,\mu,E) $ 在 $ \mathcal{L}_\infty(X,\mu,E) $ 中一般不是稠密的。
（提示：(i) 推论 V.4.8 与注 1.16(e)。(ii) 取 $ A \subset C_c(X,\mathbb{K}) $，令 $ B $ 为 $ E $ 中的可数稠密集。对 $ a \in A $ 与 $ b \in B $，置 $ (a\otimes b)(x) := a(x)b $，$ x \in X $，并考虑

\[
\left\{ \sum_{j=0}^{m} a_j \otimes b_j ; m \in \mathbb{N}, \ (a_j,b_j) \in A \times B, \ j = 0,\dots,m \right\} .
\]
\end{subitems}

(iii) 定理 4.14。(iv) 找出 $ L_\infty $ 的不可数子集 $ A $，使得对 $ A $ 中所有互不相同的 $ f,g $ 都有 $ \|f-g\|_\infty \ge 1 $。）

\item

若 $ \mu $ 有限且 $ E $ 有限维，证明 $ \mathcal{S}(X,\mu,E) $ 在 $ \mathcal{L}_\infty(X,\mu,E) $ 中稠密。

\item

证明注 4.9(c) 的陈述。

\item

证明：

\begin{subitems}
\item

$ \ell_p = \mathcal{L}_p(\mathbb{N},\mathcal{H}^0,\mathbb{K}) $，$ 1 \le p \le \infty $。

\item

若 $ 1 \le p \le q \le \infty $，则 $ \ell_p \hookrightarrow \ell_q $，且 $ \|\cdot\|_q \le \|\cdot\|_p $。

\item

若 $ 1 \le p \le q < \infty $，则 $ \ell_p \xrightarrow{d} \ell_q \xrightarrow{d} c_0 \hookrightarrow \ell_\infty $（见第 II.2 节）。

\end{subitems}

\item

设 $ p,q \in [1,\infty] $ 且 $ 1 \le p \le q \le \infty $。证明：

\begin{subitems}
\item

$ L_\infty(X,\mu,E) \subset L_1(X,\mu,E) \Rightarrow L_q(X,\mu,E) \hookrightarrow L_p(X,\mu,E) $。

\item

$ L_1(X,\mu,E) \subset L_\infty(X,\mu,E) \Rightarrow L_p(X,\mu,E) \hookrightarrow L_q(X,\mu,E) $。

\item

存在完备的 $ \sigma $ 有限测度空间 $ (X,\mathcal{A},\mu) $ [或 $ (Y,\mathcal{B},\nu) $]，它实现嵌入 $ L_\infty(X,\mu,\mathbb{R}) \hookrightarrow L_1(X,\mu,\mathbb{R}) $ [或 $ L_1(Y,\nu,\mathbb{R}) \hookrightarrow L_\infty(Y,\nu,\mathbb{R}) $]。
（提示：(i) 赫尔德不等式。(ii) 证明 $ L_p \hookrightarrow L_\infty $，并应用习题 4(i)。）

\end{subitems}

\item

对 $ p \in (0,1) $，证明：

\begin{subitems}
\item

$ \|f+g\|_p^p \le \|f\|_p^p + \|g\|_p^p $，$ f,g \in \mathcal{L}_0(X,\mu,E) $。

\item

$ \|f+g\|_p \le 2^{1/p-1}(\|f\|_p + \|g\|_p) $，$ f,g \in \mathcal{L}_0(X,\mu,E) $。

\item

$ \mathcal{L}_p(X,\mu,E) $ 是 $ \mathcal{L}_0(X,\mu,E) $ 的向量子空间。

\item

$ \mathcal{N} := \{ f \in \mathcal{L}_0(X,\mu,E) ; f = 0 \ \mu\text{-几乎处处} \} $ 是 $ \mathcal{L}_p(X,\mu,E) $ 的向量子空间，且

\[
\mathcal{N} = \{ f \in \mathcal{L}_p(X,\mu,E) ; \|f\|_p = 0 \} .
\]
\item

置 $ \rho(f,g) := \|f-g\|_p^p $，它在

\[
L_p(X,\mu,E) := \mathcal{L}_p(X,\mu,E)/\mathcal{N}
\]

上诱导出一个度量。

\item

$ (L_p(X,\mu,E),\rho) $ 是完备的。

\item

对 $ f,g \in \mathcal{L}_p(X,\mu,\mathbb{R}) $，若 $ f \ge 0 $ 且 $ g \ge 0 $，则 $ \|f+g\|_p \ge \|f\|_p + \|g\|_p $。

\item

映射

\[
L_p(X,\mu,\mathbb{R}) \to \mathbb{R}^+ , \quad [f] \mapsto \|f\|_p
\]

不是范数。
（提示：(i) 对 $ a > 0 $，映射 $ [t \mapsto a^p + t^p - (a+t)^p] $ 在 $ \mathbb{R}^+ $ 上递增。(ii) 对 $ a > 0 $，考察 $ [t \mapsto (a^{1/p} + t^{1/p})/(a+t)^{1/p}] $。(vi) 仿照引理 4.5 与定理 4.6 的证明。(vii) 定理 4.2。）

\end{subitems}

\item

设 $ p_j \in [1,\infty] $，$ j = 1,\dots,m $；令 $ 1/r := \sum_{j=1}^{m} 1/p_j $。对 $ f_j \in L_{p_j}(X,\mu,\mathbb{K}) $，证明 $ \prod_{j=1}^{m} f_j $ 属于 $ L_r(X,\mu,\mathbb{K}) $，且

\[
\left\| \prod_{j=1}^{m} f_j \right\|_r \le \prod_{j=1}^{m} \|f_j\|_{p_j} .
\]

（提示：赫尔德不等式。）

\item

设 $ X $ 是度量空间。函数 $ f \in E^X $ \textbf{在无穷远处消失}，是指对每个 $ \varepsilon > 0 $ 都存在 $ X $ 的紧子集 $ K $，使得对一切 $ x \in K^c $ 都有 $ |f(x)| < \varepsilon $。验证

\[
C_0(X,E) := \{ f \in C(X,E) ; f \text{ 在无穷远处消失} \}
\]

是 $ C_c(X,E) $ 在 $ BUC(X,E) $ 中的闭包。

\item

对 $ f \in \mathcal{L}_0(X,\mu,E) $，置

\[
\lambda_f(t) := \mu([|f| > t]) \quad \text{与} \quad f^*(t) := \inf\{ s \ge 0 ; \lambda_f(s) \le t \} \quad \text{对 } t \in [0,\infty) .
\]

我们称 $ f^*: [0,\infty) \to [0,\infty] $ 为 $ f $ 的\textbf{递减重排}。证明：

\begin{subitems}
\item

$ \lambda_f $ 与 $ f^* $ 递减、右连续且勒贝格可测。

\item

若对 $ g \in \mathcal{L}_0(X,\mu,E) $ 有 $ |f| \le |g| $，则 $ \lambda_f \le \lambda_g $ 且 $ f^* \le g^* $。

\item

若 $ (f_j) $ 是递增序列且 $ |f_j| \uparrow |f| $，则 $ \lambda_{f_j} \uparrow \lambda_f $ 且 $ f_j^* \uparrow f^* $。

\item

对 $ p \in (0,\infty) $，

\[
\int_X |f|^p \, d\mu = p \int_{\mathbb{R}^+} t^{p-1} \lambda_f(t) \, \lambda_1(dt) = \int_{\mathbb{R}^+} (f^*)^p \, d\lambda_1 .
\]
\item

$ \|f\|_\infty = f^*(0) $。

\item

$ \lambda_f = \lambda_{f^*} $。
（关于 (iv) 的提示：先考虑简单函数，再应用 (iii) 以及定理 1.12 与 3.4。）

\end{subitems}

\item

对 $ j \in \mathbb{N} $，令 $ I_{j,k} := [k2^{-j},(k+1)2^{-j}] $，$ k = 0,\dots,2^j-1 $。又设 $ \{ J_n ; n \in \mathbb{N} \} $ 是 $ \{ I_{j,k} ; j \in \mathbb{N}, k = 0,\dots,2^j-1 \} $ 的一个重编号，并对 $ j \in \mathbb{N} $ 置 $ f_n := \chi_{J_n} $。证明 $ (f_n) $ 对每个 $ p \in [1,\infty) $ 都是 $ \mathcal{L}_p([0,1]) $ 中的零序列，尽管 $ (f_n(x)) $ 对每个 $ x \in [0,1] $ 都发散。

\item

设 $ (f_k) $ 是 $ L_p(X) $ 中的序列，其中 $ 1 \le p < \infty $。若

\[
\int_X f_k\varphi \, dx \to \int_X f\varphi \, dx \quad \text{对 } \varphi \in L_{p'}(X) ,
\]

则称 $ (f_k) $ 在 $ L_p(X) $ 中\textbf{弱收敛}于 $ f \in L_p(X) $。此时称 $ f $ 为 $ (f_k) $ 在 $ L_p(X) $ 中的\textbf{弱极限}。

证明：

\begin{subitems}
\item

$ L_p(X) $ 中的弱极限是唯一的。

\item

$ L_p(X) $ 中每个收敛序列都在 $ L_p(X) $ 中弱收敛。

\item

若 $ (f_k) $ 在 $ L_p(X) $ 中弱收敛于 $ f $，且 $ \mu $-几乎处处收敛于 $ g \in L_p(X) $，则 $ f = g $。

\item

若 $ (f_k) $ 在 $ L_2(X) $ 中弱收敛于 $ f $，且 $ \|f_k\|_2 \to \|f\|_2 $，则 $ (f_k) $ 在 $ L_2(X) $ 中收敛于 $ f $。

\item

令 $ \mathrm{e}_k(t) := (2\pi)^{-1/2} e^{ikt} $，$ 0 < t < 2\pi $，$ k \in \mathbb{N} $。则序列 $ (\mathrm{e}_k) $ 在 $ L_2((0,2\pi)) $ 中弱收敛于 0，尽管它在 $ L_2((0,2\pi)) $ 中发散。
（提示：(i) 对 $ f \in L_p(X) $ 考察 $ \varphi(x) := \overline{f(x)}\,|f(x)|^{p/p'-1} $。(ii) 赫尔德不等式。(iii) 证明 $ g \in L_p(X) $，从而 $ [|g| = \infty] $ 是 $ \mu $-零集。若 $ X_n := [\sup_{k\ge n}|f_k(x)| \ge n] $，则 $ \bigcap X_n $ 也是 $ \mu $-零集。现在对 $ \varphi \in L_{p'}(X) $ 考察 $ \lim\int_{X_n^c} f_n\varphi \, dx $。(iv) 应用 $ L_2(X) $ 中的平行四边形恒等式。(v) 第一个结论由贝塞尔不等式推出，第二个由 (ii) 推出。）

\end{subitems}

\end{exercises}

\section{n 维 Bochner–Lebesgue 积分}

在这一简短的节中，我们讨论 Bochner–Lebesgue 积分与第 VI 章定义的柯西–黎曼积分之间的关系。我们将证明每个跳跃连续函数都是勒贝格可测的，并且相应的积分相等。这一联系将使我们能够把为柯西–黎曼积分发展出来的方法引入勒贝格积分理论。

我们还证明：紧区间上的有界标量值函数黎曼可积，当且仅当其不连续点集具有测度零。由此可知存在勒贝格可积但不是黎曼可积的函数。因此勒贝格积分是黎曼积分的真扩充——从而也是柯西–黎曼积分的真扩充。

在本节全文中，设

- $ X \subset \mathbb{R}^{n} $ 是正测度的 $ \lambda_{n} $-可测集；$ E = (E,|\cdot|) $ 是 Banach 空间。

\needspace{3\baselineskip}\subsection*{勒贝格测度空间}

由习题 IX.1.7 可知，$ \mathcal{L}_{X} := \mathcal{L}(n)|X $ 是 $ X $ 上的 $ \sigma $-代数。因此限制 $ \lambda_{n}|X := \lambda_{n}|\mathcal{L}_{X} $ 是 $ X $ 上的测度，称为 $ X $ 上的 \textbf{n 维勒贝格测度}（或 \textbf{勒贝格 n-测度}）。我们也用 $ \lambda_{n} $ 来记这个限制。容易验证 $ (X,\mathcal{L}_{X},\lambda_{n}) $ 是完备的 $ \sigma $-有限测度空间。如果不存在误解的危险，我们就从“可测的”“测度”“可积的”等词中去掉限定词“勒贝格”（或“$ \lambda_{n} $”）。

若 $ f \in E^{X} $ 可积，则称

\[
\int_{X} f d\lambda_{n} := \int_{X} f d(\lambda_{n}|X) = \int_{\mathbb{R}^{n}} \widetilde{f}\chi_{X} d\lambda_{n}
\]

为 $ f $ 在 $ X $ 上的（n 维）（Bochner–Lebesgue）积分。记号

\[
\int_{X} f(x) d\lambda_{n}(x) \quad \text{和} \quad \int_{X} f(x) \lambda_{n}(dx)
\]

也很常见。

简言之，我们令

\[
\mathcal{L}_{p}(X,E) := \mathcal{L}_{p}(X,\lambda_{n},E) \quad \text{和} \quad L_{p}(X,E) := L_{p}(X,\lambda_{n},E) .
\]

我们还令 $ \mathcal{L}_{p}(X) := \mathcal{L}_{p}(X,\mathbb{K}) $ 与 $ L_{p}(X) := L_{p}(X,\mathbb{K}) $，其中 $ p \in [1,\infty] \cup \{0\} $。

下一个定理列出 n 维积分的重要性质。

\begin{theorem}

设 $ X $ 在 $ \mathbb{R}^{n} $ 中是开的，或者当 $ n = 1 $ 时 $ X $ 是完全区间。则：

\begin{subitems}
\item

$ \lambda_{n} $ 是 $ X $ 上的 massive 拉东测度。

\item

$ C(X,E) $ 是 $ L_{0}(X,E) $ 的向量子空间。

\item

$ BC(X,E) $ 是 $ L_{\infty}(X,E) $ 的闭向量子空间。

\item

对 $ p \in [1,\infty) $，$ C_{c}(X,E) $ 是 $ L_{p}(X,E) $ 的稠密向量子空间。若 $ K $ 是 $ X $ 的紧子集，则

\[
\|f\|_{p} \le \lambda_{n}(K)^{1/p}\|f\|_{\infty} \quad \text{对满足 } \operatorname{supp}(f) \subset K \text{ 的 } f \in C_{c}(X,E) .
\]
\item

若 $ X $ 具有有限测度且 $ 1 \le p < q \le \infty $，则

\[
L_{q}(X,E) \overset{d}{\hookrightarrow} L_{p}(X,E)
\]

以及

\[
\|f\|_{p} \le \lambda_{n}(X)^{1/p-1/q}\|f\|_{q} \quad \text{对 } f \in L_{q}(X,E) .
\]

\proofhead
(i) $ X $ 是 $ \sigma $-紧的度量空间——若 $ X $ 是开集，由注 1.16(e)；若 $ X $ 是区间，理由显然。于是论断由注 1.16(h) 与习题 IX.5.21 推出。
\end{subitems}

(ii) 与 (iii) 分别由命题 4.17 与定理 4.18(ii) 覆盖。

\begin{subitems}
\item[（iv）]

第一个论断是定理 4.18(i) 的推论，第二个论断显然。

\item[（v）]

是定理 4.16 的特例。\nobreak\hspace{0.35em}■

\end{subitems}

\end{theorem}

\begin{remark}

设 $ X $ 可测且其边界 $ \partial X $ 是 $ \lambda_{n} $-零集。则博雷尔集 $ \mathring{X} $ 属于 $ \mathcal{L}(n) $，且 $ \lambda_{n}(\mathring{X}) = \lambda_{n}(X) $。进一步，容易验证映射

\[
L_{p}(X,E) \to L_{p}(\mathring{X},E) , \quad [f] \mapsto [f|\mathring{X}]
\]

对 $ p \in [1,\infty] \cup \{0\} $ 是向量空间同构。若 $ p \in [1,\infty] $，它是等距。因此对 $ p \in [1,\infty] \cup \{0\} $ 我们可以把 $ L_{p}(X,E) $ 与 $ L_{p}(\mathring{X},E) $ 等同。特别地，对 $ \mathbb{R} $ 中以 $ a := \inf X $ 与 $ b := \sup X $ 为端点的区间 $ X $，有

\[
L_{p}(X,E) = L_{p}([a,b],E) = L_{p}([a,b),E) = L_{p}((a,b],E) = L_{p}((a,b),E)
\]

其中 $ p \in [1,\infty] \cup \{0\} $。

\end{remark}

\needspace{3\baselineskip}\subsection*{绝对可积函数的勒贝格积分}

我们现在证明每个绝对可积函数都是勒贝格可积的，且它在第 VI.8 节意义下的积分等于勒贝格积分。

\qyxffntextdef{1}{此处及类似情形中，我们把 $ \chi_{[c,b)}f $ 视为 $ \mathbb{R} $ 上的函数。写成 $ \chi_{[c,b)}\widetilde{f} $ 会更精确，但很繁琐。}
\begin{theorem}

设 $ f: (a,b) \to E $ 绝对可积，其中 $ a,b \in \overline{\mathbb{R}} $ 且 $ a < b $。则 $ f $ 属于 $ \mathcal{L}_{1}((a,b),E) $，且

\[
\int_{(a,b)} f d\lambda_{1} = \int_{a}^{b} f .
\]

证明

\begin{subitems}
\item

设 $ a < \alpha < \beta < b $。若 $ g: [\alpha,\beta] \to E $ 是阶梯函数，则 $ g $ 显然是 $ \lambda_{1} $-简单的，且

\begin{equation}
\int_{(\alpha,\beta)} g d\lambda_{1} = \int_{\alpha}^{\beta} g .\tag{5.1}
\end{equation}

现在设 $ g: [\alpha,\beta] \to E $ 跳跃连续。则存在阶梯函数序列 $ (g_{j}) $ 一致收敛于 $ g $。因此 $ g $ 可测，且注 VI.1.1(d) 与推论 3.15(ii) 表明 $ g $ 属于 $ \mathcal{L}_{1}((\alpha,\beta),E) $。由于 $ g $ 有界且序列 $ (g_{j}) $ 一致收敛，存在 $ M \ge 0 $ 使得对所有 $ j \in \mathbb{N} $ 有 $ |g_{j}| \le M $。因此由勒贝格控制收敛定理，

\[
\lim_{j \to \infty} \int_{(\alpha,\beta)} g_{j} d\lambda_{1} = \int_{(\alpha,\beta)} g d\lambda_{1}
\]

在 $ E $ 中成立，并利用 (5.1) 与柯西–黎曼积分的定义得出

\[
\int_{\alpha}^{\beta} g = \lim_{j \to \infty} \int_{\alpha}^{\beta} g_{j} = \lim_{j \to \infty} \int_{(\alpha,\beta)} g_{j} d\lambda_{1} = \int_{(\alpha,\beta)} g d\lambda_{1} .
\]
\item

我们固定 $ c \in (a,b) $ 并选取 $ (c,b) $ 中的序列 $ (\beta_{j}) $ 使得 $ \beta_{j} \uparrow b $。我们还设\qyxffnmark{1}

\[
g := \chi_{[c,b)}f , \quad g_{j} := \chi_{[c,\beta_{j}]}f \quad \text{对 } j \in \mathbb{N} .
\]

由 (i)，$ (g_{j}) $ 是 $ \mathcal{L}_{1}(\mathbb{R},E) $ 中的序列。显然 $ (g_{j}) $ 逐点收敛于 $ g $，且 $ (|g_{j}|) $ 是递增序列收敛于 $ |g| $。因此 $ g $ 可测。由 (i) 可得

\[
\int_{\mathbb{R}} |g_{j}| d\lambda_{1} = \int_{(c,\beta_{j})} |f| d\lambda_{1} = \int_{c}^{\beta_{j}} |f| ,
\]

而 $ \int_{c}^{b} f $ 的绝对收敛性蕴含

\begin{equation}
\lim_{j \to \infty} \int_{\mathbb{R}} |g_{j}| d\lambda_{1} = \lim_{j \to \infty} \int_{c}^{\beta_{j}} |f| = \int_{c}^{b} |f| .\tag{5.2}
\end{equation}

另一方面，单调收敛定理表明

\[
\int_{\mathbb{R}} |g| d\lambda_{1} = \lim_{j \to \infty} \int_{\mathbb{R}} |g_{j}| d\lambda_{1} ,
\]

且由 (5.2) 看出 $ g $ 属于 $ \mathcal{L}_{1}(\mathbb{R},E) $。因此我们可以对序列 $ (g_{j}) $ 应用控制收敛定理，得到

\[
\lim_{j \to \infty} \int_{\mathbb{R}} g_{j} d\lambda_{1} = \int_{\mathbb{R}} g d\lambda_{1} = \int_{[c,b)} f d\lambda_{1}
\]

在 $ E $ 中成立。进一步地，由 (i) 得

\[
\int_{\mathbb{R}} g_{j} d\lambda_{1} = \int_{[c,\beta_{j})} f d\lambda_{1} = \int_{c}^{\beta_{j}} f ,
\]

从而由命题 VI.8.7，

\[
\lim_{j \to \infty} \int_{\mathbb{R}} g_{j} d\lambda_{1} = \lim_{j \to \infty} \int_{c}^{\beta_{j}} f = \int_{c}^{b} f
\]

在 $ E $ 中成立。于是极限 $ \int_{[c,b)} f d\lambda_{1} $ 与 $ \int_{c}^{b} f $ 相等。以类似方式，我们证明 $ \chi_{(a,c]}f $ 属于 $ \mathcal{L}_{1}(\mathbb{R},E) $ 且 $ \int_{(a,c]} f d\lambda_{1} = \int_{a}^{c} f $。这表明 $ f $ 勒贝格可积，且 $ \int_{(a,b)} f d\lambda_{1} = \int_{a}^{b} f $。\nobreak\hspace{0.35em}■

\end{subitems}

\end{theorem}

\begin{corollary}

对 $ -\infty < a < b < \infty $，有 $ \mathcal{S}([a,b],E) \hookrightarrow \mathcal{L}_{1}([a,b],E) $ 且

\[
\int_{[a,b]} f d\lambda_{1} = \int_{a}^{b} f \quad \text{对 } f \in \mathcal{S}([a,b],E) .
\]

\proofhead
这由定理 5.3 与命题 VI.8.3 推出。\nobreak\hspace{0.35em}■
\end{corollary}

\begin{remark}

固定 $ a,b \in \mathbb{R} $ 且 $ a < b $。

\begin{items}
\item

设 $ f: (a,b) \to E $ 可容许且 $ \int_{a}^{b} f $ 作为反常积分存在。则 $ f $ 不必属于 $ \mathcal{L}_{1}((a,b),E) $。

\proofhead
我们通过下式定义 $ f: \mathbb{R} \to \mathbb{R} $：
\begin{equation*}
f(x) := \begin{cases} 0 & \text{若 } x \in (-\infty,0), \\ (-1)^{j}/j & \text{若 } x \in [j-1,j), \text{ 其中 } j \in \mathbb{N}^{\times} . \end{cases}
\end{equation*}

显然 $ f $ 可容许，且 $ \int_{-\infty}^{\infty} f $ 在 $ \mathbb{R} $ 中存在，因为

\[
\int_{-\infty}^{\infty} f = \sum_{j=1}^{\infty} (-1)^{j}/j .
\]

若 $ f $ 属于 $ \mathcal{L}_{1}(\mathbb{R}) $，则应有 $ \int_{\mathbb{R}} |f| d\lambda_{1} < \infty $，这与单调收敛定理矛盾，该定理给出

\[
\int_{\mathbb{R}} |f| d\lambda_{1} = \lim_{k \to \infty} \int_{\mathbb{R}} \chi_{[0,k]}|f| d\lambda_{1} = \lim_{k \to \infty} \sum_{j=1}^{k} 1/j = \infty . \quad \text{■}
\]
\item

设 $ f: (a,b) \to E $ 可容许且 $ f $ 属于 $ \mathcal{L}_{1}((a,b),E) $。则 $ f $ 绝对可积，且

\[
\int_{(a,b)} f d\lambda_{1} = \int_{a}^{b} f \quad \text{在 } E \text{ 中} .
\]

\proofhead
取 $ c \in (a,b) $，并设 $ (a,c) $ 中的序列 $ (\alpha_{j}) $ 满足 $ \alpha_{j} \to a $。又设 $ f_{j} := \chi_{[\alpha_{j},c]}f $。则 $ (f_{j}) $ 逐点收敛于 $ \chi_{(a,c]}f $，且对 $ j \in \mathbb{N} $ 有 $ |f_{j}| \le |f| $。由于 $ f $ 可容许，命题 VI.4.3 表明 $ |f|[\alpha_{j},c] $ 属于 $ \mathcal{S}([\alpha_{j},c],\mathbb{R}) $。因此由推论 5.4 与控制收敛定理，
\[
\int_{\alpha_{j}}^{c} |f| = \int_{[\alpha_{j},c]} |f| d\lambda_{1} \to \int_{(a,c]} |f| d\lambda_{1} .
\]

因此 $ \int_{a}^{c} |f| $ 存在。类似地，我们证明 $ \int_{c}^{b} |f| $ 存在，从而 $ \int_{a}^{b} f $ 绝对收敛。第二个论断现在由定理 5.3 推出。\nobreak\hspace{0.35em}■

\end{items}

设 $ f \in \mathcal{L}_{1}((a,b),E) $。注 5.5(b) 表明，此时若把 $ \int_{(a,b)} f d\lambda_{1} $ 记为 $ \int_{a}^{b} f $ 或 $ \int_{a}^{b} f(x) dx $，不会产生误解。今后在 n 维情形我们通常写

\[
\int_{X} f dx := \int_{X} f d\lambda_{n} .
\]

定理 5.3 及其推论使我们能够把第二卷中发展出来的积分方法移植到勒贝格理论的框架中。结合定理 3.14 的可积性判别准则与控制收敛定理，这些方法为证明积分的存在性提供了十分有效的工具。在本章余下的各节中，当我们建立“多维”积分的具体计算程序时，这一点将变得清楚。

\end{remark}

\needspace{3\baselineskip}\subsection*{黎曼可积函数的一个刻画}

定理 5.3 表明勒贝格积分是柯西–黎曼积分的扩充。我们现在刻画黎曼可积函数，并证明这一扩充是真扩充。

\begin{theorem}

设 $ I $ 是紧区间，$ f: I \to \mathbb{K} $ 有界。则 $ f $ 黎曼可积当且仅当它 $ \lambda_{1} $-几乎处处连续。此时 $ f $ 勒贝格可积，且黎曼积分与勒贝格积分相等。

证明

\begin{subitems}
\item

不失一般性，可以假设 $ \mathbb{K} = \mathbb{R} $ 且 $ I := [0,1] $。对 $ k \in \mathbb{N} $，设 $ \mathfrak{Z}_{k} := (\xi_{0,k},\ldots,\xi_{2^{k},k}) $ 是 $ [0,1] $ 的划分，其中 $ \xi_{j,k} := j2^{-k} $，$ j = 0,\ldots,2^{k} $。又设

\[
I_{0,k} := [\xi_{0,k},\xi_{1,k}] , \quad I_{j,k} := (\xi_{j,k},\xi_{j+1,k}] \quad \text{对 } j = 1,\ldots,2^{k}-1 .
\]

最后，令 $ \alpha_{j,k} := \inf_{x \in \overline{I}_{j,k}} f(x) $，$ \beta_{j,k} := \sup_{x \in \overline{I}_{j,k}} f(x) $，以及

\[
g_{k} := \sum_{j=0}^{2^{k}-1} \alpha_{j,k}\chi_{I_{j,k}} , \quad h_{k} := \sum_{j=0}^{2^{k}-1} \beta_{j,k}\chi_{I_{j,k}} \quad \text{对 } k \in \mathbb{N} .
\]

则 $ (g_{k}) $ 是 $ \lambda_{1} $-简单函数的递增序列，$ (h_{k}) $ 是 $ \lambda_{1} $-简单函数的递减序列。因此它们的逐点极限 $ g := \lim_{k} g_{k} $ 与 $ h := \lim_{k} h_{k} $ 有定义且 $ \lambda_{1} $-可测，并且 $ g \le f \le h $。此外有

\[
\int_{[0,1]} g_{k} d\lambda_{1} = \underline{S}(f,k) \quad \text{和} \quad \int_{[0,1]} h_{k} d\lambda_{1} = \overline{S}(f,k) ,
\]

其中 $ \underline{S}(f,k) $ 与 $ \overline{S}(f,k) $ 表示 $ f $ 在 $ [0,1] $ 上关于划分 $ \mathfrak{Z}_{k} $ 的下和与上和（见习题 VI.3.7）。用 $ \underline{\int} f $ 与 $ \overline{\int} f $ 记 $ f $ 在 $ [0,1] $ 上的下、上黎曼积分，则由单调收敛定理得

\begin{equation}
\int_{[0,1]} (h-g) d\lambda_{1} = \overline{\int} f - \underline{\int} f .\tag{5.3}
\end{equation}
\item

设 $ R := \bigcup_{k \in \mathbb{N}} \{\xi_{0,k},\ldots,\xi_{2^{k},k}\} $ 是区间 $ I_{j,k} $ 的端点集。设 $ C $ 是 $ f $ 的连续点集。则

\begin{equation}
[g = h] \cap R^{c} \subset C \subset [g = h] .\tag{5.4}
\end{equation}

为看出这一点，取 $ \varepsilon > 0 $。先设 $ x_{0} \in R^{c} $ 且 $ g(x_{0}) = h(x_{0}) $。可以找到 $ k \in \mathbb{N} $ 使 $ h_{k}(x_{0}) - g_{k}(x_{0}) < \varepsilon $，以及 $ j \in \{0,\ldots,2^{k}-1\} $ 使 $ x_{0} $ 落在区间 $ (\xi_{j,k},\xi_{j+1,k}) $ 内。于是对 $ x \in I_{j,k} $ 有

\[
|f(x) - f(x_{0})| \le \sup_{y \in \overline{I}_{j,k}} f(y) - \inf_{y \in \overline{I}_{j,k}} f(y) = h_{k}(x_{0}) - g_{k}(x_{0}) < \varepsilon ,
\]

这证明了 $ f $ 在 $ x_{0} $ 处的连续性。
现设 $ x_{0} \in C $。取 $ \delta > 0 $ 使得对 $ x \in [x_{0}-\delta,x_{0}+\delta] \cap [0,1] $ 有 $ |f(x) - f(x_{0})| < \varepsilon/2 $。选取 $ k_{0} \in \mathbb{N} $ 使 $ 2^{-k_{0}} \le \delta $，并对每个 $ k \ge k_{0} $ 取 $ j \in \{0,\ldots,2^{k}-1\} $ 使 $ x_{0} \in I_{j,k} \subset [x_{0}-\delta,x_{0}+\delta] $。则

\[
0 \le h_{k}(x_{0}) - g_{k}(x_{0}) = \sup_{x \in \overline{I}_{j,k}} (f(x) - f(x_{0})) - \inf_{x \in \overline{I}_{j,k}} (f(x) - f(x_{0})) < \varepsilon .
\]

由此得 $ h(x_{0}) - g(x_{0}) = \lim_{k} (h_{k}(x_{0}) - g_{k}(x_{0})) = 0 $。这证明了 (5.4)。

\item

若 $ f $ 是黎曼可积函数，则 $ \underline{\int} f = \overline{\int} f = \int f $（习题 VI.3.10）。因此 (5.3) 表明

\begin{equation}
h = g = f \qquad \lambda_1\text{-a.e.} ,\tag{5.5}
\end{equation}

这蕴涵 $ f $ 的 $ \lambda_1 $-可测性。由于 $ f $ 有界，故 $ f \in \mathcal{L}_1([0,1]) $。我们还有 $ |g_k| \le \|f\|_\infty $ 对 $ k \in \mathbb{N} $ 是 $ \lambda_1 $-几乎处处的。于是勒贝格控制收敛定理给出

\[
\int_{[0,1]} g \, d\lambda_1 = \lim_{k} \int_{[0,1]} g_k \, d\lambda_1 = \lim_{k} \underline{S}(f,k) = \int_0^1 f ,
\]

其中在最后一个等式中我们又一次用到了习题 VI.3.10。由 (5.5) 与引理 2.15 可得 $ \int_{[0,1]} f \, d\lambda_1 = \int_0^1 f $。最后，(5.4)、(5.5) 以及 $ R $ 的可数性蕴涵：$ f $ 的不连续点构成一个勒贝格测度为零的集合。

\item

反之，设 $ C^c $ 的测度为零。由 (5.4)，$ [g \ne h] $ 的测度也为零，而 $ f $ 的黎曼可积性由 (5.3) 推出。这就完成了证明。\nobreak\hspace{0.35em}■

\end{subitems}

\end{theorem}

\begin{corollary}

有些勒贝格可积函数不是黎曼可积的。因此勒贝格积分是黎曼积分的真扩张。

\proofhead
考虑狄利克雷函数
\begin{equation*}
f: [0,1] \to \mathbb{R}, \quad f(x) := \begin{cases} 1 & \text{若 } x \in \mathbb{Q}, \\ 0 & \text{若 } x \notin \mathbb{Q}, \end{cases}
\end{equation*}

在 $ [0,1] $ 上。由引理 2.15，$ f $ 属于 $ \mathcal{L}_1([0,1]) $，因为 $ f $ 几乎处处为零。但我们由例 III.1.3(c) 已知 $ f $ 处处不连续，故由定理 5.6 它不是黎曼可积的。\nobreak\hspace{0.35em}■

\end{corollary}

与狄利克雷函数几乎处处相等的映射构成的等价类中含有黎曼可积函数——例如零函数。因此从 $ L_1 $-空间的观点看，这个例子是无趣的。然而，习题 13 将证明存在 $ f \in \mathcal{L}_1([0,1],\mathbb{R}) $，使得任何 $ g \in [f] $ 都不是黎曼可积的。这意味着黎曼积分对 $ L_p $-空间的理论而言是不充分的。

\exerciseshead

\begin{exercises}
\item

设 $ p, q \in [1,\infty] $ 且 $ p \ne q $，证明 $ L_p(\mathbb{R},E) \not\subset L_q(\mathbb{R},E) $。

\item

设 $ J $ 是开区间，且 $ f \in C^1(J,E) $ 具有紧支集。则 $ \int_J f' = 0 $。

\item

设 $ f \in \mathcal{L}_0([0,1],\mathbb{R}^+) $ 有界。证明

\[
\underline{\int} f \le \int_{[0,1]} f \, d\lambda_1 \le \overline{\int} f .
\]
\item

设 $ I $ 是紧区间，并如下定义 $ I $ 上的\textbf{有界变差函数空间}：

\[
BV(I,E) := \{ f: I \to E ;\ \operatorname{Var}(f,I) < \infty \} .
\]

\begin{items}
\item

作为向量子空间，我们有包含关系

\[
C^{1-}(I,E) \subset BV(I,E) \subset B(I,E) .
\]
\item

令 $ \alpha := \inf I $ 且 $ f \in \mathcal{L}_1(I,E) $。则 $ F: I \to E $，$ x \mapsto \int_\alpha^x f(t) \, dt $ 属于 $ BV(I,E) $，且 $ \operatorname{Var}(F,I) \le \|f\|_1 $。

\item

对每个 $ f \in BV(I,\mathbb{R}) $，存在递增映射 $ s^\pm: I \to \mathbb{R} $ 使得 $ f = s^+ - s^- $。

\item

$ BV(I,\mathbb{R}) $ 是跳跃连续函数 $ I \to \mathbb{R} $ 构成的空间 $ \mathcal{S}(I,\mathbb{R}) $ 的向量子空间。

\item

每个单调函数都属于 $ BV(I,\mathbb{R}) $。
（(c) 的提示：令 $ \alpha := \inf I $，考虑函数 $ s^+ := (x \mapsto \operatorname{Var}(f^+,[\alpha,x])) $ 与 $ s^- := s - f $。）

\end{items}

\item

设 $ H $ 是可分希尔伯特空间。证明\footnote{可以证明，若把 $ H $ 换成任意巴拿赫空间，习题 5 的结论仍然成立。} $ BV([a,b],H) $ 是 $ \mathcal{L}_\infty([a,b],H) $ 的向量子空间，并且

\[
\int_a^{b-h} \|f(t+h) - f(t)\| \, dt \le h \operatorname{Var}(f,[a,b]) \qquad \text{对 } 0 < h < b - a .
\]

（提示：注意习题 1.1 与 4(d)。对 $ 0 < h < b - a $ 及 $ t \in [a,b-h] $，证明 $ \|f(t+h) - f(t)\| \le \operatorname{Var}(f,[a,t+h]) - \operatorname{Var}(f,[a,t]) $。）

\item

设 $ J \subset \mathbb{R} $ 是完全区间。函数 $ f: J \to E $ 称为\textbf{绝对连续}的，若对每个 $ \varepsilon > 0 $ 都存在 $ \delta > 0 $，使得

\[
\sum_{k=0}^{m} |f(\beta_k) - f(\alpha_k)| < \varepsilon
\]

对 $ J $ 的每个由两两不相交子区间构成、且满足 $ \sum_{k=0}^{m} (\beta_k - \alpha_k) < \delta $ 的有限族 $ \{ (\alpha_k, \beta_k) ;\ k = 0, \ldots, m \} $ 都成立。我们用 $ W_1^1(J,E) $ 记 $ E^J $ 中全体绝对连续函数构成的集合。证明：

\begin{items}
\item

作为向量子空间，我们有包含关系

\[
BC^1(J,E) \subset W_1^1(J,E) \subset C(J,E) .
\]
\item

若 $ J $ 紧，则 $ W_1^1(J,E) \subset BV(J,E) $。

\item

康托尔函数（习题 III.3.8）是连续的，但不是绝对连续的。

\item

令 $ \alpha := \inf J $ 并取 $ f \in \mathcal{L}_1(J,E) $。则 $ F: J \to E $，$ x \mapsto \int_\alpha^x f(t) \, dt $ 是绝对连续的。

\end{items}

\item

对 $ j = 1, 2 $，定义 $ f_j: [0,1] \to \mathbb{R} $ 为

\begin{equation*}
f_j(x) := \begin{cases} x^2 \sin(1/x^j) & \text{若 } x \in (0,1], \\ 0 & \text{若 } x = 0 ; \end{cases}
\end{equation*}

与习题 IV.1.2 比较。证明：

\begin{items}
\item

$ f_1 \in BV([0,1],\mathbb{R}) $。

\item

$ f_2 \notin BV([0,1],\mathbb{R}) $。

\end{items}

\item

设 $ \mu $ 与 $ \nu $ 是可测空间 $ (X,\mathcal{A}) $ 上的测度。我们称 $ \nu $ 是\textbf{$ \mu $-绝对连续}的，若每个 $ \mu $-零集也是 $ \nu $-零集。此时我们记 $ \nu \ll \mu $。

\begin{items}
\item

设 $ (X,\mathcal{A},\mu) $ 是 $ \sigma $-有限的完备测度空间。对 $ f \in \mathcal{L}_0(X,\mu,\overline{\mathbb{R}}^+) $，定义

\[
f \cdot \mu: \mathcal{A} \to [0,\infty], \quad A \mapsto \int_A f \, d\mu .
\]

\proofhead
$ f \cdot \mu $ 是 $ (X,\mathcal{A}) $ 上的完备测度，且 $ f \cdot \mu \ll \mu $。
\item

设 $ \mathcal{A} := \mathcal{L}_{[0,1]} $，$ \nu := \lambda_1 $，$ \mu := \mathcal{H}^0 $。验证：

\end{items}

\begin{subitems}
\item

$ \nu \ll \mu $。

\item

不存在 $ f \in \mathcal{L}_0([0,1],\mathcal{A},\mu) $ 使得 $ \nu = f \cdot \mu $。

\end{subitems}

\item

设 $ (X,\mathcal{A},\nu) $ 是有限测度空间，$ \mu $ 是 $ (X,\mathcal{A}) $ 上的测度。下列陈述等价：

\begin{subitems}
\item

$ \nu \ll \mu $。

\item

对每个 $ \varepsilon > 0 $ 存在 $ \delta > 0 $，使得对一切满足 $ \mu(A) < \delta $ 的 $ A \in \mathcal{A} $ 都有 $ \nu(A) < \varepsilon $。

\end{subitems}

\item

对 $ f \in \mathcal{L}_0(\mathbb{R},\lambda_1,\overline{\mathbb{R}}^+) $，令 $ F(x) := \int_{-\infty}^{x} f(t) \, dt $（$ x \in \mathbb{R} $），并用 $ \mu_F $ 记由 $ F $ 生成的 $ \mathbb{R} $ 上的勒贝格–斯蒂尔杰斯测度。证明：

\begin{items}
\item

$ F \in W_1^1(\mathbb{R},\mathbb{R}) $ 蕴涵 $ \mu_F \ll \lambda_1 $。

\item

若 $ \mu_F $ 有限，则 $ \mu_F \ll \beta_1 $ 蕴涵 $ F \in W_1^1(\mathbb{R},\mathbb{R}) $。

\end{items}

\item

设 $ I $ 是区间，取 $ f \in \mathcal{L}_1(I,\mathbb{R}^n) $。对固定的 $ a \in I $，假设对 $ x \in I $ 有 $ \int_a^x f(t) \, dt = 0 $。证明对几乎每个 $ x \in I $ 都有 $ f(x) = 0 $。

\item

设 $ 0 \le a < b < \infty $，$ I := (-b,-a) \cup (a,b) $，并设 $ f \in \mathcal{L}_1(I,E) $。证明：若 $ f $ 为奇函数，则 $ \int_I f \, dx = 0 $；若 $ f $ 为偶函数，则 $ \int_I f \, dx = 2 \int_a^b f \, dx $。

\item

定义

\begin{equation*}
\begin{aligned}
K_0 &:= [0,1] , \\ K_1 &:= K_0 \setminus (3/8, 5/8) , \\ K_2 &:= K_1 \setminus ((5/32, 7/32) \cup (25/32, 27/32)) , \ldots
\end{aligned}
\end{equation*}

一般地，$ K_{n+1} $ 由 $ K_n $ 挖去长度为 $ (1/4)^{n+1} $ 的开“中间四分之一”区间得到，而不是像传统康托尔集（习题 III.3.8）的构造那样挖去中间三分之一。令 $ K := \bigcap K_n $，$ f := \chi_K $。证明 $ f $ 属于 $ \mathcal{L}_1([0,1]) $，且任何 $ g \in [f] $ 都不是黎曼可积的。

\end{exercises}

\section{富比尼定理}

本节的核心是证明：多变量函数的勒贝格积分可以迭代地计算，并且这一系列一元积分可以按任意顺序进行。因此多变量积分归结为对一元函数的积分。借助上一节的结果与第 II 卷中发展的方法，多维积分在许多情形下可以显式计算。

积分迭代求值的方法有广泛的理论应用，我们将介绍其中少数几个。

本节自始至终假设 $ m, n $ 是正整数，且 $ E $ 是 Banach 空间。

此外，我们一般把 $ \mathbb{R}^{m+n} $ 与 $ \mathbb{R}^m \times \mathbb{R}^n $ 等同起来。

\needspace{3\baselineskip}\subsection*{几乎处处有定义的映射}

设 $ (X, \mathcal{A}, \mu) $ 是测度空间。我们常常考虑仅在 $ \mu $-几乎处处有定义的非负 $ \overline{\mathbb{R}} $ 值函数。对于这些函数，我们简单地写作 $ x \mapsto f(x) $，而不具体指明其定义域。我们说这样的函数 $ x \mapsto f(x) $ 是\textbf{可测}的，如果存在 $ \mu $-零集 $ N $，使得 $ f \mid N^{c} : N^{c} \to \overline{\mathbb{R}}^{+} $ 有定义且 $ \mu $-可测。因此 $ \int_{N^{c}} f \, d\mu $ 有定义。若 $ M $ 是另一个 $ \mu $-零集，使得 $ f \mid M^{c} : M^{c} \to \overline{\mathbb{R}}^{+} $ 有定义且 $ \mu $-可测，则等式 $ \mu(N) = \mu(M) = \mu(M \cup N) = 0 $ 以及注 3.3(a) 与 (b) 蕴含

\[
\int_{N^{c}} f \, d\mu = \int_{M^{c} \cap N^{c}} f \, d\mu = \int_{M^{c}} f \, d\mu .
\]

因此

\begin{equation}
\int_{X} f \, d\mu := \int_{N^{c}} f \, d\mu\tag{6.1}
\end{equation}

是有良好定义的，并且与所选零集 $ N $ 无关。

对于 $ \mu $-几乎处处有定义的 $ E $ 值函数 $ x \mapsto f(x) $，我们完全像上面那样定义可测性。我们说这样的 $ f $ 是\textbf{可积}的，如果 $ f \mid N^{c} $ 属于 $ \mathcal{L}_{1}(N^{c}, \mu, E) $。此时 $ \int_{X} f \, d\mu $ 也通过 (6.1) 定义，并且引理 2.15 表明这个定义是有意义的。

例如考虑 $ A \in \mathcal{L}(m+n) $，并假设截面 $ A_{[x]} $ 对 $ \lambda_m $-几乎每个 $ x \in \mathbb{R}^m $ 是 $ \lambda_n $-可测的。则 $ x \mapsto \lambda_n(A_{[x]}) $ 是 $ \lambda_m $-几乎处处有定义的非负 $ \overline{\mathbb{R}} $ 值函数。若 $ x \mapsto \lambda_n(A_{[x]}) $ 可测，则积分 $ \int_{\mathbb{R}^m} \lambda_n(A_{[x]}) \, dx $ 有良好定义。

\needspace{3\baselineskip}\subsection*{卡瓦列里原理}

我们用 $ \mathcal{C}(m, n) $ 表示满足以下条件的全体 $ A \in \mathcal{L}(m+n) $ 构成的集合：

\begin{subitems}
\item

$ A_{[x]} \in \mathcal{L}(n) $ 对 $ \lambda_m $-几乎每个 $ x \in \mathbb{R}^m $ 成立；

\item

$ x \mapsto \lambda_n(A_{[x]}) $ 是 $ \lambda_m $-可测的；

\item

$ \lambda_{m+n}(A) = \int_{\mathbb{R}^m} \lambda_n(A_{[x]}) \, dx $。
我们想证明 $ \mathcal{C}(m, n) $ 与 $ \mathcal{L}(m+n) $ 一致，但这需要一些准备工作。

\end{subitems}

\begin{remark}

\begin{items}
\item

设 $ A \in \mathcal{C}(1, n) $ 有界，且 $ \mathrm{pr}_1(A) $ 是以 $ a $ 与 $ b $ 为端点的区间。则

\[
\lambda_{n+1}(A) = \int_{a}^{b} \lambda_n(A_{[x]}) \, dx .
\]

这一结论称为\textbf{卡瓦列里原理}，它精确表达了如下几何思想：$ A $ 的测度（体积）可以通过把 $ A $ 分割成薄的平行切片、并连续地对这些切片的体积求和（积分）来确定。

\end{items}

\bookfigure[0.95]{figures/d325120214151d5f.png}

\begin{items}
\item[（b）]

$ \mathcal{L}(m) \boxtimes \mathcal{L}(n) \subset \mathcal{C}(m, n) $。

\item[（c）]

对 $ \mathcal{C}(m, n) $ 中的每个递增序列 $ (A_j) $，并集 $ \bigcup_j A_j $ 属于 $ \mathcal{C}(m, n) $。

\begin{subitems}
\item

\proofhead
(i) 对 $ j \in \mathbb{N} $，设 $ M_j $ 是 $ \lambda_m $-零集，使得对 $ x \in M_j^{c} $ 有 $ A_{j,[x]} := (A_j)_{[x]} \in \mathcal{L}(n) $。令 $ A := \bigcup_j A_j $ 且 $ M := \bigcup_j M_j $，则对 $ x \in M^{c} $ 有 $ A_{[x]} = \bigcup_j A_{j,[x]} \in \mathcal{L}(n) $。$ \lambda_n $ 的下连续性蕴含，对 $ x \in M^{c} $ 有 $ \lambda_n(A_{[x]}) = \lim_j \lambda_n(A_{j,[x]}) $；借助命题 1.11 我们得出 $ x \mapsto \lambda_n(A_{[x]}) $ 是 $ \lambda_m $-可测的。
\item

由于 $ A_j \in \mathcal{C}(m, n) $，我们有

\[
\int_{\mathbb{R}^m} \lambda_n(A_{j,[x]}) \, dx = \lambda_{m+n}(A_j) \quad \text{对 } j \in \mathbb{N} ,
\]

由单调收敛定理推出

\begin{equation}
\lim_j \int_{\mathbb{R}^m} \lambda_n(A_{j,[x]}) \, dx = \int_{\mathbb{R}^m} \lambda_n(A_{[x]}) \, dx .\tag{6.2}
\end{equation}

$ \lambda_{m+n} $ 的下连续性于是表明

\[
\lambda_{m+n}(A) = \lim_j \lambda_{m+n}(A_j) = \lim_j \int_{\mathbb{R}^m} \lambda_n(A_{j,[x]}) \, dx = \int_{\mathbb{R}^m} \lambda_n(A_{[x]}) \, dx .
\]

因此 $ A $ 属于 $ \mathcal{C}(m, n) $。\nobreak\hspace{0.35em}■

\end{subitems}

\item[（d）]

设 $ (A_j) $ 是 $ \mathcal{C}(m, n) $ 中的递减序列，且存在 $ k \in \mathbb{N} $ 使 $ \lambda_{m+n}(A_k) < \infty $。则 $ \bigcap_j A_j $ 属于 $ \mathcal{C}(m, n) $。

\proofhead
我们令 $ A := \bigcap_j A_j $。对 $ \lambda_m $-几乎所有截面 $ A_{[x]} $ 的可测性以及 $ x \mapsto \lambda_n(A_{[x]}) $ 的可测性，都如同 (c) 中那样得到。其次，勒贝格控制收敛定理表明 (6.2) 在此情形下成立。结论现在如同 (c) 中那样推出。\nobreak\hspace{0.35em}■
\item[（e）]

设 $ (A_j) $ 是 $ \mathcal{C}(m, n) $ 中的不交序列。则 $ \bigcup_j A_j $ 也属于 $ \mathcal{C}(m, n) $。

\proofhead
由于 (c)，只需对有限不交序列证明该结论。我们把这个练习留给读者。\nobreak\hspace{0.35em}■
\item[（f）]

$ \mathbb{R}^{m+n} $ 中的每个开集都属于 $ \mathcal{C}(m, n) $。

\proofhead
这由命题 IX.5.6、(e) 与 (b) 推出。\nobreak\hspace{0.35em}■
\item[（g）]

$ \mathbb{R}^{m+n} $ 中的每个有界 $ G_\delta $ 集都属于 $ \mathcal{C}(m, n) $。

\proofhead
这由 (f) 与 (d) 推出。\nobreak\hspace{0.35em}■
\item[（h）]

设 $ A $ 是 $ \lambda_{m+n} $-零集。则 $ A $ 属于 $ \mathcal{C}(m, n) $，且存在 $ \lambda_m $-零集 $ M $，使得对每个 $ x \in M^{c} $，$ A_{[x]} $ 是 $ \lambda_n $-零集。

\proofhead
只需验证存在 $ \lambda_m $-零集 $ M $，使得对 $ x \in M^{c} $ 有 $ \lambda_n(A_{[x]}) = 0 $。于是令 $ A_j := A \cap (j \mathbb{B}^{m+n}) $ 对 $ j \in \mathbb{N} $。则 $ (A_j) $ 是有界 $ \lambda_{m+n} $-零集的递增序列，且 $ \bigcup_j A_j = A $。由推论 IX.5.5，存在有界 $ G_\delta $ 集序列 $ (G_j) $，使得 $ G_j \supset A_j $ 且 $ \lambda_{m+n}(G_j) = 0 $ 对 $ j \in \mathbb{N} $。由 (g) 因此推出
\[
0 = \lambda_{m+n}(G_j) = \int_{\mathbb{R}^m} \lambda_n(G_{j,[x]}) \, dx .
\]

于是，对每个 $ j \in \mathbb{N} $ 存在 $ \lambda_m $-零集 $ M_j $，使得对 $ x \in M_j^{c} $ 有 $ \lambda_n(G_{j,[x]}) = 0 $（见注 3.3(c)）。因为

\[
\bigcup_j G_{j,[x]} \supset \bigcup_j A_{j,[x]} = \left( \bigcup_j A_j \right)_{[x]} = A_{[x]} \quad \text{对 } x \in \mathbb{R}^m ,
\]

故 $ M := \bigcup_j M_j $ 具有所需的性质。\nobreak\hspace{0.35em}■

\end{items}

在这些注之后，我们现在可以证明 $ \mathcal{C}(m, n) $ 与 $ \mathcal{L}(m+n) $ 相等。

\end{remark}

\begin{proposition}

$ \mathcal{C}(m, n) = \mathcal{L}(m+n) $。

\proofhead
我们只需验证包含关系 $ \mathcal{L}(m+n) \subset \mathcal{C}(m, n) $。

\begin{subitems}
\item

设 $ A \in \mathcal{L}(m+n) $ 有界。由推论 IX.5.5，存在有界 $ G_\delta $ 集 $ G $，使得 $ G \supset A $ 且 $ \lambda_{m+n}(G) = \lambda_{m+n}(A) $。由于 $ A $ 的测度有限，$ G \backslash A $ 是有界 $ \lambda_{m+n} $-零集（命题 IX.2.3(ii)），并且利用注 6.1(h) 得出 $ (G \backslash A)_{[x]} = G_{[x]} \backslash A_{[x]} $ 对 $ \lambda_m $-几乎每个 $ x \in \mathbb{R}^m $ 是 $ \lambda_n $-零集。由注 6.1(g)，$ G_{[x]} $ 对 $ \lambda_m $-几乎每个 $ x \in \mathbb{R}^m $ 属于 $ \mathcal{L}(n) $。因为

\[
A_{[x]} = G_{[x]} \cap (G_{[x]} \backslash A_{[x]})^{c} \quad \text{对 } x \in \mathbb{R}^m ,
\]

这对 $ \lambda_m $-几乎每个交集 $ A_{[x]} $ 也成立。此外，$ \lambda_n(A_{[x]}) = \lambda_n(G_{[x]}) $ 对 $ \lambda_m $-几乎每个 $ x \in \mathbb{R}^m $ 成立。由注 6.1(g) 我们知道 $ G $ 属于 $ \mathcal{C}(m, n) $。因此 $ x \mapsto \lambda_n(A_{[x]}) $ 可测，且

\[
\lambda_{m+n}(G) = \int_{\mathbb{R}^m} \lambda_n(G_{[x]}) \, dx = \int_{\mathbb{R}^m} \lambda_n(A_{[x]}) \, dx .
\]

因此 $ A $ 属于 $ \mathcal{C}(m, n) $。

\item

若 $ A $ 无界，我们令 $ A_j := A \cap (j \mathbb{B}^{m+n}) $ 对 $ j \in \mathbb{N} $。则 $ (A_j) $ 是 $ \mathcal{L}(m+n) $ 中的递增序列，且 $ \bigcup_j A_j = A $。结论现在由 (i) 与注 6.1(c) 推出。\nobreak\hspace{0.35em}■

\end{subitems}

\end{proposition}

\begin{corollary}

若 $ A \in \mathcal{L}(m+n) $ 的测度有限，则 $ \lambda_n(A_{[x]}) < \infty $ 对 $ \lambda_m $-几乎每个 $ x \in \mathbb{R}^m $ 成立。

\proofhead
由于命题 6.2 蕴含
\[
\int_{\mathbb{R}^m} \lambda_n(A_{[x]}) \, dx = \lambda_{m+n}(A) < \infty ,
\]

结论由注 3.11(c) 推出。\nobreak\hspace{0.35em}■

\end{corollary}

对 $ A \in \mathcal{L}(m+n) $ 与 $ x \in \mathbb{R}^m $，我们有 $ \chi_A(x, \cdot) = \chi_{A_{[x]}} $，因此命题 6.2 也可以用特征函数来表述。于是很容易把该结论应用到特征函数的线性组合上，从而应用到简单函数上。

\begin{lemma}

设 $ f \in \mathcal{S}(\mathbb{R}^{m+n}, E) $。

\begin{subitems}
\item

$ f(x, \cdot) \in \mathcal{S}(\mathbb{R}^n, E) $ 对 $ \lambda_m $-几乎每个 $ x \in \mathbb{R}^m $ 成立。

\item

$ E $ 值函数 $ x \mapsto \int_{\mathbb{R}^n} f(x, y) \, dy $ 是 $ \lambda_m $-可积的。

\item

$ \int_{\mathbb{R}^{m+n}} f \, d(x, y) = \int_{\mathbb{R}^m} \left[ \int_{\mathbb{R}^n} f(x, y) \, dy \right] dx $。

\proofhead
(i) 取 $ f = \sum_{j=0}^{k} e_j \chi_{A_j} $，则对 $ x \in \mathbb{R}^m $ 有 $ f(x, \cdot) = \sum_{j=0}^{k} e_j \chi_{A_{j,[x]}} $。于是由命题 6.2 与推论 6.3 容易得出：存在 $ \lambda_m $-零集 $ M $，使得对每个 $ x \in M^{c} $，$ f(x, \cdot) $ 属于 $ \mathcal{S}(\mathbb{R}^n, E) $。
\end{subitems}

\begin{subitems}
\item[（ii）]

我们令

\begin{equation}
g(x) := \int_{\mathbb{R}^n} f(x, y) \, dy = \sum_{j=0}^{k} e_j \lambda_n(A_{j,[x]}) \quad \text{对 } x \in M^{c} .\tag{6.3}
\end{equation}

则命题 6.2 与注 1.2(d) 表明 $ x \mapsto g(x) $ 是 $ \lambda_m $-可测的。此外，我们有

\[
\int_{\mathbb{R}^m} |g| \, dx \le \sum_{j=0}^{k} |e_j| \int_{\mathbb{R}^m} \lambda_n(A_{j,[x]}) \, dx = \sum_{j=0}^{k} |e_j| \lambda_{m+n}(A_j) < \infty .
\]

因此 $ x \mapsto g(x) $ 是 $ \lambda_m $-可积的。

\item[（iii）]

最后，由命题 6.2 与 (6.3) 推出

\begin{equation*}
\begin{aligned}
\int_{\mathbb{R}^{m+n}} f \, d(x, y) &= \sum_{j=0}^{k} e_j \lambda_{m+n}(A_j) = \sum_{j=0}^{k} e_j \int_{\mathbb{R}^m} \lambda_n(A_{j,[x]}) \, dx = \int_{\mathbb{R}^m} g \, dx \\ &= \int_{\mathbb{R}^m} \left[ \int_{\mathbb{R}^n} f(x, y) \, dy \right] dx ,
\end{aligned}
\end{equation*}

这就完成了证明。\nobreak\hspace{0.35em}■

\end{subitems}

\end{lemma}

\begin{remark}

在集合 $ \mathcal{C}(m, n) $ 的定义中，我们选择标出 $ \mathbb{R}^{m+n} $ 的前 $ m $ 个坐标。我们本来也完全可以选取后 $ n $ 个坐标，并用 $ \lambda_m(A^{[y]}) $ 而不是 $ \lambda_n(A_{[x]}) $ 对 $ \lambda_n $-几乎每个 $ y \in \mathbb{R}^n $ 作同样的论证。按这种 $ \mathcal{C}(m, n) $ 的定义，我们显然同样会得到 $ \mathcal{C}(m, n) = \mathcal{L}(m+n) $。因此引理 6.4 中 $ x $ 与 $ y $ 的地位可以互换，于是我们得出：对 $ f \in \mathcal{S}(\mathbb{R}^{m+n}, E) $，

\begin{subitems}
\item

$ f(\cdot, y) \in \mathcal{S}(\mathbb{R}^m, E) $ 对 $ \lambda_n $-几乎每个 $ y \in \mathbb{R}^n $ 成立；

\item

$ E $ 值函数 $ y \mapsto \int_{\mathbb{R}^m} f(x, y) \, dx $ 是 $ \lambda_n $-可积的；

\item

$ \int_{\mathbb{R}^{m+n}} f \, d(x, y) = \int_{\mathbb{R}^n} \left[ \int_{\mathbb{R}^m} f(x, y) \, dx \right] dy $。
特别地，我们得到

\[
\int_{\mathbb{R}^m} \left[ \int_{\mathbb{R}^n} f(x, y) \, dy \right] dx = \int_{\mathbb{R}^n} \left[ \int_{\mathbb{R}^m} f(x, y) \, dx \right] dy
\]

对 $ f \in \mathcal{S}(\mathbb{R}^{m+n}, E) $ 成立。换句话说，对简单函数而言，积分 $ \int_{\mathbb{R}^{m+n}} f \, d(x, y) $ 可以迭代计算，而且积分进行的顺序无关紧要。\nobreak\hspace{0.35em}■

\end{subitems}

\end{remark}

\needspace{3\baselineskip}\subsection*{卡瓦列里原理的应用}

本节的主要结果是：注 6.5 关于积分迭代计算的结论对任意可积函数 $ f $ 仍然成立。在证明这个定理之前，我们先给出卡瓦列里原理的几个应用，也就是说我们在 $ f = \chi_A $ 的情形下工作。

\begin{example}

\begin{items}
\item

（积分的几何解释）对 $ M \in \mathcal{L}(m) $ 与 $ f \in \mathcal{L}_{0}(M, \mathbb{R}^{+}) $，集合

\[
S_f := S_{f,M} := \{ (x, y) \in \mathbb{R}^m \times \mathbb{R} ; 0 \le y \le f(x), x \in M \}
\]

属于 $ \mathcal{L}(m+1) $，且

\[
\int_{M} f \, dx = \lambda_{m+1}(S_f) ,
\]

也就是说，积分 $ \int_{M} f \, dx $ 等于 $ f $ 的图像下方的点集的 $ (m+1) $ 维勒贝格测度。\footnote{比较第 VI.3 节的引言性评注。}

\end{items}

\bookfigure[0.46]{figures/4050ed19dc01989c.png}

\proofhead
令 $ f_1 := \mathrm{pr}_{\mathbb{R}} $ 与 $ f_2 := f \circ \mathrm{pr}_{\mathbb{R}^m} $。则 $ f_1 $ 与 $ f_2 $ 属于 $ \mathcal{L}_{0}(M \times \mathbb{R}, \overline{\mathbb{R}}^{+}) $，且 $ S_f = [0 \le f_1 \le f_2] $。因此命题 1.9 蕴含 $ S_f $ 的 $ \lambda_{m+1} $-可测性。因为 $ (S_f)_{[x]} = [0, f(x)] $ 对 $ x \in M $ 成立，所以 $ \lambda_1((S_f)_{[x]}) = f(x) $，从而
\[
\lambda_{m+1}(S_f) = \int_{\mathbb{R}^m} \lambda_1((S_f)_{[x]}) \, dx = \int_{M} f \, dx ,
\]

由命题 6.2 即得。\nobreak\hspace{0.35em}■

\begin{items}
\item[（b）]

（线性映射的换元法则）设 $ T \in \mathcal{L}(\mathbb{R}^m) $、$ a \in \mathbb{R}^m $ 与 $ M \in \mathcal{L}(m) $。又设 $ \varphi(x) := a + Tx $ 对 $ x \in \mathbb{R}^m $，且 $ f \in \mathcal{L}_{1}(\varphi(M)) $。则 $ f \circ \varphi $ 属于 $ \mathcal{L}_{1}(M) $，且

\begin{equation}
\int_{\varphi(M)} f \, dy = |\det T| \int_{M} (f \circ \varphi) \, dx .\tag{6.4}
\end{equation}

特别地，勒贝格积分是\textbf{仿射等距不变}的，即对 $ \mathbb{R}^m $ 的每个仿射等距 $ \varphi $，我们有

\[
\int_{\mathbb{R}^m} f = \int_{\mathbb{R}^m} f \circ \varphi \quad \text{对 } f \in \mathcal{L}_{1}(\mathbb{R}^m) .
\]

\proofhead
(i) 由定理 IX.5.12，$ \varphi $ 把 $ \sigma $-代数 $ \mathcal{L}(m) $ 映入自身。因此 $ \varphi(M) $ 属于 $ \mathcal{L}(m) $，且定理 1.4 蕴含 $ f \circ \varphi $ 位于 $ \mathcal{L}_{0}(M) $ 中。分解 $ f = f_1 - f_2 + i(f_3 - f_4) $，其中 $ f_j \in \mathcal{L}_{1}(\varphi(M), \mathbb{R}^{+}) $，表明我们可以限于 $ f \in \mathcal{L}_{1}(\varphi(M), \mathbb{R}^{+}) $ 的情形。此时 (a) 说明
\begin{equation}
\int_{\varphi(M)} f = \lambda_{m+1}(S_{f,\varphi(M)}) , \quad \int_{M} f \circ \varphi = \lambda_{m+1}(S_{f \circ \varphi, M}) .\tag{6.5}
\end{equation}

\begin{subitems}
\item

我们令 $ \hat{a} := (a, 0) \in \mathbb{R}^m \times \mathbb{R} $ 以及 $ \hat{T}(x, t) := (Tx, t) $ 对 $ (x, t) \in \mathbb{R}^m \times \mathbb{R} $。则 $ \hat{a} + \hat{T}(S_{f \circ \varphi}) = S_f $ 且 $ \det T = \det \hat{T} $，因为表示矩阵 $ \hat{T} $ 具有分块结构

\begin{equation*}
[\hat{T}] = \begin{bmatrix} [T] & 0 \\ 0 & 1 \end{bmatrix} .
\end{equation*}

由推论 IX.5.23 与定理 IX.5.25 因此推出

\[
\lambda_{m+1}(S_f) = \lambda_{m+1}(\hat{T}(S_{f \circ \varphi})) = |\det T| \lambda_{m+1}(S_{f \circ \varphi}) ,
\]

结合 (6.5) 即证得 (6.4)。$ f \circ \varphi $ 的可积性由注 3.11(a) 推出。\nobreak\hspace{0.35em}■

\end{subitems}

\item[（c）]

（$ \mathbb{R}^m $ 中单位球的体积）对 $ m \in \mathbb{N}^{\times} $，我们有

\[
\lambda_m(\mathbb{B}^m) = \frac{\pi^{m/2}}{\Gamma(1 + m/2)} ;
\]

特别地，$ \lambda_1(\mathbb{B}^1) = 2 $、$ \lambda_2(\mathbb{B}^2) = \pi $，且 $ \lambda_3(\mathbb{B}^3) = 4\pi/3 $。

\proofhead
令 $ \omega_m := \lambda_m(\mathbb{B}^m) $，由卡瓦列里原理以及注 IX.5.26(b) 与 6.5 得到
\begin{equation*}
\begin{aligned}
\omega_m &= \int_{-1}^{1} \lambda_{m-1}((\mathbb{B}^m)^{[y]}) \, dy \\ &= \int_{-1}^{1} \lambda_{m-1}(\sqrt{1 - y^2} \mathbb{B}^{m-1}) \, dy \\ &= \omega_{m-1} \int_{-1}^{1} (\sqrt{1 - y^2})^{m-1} \, dy .
\end{aligned}
\end{equation*}
\end{items}

\bookfigure[0.32]{figures/4372387e0b3742fd.png}

为了计算积分

\[
B_m := \int_{-1}^{1} (1 - y^2)^{(m-1)/2} \, dy = 2 \int_{0}^{1} (1 - y^2)^{(m-1)/2} \, dy \quad \text{对 } m \in \mathbb{N}^{\times} ,
\]

我们令 $ y = -\cos x $，于是 $ dy = \sin x \, dx $。由此得到 $ B_m = 2 \int_{0}^{\pi/2} \sin^m x \, dx $。由例 VI.5.5(d) 的证明可得

\[
B_{2m} = \frac{(2m-1)(2m-3) \cdots 1}{2m(2m-2) \cdots 2} \cdot \pi , \quad B_{2m+1} = \frac{2m(2m-2) \cdots 2}{(2m+1)(2m-1) \cdots 1} \cdot 2 .
\]

于是得到 $ B_m B_{m-1} = 2\pi/m $，并且

\begin{equation}
\omega_m = B_m \omega_{m-1} = B_m B_{m-1} \omega_{m-2} = \frac{2\pi}{m} \omega_{m-2} .\tag{6.6}
\end{equation}

由于 $ \omega_1 = 2 $，我们得到 $ \omega_2 = B_2 \omega_1 = 2B_2 = \pi $，于是由 (6.6) 得

\[
\omega_{2m} = \frac{\pi^m}{m!} , \quad \omega_{2m+1} = \frac{(2\pi)^m}{1 \cdot 3 \cdot 5 \cdots (2m+1)} \cdot 2 .
\]

这两个表达式可以借助 Gamma 函数统一起来，因为

\[
\Gamma(m+1) = m! , \quad \Gamma\left(m + \frac{3}{2}\right) = \frac{\sqrt{\pi}}{2^{m+1}} \cdot 1 \cdot 3 \cdots (2m+1)
\]

（见定理 VI.9.2 与习题 VI.9.1）。\nobreak\hspace{0.35em}■

\end{example}

\needspace{3\baselineskip}\subsection*{托内利定理}

我们现在证明前面预告的定理，它为非负 $ \overline{\mathbb{R}} $ 值函数积分的迭代计算提供了依据。这个版本，即托内利定理，将为我们提供关于 $ E $ 值函数的一个重要的可积性判别法。

\begin{theorem}

（托内利）设 $ f \in \mathcal{L}_{0}(\mathbb{R}^{m+n}, \overline{\mathbb{R}}^{+}) $，

\begin{subitems}
\item

$ f(x, \cdot) \in \mathcal{L}_{0}(\mathbb{R}^n, \overline{\mathbb{R}}^{+}) $ 对 $ \lambda_m $-几乎所有 $ x \in \mathbb{R}^m $ 成立，$ f(\cdot, y) \in \mathcal{L}_{0}(\mathbb{R}^m, \overline{\mathbb{R}}^{+}) $ 对 $ \lambda_n $-几乎所有 $ y \in \mathbb{R}^n $ 成立；

\item

$ x \mapsto \int_{\mathbb{R}^n} f(x, y) \, dy $ 是 $ \lambda_m $-可测的，$ y \mapsto \int_{\mathbb{R}^m} f(x, y) \, dx $ 是 $ \lambda_n $-可测的；

\item

$ \int_{\mathbb{R}^{m+n}} f \, d(x, y) = \int_{\mathbb{R}^m} \left[ \int_{\mathbb{R}^n} f(x, y) \, dy \right] dx = \int_{\mathbb{R}^n} \left[ \int_{\mathbb{R}^m} f(x, y) \, dx \right] dy $。

\proofhead
(i) 由定理 1.12，存在 $ \mathcal{S}(\mathbb{R}^{m+n}, \mathbb{R}^{+}) $ 中的递增序列 $ (f_j) $ 收敛于 $ f $。单调收敛定理于是给出
\begin{equation}
\lim_j \int_{\mathbb{R}^{m+n}} f_j \, d(x, y) = \int_{\mathbb{R}^{m+n}} f \, d(x, y) \quad \text{在 } \overline{\mathbb{R}}^{+} \text{ 中} .\tag{6.7}
\end{equation}

其次，由引理 6.4，对每个 $ j \in \mathbb{N} $ 存在 $ \lambda_m $-零集 $ M_j $，使得对 $ x \in M_j^{c} $ 有 $ f_j(x, \cdot) \in \mathcal{S}(\mathbb{R}^n, \mathbb{R}^{+}) $。若令 $ M := \bigcup_j M_j $，则由单调收敛定理可见

\begin{equation}
\int_{\mathbb{R}^n} f_j(x, y) \, dy \uparrow \int_{\mathbb{R}^n} f(x, y) \, dy \quad \text{对 } x \in M^{c} .\tag{6.8}
\end{equation}

引理 6.4(ii)、命题 1.11、$ M $ 的测度为零这一事实以及 (6.8) 蕴含 $ \overline{\mathbb{R}} $ 值函数 $ x \mapsto \int_{\mathbb{R}^n} f(x, y) \, dy $ 是 $ \lambda_m $-可测的。其次，由 (6.7)、引理 6.4(iii)、(6.8) 以及单调收敛定理推出

\begin{equation*}
\begin{aligned}
\int_{\mathbb{R}^{m+n}} f \, d(x, y) &= \lim_j \int_{\mathbb{R}^{m+n}} f_j \, d(x, y) = \lim_j \int_{\mathbb{R}^m} \left[ \int_{\mathbb{R}^n} f_j(x, y) \, dy \right] dx \\ &= \int_{\mathbb{R}^m} \left[ \int_{\mathbb{R}^n} f(x, y) \, dy \right] dx .
\end{aligned}
\end{equation*}

其余结论可类似地证明（注意注 6.5）。\nobreak\hspace{0.35em}■

\end{subitems}

\end{theorem}

\begin{corollary}

设 $ f \in \mathcal{L}_{0}(\mathbb{R}^{m+n}, E) $ 满足 $ f = 0 $ 关于 $ \lambda_{m+n} $ 几乎处处。则存在 $ \lambda_m $-零集 $ M $，使得对每个 $ x \in M^{c} $，$ f(x, \cdot) $ 关于 $ \lambda_n $ 几乎处处为零；并且存在 $ \lambda_n $-零集 $ N $，使得对每个 $ y \in N^{c} $，$ f(\cdot, y) = 0 $ 关于 $ \lambda_m $ 几乎处处成立。

\proofhead
显然只需证明 $ M $ 的存在性（比较注 6.5）。托内利定理给出
\[
\int_{\mathbb{R}^m} \left[ \int_{\mathbb{R}^n} |f(x, y)| \, dy \right] dx = \int_{\mathbb{R}^{m+n}} |f| \, d(x, y) = 0 .
\]

于是由注 3.3(c) 存在 $ \lambda_m $-零集 $ M $，使得

\[
\int_{\mathbb{R}^n} |f(x, y)| \, dy = 0 \quad \text{对 } x \in M^{c} ,
\]

再次由注 3.3(c) 即得结论。\nobreak\hspace{0.35em}■

\end{corollary}

\needspace{3\baselineskip}\subsection*{标量函数的富比尼定理}

现在很容易把托内利定理推广到可积的 $ \mathbb{K} $ 值函数的情形，这在应用上特别有意义。

\begin{theorem}

（富比尼）设 $ f \in \mathcal{L}_{1}(\mathbb{R}^{m+n}) $，

\begin{subitems}
\item

$ f(x, \cdot) \in \mathcal{L}_{1}(\mathbb{R}^n) $ 对 $ \lambda_m $-几乎所有 $ x \in \mathbb{R}^m $ 成立，$ f(\cdot, y) \in \mathcal{L}_{1}(\mathbb{R}^m) $ 对 $ \lambda_n $-几乎所有 $ y \in \mathbb{R}^n $ 成立；

\item

$ x \mapsto \int_{\mathbb{R}^n} f(x, y) \, dy $ 是 $ \lambda_m $-可积的，$ y \mapsto \int_{\mathbb{R}^m} f(x, y) \, dx $ 是 $ \lambda_n $-可积的；

\item

$ \int_{\mathbb{R}^{m+n}} f \, d(x, y) = \int_{\mathbb{R}^m} \left[ \int_{\mathbb{R}^n} f(x, y) \, dy \right] dx = \int_{\mathbb{R}^n} \left[ \int_{\mathbb{R}^m} f(x, y) \, dx \right] dy $。

\end{subitems}

证明

\begin{items}
\item

对 $ f \in \mathcal{L}_{1}(\mathbb{R}^{m+n}, \mathbb{R}^{+}) $，结论由托内利定理与注 3.3(e) 推出。

\item

给定表示 $ f = f_1 - f_2 + i(f_3 - f_4) $，其中 $ f_j \in \mathcal{L}_{1}(\mathbb{R}^{m+n}, \mathbb{R}^{+}) $，一般情形现在由推论 2.12 与积分的线性性推出。\nobreak\hspace{0.35em}■

\end{items}

\end{theorem}

\begin{corollary}

设 $ A \in \mathcal{L}(m) $ 且 $ f \in \mathcal{L}_{1}(A) $。设 $ (j_1, \ldots, j_m) $ 是 $ (1, \ldots, m) $ 的一个排列。则

\[
\int_{A} f \, dx = \int_{\mathbb{R}} \left( \int_{\mathbb{R}} \cdots \left( \int_{\mathbb{R}} \tilde{f}(x_1 \ldots, x_m) \, dx_{j_1} \right) \cdots dx_{j_{m-1}} \right) dx_{j_m} .
\]

富比尼定理保证\textbf{可积}函数可以按任意顺序积分。与托内利定理相结合，我们得到多变量函数可积性的一个简单、通用且极其重要的判别法，以及一种显式计算积分的方法。

\end{corollary}

\qyxffntextdef{1}{正如注 IX.5.14(b) 所示，情形一般并非如此。}
\begin{theorem}

（富比尼–托内利）设 $ A \in \mathcal{L}(m+n) $ 且 $ f \in \mathcal{L}_{0}(A) $。

\begin{subitems}
\item

若下列积分之一

\[
\int_{\mathbb{R}^m} \left[ \int_{\mathbb{R}^n} |\tilde{f}(x, y)| \, dy \right] dx , \quad \int_{\mathbb{R}^n} \left[ \int_{\mathbb{R}^m} |\tilde{f}(x, y)| \, dx \right] dy , \quad \int_{A} |f| \, d(x, y)
\]

为有限，则其余每个也为有限，且它们全都相等。此时 $ f $ 可积，并且定理 6.9 的结论对 $ \tilde{f} $ 成立。

\item

若 $ \operatorname{pr}_{\mathbb{R}^m}(A) $ 可测\qyxffnmark{1}且 $ f $ 可积，则

\[
\int_A f\,d(x,y) = \int_{\operatorname{pr}_{\mathbb{R}^m}(A)}\left[\int_{A_{[x]}} f(x,y)\,dy\right]dx .
\]

\proofhead
因为 $ \widetilde{f} $ 属于 $ \mathcal{L}_0(\mathbb{R}^{m+n}) $，第一个结论由托内利定理立即得到。再由定理 3.9，$ \widetilde{f} $ 可积，从而 $ f $ 也可积。结论至此显然。\nobreak\hspace{0.35em}■
\end{subitems}

\end{theorem}

\begin{remark}

\begin{items}
\item

通过选取前 $ m $ 个坐标，我们没有损失一般性，因为由推论 6.10，这种次序总可以通过一个置换达到。

\item

通常我们省略 $ \int_{\mathbb{R}^n}\left[\int_{\mathbb{R}^m} f(x,y)\,dx\right]dy $ 中的括号，而写成比如说

\begin{equation}
\int_{\mathbb{R}^n}\int_{\mathbb{R}^m} f(x,y)\,dx\,dy .\tag{6.9}
\end{equation}

在这种记法下，约定积分按从内到外的次序计算。\footnote{也就是说，积分 $ \int_{\mathbb{R}^m} f(x,y)\,dx $ 是对固定的 $ y $ 计算的，然后把结果在 $ \mathbb{R}^n $ 上对 $ y $ 积分。} \textbf{累次积分} (6.9) 要与 $ (m+n) $ 维积分

\[
\int_{\mathbb{R}^{m+n}} f\,d(x,y) = \int_{\mathbb{R}^{m+n}} f\,d\lambda_{m+n} .
\]

区分开来。

\item

存在 $ f \in \mathcal{L}_0(\mathbb{R}^2)\setminus\mathcal{L}_1(\mathbb{R}^2) $，使得

\[
\int_{\mathbb{R}}\int_{\mathbb{R}} f(x,y)\,dx\,dy = \int_{\mathbb{R}}\int_{\mathbb{R}} f(x,y)\,dy\,dx = 0 .
\]

因此，累次积分的存在与相等并不蕴含 $ f $ 可积。

\proofhead
定义 $ f : \mathbb{R}^2 \to \mathbb{R} $ 为
\begin{equation}
f(x,y) := \begin{cases} \frac{xy}{(x^2+y^2)^2} & \text{若 } (x,y) \ne (0,0) , \\ 0 & \text{若 } (x,y) = (0,0) . \end{cases}\tag{6.10}
\end{equation}

于是 $ f $ 是 $ \lambda_2 $-可测的。对每个 $ y \in \mathbb{R} $，广义黎曼积分 $ \int_{\mathbb{R}} f(x,y)\,dx $ 绝对收敛。又 $ f(\cdot,y) $ 为奇函数。故 $ \int_{\mathbb{R}} f(x,y)\,dx = 0 $ 对每个\footnote{情形 $ y = 0 $ 已包含在所给论证之中，尽管它可由 $ f(\cdot,0) = 0 $ 更简单地推出。} $ y \in \mathbb{R} $ 成立，因此，由于 $ f(x,y) = f(y,x) $，我们有

\[
\int_{\mathbb{R}}\int_{\mathbb{R}} f(x,y)\,dx\,dy = \int_{\mathbb{R}}\int_{\mathbb{R}} f(x,y)\,dy\,dx = 0 .
\]

现在假设 $ f $ 可积。那么由富比尼定理，$ x \mapsto \int_{\mathbb{R}}|f(x,y)|\,dy $ 也将可积，而这由于

\[
\int_{\mathbb{R}} \frac{|xy|}{(x^2+y^2)^2}\,dy = \frac{1}{|x|} \quad \text{对 } x \ne 0 ,
\]

不可能成立。\nobreak\hspace{0.35em}■

\item

存在 $ g \in \mathcal{L}_0(\mathbb{R}^2)\setminus\mathcal{L}_1(\mathbb{R}^2) $，使得

\[
0 < \left|\int_{\mathbb{R}}\int_{\mathbb{R}} g(x,y)\,dx\,dy\right| = \left|\int_{\mathbb{R}}\int_{\mathbb{R}} g(x,y)\,dy\,dx\right| < \infty .
\]

\proofhead
取 $ f $ 为 (6.10) 中的函数，并取 $ h \in \mathcal{L}_1(\mathbb{R}^2) $ 使 $ \int h\,d(x,y) > 0 $。则 $ g := f + h $ 具有所述性质。\nobreak\hspace{0.35em}■
\end{items}

\end{remark}

\begin{example}

\begin{items}
\item

（多维高斯积分）对 $ n \in \mathbb{N}^\times $，有

\[
\int_{\mathbb{R}^n} e^{-|x|^2} dx = \pi^{n/2} .
\]

\proofhead
利用 $ |x|^2 = x_1^2 + \cdots + x_n^2 $ 以及指数函数的性质，由托内利定理可得
\[
\int_{\mathbb{R}^n} e^{-|x|^2} dx = \int_{\mathbb{R}} \cdots \int_{\mathbb{R}} e^{-x_1^2} e^{-x_2^2} \cdot \cdots \cdot e^{-x_n^2} \, dx_1 \cdots dx_n = \prod_{j=1}^{n}\int_{\mathbb{R}} e^{-x_j^2} dx_j = \left(\int_{\mathbb{R}} e^{-t^2} dt\right)^n .
\]

现在结论由应用 VI.9.7 得到。\nobreak\hspace{0.35em}■

\item

（贝塔函数\footnote{比较注 VI.9.12(a)。}的一个表示）对 $ v,w \in [\operatorname{Re} z > 0] $，

\[
\mathrm{B}(v,w) = \frac{\Gamma(v)\Gamma(w)}{\Gamma(v+w)} .
\]

\proofhead
令 $ A := \{(s,t)\in\mathbb{R}^2 ; 0 < t < s\} $，并定义 $ \gamma_{v,w} : A \to \mathbb{C} $ 为 $ \gamma_{v,w}(s,t) := t^{v-1}(s-t)^{w-1}e^{-s} $，其中 $ v,w \in [\operatorname{Re} z > 0] $。再令 $ \gamma_v(t) := t^{v-1}e^{-t} $（$ t > 0 $），由托内利定理可得
\[
\int_A |\gamma_{v,w}(s,t)|\,d(s,t) = \int_0^\infty\int_t^\infty |\gamma_{v,w}(s,t)|\,ds\,dt = \left(\int_0^\infty \gamma_{\operatorname{Re} v}(t)\,dt\right)\left(\int_0^\infty \gamma_{\operatorname{Re} w}(s)\,ds\right) = \Gamma(\operatorname{Re} v)\Gamma(\operatorname{Re} w) < \infty .
\]

因此 $ \gamma_{v,w} $ 可积，而富比尼定理类似地给出

\begin{equation}
\int_A \gamma_{v,w}(s,t)\,d(s,t) = \int_0^\infty\int_t^\infty \gamma_{v,w}(s,t)\,ds\,dt = \Gamma(v)\Gamma(w) .\tag{6.11}
\end{equation}

由于 $ \operatorname{pr}_1(A) = \mathbb{R}^+ $ 且 $ A_{[s]} = [0,s] $ 对 $ s > 0 $ 成立，由 (6.11) 与定理 6.11(ii) 我们得到

\[
\Gamma(v)\Gamma(w) = \int_0^\infty\left(\int_0^s t^{v-1}(s-t)^{w-1} dt\right) e^{-s} ds .
\]

在内层积分中作代换 $ r = t/s $，并利用贝塔函数的定义，得到

\begin{equation}
\Gamma(v)\Gamma(w) = \int_0^\infty\left(\int_0^1 r^{v-1}(1-r)^{w-1} dr\right) s^{v+w-1} e^{-s} ds = \mathrm{B}(v,w)\Gamma(v+w) ,\tag{6.12}
\end{equation}

这就完成了证明。\nobreak\hspace{0.35em}■

\end{items}

例 (b) 表明，只要巧妙地选择积分次序，复杂的积分就可以被简化。

\end{example}

\needspace{3\baselineskip}\subsection*{向量值函数的富比尼定理}

我们现在要说明富比尼定理对 $ E $ 值函数同样成立\footnote{本节初读时可以跳过。}，并给出一些应用。几条预备注记将会有用。

假设 $ A \in \mathcal{L}(m+n) $ 具有有限测度。由命题 6.2 与推论 6.3，存在一个 $ \lambda_m $-零集 $ M $，使得 $ A_{[x]} \in \mathcal{L}(n) $ 且 $ \lambda_n(A_{[x]}) < \infty $ 对 $ x \in M^c $ 成立。我们固定 $ q \in [1,\infty) $。因为 $ |\chi_{A_{[x]}}|^q = \chi_{A_{[x]}} $，我们有

\[
\int_{\mathbb{R}^n}|\chi_{A_{[x]}}(y)|^q dy = \int_{\mathbb{R}^n}\chi_{A_{[x]}}(y)\,dy = \lambda_n(A_{[x]}) < \infty .
\]

若按 §4 中约定的那样，把 $ \chi_{A_{[x]}} $ 与所有在 $ \lambda_n $-几乎处处 意义下与 $ y \mapsto \chi_{A_{[x]}}(y) $ 相重合的函数构成的等价类等同起来，我们便得到映射

\[
M^c \to F := L_q(\mathbb{R}^n) , \quad x \mapsto \chi_{A_{[x]}} .
\]

因为 $ F $ 是巴拿赫空间，我们可以研究它的可测性与可积性。

\begin{lemma}

假设 $ A \in \mathcal{L}(m+n) $ 具有有限测度。则在 $ \lambda_m $-几乎处处 意义下有定义的 $ F $ 值映射 $ x \mapsto \chi_{A_{[x]}} $ 是 $ \lambda_m $-可测的。

\proofhead
我们把 $ x \mapsto \chi_{A_{[x]}} $ 的平凡延拓记为 $ \psi_A : \mathbb{R}^m \to F $。

\begin{subitems}
\item

假设 $ A $ 是 $ \lambda_{m+n} $-零集。由注 6.1(h)，存在 $ \lambda_m $-零集 $ M $，使得对 $ x \in M^c $，$ A_{[x]} $ 是 $ \lambda_n $-零集。因此对 $ x \in M^c $，在 $ F $ 中 $ \psi_A(x) = 0 $。结论成立。

\item

现在假设 $ A $ 是形如 $ [a,b) $ 的区间，其中 $ a,b \in \mathbb{R}^{m+n} $。令 $ J_1 := \prod_{j=1}^{m}[a_j,b_j) $、$ J_2 := \prod_{j=m+1}^{m+n}[a_j,b_j) $。因为 $ A = J_1 \times J_2 $，我们有

\[
\chi_{A_{[x]}} = \chi_{J_1}(x)\chi_{J_2} \quad \text{对 } x \in \mathbb{R}^m ,
\]

于是可见在这种情形下 $ \psi_A $ 属于 $ \mathcal{S}(\mathbb{R}^m,F) $。

\item

假设 $ A \subset \mathbb{R}^{m+n} $ 是开集，且 $ (I_j) $ 是形如 $ [a,b) $ 的互不相交区间序列，满足 $ A = \bigcup_j I_j $（见命题 IX.5.6）。我们令

\[
f_k := \sum_{j=0}^{k}\psi_{I_j} \quad \text{对 } k \in \mathbb{N} .
\]

由 (ii) 与注 1.2(a)，$ (f_k) $ 是 $ \mathcal{S}(\mathbb{R}^m,F) $ 中的一个序列。此外，存在一个勒贝格测度为零的集合 $ M $，使得

\[
\|\psi_A(x) - f_k(x)\|_F^q = \int_{\mathbb{R}^n}\left|\chi_{A_{[x]}}(y) - \left(\sum_{j=0}^{k}\chi_{(I_j)_{[x]}}(y)\right)\right|^q dy = \lambda_n\left(\bigcup_{j=k+1}^{\infty}(I_j)_{[x]}\right) = \sum_{j=k+1}^{\infty}\lambda_n((I_j)_{[x]})
\]

对 $ x \in M^c $ 成立。此外，由推论 6.3，$ A_{[x]} $ 具有有限测度，且 $ \lambda_n(A_{[x]}) = \sum_{j=0}^{\infty}\lambda_n((I_j)_{[x]}) $ 对 $ x \in M^c $ 成立。因此 $ (f_k) $ 在 $ F $ 中 $ \lambda_m $-几乎处处 收敛于 $ \psi_A $，于是可见 $ \psi_A $ 属于 $ \mathcal{L}_0(\mathbb{R}^m,F) $。

\item

假设 $ A $ 是 $ G_\delta $ 集。推论 IX.5.5 的证明表明，存在开集序列 $ (O_j) $，使得 $ \lambda_{m+n}(O_j) < \infty $ 且 $ A = \bigcap O_j $。令

\[
f_k := \psi_{\bigcap_{j=0}^{k}O_j} , \quad R_k := \bigcap_{j=0}^{k}O_j\setminus A \quad \text{对 } k \in \mathbb{N} .
\]

则 $ (f_k) $ 由 (iii) 是 $ \mathcal{L}_0(\mathbb{R}^m,F) $ 中的一个序列，而 $ (R_k) $ 是递减序列，满足 $ \bigcap_{k=0}^{\infty} R_k = \varnothing $ 且 $ \lambda_{m+n}(R_0) < \infty $。又我们有

\[
\|f_k(x) - \psi_A(x)\|_F^q = \int_{\mathbb{R}^n}\left|\chi_{(\bigcap_{j=0}^{k}O_j)_{[x]}}(y) - \chi_{A_{[x]}}(y)\right|^q dy = \lambda_n((R_k)_{[x]})
\]

对 $ \lambda_m $-几乎每个 $ x \in \mathbb{R}^m $ 成立。因此 $ \lambda_n $ 的从上连续性蕴含 $ (f_k) $ 在 $ \lambda_m $-几乎处处 意义下收敛于 $ \psi_A $。由定理 1.14 即得 $ \psi_A $ 属于 $ \mathcal{L}_0(\mathbb{R}^m,F) $。

\item

最后，考虑 $ A \in \mathcal{L}(m+n) $ 使 $ \lambda_{m+n}(A) $ 有限。由推论 IX.5.5，存在 $ G_\delta $ 集 $ G $ 包含 $ A $ 且与 $ A $ 具有相同的测度。由命题 IX.2.3(ii)，$ N := G\setminus A $ 是 $ \lambda_{m+n} $-零集，且 $ \psi_A = \psi_G - \psi_N $ 在 $ \lambda_m $-几乎处处 意义下成立。现在结论由 (i) 与 (iv) 得到。\nobreak\hspace{0.35em}■

\end{subitems}

\end{lemma}

\begin{corollary}

设 $ p,q \in [1,\infty) $，并设 $ \varphi \in \mathcal{S}(\mathbb{R}^{m+n},E) $ 具有紧支集。则 $ L_q(\mathbb{R}^n,E) $ 值函数 $ x \mapsto [\varphi(x,\cdot)] $ 是 $ \lambda_m $-几乎处处 有定义的，且是 $ L_p $-可积的，即

\[
\int_{\mathbb{R}^m}\|\varphi(x,\cdot)\|_{L_q(\mathbb{R}^n,E)}^p dx < \infty .
\]

若 $ p = q $，则对每个 $ \varphi \in \mathcal{S}(\mathbb{R}^{m+n},E) $ 此结论都成立。

\proofhead
由闵可夫斯基不等式，只需对 $ \varphi := e\chi_A $ 证明本结论，其中 $ e \in E $、$ A \in \mathcal{L}(m+n) $，并且当 $ p = q $ 时 $ A $ 具有有限测度，当 $ p \ne q $ 时 $ A $ 有界。

由引理 6.14，存在 $ \lambda_m $-零集 $ M $，使得函数
\[
M^c \to L_q(\mathbb{R}^n) , \quad x \mapsto \chi_{A_{[x]}}
\]

是 $ \lambda_m $-可测的。因为 $ \varphi(x,\cdot) = e\chi_{A_{[x]}} $，故 $ (x \mapsto \varphi(x,\cdot)) \in \mathcal{L}_0(M^c, L_q(\mathbb{R}^n,E)) $。由

\[
\|\varphi(x,\cdot)\|_{L_q(\mathbb{R}^n,E)} = \left(\int_{\mathbb{R}^n}|e|^q\chi_{A_{[x]}}(y)\,dy\right)^{1/q} = |e|\,[\lambda_n(A_{[x]})]^{1/q} \quad \text{对 } x \in M^c ,
\]

我们得到

\[
\int_{\mathbb{R}^m}\|\varphi(x,\cdot)\|_{L_q(\mathbb{R}^n,E)}^p dx = |e|^p\int_{\mathbb{R}^m}\lambda_n(A_{[x]})^{p/q} dx .
\]

在 $ p = q $ 的情形，命题 6.2 蕴含

\[
\int_{\mathbb{R}^n}\lambda_n(A_{[x]})\,dx = \lambda_{m+n}(A) < \infty .
\]

因此假设 $ p \ne q $。因为 $ \varphi $ 具有紧支集，存在紧子集 $ K \subset \mathbb{R}^m $ 与 $ L \subset \mathbb{R}^n $，使得 $ A \subset K \times L $。于是 $ A_{[x]} \subset L $，这蕴含 $ \lambda_n(A_{[x]}) \le \lambda_n(L) $ 对 $ \lambda_m $-几乎每个 $ x \in \mathbb{R}^m $ 成立。由此我们推出

\[
\int_{\mathbb{R}^m}\lambda_n(A_{[x]})^{p/q} dx = \int_K \lambda_n(A_{[x]})^{p/q} dx \le \lambda_n(L)^{p/q}\lambda_m(K) < \infty .\quad\blacksquare
\]
\end{corollary}

这些准备工作比我们当前目的所需更为一般，但对进一步的应用会有用。我们现在准备证明 $ E $ 值情形的富比尼定理。

\begin{theorem}

（富比尼）对 $ f \in \mathcal{L}_1(\mathbb{R}^{m+n},E) $，

\begin{subitems}
\item

$ f(x,\cdot) \in \mathcal{L}_1(\mathbb{R}^n,E) $ 对 $ \lambda_m $-几乎所有 $ x \in \mathbb{R}^m $ 成立；$ f(\cdot,y) \in \mathcal{L}_1(\mathbb{R}^m,E) $ 对 $ \lambda_n $-几乎所有 $ y \in \mathbb{R}^n $ 成立；

\item

$ x \mapsto \int_{\mathbb{R}^n} f(x,y)\,dy $ 是 $ \lambda_m $-可积的；$ y \mapsto \int_{\mathbb{R}^m} f(x,y)\,dx $ 是 $ \lambda_n $-可积的；

\item

$ \int_{\mathbb{R}^{m+n}} f\,d(x,y) = \int_{\mathbb{R}^m}\left[\int_{\mathbb{R}^n} f(x,y)\,dy\right]dx = \int_{\mathbb{R}^n}\left[\int_{\mathbb{R}^m} f(x,y)\,dx\right]dy $。

\end{subitems}

证明

\begin{items}
\item

设 $ f \in \mathcal{L}_1(\mathbb{R}^{m+n},E) $。则存在 $ \mathcal{S}(\mathbb{R}^{m+n},E) $ 中的 $ L_1 $-柯西序列 $ (f_j) $ 与一个 $ \lambda_{m+n} $-零集 $ L $，使得 $ f_j(x,y) \to f(x,y) $ 对 $ (x,y) \in L^c $ 成立。由注 6.1(h)，存在测度为零的集合 $ M_1 $，使得

\begin{equation}
f_j(x,\cdot) \to f(x,\cdot) \quad \lambda_n\text{-a.e.} ,\tag{6.13}
\end{equation}

对 $ x \in M_1^c $ 成立。我们令 $ F := L_1(\mathbb{R}^n,E) $，并把 $ x \mapsto f_j(x,\cdot) $ 的平凡延拓记为 $ \varphi_j $。根据推论 6.15，$ (\varphi_j) $ 是 $ \mathcal{L}_1(\mathbb{R}^m,F) $ 中的一个序列，满足

\[
\|\varphi_j - \varphi_k\|_1 = \int_{\mathbb{R}^m}\|\varphi_j(x) - \varphi_k(x)\|_F\,dx = \int_{\mathbb{R}^m}\int_{\mathbb{R}^n}|f_j(x,y) - f_k(x,y)|\,dy\,dx .
\]

进一步，引理 6.4 表明

\[
\int_{\mathbb{R}^m}\int_{\mathbb{R}^n}|f_j(x,y) - f_k(x,y)|\,dy\,dx = \int_{\mathbb{R}^{m+n}}|f_j - f_k|\,d(x,y) = \|f_j - f_k\|_1 ,
\]

于是可见 $ (\varphi_j) $ 是 $ \mathcal{L}_1(\mathbb{R}^m,F) $ 中的柯西序列。由定理 2.10 与 2.18，存在 $ \widehat{g} \in \mathcal{L}_1(\mathbb{R}^m,F) $、$ \lambda_m $-零集 $ M_2 $ 以及 $ (\varphi_j) $ 的一个子序列（为简单起见，我们仍记为 $ (\varphi_j) $），使得

\begin{equation}
\lim_{j\to\infty}\varphi_j(x) = \widehat{g}(x) \quad \text{对 } x \in M_2^c\tag{6.14}
\end{equation}

在 $ F $ 中成立，且 $ \varphi_j \to \widehat{g} $ 在 $ \mathcal{L}_1(\mathbb{R}^m,F) $ 中成立。对 $ x \in M_2^c $，取 $ \widehat{g}(x) $ 的代表元 $ g(x) \in \mathcal{L}_1(\mathbb{R}^n,E) $。则存在勒贝格测度为零的集合 $ N(x) $ 以及 $ (\varphi_j(x)) $ 的一个子序列（我们也写作 $ (\varphi_j(x)) $），使得在 $ E $ 中

\[
\lim_{j\to\infty} f_j(x,y) = \lim_{j\to\infty}\varphi_j(x)(y) = g(x)(y) \quad \text{对 } x \in M_2^c \text{ 与 } y \in (N(x))^c .
\]

因此 (6.13) 蕴含：对每个 $ x \in M_1^c \cap M_2^c $，映射 $ f(x,\cdot), g(x) : \mathbb{R}^n \to E $ 在 $ \lambda_n $-几乎处处 意义下相等。引理 2.15 现在表明 $ f(x,\cdot) $ 属于 $ \mathcal{L}_1(\mathbb{R}^n,E) $，且

\begin{equation}
\int_{\mathbb{R}^n} g(x)(y)\,dy = \int_{\mathbb{R}^n} f(x,y)\,dy \quad \text{对 } x \in M_1^c \cap M_2^c .\tag{6.15}
\end{equation}

此外，由 (6.13)、(6.14) 与定理 2.18(ii) 可得

\begin{equation}
\int_{\mathbb{R}^n} f_j(x,y)\,dy = \int_{\mathbb{R}^n}\varphi_j(x)(y)\,dy \to \int_{\mathbb{R}^n} g(x)(y)\,dy = \int_{\mathbb{R}^n} f(x,y)\,dy\tag{6.16}
\end{equation}

对 $ x \in M_1^c \cap M_2^c $ 成立。

\item

对 $ \varphi \in F = L_1(\mathbb{R}^n,E) $，令 $ A\varphi := \int_{\mathbb{R}^n}\varphi\,dy $。由定理 2.11(i)，$ A $ 属于 $ \mathcal{L}(F,E) $。定理 2.11(iii) 蕴含 $ g_j := A\varphi_j $ 定义了 $ \mathcal{L}_1(\mathbb{R}^m,E) $ 中的一个序列。
因为

\begin{equation}
g_j(x) = \int_{\mathbb{R}^n}\varphi_j(x)(y)\,dy = \int_{\mathbb{R}^n} f_j(x,y)\,dy ,\tag{6.17}
\end{equation}

由定理 2.11(i) 我们知道

\[
|g_j(x) - g_k(x)| = \left|\int_{\mathbb{R}^n}\left(f_j(x,y) - f_k(x,y)\right)dy\right| \le \int_{\mathbb{R}^n}|f_j(x,y) - f_k(x,y)|\,dy .
\]

因此定理 2.11(ii) 给出

\[
\int_{\mathbb{R}^m}|g_j - g_k|\,dx \le \int_{\mathbb{R}^m}\int_{\mathbb{R}^n}|(f_j(x,y) - f_k(x,y))|\,dy\,dx = \|f_j - f_k\|_1 ,
\]

其中最后一个等式由托内利定理得到。因此 $ (g_j) $ 是 $ \mathcal{L}_1(\mathbb{R}^m,E) $ 中的柯西序列，由完备性存在 $ h \in \mathcal{L}_1(\mathbb{R}^m,E) $，使得 $ g_j \to h $ 在 $ \mathcal{L}_1(\mathbb{R}^m,E) $ 中成立。于是我们可以找到 $ \lambda_m $-零集 $ M_3 $ 以及 $ (g_j) $ 的一个子序列（我们也记为 $ (g_j) $），使得 $ g_j(x) \to h(x) $ 对 $ x \in M_3^c $ 且 $ j \to \infty $ 成立。鉴于 (6.17)，由 (6.16) 可得

\begin{equation}
h(x) = \int_{\mathbb{R}^n} f(x,y)\,dy \quad \text{对 } x \in M_1^c \cap M_2^c \cap M_3^c ,\tag{6.18}
\end{equation}

这就证明了 (ii) 的第一个结论。

\item

由于 $ g_j \to h $ 在 $ \mathcal{L}_1(\mathbb{R}^m,E) $ 中成立，并且由 (6.17) 与 (6.18)，定理 2.18(ii) 蕴含

\[
\int_{\mathbb{R}^m}\int_{\mathbb{R}^n} f_j(x,y)\,dy\,dx \to \int_{\mathbb{R}^m}\int_{\mathbb{R}^n} f(x,y)\,dy\,dx .
\]

最后，由引理 6.4 可得

\[
\int_{\mathbb{R}^m}\int_{\mathbb{R}^n} f_j(x,y)\,dy\,dx = \int_{\mathbb{R}^{m+n}} f_j\,d(x,y) ,
\]

而由 $ \int_{\mathbb{R}^{m+n}} f\,d(x,y) = \lim_j \int_{\mathbb{R}^{m+n}} f_j\,d(x,y) $，我们有

\[
\int_{\mathbb{R}^m}\int_{\mathbb{R}^n} f(x,y)\,dy\,dx = \int_{\mathbb{R}^{m+n}} f\,d(x,y) .
\]

我们已经证明了 (i) 与 (ii) 各自的第一部分，以及 (iii) 中的第一个等式。其余的结论通过交换 $ x $ 与 $ y $ 的角色得到。\nobreak\hspace{0.35em}■

\end{items}

\end{theorem}

\begin{remark}

富比尼–托内利定理与推论 6.10 的类似结论显然在 $ E $ 值情形也成立。\nobreak\hspace{0.35em}■

\end{remark}

\needspace{3\baselineskip}\subsection*{积分的闵可夫斯基不等式}

作为应用，我们现在证明闵可夫斯基不等式的一个连续版本。

固定 $ p,q \in [1,\infty) $。对 $ f \in \mathcal{L}_0(\mathbb{R}^{m+n},E) $，定理 1.7(i) 表明 $ |f|^q $ 属于 $ \mathcal{L}_0(\mathbb{R}^{m+n},\mathbb{R}^+) $。因此托内利定理蕴含：对 $ \lambda_m $-几乎所有 $ x \in \mathbb{R}^m $，$ |f(x,\cdot)|^q $ 落在 $ \mathcal{L}_0(\mathbb{R}^n,\mathbb{R}^+) $ 中，并且函数

\[
x \mapsto \int_{\mathbb{R}^n}|f(x,y)|^q dy
\]

在 $ \lambda_m $-几乎处处 意义下有定义、取值于 $ \overline{\mathbb{R}}^+ $，且是 $ \lambda_m $-可测的。因此

\[
\|f\|_{(p,q)} := \left(\int_{\mathbb{R}^m}\left[\int_{\mathbb{R}^n}|f(x,y)|^q dy\right]^{p/q} dx\right)^{1/p}
\]

在 $ \overline{\mathbb{R}}^+ $ 中有定义。我们容易验证

\[
\mathcal{L}_{(p,q)}(\mathbb{R}^{m+n},E) := \left\{ f \in \mathcal{L}_0(\mathbb{R}^{m+n},E) ; \|f\|_{(p,q)} < \infty \right\}
\]

是 $ \mathcal{L}_{0}(\mathbb{R}^{m+n},E) $ 的向量子空间，且 $ \|\cdot\|_{(p,q)} $ 在 $ \mathcal{L}_{(p,q)}(\mathbb{R}^{m+n},E) $ 上定义一个半范数。最后，我们置

\[
\mathcal{S}_{c}(\mathbb{R}^{m+n},E) := \left\{\, f \in \mathcal{S}(\mathbb{R}^{m+n},E) ;\ \operatorname{supp}(f) \text{ 是紧的} \,\right\} .
\]

\begin{lemma}

$ \mathcal{S}_{c}(\mathbb{R}^{m+n},E) $ 是 $ \mathcal{L}_{(p,q)}(\mathbb{R}^{m+n},E) $ 的稠密向量子空间。

证明

\begin{subitems}
\item

取 $ f \in \mathcal{L}_{(p,q)}(\mathbb{R}^{m+n},E) $，并设 $ (g_{k}) $ 是 $ \mathcal{S}(\mathbb{R}^{m+n},E) $ 中的序列，使得 $ g_{k} \to f $ 几乎处处。置 $ A_{k} := [|g_{k}| \le 2|f|] \cap k\mathbb{B}^{m+n} $ 以及 $ f_{k} := \chi_{A_{k}}g_{k} $。则 $ (f_{k}) $ 是 $ \mathcal{S}_{c}(\mathbb{R}^{m+n},E) $ 中的序列，且存在 $ \lambda_{m+n} $-零集 $ L $，使得

\begin{equation}
f_{k}(x,y) \to f(x,y) \quad \text{对 } (x,y) \in L^{c} .\tag{6.19}
\end{equation}

此外

\begin{equation}
|f_{k}-f| \le |f_{k}|+|f| \le 3|f| \quad \text{对 } k \in \mathbb{N} .\tag{6.20}
\end{equation}
\item

由 (6.20)，从托内利定理与定理 3.9 可得，存在 $ \lambda_{m} $-零集 $ M_{0} $，使得

\begin{equation}
|f(x,\cdot)-f_{k}(x,\cdot)|^{q},\ |f(x,\cdot)|^{q} \in \mathcal{L}_{1}(\mathbb{R}^{n}) \quad \text{对 } x \in M_{0}^{c} \text{ 与 } k \in \mathbb{N} .\tag{6.21}
\end{equation}

注 6.1(h) 说明，存在 $ \lambda_{m} $-零集 $ M_{1} $，使得对每个 $ x \in M_{1}^{c} $，$ L_{[x]} $ 都是 $ \lambda_{n} $-零集。置 $ M := M_{0} \cup M_{1} $ 并取 $ x \in M^{c} $。由 (6.19) 可读出，对 $ y \in (L_{[x]})^{c} $ 有 $ f_{k}(x,y) \to f(x,y) $。由 (6.20) 与 (6.21)，我们可以对序列 $ \left(|f(x,\cdot)-f_{k}(x,\cdot)|^{p}\right)_{k \in \mathbb{N}} $ 应用勒贝格控制收敛定理，并得到

\[
\lim_{k\to\infty}\int_{\mathbb{R}^{n}}|f(x,y)-f_{k}(x,y)|^{q}\,dy = 0 \quad \text{对 } x \in M^{c} .
\]

现在定义

\[
\varphi_{k} := \left(x \mapsto \left(\int_{\mathbb{R}^{n}}|f(x,y)-f_{k}(x,y)|^{q}\,dy\right)^{p/q}\right)^{\sim} \quad \text{对 } k \in \mathbb{N} .
\]

则序列 $ (\varphi_{k}) $ 关于 $ \lambda_{m} $ 几乎处处收敛到 0。

\item

最后，置

\[
\varphi := \left(x \mapsto 3^{p}\left(\int_{\mathbb{R}^{n}}|f(x,y)|^{q}\,dy\right)^{p/q}\right)^{\sim} .
\]

因为 $ f \in \mathcal{L}_{(p,q)}(\mathbb{R}^{m+n},E) $，我们知道 $ \varphi $ 属于 $ \mathcal{L}_{1}(\mathbb{R}^{m}) $，且 (6.20) 蕴含：对 $ k \in \mathbb{N} $，$ 0 \le \varphi_{k} \le \varphi $ 关于 $ \lambda_{m} $ 几乎处处成立。因此我们可以对 $ (\varphi_{k}) $ 应用控制收敛定理，看出 $ (\int_{\mathbb{R}^{m}}\varphi_{k})_{k \in \mathbb{N}} $ 是 $ \mathbb{R}^{+} $ 中的零序列。断言现在成立，因为

\begin{equation*}
\begin{aligned}
\int_{\mathbb{R}^{m}}\varphi_{k} = \int_{\mathbb{R}^{m}}\left[\int_{\mathbb{R}^{n}}|f(x,y)-f_{k}(x,y)|^{q}\,dy\right]^{p/q}dx = \|f-f_{k}\|^{p}_{(p,q)} . \\quad\blacksquare
\end{aligned}
\end{equation*}
\end{subitems}

\end{lemma}

容易验证 $ \mathcal{N} := \left\{\, f \in \mathcal{L}_{0}(\mathbb{R}^{m+n},E) ;\ f = 0 \text{ 几乎处处} \,\right\} $ 是 $ \mathcal{L}_{(p,q)}(\mathbb{R}^{m+n},E) $ 的向量子空间，且 $ f $ 属于 $ \mathcal{N} $ 当且仅当 $ \|f\|_{(p,q)} = 0 $。因此

\[
L_{(p,q)}(\mathbb{R}^{m+n},E) := \mathcal{L}_{(p,q)}(\mathbb{R}^{m+n},E)/\mathcal{N}
\]

是良定义的向量空间，且对应 $ [f] \mapsto \|f\|_{(p,q)} $ 在 $ L_{(p,q)}(\mathbb{R}^{m+n},E) $ 上定义一个范数，我们仍将其记为 $ \|\cdot\|_{(p,q)} $。在下文中，我们总是给空间 $ L_{(p,q)}(\mathbb{R}^{m+n},E) $ 赋予由 $ \|\cdot\|_{(p,q)} $ 诱导的拓扑。

我们置

\[
S_{c}(\mathbb{R}^{m+n},E) := \left\{\, [f] \in L_{0}(\mathbb{R}^{m+n},E) ;\ [f] \cap \mathcal{S}_{c}(\mathbb{R}^{m+n},E) \ne \varnothing \,\right\} .
\]

\begin{remark}

\begin{items}
\item

$ S_{c}(\mathbb{R}^{m+n},E) $ 是 $ L_{(p,q)}(\mathbb{R}^{m+n},E) $ 的稠密向量子空间。

\proofhead
这由引理 6.18 推出。 \nobreak\hspace{0.35em}■
\item

设 $ f \in \mathcal{L}_{0}(\mathbb{R}^{m+n},E) $。若对几乎每个 $ x \in \mathbb{R}^{m} $，$ f(x,\cdot) $ 属于 $ \mathcal{L}_{q}(\mathbb{R}^{n},E) $，且

\[
\left[x \mapsto \left(\int_{\mathbb{R}^{n}}|f(x,y)|^{q}dx\right)^{1/q}\right]^{\sim} \in \mathcal{L}_{p}(\mathbb{R}^{n}) ,
\]

则 $ [f] $ 属于 $ L_{(p,q)}(\mathbb{R}^{m+n},E) $。

\item

$ L_{(p,p)}(\mathbb{R}^{m+n},E) = L_{p}(\mathbb{R}^{m+n},E) $。

\proofhead
这由注 4.9(b) 与富比尼–托内利定理推出。 \nobreak\hspace{0.35em}■
\item

$ S_{c}(\mathbb{R}^{n},E) $ 是 $ L_{p}(\mathbb{R}^{n},E) $ 的稠密向量子空间。

\proofhead
这是 (a) 与 (c) 的推论。 \nobreak\hspace{0.35em}■
\end{items}

设 $ g \in \mathcal{S}_{c}(\mathbb{R}^{m+n},E) $。由推论 6.15，$ T_{0}g := (x \mapsto [g(x,\cdot)])^{\sim} $ 属于 $ \mathcal{L}_{p}\left(\mathbb{R}^{m},L_{q}(\mathbb{R}^{n},E)\right) $。将 $ T_{0}g $ 关于 $ \mathcal{L}_{0}(\mathbb{R}^{m},L_{q}(\mathbb{R}^{n},E)) $ 中所有 $ \lambda_{m} $-几乎处处为零的元素构成的向量子空间的等价类记为 $ [T_{0}g] $，则 $ [T_{0}g] \in L_{p}\left(\mathbb{R}^{m},L_{q}(\mathbb{R}^{n},E)\right) $。此外，由推论 6.8 可得，若 $ g,h \in \mathcal{S}_{c}(\mathbb{R}^{m+n},E) $ 在 $ \lambda_{m+n} $-几乎处处相等，则 $ [T_{0}g] = [T_{0}h] $。于是

\[
T:\ S_{c}(\mathbb{R}^{m+n},E) \to L_{p}\left(\mathbb{R}^{m},L_{q}(\mathbb{R}^{n},E)\right) ,\quad [g] \mapsto [T_{0}g]
\]

是良定义的线性映射。

\end{remark}

\begin{lemma}

存在 $ T $ 的唯一延拓

\[
\overline{T} \in \mathcal{L}\left(L_{(p,q)}(\mathbb{R}^{m+n},E),L_{p}\left(\mathbb{R}^{m},L_{q}(\mathbb{R}^{n},E)\right)\right)
\]

且 $ \overline{T} $ 是像稠密的等距。

证明

\begin{subitems}
\item

对 $ f \in S_{c}(\mathbb{R}^{m+n},E) $，取 $ g \in f \cap \mathcal{S}_{c}(\mathbb{R}^{m+n},E) $。则

\[
\int_{\mathbb{R}^{m}}\|Tf\|^{p}_{L_{q}(\mathbb{R}^{n},E)}\,dx = \int_{\mathbb{R}^{m}}\left(\int_{\mathbb{R}^{n}}|g(x,y)|^{q}\,dy\right)^{p/q}dx = \|g\|^{p}_{(p,q)} = \|f\|^{p}_{(p,q)} .
\]

因此 $ T \in \mathcal{L}\left(S_{c}(\mathbb{R}^{m+n},E),L_{p}\left(\mathbb{R}^{m},L_{q}(\mathbb{R}^{n},E)\right)\right) $ 是等距。现在由定理 VI.2.6 与注 6.19(a) 可得，存在唯一确定的 $ T $ 的等距延拓 $ \overline{T} $。

\item

置 $ F := L_{q}(\mathbb{R}^{n},E) $，并取 $ w \in \mathcal{L}_{p}(\mathbb{R}^{m},F) $ 与 $ \varepsilon > 0 $。由注 6.19(d) 可得，存在 $ \varphi \in \mathcal{S}_{c}(\mathbb{R}^{m},F) $，使得 $ \|w-\varphi\|_{p} < \varepsilon/2 $。设 $ \sum_{j=0}^{r}\chi_{A_{j}}\hat{f}_{j} $ 是 $ \varphi $ 的标准形。则 $ \bigcup_{j=0}^{r}A_{j} $ 在 $ \mathbb{R}^{m} $ 中有界，且 $ \alpha := \sum_{j=0}^{r}\lambda_{m}(A_{j}) $ 有限。在 $ \alpha = 0 $ 的情形，我们有

\[
\|w\|_{p} = \|w-T0\|_{p} < \varepsilon/2 .
\]

在 $ \alpha > 0 $ 的情形，对每个 $ j \in \{0,\dots,r\} $，我们选取 $ \hat{f}_{j} $ 的一个代表元 $ f_{j} $ 以及 $ \psi_{j} \in \mathcal{S}_{c}(\mathbb{R}^{n},E) $，使得

\[
\|\psi_{j}-f_{j}\|_{q} < \alpha^{-1/p}(r+1)^{-1/q'}\varepsilon .
\]

此外，置

\[
h(x,y) := \sum_{j=0}^{r}\chi_{A_{j}}(x)\psi_{j}(y) \quad \text{对 } (x,y) \in \mathbb{R}^{m+n} .
\]

若对 $ j \in \{0,\dots,r\} $ 有 $ \psi_{j} = \sum_{k_{j}=0}^{s_{j}}\chi_{B_{k_{j}}}e_{k_{j}} $，则我们有

\[
h = \sum_{j=0}^{r}\sum_{k_{j}=0}^{s_{j}}\chi_{A_{j}}\chi_{B_{k_{j}}}e_{k_{j}} = \sum_{j=0}^{r}\sum_{k_{j}=0}^{s_{j}}\chi_{A_{j}\times B_{k_{j}}}e_{k_{j}} ,
\]

并且看出 $ h $ 属于 $ \mathcal{S}_{c}(\mathbb{R}^{m+n},E) $。最后，设 $ g $ 是 $ h $ 在 $ L_{0}(\mathbb{R}^{m+n},E) $ 中的等价类。则 $ g $ 属于 $ S_{c}(\mathbb{R}^{m+n},E) $，且 $ Tg = \sum_{j=0}^{r}[\chi_{A_{j}}\psi_{j}] $。由（和的）赫尔德不等式与等式 $ \chi_{A}^{2} = \chi_{A} $ 可得

\begin{equation*}
\begin{aligned}
\int_{\mathbb{R}^{m}}\|Tg-\varphi\|^{p}_{F} &= \int_{\mathbb{R}^{m}}\left[\int_{\mathbb{R}^{n}}\left|\sum_{j=0}^{r}\chi_{A_{j}}(x)(\psi_{j}(y)-f_{j}(y))\right|^{q}dy\right]^{p/q}dx \\ &\le (r+1)^{p/q'}\int_{\mathbb{R}^{m}}\left[\int_{\mathbb{R}^{n}}\sum_{j=0}^{r}\chi_{A_{j}}(x)|\psi_{j}(y)-f_{j}(y)|^{q}dy\right]^{p/q}dx \\ &= (r+1)^{p/q'}\int_{\mathbb{R}^{m}}\left[\sum_{j=0}^{r}\chi_{A_{j}}(x)\|\psi_{j}-f_{j}\|^{q}_{F}\right]^{p/q}dx \\ &\le (r+1)^{p/q'}\sum_{j=0}^{r}\lambda_{m}(A_{j})\|\psi_{j}-f_{j}\|^{p}_{F} .
\end{aligned}
\end{equation*}

因此，

\[
\|Tg-\varphi\|_{p} \le \alpha^{1/p}(r+1)^{1/q'}\max_{j}\|\psi_{j}-f_{j}\|_{F} < \varepsilon/2 ,
\]

从而 $ \|Tg-w\|_{p} < \varepsilon $。因为这对 $ w $ 与 $ \varepsilon $ 的每个选择都成立，所以我们看出 $ T $ 的像，因而 $ \overline{T} $ 的像，是稠密的。 \nobreak\hspace{0.35em}■

\end{subitems}

\end{lemma}

按惯例，我们通过把 $ \overline{T} $ 写作 $ T $ 来简化记号。此外，正如第 4 节所述，我们对勒贝格空间中元素的记号不区分陪集与其代表元。这意味着对 $ f \in L_{(p,q)}(\mathbb{R}^{m+n},E) $，我们可以把 $ Tf(x) $ 写作 $ f(x,\cdot) $。在这些约定下，引理 6.20 说明

\begin{equation}
T:\ L_{(p,q)}(\mathbb{R}^{m+n},E) \to L_{p}\left(\mathbb{R}^{m},L_{q}(\mathbb{R}^{n},E)\right) ,\quad f \mapsto (x \mapsto f(x,\cdot))\tag{6.22}
\end{equation}

是像稠密的线性等距。

现在容易证明我们的连续型闵可夫斯基不等式。

\begin{proposition}

（积分的闵可夫斯基不等式） 对 $ 1 \le q < \infty $，有：

\begin{subitems}
\item

对 $ f \in \mathcal{L}_{0}(\mathbb{R}^{m+n},E) $，有

\[
\left(\int_{\mathbb{R}^{n}}\left[\int_{\mathbb{R}^{m}}|f(x,y)|\,dx\right]^{q}dy\right)^{1/q} \le \int_{\mathbb{R}^{m}}\left[\int_{\mathbb{R}^{n}}|f(x,y)|^{q}dy\right]^{1/q}dx
\]
\item

对 $ f \in \mathcal{L}_{(1,q)}(\mathbb{R}^{m+n},E) $，有

\[
\left(\int_{\mathbb{R}^{n}}\left|\int_{\mathbb{R}^{m}}f(x,y)\,dx\right|^{q}dy\right)^{1/q} \le \int_{\mathbb{R}^{m}}\left[\int_{\mathbb{R}^{n}}|f(x,y)|^{q}dy\right]^{1/q}dx < \infty
\]

\proofhead
在 (i) 的情形，我们可以不失一般性假设
\[
\int_{\mathbb{R}^{m}}\left[\int_{\mathbb{R}^{n}}|f(x,y)|^{q}dy\right]^{1/q}dx < \infty .
\]

则 $ |f| $ 属于 $ \mathcal{L}_{(1,q)}(\mathbb{R}^{m+n},\mathbb{R}) $，且该断言是 (ii) 的特例，其中 $ f $ 换成 $ |f| $、$ E $ 换成 $ \mathbb{R} $。因此假设 $ f \in \mathcal{L}_{(1,q)}(\mathbb{R}^{m+n},E) $。由引理 6.20 与定理 2.11(i)（其中 $ E $ 换成 $ L_{q}(\mathbb{R}^{n},E) $）可得

\[
\int_{\mathbb{R}^{m}}Tf\,dx = \int_{\mathbb{R}^{m}}f(x,\cdot)\,dx \in L_{q}(\mathbb{R}^{n},E)
\]

以及

\begin{equation*}
\begin{aligned}
\left(\int_{\mathbb{R}^{n}}\left|\int_{\mathbb{R}^{m}}f(x,y)\,dx\right|^{q}dy\right)^{1/q} &= \left\|\int_{\mathbb{R}^{m}}Tf\,dx\right\|_{L_{q}(\mathbb{R}^{n},E)} \le \int_{\mathbb{R}^{m}}\|Tf\|_{L_{q}(\mathbb{R}^{n},E)}\,dx \\ &= \int_{\mathbb{R}^{m}}\left(\int_{\mathbb{R}^{n}}|f(x,y)|^{q}dy\right)^{1/q}dx .
\end{aligned}
\end{equation*}

这就完成了证明。 \nobreak\hspace{0.35em}■

\end{subitems}

\end{proposition}

\needspace{3\baselineskip}\subsection*{$ L_p(\mathbb{R}^{m+n},E) $ 的一个刻画}

作为引理 6.20 的另一个推论，我们得到富比尼定理的一个常常有用的推广与精细化。

\begin{theorem}

对 $ 1 \le p < \infty $，

\[
L_{p}(\mathbb{R}^{m+n},E) \to L_{p}\left(\mathbb{R}^{m},L_{p}(\mathbb{R}^{n},E)\right) ,\quad f \mapsto (x \mapsto f(x,\cdot))
\]

是等距同构。

\proofhead
设 $ v \in L_{p}\left(\mathbb{R}^{m},L_{p}(\mathbb{R}^{n},E)\right) $。由引理 6.20，存在 $ L_{p}(\mathbb{R}^{m+n},E) $ 中的序列 $ (f_{j}) $，使得在 $ L_{p}\left(\mathbb{R}^{m},L_{p}(\mathbb{R}^{n},E)\right) $ 中 $ \lim_{j}Tf_{j} = v $。因为 $ T $ 是线性等距，容易推出 $ (f_{j}) $ 是 $ L_{p}(\mathbb{R}^{m+n},E) $ 中的柯西序列。把它的极限记为 $ f \in L_{p}(\mathbb{R}^{m+n},E) $，我们有 $ Tf = v $。因此 $ T $ 是满射。这证明了断言。 \nobreak\hspace{0.35em}■
\end{theorem}

借助这个等距同构，我们可以把巴拿赫空间 $ L_{p}(\mathbb{R}^{m+n},E) $ 与 $ L_{p}\left(\mathbb{R}^{m},L_{p}(\mathbb{R}^{n},E)\right) $ 等同起来：

\[
L_{p}(\mathbb{R}^{m+n},E) = L_{p}\left(\mathbb{R}^{m},L_{p}(\mathbb{R}^{n},E)\right) .
\]

\begin{remark}

\begin{items}
\item

定理 6.22 的陈述对 $ p = \infty $ 不成立，即

\[
L_{\infty}(\mathbb{R}^{m+n},E) \ne L_{\infty}\left(\mathbb{R}^{m},L_{\infty}(\mathbb{R}^{n},E)\right) .
\]

\proofhead
取 $ A := \left\{\, (x,y) \in \mathbb{R}^{2} ;\ 0 \le y \le x \le 1 \,\right\} $ 以及 $ f := \chi_{A} $。因为 $ A $ 是勒贝格可测的，所以 $ f $ 属于 $ L_{\infty}(\mathbb{R}^{2}) $。若我们置
\begin{equation*}
g(x) := f(x,\cdot) = \begin{cases} \chi_{[0,x]} & \text{若 } 0 \le x \le 1 , \\ 0 & \text{其他情形} , \end{cases}
\end{equation*}

则 $ g(x) $ 属于 $ L_{\infty}(\mathbb{R}) $，且对 $ x \in \mathbb{R} $ 有 $ \|g(x)\|_{\infty} \le 1 $。但 $ g $ 仍然不属于 $ L_{\infty}\left(\mathbb{R},L_{\infty}(\mathbb{R})\right) $，因为映射 $ g: \mathbb{R} \to L_{\infty}(\mathbb{R}) $ 不是 $ \lambda_{1} $-可测的。为看出这一点，由定理 1.4，只需说明 $ g $ 不是 $ \lambda_{1} $-几乎可分值的。为验证这一点，注意

\begin{equation}
\|g(x)-g(r)\|_{L_{\infty}(\mathbb{R})} = 1 \quad \text{对 } r \in \mathbb{R}\setminus\{x\}\tag{6.23}
\end{equation}

对 $ x \in (0,1] $ 成立。假如 $ g $ 是 $ \lambda_{1} $-几乎可分值的，则存在 $ \lambda_{1} $-零集 $ N \subset \mathbb{R} $ 以及 $ \mathbb{R} $ 中的序列 $ (r_{j}) $，使得

\begin{equation}
\inf_{j \in \mathbb{N}}\|g(x)-g(r_{j})\|_{\infty} < 1/2 \quad \text{对 } x \in N^{c} .\tag{6.24}
\end{equation}

因为 $ \lambda_{1}((0,1]\setminus N) = 1 $，所以集合 $ (0,1]\setminus N $ 不可数。于是由 (6.23) 可得 (6.24) 不可能成立，且 $ g $ 不是 $ \lambda_{1} $-可测的。 \nobreak\hspace{0.35em}■

\item

推广定理 6.22，可以证明：对任意 $ p,q \in [1,\infty) $，映射

\[
L_{(p,q)}(\mathbb{R}^{m+n},E) \to L_{p}\left(\mathbb{R}^{m},L_{q}(\mathbb{R}^{n},E)\right) ,\quad f \mapsto (x \mapsto f(x,\cdot))
\]

是等距同构。因此 $ L_{(p,q)}(\mathbb{R}^{m+n},E) $ 是完备的。 \nobreak\hspace{0.35em}■

\end{items}

\end{remark}

\needspace{3\baselineskip}\subsection*{一个迹定理}

由例 IX.5.2 以及勒贝格测度在等距下的不变性可得，$ \mathbb{R}^{n} $ 中的每个超平面 $ \Gamma $ 都是 $ \lambda_{n} $-零集。因此对 $ u \in L_{p}(\mathbb{R}^{n}) $，限制 $ u|_{\Gamma} $，即 $ u $ 在 $ \Gamma $ 上的迹，是无定义的，因为 $ u $ 可以在 $ \Gamma $ 上被“任意改变”。作为富比尼–托内利定理的另一个应用，我们现在说明，对 $ L_{p}(\mathbb{R}^{n}) $ 的某些向量子空间的元素，仍然可以在 $ \Gamma $ 上定义这样的迹。当然，对向量子空间 $ C_{c}^{1}(\mathbb{R}^{n}) $ 来说这是平凡的情形。以下结果的意义在于：赋予这个空间的不是上确界范数，而是 $ L_{p} $ 范数，并把导数也纳入其中。在下一节中，我们将更好地理解 $ L_{p}(\mathbb{R}^{n}) $ 的这些子空间的意义。

考虑坐标超平面 $ \Gamma := \mathbb{R}^{n-1} \times \{0\} $，我们把它与 $ \mathbb{R}^{n-1} $ 等同。对 $ u \in C(\mathbb{R}^{n}) $，令 $ \gamma u := u|_{\Gamma} $ 为 $ u $ 在 $ \Gamma $ 上的\textbf{迹}：

\[
(\gamma u)(x) := u(x,0) \quad \text{对 } x \in \mathbb{R}^{n-1} .
\]

则 $ \gamma: C_{c}^{1}(\mathbb{R}^{n}) \to C_{c}(\mathbb{R}^{n-1}) $，$ u \mapsto \gamma u $ 是良定义的线性映射。

现在取 $ 1 \le p < \infty $，并给 $ C_{c}^{1}(\mathbb{R}^{n}) $ 赋予范数

\[
\|u\|_{1,p} := \left(\|u\|^{p}_{p}+\sum_{j=1}^{n}\|\partial_{j}u\|^{p}_{p}\right)^{1/p} .
\]

此外，置

\[
\widehat{H}^{1}_{p}(\mathbb{R}^{n}) := \left(C^{1}_{c}(\mathbb{R}^{n}),\|\cdot\|_{1,p}\right) .
\]

因为 $ C_{c}(\mathbb{R}^{n-1}) $ 是 $ L_{p}(\mathbb{R}^{n-1}) $ 的向量子空间，

\[
\gamma:\ \widehat{H}^{1}_{p}(\mathbb{R}^{n}) \to L_{p}(\mathbb{R}^{n-1}) ,\quad u \mapsto \gamma u
\]

是良定义的线性映射，即关于 $ \Gamma = \mathbb{R}^{n-1} $ 的\textbf{迹算子}。下面的迹定理说明 $ \gamma $ 是连续的。

\begin{proposition}

对 $ 1 \le p < \infty $，有 $ \gamma \in \mathcal{L}\left(\widehat{H}^{1}_{p}(\mathbb{R}^{n}),L_{p}(\mathbb{R}^{n-1})\right) $。

\proofhead
定义 $ h \in C^{1}(\mathbb{R}) $ 为 $ h(t) := |t|^{p-1}t $。对 $ v \in C_{c}^{1}(\mathbb{R}^{n}) $，由链式法则可得 $ \partial_{n}h(v) = h'(v)\partial_{n}v $。因为 $ v $ 具有紧支集，微积分基本定理蕴含
\[
-h(v(x,0)) = \int_{0}^{\infty}\partial_{n}h(v)(x,y)\,dy = \int_{0}^{\infty}h'(v(x,y))\partial_{n}v(x,y)\,dy \quad \text{对 } x \in \mathbb{R}^{n-1} .
\]

因为 $ h'(t) = p|t|^{p-1} $，我们得到

\begin{equation*}
\begin{aligned}
|v(x,0)|^{p} = |h(v(x,0))| &\le \int_{0}^{\infty}|h'(v(x,y))|\,|\partial_{n}v(x,y)|\,dy \\ &= p\int_{0}^{\infty}|v(x,y)|^{p-1}|\partial_{n}v(x,y)|\,dy .
\end{aligned}
\end{equation*}

又，对 $ \xi,\eta \in [0,\infty) $，Young 不等式给出 $ \xi^{p-1}\eta \le \frac{p-1}{p}\xi^p + \frac{1}{p}\eta^p $，故

\[
|v(x,0)|^p \le (p-1)\int_0^\infty |v(x,y)|^p dy + \int_0^\infty |\partial_n v(x,y)|^p dy .
\]

令 $ c_p := \max\{p-1,1\} $，则由富比尼–托内利定理可得

\begin{equation}
\int_{\mathbb{R}^{n-1}} |v(x,0)|^p dx \le c_p \left( \int_{\mathbb{R}^{n-1}\times\mathbb{R}} |v(x,y)|^p d(x,y) + \int_{\mathbb{R}^{n-1}\times\mathbb{R}} |\partial_n v(x,y)|^p d(x,y) \right) .\tag{6.25}
\end{equation}

因此

\[
\|\gamma v\|_{L_p(\mathbb{R}^{n-1})} \le c \|v\|_{\widehat{H}_p^1(\mathbb{R}^n)} \quad \text{对 } v \in \widehat{H}_p^1(\mathbb{R}^n) ,
\]

其中 $ c := c_p^{1/p} $。这证明了定理。\nobreak\hspace{0.35em}■

\end{proposition}

\begin{remark}

记 $ \mathbb{H}^n $ 为 $ \mathbb{R}^n $ 的上半空间：

\[
\mathbb{H}^n := \mathbb{R}^{n-1}\times(0,\infty) = \{ (x,y) \in \mathbb{R}^{n-1}\times\mathbb{R} ; y > 0 \} .
\]

则 $ \Gamma = \mathbb{R}^{n-1}\times\{0\} = \partial\mathbb{H}^n $。若令

\[
\widehat{H}_p^1(\mathbb{H}^n) := ( \{ u|_{\mathbb{H}^n} ; u \in C_c^1(\mathbb{R}^n) \}, \|\cdot\|_{1,p} ) ,
\]

则 $ \widehat{H}_p^1(\mathbb{H}^n) $ 是 $ L_p(\mathbb{H}^n) $ 的向量子空间，且由一条与 (6.25) 类似的陈述可得

\[
\gamma \in \mathcal{L}(\widehat{H}_p^1(\mathbb{H}^n), L_p(\mathbb{R}^{n-1})) .
\]

此时，对 $ u \in \widehat{H}_p^1(\mathbb{R}^n) $，$ \gamma u $ 是 $ u $ 在\textbf{边界} $ \partial\mathbb{H}^n $ 上的\textbf{迹}。\nobreak\hspace{0.35em}■

\end{remark}

\exerciseshead

\begin{exercises}
\item

设 $ B \in \mathcal{L}(n) $，$ a \in \mathbb{R}^{n+1} $。记

\[
Z_a(B) := \{ (x,0) + ta \in \mathbb{R}^{n+1} ; x \in B,\ t \in [0,1] \}
\]

为以 $ B $ 为底、$ a $ 为棱的柱体，并令

\[
K_a(B) := \{ (1-t)(x,0) + ta \in \mathbb{R}^{n+1} ; x \in B,\ t \in [0,1] \}
\]

为以 $ B $ 为底、$ a $ 为顶点的锥体。证明：

\begin{items}
\item

$ \lambda_{n+1}(Z_a(B)) = |a_{n+1}|\lambda_n(B) $；

\item

$ \lambda_{n+1}(K_a(B)) = |a_{n+1}|\lambda_n(B)/(n+1) $。
若把 $ |a_{n+1}| $ 解释为柱体 $ Z_a(B) $ 或锥体 $ K_a(B) $ 的高，则 (b) 说明 $ n $ 维锥体的体积等于 $ n $ 个具有相同底与高的柱体的体积之和。

\end{items}

\item

设 $ 0 < r < a $，$ V_{a,r} $ 是 $ \mathbb{R}^3 $ 中由 2-环面 $ \mathbb{T}^2_{a,r} $ 围成的区域。证明 $ V_{a,r} = 2\pi^2 ar^2 $。

\item

设 $ J \subset \mathbb{R} $ 是以 $ a := \inf J $ 与 $ b := \sup J $ 为端点的区间。又设 $ f \in \mathcal{L}_0(J, \mathbb{R}^+) $，并记

\[
R_f := \{ (x,t) \in \mathbb{R}^n \times J ; |x| \le f(t) \}
\]

为 $ f $ 的图像绕 $ t $ 轴旋转所得的\textbf{旋转体}。证明

\[
\lambda_{n+1}(R_f) = \omega_n \int_a^b (f(t))^n dt ,
\]

其中 $ \omega_n $ 是 $ \mathbb{B}^n $ 的体积。在 $ n = 2 $ 的情形从几何上解释这个公式。

\item

设 $ K $ 是 $ \mathbb{R}^n $ 中的紧集，且对 $ \rho \in \mathcal{L}_1(K, \mathbb{R}^+) $ 有 $ \rho_K := \int_K \rho(x) dx > 0 $。则

\[
S(K,\rho) := \frac{1}{\rho_K} \int_K x\rho(x) dx \in \mathbb{R}^n
\]

是 $ K $ 关于密度 $ \rho $ 的\textbf{形心}。我们令 $ S(K) := S(K,1) $。现在设 $ J := [a,b] $ 是 $ \mathbb{R} $ 中的紧致完全区间，并令 $ f \in \mathcal{L}_0(J, \mathbb{R}^+) $。再令

\[
A_f := \{ (x,y) \in \mathbb{R}^2 ; 0 \le y \le f(x),\ x \in J \} ,
\]

并记 $ R_f $ 为 $ f $ 生成的 $ \mathbb{R}^3 $ 中的旋转体（绕 $ x $ 轴旋转）。证明：

\begin{items}
\item

对 $ f \in \mathcal{L}_1(J, \mathbb{R}^+) $，

\[
S(A_f) = (S_1(A_f), S_2(A_f)) = \frac{1}{\|f\|_1} \left( \int_a^b xf(x) dx, \frac{1}{2}\int_a^b (f(x))^2 dx \right) .
\]
\item

对 $ f \in \mathcal{L}_2(J, \mathbb{R}^+) $，

\[
S(R_f) = \left( \frac{1}{\|f\|_2^2} \int_a^b t(f(t))^2 dt, 0, 0 \right) .
\]
\item

对 $ f \in \mathcal{L}_1(J, \mathbb{R}^+) $，有\textbf{古尔丁第一法则}

\[
\lambda_3(R_f) = \pi \int_a^b (f(x))^2 dx = 2\pi S_2(A_f)\lambda_2(A_f) .
\]

换言之，旋转体的体积等于一个子午截面的面积\footnote{即与包含旋转轴的平面的交。}乘以该截面的形心在整周旋转中画出的圆的周长\footnote{对并非由某个图像旋转得到的旋转体，古尔丁第一法则同样成立；见习题 XII.1.11。}。

\end{items}

\item

\begin{items}
\item

对 $ \alpha \in [0,\pi/2) $，令 $ a := (\cos\alpha, 0, \sin\alpha) $。确定柱体 $ Z_a(\mathbb{B}_2) $ 与锥体 $ K_a(\mathbb{B}_2) $ 关于密度 $ 1 $ 的形心。

\item

对 $ \lambda > 0 $，令 $ A_\lambda := \{ (x,y) \in \mathbb{R}^2 ; 0 \le y \le e^{-\lambda x},\ x \ge 0 \} $。证明 $ S(A_\lambda) \in A_\lambda $。

\item

给出一个 $ S(A_f) \notin A_f $ 的例子。

\end{items}

\item

设 $ K \subset \mathbb{R}^n $ 是凸紧集。验证对 $ \rho \in \mathcal{L}_1(K, \mathbb{R}^+) $ 有 $ S(K,\rho) \in K $。

\item

记 $ \Delta_n := \{ x \in \mathbb{R}^n ; x_j \ge 0,\ \sum_{j=1}^n x_j \le 1 \} $ 为 $ \mathbb{R}^n $ 中的\textbf{标准单纯形}。证明：

\begin{items}
\item

$ \lambda_n(\Delta_n) = 1/n! $。

\item

$ S(\Delta_n) = (1/(n+1), 1/(n+1), \ldots, 1/(n+1)) $。

\end{items}

\item

给定 $ f \in \mathcal{L}_1(\mathbb{R}^m, \mathbb{K}) $、$ g \in \mathcal{L}_1(\mathbb{R}^n, E) $，定义 $ F(x,y) := f(x)g(y) $，$ (x,y) \in \mathbb{R}^m\times\mathbb{R}^n $。证明 $ F $ 属于 $ \mathcal{L}_1(\mathbb{R}^{m+n}, E) $，且

\[
\int_{\mathbb{R}^{m+n}} F(x,y) d(x,y) = \int_{\mathbb{R}^m} f(x) dx \int_{\mathbb{R}^n} g(y) dy .
\]
\item

对 $ D := \{ (x,y) \in \mathbb{R}^2 ; x,y \ge 0,\ x+y \le 1 \} $，证明

\[
\int_D x^m y^n d(x,y) = \frac{1}{n+1} \mathrm{B}(m+1, n+2) \quad \text{对 } m,n \in \mathbb{N} .
\]
\item

证明 $ \int_{[0,1]\times[0,1]} y/\sqrt{x} d(x,y) = 1 $。

\item

证明对 $ \varphi \in C_c^1(\mathbb{R}^n, E) $ 与 $ j \in \{1,\ldots,n\} $ 有 $ \int_{\mathbb{R}^n} \partial_j\varphi dx = 0 $。

\item

对下列每个映射 $ f: (0,1)\times(0,1) \to \mathbb{R} $，计算

\[
\int_0^1\int_0^1 f(x,y) dx dy , \quad \int_0^1\int_0^1 f(x,y) dy dx , \quad \int_0^1\int_0^1 |f(x,y)| dx dy , \quad \int_0^1\int_0^1 |f(x,y)| dy dx .
\]

\begin{items}
\item

$ f(x,y) := (x-y)/(x^2+y^2)^{3/2} $。

\item

$ f(x,y) := 1/(1-xy)^\alpha $，$ \alpha > 0 $。

\end{items}

\item

设 $ p,q \in [1,\infty] $。证明：

\begin{items}
\item

当 $ p \ne q $ 时 $ L_p(\mathbb{R}^n) \not\subset L_q(\mathbb{R}^n) $。

\item

若 $ X \subset \mathbb{R}^n $ 是开集且有界，则当 $ p > q $ 时 $ L_p(X) \subseteq L_q(X) $。

\end{items}

\end{exercises}

\section{卷积}

本节中，我们利用勒贝格测度的平移不变性在 $ L_{1}(\mathbb{R}^{n}) $ 上引进一个新的乘法——卷积，它依赖于勒贝格积分。我们将证明，这一运算不仅在 $ L_{1}(\mathbb{R}^{n}) $ 上有定义，而且在其他函数空间上也有定义，并且它具有重要的磨光性质。它的应用之一是某些逼近定理，我们在此证明它们，因为它们在后续的构造中极为有用。

我们将主要考虑定义在整个 $ \mathbb{R}^{n} $ 上的 $ \mathbb{K} $ 值函数空间。对于这类空间，我们在记号中省略定义域与陪域。换言之，若 $ \mathcal{F}(\mathbb{R}^{n}) = \mathcal{F}(\mathbb{R}^{n},\mathbb{K}) $ 是 $ \mathbb{R}^{n} $ 上 $ \mathbb{K} $ 值函数构成的向量空间，在没有混淆风险时我们就简写为 $ \mathcal{F} $。于是 $ L_{p} $ 表示 $ L_{p}(\mathbb{R}^{n}) = L_{p}(\mathbb{R}^{n},\mathbb{K}) $，依此类推。此外，$ \int f\,dx $ 将始终表示 $ \int_{\mathbb{R}^{n}} f\,dx $。

\needspace{3\baselineskip}\subsection*{卷积的定义}

设 $ F $ 是 $ \mathbb{K} $-向量空间。对 $ f \in \mathcal{F}(\mathbb{R}^{n},F) $，我们定义另一个函数 $ \tilde{f} \in \mathcal{F}(\mathbb{R}^{n},F) $ 为 $ \tilde{f}(x) := f(-x) $，其中 $ x \in \mathbb{R}^{n} $。映射 $ f \mapsto \tilde{f} $ 称为（关于原点的）\textbf{反演}。

回忆 IX.5.15 中平移群 $ \mathcal{T} := \{\, \tau_{a} ; a \in \mathbb{R}^{n} \,\} $ 的定义。现在我们定义该群在 $ \mathcal{F}(\mathbb{R}^{n},F) $ 上的一个作用\footnote{见习题 I.7.6。}：

\begin{equation}
\mathcal{T} \times \mathcal{F}(\mathbb{R}^{n},F) \to \mathcal{F}(\mathbb{R}^{n},F) , \quad (\tau_{a},f) \mapsto \tau_{a}f ,\tag{7.1}
\end{equation}

其中

\begin{equation}
\tau_{a}f(x) := f(x-a) \quad \text{对 } a, x \in \mathbb{R}^{n} .\tag{7.2}
\end{equation}

因此

\[
\tau_{a}f = f \circ \tau_{-a} = (\tau_{-a})^{\ast}f ,
\]

其中 $ (\tau_{-a})^{\ast} $ 是 VIII.3 节中定义的拉回。

\begin{remark}

\begin{items}
\item

对 $ f \in \mathcal{F} := \mathcal{F}(\mathbb{R}^{n},\mathbb{K}) $，我们有 $ \tilde{f} = (-\mathrm{id}_{\mathbb{R}^{n}})^{\ast}f $。

\item

反演是 $ \mathcal{F} $ 上以及 $ p \in [1,\infty] \cup \{0\} $ 时的 $ \mathcal{L}_{p} $ 上的对合\footnote{映射 $ f \in X^{X} $ 称为\textbf{对合的}，若 $ f \circ f = \mathrm{id}_{X} $。}向量空间同构。

\item

设 $ E \in \{ BC^{k}, BUC^{k}, C_{0} ; k \in \mathbb{N} \} $。则反演属于 $ \mathcal{L}\mathrm{aut}(E) $。

\item

对 $ f \in \mathcal{F} $ 与 $ x \in \mathbb{R}^{n} $，我们有

\[
(\tau_{-x}f)^{\sim}(y) = \tau_{x}\tilde{f}(y) = f(x-y) \quad \text{对 } y \in \mathbb{R}^{n} .
\]
\item

设 $ n = 1 $ 且 $ a > 0 $。则 $ \tau_{a} : \mathbb{R} \to \mathbb{R} $，$ x \mapsto x+a $ 是 $ \mathbb{R} $ 上由 $ a $ 给出的\textbf{右平移}。定义 (7.2) 意味着 $ \tau_{a} $ 也把 $ f $ 的图像向右平移 $ a $。

\end{items}

\bookfigure[0.61]{figures/3c7e9f3b05967fd8.png}

因此 $ \mathcal{T} $ 在 $ \mathcal{F}(\mathbb{R},F) $ 上作为\textbf{右平移}起作用，这就澄清了把 $ \tau_{a}f $ 定义为 $ \mathbb{R} $ 的左平移 $ \tau_{-a} $ 的拉回这一做法。 \nobreak\hspace{0.35em}■

取 $ f, g \in \mathcal{L}_{1} $ 与 $ x \in \mathbb{R}^{n} $，并设 $ O $ 在 $ \mathbb{K} $ 中开。

\[
(\tau_{-x}f)^{-1}(O) = (f \circ \tau_{x})^{-1}(O) = \tau_{-x}(f^{-1}(O)) .
\]

因此由推论 1.5 与引理 IX.5.16 可知 $ (\tau_{-x}f)^{-1}(O) $ 可测。于是再由推论 1.5，$ \tau_{-x}f $ 属于 $ \mathcal{L}_{0} $。现在由注 1.2(d) 以及注 7.1 的 (b) 与 (d) 推得：对每个 $ x \in \mathbb{R}^{n} $，$ y \mapsto f(x-y)g(y) $ 属于 $ \mathcal{L}_{0} $。若此函数可积，我们定义 $ f $ 与 $ g $ 在 $ x $ 处的卷积为

\[
f * g(x) := \int f(x-y)g(y)\,dy .
\]

若 $ f * g(x) $ 对几乎每个 $ x \in \mathbb{R}^{n} $ 有定义，我们就说 $ f $ 与 $ g $ 是\textbf{可卷积的}。此时几乎处处有定义的函数

\[
f * g := (x \mapsto f * g(x))
\]

称为 $ f $ 与 $ g $ 的\textbf{卷积}。若 $ f $ 与 $ g $ 可卷积，且 $ (f * g)^{p} $ 可积（或当 $ p = \infty $ 时 $ f * g $ 本质有界），我们就写 $ f * g \in \mathcal{L}_{p} $，这略有滥用记号之嫌。\footnote{我们的确切意思是 $ f * g $ 的平凡延拓属于 $ \mathcal{L}_{p} $。})

现在我们要证明：对 $ p \in [1,\infty] $，每一对 $ (f,g) \in \mathcal{L}_{p} \times \mathcal{L}_{1} $ 都是可卷积的。下面的观察将有用。

\end{remark}

\begin{lemma}

设 $ f \in \mathcal{L}_{0} $，$ (x,y) \in \mathbb{R}^{n} \times \mathbb{R}^{n} = \mathbb{R}^{2n} $，令

\[
F_{1}(x,y) := f(x) \quad \text{与} \quad F_{2}(x,y) := f(x-y) .
\]

则 $ F_{1} $ 与 $ F_{2} $ 属于 $ \mathcal{L}_{0}(\mathbb{R}^{2n}) $。

证明

\begin{subitems}
\item

设 $ O $ 在 $ \mathbb{K} $ 中开，$ A := f^{-1}(O) $。则 $ A $ 属于 $ \mathcal{L}(n) $。因此注 6.1(b) 与命题 6.2 表明 $ F_{1}^{-1}(O) = A \times \mathbb{R}^{n} $ 是 $ \lambda_{2n} $-可测的。现在关于 $ F_{1} $ 的断言由推论 1.5 推出。

\item

令 $ \varphi(x,y) := (x-y,y) $，$ (x,y) \in \mathbb{R}^{n} \times \mathbb{R}^{n} $。则 $ \varphi \in \mathcal{L}\mathrm{aut}(\mathbb{R}^{2n}) $ 且 $ F_{2} = F_{1} \circ \varphi $。断言于是由 (i) 与定理 IX.5.12 推出。 \nobreak\hspace{0.35em}■

\end{subitems}

\end{lemma}

\qyxffntextdef{1}{通读上一节最后一部分的读者会看出，这个界也可以由积分的闵可夫斯基不等式容易地推出。}
\begin{theorem}

设 $ p \in [1,\infty] $ 且 $ (f,g) \in \mathcal{L}_{p} \times \mathcal{L}_{1} $。

\begin{subitems}
\item

$ f $ 与 $ g $ 可卷积。

\item

（杨不等式）$ f * g \in \mathcal{L}_{p} $，且 $ \|f * g\|_{p} \le \|f\|_{p}\|g\|_{1} $。

\end{subitems}

证明

\begin{items}
\item

先设 $ p \in [1,\infty) $。由引理 7.2 与注 1.2(d)，映射 $ (x,y) \mapsto f(x-y)g(y) $ 属于 $ \mathcal{L}_{0}(\mathbb{R}^{2n}) $。利用赫尔德不等式推得

\begin{equation*}
\begin{aligned}
\int |f(x-y)g(y)|\,dy &= \int |f(x-y)|\,|g(y)|^{1/p}|g(y)|^{1/p'}\,dy \\ &\le \left(\int |f(x-y)|^{p}|g(y)|\,dy\right)^{1/p}\left(\int |g(y)|\,dy\right)^{1/p'} .
\end{aligned}
\end{equation*}

由此与托内利定理，我们得到

\begin{equation*}
\begin{aligned}
\int \left(\int |f(x-y)g(y)|\,dy\right)^{p} dx &\le \|g\|_{1}^{p/p'}\iint |f(x-y)|^{p}|g(y)|\,dy\,dx \\ &= \|g\|_{1}^{p/p'}\iint |f(x-y)|^{p}\,dx\,|g(y)|\,dy \\ &= \|g\|_{1}^{1+p/p'}\|f\|_{p}^{p} < \infty ,
\end{aligned}
\end{equation*}

其中最后一步我们再次用到了勒贝格积分的平移不变性。于是得到\qyxffnmark{1}

\begin{equation}
\left(\int \left[\int |f(x-y)g(y)|\,dy\right]^{p} dx\right)^{1/p} \le \|f\|_{p}\|g\|_{1} < \infty .\tag{7.3}
\end{equation}

现在由注 3.11(c) 断定：对几乎每个 $ x \in \mathbb{R}^{n} $ 有 $ \int |f(x-y)g(y)|\,dy < \infty $；由注 3.11(a)，这已足以说明 $ f $ 与 $ g $ 可卷积。定理的 (ii) 部分于是由 (7.3) 推出。

\item

在 $ p = \infty $ 的情形，我们有

\[
\int |f(x-y)g(y)|\,dy \le \|f\|_{\infty}\|g\|_{1} < \infty \quad \text{对几乎每个 } x \in \mathbb{R}^{n} ,
\]

这立即蕴含 (i) 与 (ii)。 \nobreak\hspace{0.35em}■

\end{items}

\end{theorem}

\begin{corollary}

设 $ ([f],[g]) \in L_{p} \times L_{1} $，$ p \in [1,\infty] $。则

\[
f * g = f^{*} * g^{*} \quad \text{几乎处处于 } \mathbb{R}^{n}
\]

对 $ (f^{*},g^{*}) \in ([f],[g]) $ 成立。

\proofhead
由定理 7.3，$ f * g $、$ f^{*} * g^{*} $ 与 $ f * g^{*} $ 均几乎处处有定义且属于 $ \mathcal{L}_{p} $。因为
\[
f * g - f^{*} * g^{*} = f * (g - g^{*}) + (f - f^{*}) * g^{*} ,
\]

由杨不等式得到

\[
\|f * g - f^{*} * g^{*}\|_{p} \le \|f\|_{p}\|g - g^{*}\|_{1} + \|f - f^{*}\|_{p}\|g^{*}\|_{1} = 0 ,
\]

由此即得断言。 \nobreak\hspace{0.35em}■

\end{corollary}

现在我们可以对 $ p \in [1,\infty) $ 的 $ L_{p} \times L_{1} $ 的元素定义卷积：事实上，推论 7.4 保证映射

\[
* : L_{p} \times L_{1} \to L_{p} , \quad ([f],[g]) \mapsto [f * g]
\]

是良定义的。我们把这个 $ * $ 称为 $ L_{p} \times L_{1} $ 上的\textbf{卷积乘积}，并称 $ [f] * [g] := [f * g] $ 为 $ [f] $ 与 $ [g] $ 的\textbf{卷积}。显然卷积也可以定义在 $ L_{1} \times L_{p} $ 上，对此我们也使用符号 $ * $。

\needspace{3\baselineskip}\subsection*{平移群}

为了能更好地探讨卷积的更多性质，我们先汇集关于平移群 $ (\mathbb{R}^{n},+) $ 在函数空间上的表示的一些重要定义与事实。

设 $ F $ 是 $ \mathbb{K} $-向量空间，$ V $ 是 $ \mathcal{F}(\mathbb{R}^{n},F) $ 的向量子空间，且它在 $ \mathbb{R}^{n} $ 的平移群 $ \mathcal{T} $ 的作用 (7.1) 下是\textbf{不变的}，即对一切 $ a \in \mathbb{R}^{n} $ 有 $ \tau_{a}(V) \subset V $。通过\textbf{限制}，(7.1) 诱导出平移群 $ \mathcal{T} $ 在 $ V $ 上的一个作用

\[
\mathcal{T} \times V \to V , \quad (\tau_{a},v) \mapsto \tau_{a}v
\]

对每个 $ a \in \mathbb{R}^{n} $，映射 $ T_{a} := (v \mapsto \tau_{a}v) $ 是 $ V $ 到自身的线性映射。因为

\[
\tau_{a}\tau_{b}v = \tau_{a+b}v \quad \text{与} \quad \tau_{0}v = v ,
\]

$ T_{a} $ 是 $ V $ 的向量空间自同构，且 $ (T_{a})^{-1} = T_{-a} $。因此\footnote{见注 I.12.2(d) 与 I.7.6(e)。}

\[
(\mathbb{R}^{n},+) \to \mathrm{Aut}(V) , \quad a \mapsto T_{a}
\]

是一个群同态，即群 $ (\mathbb{R}^{n},+) $ 在 $ V $ 上的一个\textbf{线性表示}。特别地，

\[
\mathcal{T}_{V} := \{\, T_{a} \in \mathrm{Aut}(V) ; a \in \mathbb{R}^{n} \,\}
\]

是 $ \mathrm{Aut}(V) $ 的子群，称为 $ V $ 上的\textbf{平移群}。我们倾向于在不致误解时用同一个符号 $ \tau_{a} $ 代替 $ T_{a} $。$ V $ 在 (7.1) 下的不变性也说成 $ (\mathbb{R}^{n},+) $ 在 $ V $ 上\textbf{线性可表示}。

若 $ V $ 是（半）赋范向量空间，则称群 $ \mathcal{T}_{V} $ 是\textbf{强连续的}，若对每个 $ v \in V $ 有 $ \lim_{a \to 0}\tau_{a}v = v $。

\begin{remark}

\begin{items}
\item

$ (\mathbb{R}^{n},+) $ 在 $ \mathcal{F} $ 上以及在 $ \mathcal{B} := B(\mathbb{R}^{n}) $ 上线性可表示。

\item

$ (\mathbb{R}^{n},+) $ 在 $ \mathcal{L}_{\infty} $ 上线性可表示，且对 $ f \in \mathcal{L}_{\infty} $ 有 $ \|\tau_{a}f\|_{\infty} = \|f\|_{\infty} $。

\proofhead
取 $ f \in \mathcal{L}_{\infty} $。对每个 $ \alpha > \|f\|_{\infty} $，存在勒贝格测度为零的集合 $ N $，使得 $ |f(x)| \le \alpha $ 对 $ x \in N^{c} $ 成立。由平移不变性（定理 IX.5.17），$ N_{a} := \tau_{a}(N) $ 的测度也为零，且
\[
|\tau_{a}f(x)| = |f(x-a)| \le \alpha \quad \text{对 } x \in N_{a}^{c} .
\]

因此 $ \tau_{a}f $ 本质有界，且 $ \|\tau_{a}f\|_{\infty} \le \|f\|_{\infty} $。断言成立，因为

\[
\|f\|_{\infty} = \|\tau_{-a}(\tau_{a}f)\|_{\infty} \le \|\tau_{a}f\|_{\infty} .\quad\blacksquare
\]
\item

平移群 $ \mathcal{T}_{\mathcal{B}} $ 与 $ \mathcal{T}_{\mathcal{L}_{\infty}} $ 不强连续。

\proofhead
$ \|\tau_{a}\chi_{\mathbb{B}^{n}} - \chi_{\mathbb{B}^{n}}\|_{\infty} = 1 $，$ a \in \mathbb{R}^{n} \setminus \{0\} $。 \nobreak\hspace{0.35em}■
\item

若 $ \mathcal{T}_{V} $ 强连续，则

\[
(a \mapsto \tau_{a}f) \in C(\mathbb{R}^{n},V) \quad \text{对 } f \in V .
\]

\proofhead
这由 $ \tau_{a}f - \tau_{b}f = \tau_{a-b}(\tau_{b}f) - \tau_{b}f $（$ f \in V $，$ a, b \in \mathbb{R}^{n} $）推出。 \nobreak\hspace{0.35em}■
\end{items}

\end{remark}

\begin{theorem}

设 $ V = \mathcal{L}_{p} $，$ p \in [1,\infty) $，或 $ V = BUC^{k} $，$ k \in \mathbb{N} $。则 $ (\mathbb{R}^{n},+) $ 在 $ V $ 上线性可表示，且平移群 $ \mathcal{T}_{V} $ 强连续。此外对 $ a \in \mathbb{R}^{n} $ 与 $ f \in V $ 有 $ \|\tau_{a}f\|_{V} = \|f\|_{V} $。

证明

\begin{subitems}
\item

先考虑 $ V = BUC^{k} $ 的情形。取 $ f \in BUC^{k} $、$ a \in \mathbb{R}^{n} $ 与 $ \varepsilon > 0 $。则存在 $ \delta > 0 $，使得对所有满足 $ |x-y| < \delta $ 的 $ x, y \in \mathbb{R}^{n} $ 有 $ |f(x) - f(y)| < \varepsilon $。由此

\begin{equation}
|\tau_{a}f(x) - \tau_{a}f(y)| = |f(x-a) - f(y-a)| < \varepsilon\tag{7.4}
\end{equation}

对满足 $ |x-y| < \delta $ 的 $ x, y \in \mathbb{R}^{n} $ 成立。因此 $ \tau_{a}f $ 属于 $ BUC $；又因为

\begin{equation}
\partial^{\alpha}\tau_{a}f = \tau_{a}\partial^{\alpha}f \quad \text{对 } \alpha \in \mathbb{N}^{n} \text{ 与 } |\alpha| \le k ,\tag{7.5}
\end{equation}

我们得到 $ \tau_{a}f \in BUC^{k} $。从而 $ (\mathbb{R}^{n},+) $ 在 $ BUC^{k} $ 上线性可表示。
由注 7.5(b) 与 (7.5)，我们得到 $ \|\tau_{a}f\|_{BC^{k}} = \|f\|_{BC^{k}} $。
现在取 $ x \in \mathbb{R}^{n} $。若 $ |a| < \delta $，可在 (7.4) 中令 $ y = x+a $，得到

\[
|\tau_{a}f(x) - f(x)| < \varepsilon \quad \text{对 } x \in \mathbb{R}^{n} ,
\]

即 $ \|\tau_{a}f - f\|_{\infty} < \varepsilon $ 对 $ a \in \delta\mathbb{B}^{n} $ 成立。类似地，用 (7.5) 可以证明存在 $ \delta_{1} > 0 $，使得 $ \|\tau_{a}f - f\|_{BC^{k}} < \varepsilon $ 对 $ a \in \delta_{1}\mathbb{B}^{n} $ 成立。因此 $ \mathcal{T}_{BUC^{k}} $ 强连续。

\item

设 $ p \in [1,\infty) $，$ f \in \mathcal{L}_{p} $。等式 $ \|\tau_{a}f\|_{p} = \|f\|_{p} $ 由勒贝格积分的平移不变性推出。
现在取 $ \varepsilon > 0 $。由定理 4.14，存在 $ g \in C_{c} $ 使得 $ \|f - g\|_{p} < \varepsilon/3 $。因为 $ g $ 有紧支集，存在 $ \mathbb{R}^{n} $ 的紧子集 $ K $，使得 $ \mathrm{supp}(\tau_{a}g) \subset K $ 对 $ |a| \le 1 $ 成立。又因 $ g $ 一致连续，存在 $ \delta \in (0,1] $ 使得

\[
\|\tau_{a}g - g\|_{\infty} < \varepsilon/3\lambda_{n}(K)^{1/p} \quad \text{对 } a \in \delta\mathbb{B}^{n} .
\]

设 $ a \in \delta\mathbb{B}^{n} $。因为 $ \mathrm{supp}(\tau_{a}g - g) \subset K $，定理 5.1(iv) 表明

\[
\|\tau_{a}g - g\|_{p} < \varepsilon/3 \quad \text{对 } a \in \delta\mathbb{B}^{n} .
\]

由于

\[
\|\tau_{a}f - f\|_{p} \le \|\tau_{a}f - \tau_{a}g\|_{p} + \|\tau_{a}g - g\|_{p} + \|g - f\|_{p}
\]

且 $ \|\tau_{a}f - \tau_{a}g\|_{p} = \|\tau_{a}(f-g)\|_{p} = \|f-g\|_{p} $，我们得到 $ \|\tau_{a}f - f\|_{p} < \varepsilon $ 对 $ a \in \delta\mathbb{B}^{n} $，证毕。 \nobreak\hspace{0.35em}■

\end{subitems}

\end{theorem}

现在我们对 $ p \in [1,\infty] $ 定义 $ \mathcal{T} $ 在 $ L_{p} $ 上的一个作用。由注 7.5(b) 与定理 7.6，对每个 $ a \in \mathbb{R}^{n} $，$ \tau_{a} $ 是 $ \mathcal{L}_{p} $ 的等距。因此映射

\[
L_{p} \to L_{p} , \quad [f] \mapsto [\tau_{a}f]
\]

对每个 $ a \in \mathbb{R}^{n} $ 都良定义。我们也用 $ \tau_{a} $ 表示它，即令

\[
\tau_{a}[f] := [\tau_{a}f] \quad \text{对 } f \in L_{p} \text{ 与 } a \in \mathbb{R}^{n} .
\]

于是

\begin{equation}
\|\tau_{a}[f]\|_{p} = \|[\tau_{a}f]\|_{p} = \|\tau_{a}f\|_{p} = \|f\|_{p} = \|[f]\|_{p} .\tag{7.6}
\end{equation}

显然

\[
\mathcal{T} \times L_{p} \to L_{p} , \quad (\tau_{a},f) \mapsto \tau_{a}f
\]

是 $ \mathbb{R}^{n} $ 的平移群 $ \mathcal{T} $ 在 $ L_{p} $ 上的作用。由注 7.5(b) 与定理 7.6，对每个 $ a \in \mathbb{R}^{n} $，$ T_{a} := (f \mapsto \tau_{a}f) $ 是 $ L_{p} $ 上的线性等距。再次把 $ T_{a} $ 写成 $ \tau_{a} $，我们断定

\[
(\mathbb{R}^{n},+) \to \mathcal{L}\mathrm{aut}(L_{p}) , \quad a \mapsto \tau_{a}
\]

是 $ \mathbb{R}^{n} $ 的加法群通过 $ L_{p} $ 上的线性等距给出的一个表示。特别地，$ L_{p} $ 上的\textbf{平移群}，即

\[
\mathcal{T}_{L_{p}} := \{\, \tau_{a} \in \mathcal{L}\mathrm{aut}(L_{p}) ; a \in \mathbb{R}^{n} \,\} ,
\]

是 $ \mathcal{L}\mathrm{aut}(L_{p}) $ 的由等距组成的子群。

\begin{corollary}

对 $ 1 \le p < \infty $，$ L_{p} $ 上的平移群强连续。

\proofhead
这是定理 7.6 与 (7.6) 的直接推论。 \nobreak\hspace{0.35em}■
\end{corollary}

\needspace{3\baselineskip}\subsection*{卷积的基本性质}

在关于平移群的这段题外话之后，我们回到卷积并导出它的主要性质。

\qyxffnnumdef{1}
\afterpage{\qyxffntextof{1}{见定理 4.18。}}
\begin{theorem}

考虑 $ (f,g) \in L_{p} \times L_{1} $，$ p \in [1,\infty] $。

\begin{subitems}
\item

卷积 $ f * g $ 属于 $ L_{p} $，且满足\textbf{杨不等式}

\[
\|f * g\|_{p} \le \|f\|_{p}\|g\|_{1} .
\]
\item

$ f * g = g * f $。

\item

若 $ p = \infty $，则卷积 $ f * g $ 属于\qyxffnmark{1} $ BUC $。

\item

对 $ \varphi \in BC^{k} $，我们有 $ \varphi * g \in BUC^{k} $，

\[
\partial^{\alpha}(\varphi * g) = \partial^{\alpha}\varphi * g \quad \text{对 } \alpha \in \mathbb{N}^{n} , \ |\alpha| \le k ,
\]

且 $ \|\varphi * g\|_{BC^{k}} \le \|\varphi\|_{BC^{k}}\|g\|_{1} $。

\proofhead
(i) 由定理 7.3(ii) 与推论 7.4 推出。
\end{subitems}

\begin{subitems}
\item[（ii）]

取 $ x \in \mathbb{R} $，设 $ f^{*} $ 与 $ g^{*} $ 分别是 $ f $ 与 $ g $ 的代表。又令 $ \psi(y) := x-y $，$ y \in \mathbb{R}^{n} $。则 $ \psi $ 是 $ \mathbb{R}^{n} $ 的对合等距。由定理 7.3(i) 与例 6.6(b) 可得

\begin{equation*}
\begin{aligned}
f^{*} * g^{*}(x) &= \int f^{*}(x-y)g^{*}(y)\,dy = \int (f^{*} \circ \psi)\left((g^{*} \circ \psi) \circ \psi\right) dy \\ &= \int (g^{*} \circ \psi)f^{*}\,dy = \int g^{*}(x-y)f^{*}(y)\,dy = g^{*} * f^{*}(x) .
\end{aligned}
\end{equation*}

因此由推论 7.4 得 $ f * g = g * f $。

\item[（iii）]

勒贝格积分的平移不变性给出 $ \|\tilde{g}\|_{1} = \|g\|_{1} $。由 (ii)，并因为 $ \mathcal{T}_{L_{1}} $ 的元素是等距，我们于是得到

\begin{equation*}
\begin{aligned}
|f^{*} * g^{*}(x) - f^{*} * g^{*}(y)| &\le \int |f^{*}(z)(g^{*}(x-z) - g^{*}(y-z))|\,dz \le \|f^{*}\|_{\infty}\|\tau_{x}\tilde{g} - \tau_{y}\tilde{g}\|_{1} \\ &= \|f\|_{\infty}\|\tau_{y}(\tau_{x-y}\tilde{g} - \tilde{g})\|_{1} = \|f\|_{\infty}\|\tau_{x-y}\tilde{g} - \tilde{g}\|_{1}
\end{aligned}
\end{equation*}

对 $ x, y \in \mathbb{R}^{n} $ 成立。因为 $ \tilde{g} \in L_{1} $，$ \mathcal{T}_{L_{1}} $ 的强连续性连同 (i) 表明 $ f^{*} * g^{*} \in BUC $。断言得证。

\item[（iv）]

鉴于 (iii)，只需考虑 $ k \ge 1 $ 的情形。为此我们定义 $ h(x,y) := \varphi(x-y)g(y) $，$ (x,y) \in \mathbb{R}^{n} \times \mathbb{R}^{n} $。则 $ h $ 满足定理 3.18 的假设，且由此得 $ \partial_{j}(\varphi * g) = \partial_{j}\varphi * g $，$ j \in \{1,\ldots,n\} $。由 (iii) 与定理 VII.2.10，我们有 $ \varphi * g \in BUC^{1} $。现在用归纳法可知 $ \varphi * g $ 属于 $ BUC^{k} $，且对每个满足 $ |\alpha| \le k $ 的 $ \alpha \in \mathbb{N}^{n} $ 有 $ \partial^{\alpha}(\varphi * g) = \partial^{\alpha}\varphi * g $。最后，由 (i) 我们有

\begin{equation*}
\begin{aligned}
\|\varphi * g\|_{BC^{k}} &= \max_{|\alpha| \le k}\|\partial^{\alpha}(\varphi * g)\|_{\infty} = \max_{|\alpha| \le k}\|(\partial^{\alpha}\varphi) * g\|_{\infty} \le \left(\max_{|\alpha| \le k}\|\partial^{\alpha}\varphi\|_{\infty}\right)\|g\|_{1} \\ &= \|\varphi\|_{BC^{k}}\|g\|_{1} . \quad \blacksquare
\end{aligned}
\end{equation*}
\end{subitems}

\end{theorem}

\begin{corollary}

\begin{subitems}
\item

设 $ p \in (1,\infty) $ 且 $ k \in \mathbb{N} $。卷积满足

\begin{equation*}
* \in \begin{cases} \mathcal{L}_{\mathrm{sym}}^{2}(L_{1},L_{1}) , \\ \mathcal{L}(L_{p},L_{1};L_{p}) , \\ \mathcal{L}(L_{\infty},L_{1};BUC) , \\ \mathcal{L}(BC^{k},L_{1};BUC^{k}) , \end{cases}
\end{equation*}

且所有这些映射的范数至多为 1。

\item

$ (L_{1},+,*) $ 是没有乘法单位元的交换巴拿赫代数。

\proofhead
(i) 与 (ii) 的第一个断言由定理 7.8 直接推出。现在假设存在 $ e \in L_{1} $ 使得对每个 $ f \in L_{1} $ 有 $ e * f = f $。我们取 $ e $ 的一个代表 $ e^{*} $，然后由习题 2.15 找到 $ \delta > 0 $ 使得
\[
\left|\int_{\delta\mathbb{B}^{n}} e^{*}(x-y)\,dy\right| = \left|\int_{\mathbb{B}^{n}(x,\delta)} e^{*}(z)\,dz\right| < 1 \quad \text{对 } x \in \mathbb{R}^{n} .
\]

此外，存在勒贝格测度为零的集合 $ N $，使得 $ \chi_{\delta\mathbb{B}^{n}}(x) = e^{*} * \chi_{\delta\mathbb{B}^{n}}(x) $ 对 $ x \in N^{c} $ 成立。然而，对 $ x \in \delta\mathbb{B}^{n} \cap N^{c} $，我们有

\[
1 = \chi_{\delta\mathbb{B}^{n}}(x) = e^{*} * \chi_{\delta\mathbb{B}^{n}}(x) = \int_{\mathbb{R}^{n}} e^{*}(x-y)\chi_{\delta\mathbb{B}^{n}}(y)\,dy = \int_{\delta\mathbb{B}^{n}} e^{*}(x-y)\,dy < 1 ,
\]

这是不可能的。 \nobreak\hspace{0.35em}■

\end{subitems}

\end{corollary}

\begin{theorem}

（支集的可加性） 设 $ f, g \in \mathcal{L}_{0} $ 可卷积且 $ f $ 有紧支集。则

\[
\mathrm{supp}(f * g) \subset \mathrm{supp}(f) + \mathrm{supp}(g) .
\]

证明

\begin{subitems}
\item

我们可以假设 $ f * g \neq 0 $。对 $ x \in [f * g \neq 0] $，存在 $ y \in \mathbb{R}^{n} $ 使得 $ f(x-y)g(y) \neq 0 $。由此 $ y \in \mathrm{supp}(g) $ 且 $ x \in y + \mathrm{supp}(f) $，于是 $ x $ 属于 $ \mathrm{supp}(f) + \mathrm{supp}(g) $。因此 $ [f * g \neq 0] \subset \mathrm{supp}(f) + \mathrm{supp}(g) $。

\item

我们证明 $ \mathrm{supp}(f) + \mathrm{supp}(g) $ 是闭的。设 $ (x_{k}) $ 是 $ \mathrm{supp}(f) + \mathrm{supp}(g) $ 中的序列，使得 $ x_{k} \to x $ 对某个 $ x \in \mathbb{R}^{n} $ 成立。则存在 $ \mathrm{supp}(f) $ 中的序列 $ (a_{k}) $ 与 $ \mathrm{supp}(g) $ 中的序列 $ (b_{k}) $，使得 $ x_{k} = a_{k} + b_{k} $，$ k \in \mathbb{N} $。因为 $ \mathrm{supp}(f) $ 紧，存在 $ (a_{k}) $ 的子序列 $ (a_{k_{\ell}})_{\ell \in \mathbb{N}} $ 与 $ a \in \mathrm{supp}(f) $，使得 $ a_{k_{\ell}} \to a $（$ \ell \to \infty $）。于是 $ b_{k_{\ell}} = x_{k_{\ell}} - a_{k_{\ell}} \to x - a $（$ k \to \infty $）。因为 $ \mathrm{supp}(g) $ 闭，我们知道 $ x-a $ 属于 $ \mathrm{supp}(g) $。因此存在 $ b \in \mathrm{supp}(g) $ 使得 $ x = a+b $。这说明 $ \mathrm{supp}(f) + \mathrm{supp}(g) $ 是闭的。断言由推论 III.2.13 推出。 \nobreak\hspace{0.35em}■

\end{subitems}

\end{theorem}

\needspace{3\baselineskip}\subsection*{恒等逼近}

由推论 7.9 我们看到卷积代数 $ L_{1} $ 没有乘法单位元。然而，下一定理保证了“恒等逼近”的存在性，即元素 $ \varphi \in L_{1} $ 满足对每个 $ f \in L_{1} $ 有 $ \|\varphi * f - f\|_{1} < \varepsilon $（对给定的 $ \varepsilon > 0 $）。

\qyxffnnumdef{1}
\afterpage{\qyxffntextof{1}{这也可由积分的闵可夫斯基不等式推出。}}
\begin{theorem}

（逼近定理） 给定 $ E \in \{ L_{p} ; 1 \le p < \infty \} $ 或 $ E \in \{ BUC^{k} ; k \in \mathbb{N} \} $，设 $ \varphi \in \mathcal{L}_{1} $ 且

\[
a := \int \varphi\,dx , \quad \varphi_{\varepsilon}(x) := \varepsilon^{-n}\varphi(x/\varepsilon) \quad \text{对 } x \in \mathbb{R}^{n} , \ \varepsilon > 0 .
\]

则 $ \lim_{\varepsilon \to 0}\varphi_{\varepsilon} * f = af $ 在 $ E $ 中对 $ f \in E $ 成立。

证明

\begin{subitems}
\item

固定 $ \varepsilon > 0 $。由换元法则——例 6.6(b)——我们知道 $ \varphi_{\varepsilon} \in L_{1} $ 且 $ \int \varphi_{\varepsilon}\,dx = a $。因此定理 7.8 表明 $ \varphi_{\varepsilon} * f \in E $ 对 $ f \in E $ 成立。

\item

为证明 $ \varepsilon \to 0 $ 时的极限，先考虑 $ E = L_{p} $ 的情形。取 $ f \in \mathcal{L}_{p} $ 与 $ \varepsilon > 0 $。由定理 7.3(i) 以及定理 7.8(ii) 的证明，并用例 6.6(b) 中的变换 $ y \mapsto y/\varepsilon $，我们得到

\begin{equation}
\begin{aligned}
\varphi_{\varepsilon} * f(x) - af(x) &= f * \varphi_{\varepsilon}(x) - af(x) = \int [f(x-y) - f(x)]\varphi_{\varepsilon}(y)\,dy \\ &= \int [f(x-\varepsilon z) - f(x)]\varphi(z)\,dz = \int [\tau_{\varepsilon z}f(x) - f(x)]\varphi(z)\,dz
\end{aligned}\tag{7.7}
\end{equation}

对几乎每个 $ x \in \mathbb{R}^{n} $ 成立。推论 7.7 与注 7.5(d) 表明

\begin{equation}
(z \mapsto (\tau_{\varepsilon z}f - f)) \in C(\mathbb{R}^{n},E) \quad \text{对 } \varepsilon > 0 ,\tag{7.8}
\end{equation}

且

\begin{equation}
\lim_{\varepsilon \to 0}\|\tau_{\varepsilon z}f - f\|_{E} = 0 \quad \text{对 } z \in \mathbb{R}^{n} .\tag{7.9}
\end{equation}

现在令

\[
g^{\varepsilon}(z) := (\tau_{\varepsilon z}f - f)\varphi(z) \quad \text{对 } z \in \mathbb{R}^{n} \text{ 与 } \varepsilon > 0 .
\]

则由 (7.8)、定理 1.17 与注 1.2(d) 可知，对每个 $ \varepsilon > 0 $，$ g^{\varepsilon} $ 属于 $ \mathcal{L}_{0}(\mathbb{R}^{n},E) $。因为 $ \|\tau_{\varepsilon z}f\|_{E} = \|f\|_{E} $，我们还可由
三角不等式

\[
\|g^{\varepsilon}(z)\|_{E} \le 2\|f\|_{E}|\varphi(z)| \quad \text{对 } z \in \mathbb{R}^{n} \text{ 与 } \varepsilon > 0 .
\]

因为 $ \varphi \in \mathcal{L}_{1}(\mathbb{R}^{n}) $，我们由此断定 $ g^{\varepsilon} \in \mathcal{L}_{1}(\mathbb{R}^{n},E) $。于是 (7.7) 与定理 2.11(i) 给出界\qyxffnmark{1}

\[
\|\varphi_{\varepsilon} * f - af\|_{E} = \left\|\int g^{\varepsilon}(z)\,dz\right\|_{E} \le \int\|g^{\varepsilon}(z)\|_{E}\,dz .
\]

现在控制收敛定理表明，当 $ \varepsilon \to 0 $ 时 $ \varphi_{\varepsilon} * f $ 在 $ E $ 中收敛到 $ af $，因为由 (7.9)，对几乎每个 $ z \in \mathbb{R}^{n} $ 都有 $ \lim_{\varepsilon \to 0}\|g^{\varepsilon}(z)\|_{E} = 0 $。

\item

现在设 $ f \in BUC^{k} $。若 $ \varphi = 0 $ 是 $ \lambda_{n} $-几乎处处成立的，则断言显然为真。于是设 $ m := \int|\varphi|\,dx > 0 $。由定理 7.8(ii) 与 (iv) 推出

\[
\partial^{\alpha}(\varphi_{\varepsilon} * f - af) = \varphi_{\varepsilon} * \partial^{\alpha}f - a\partial^{\alpha}f \quad \text{对 } \alpha \in \mathbb{N}^{n} \text{ 且 } |\alpha| \le k .
\]

因此只需考虑 $ k = 0 $ 的情形。
设 $ \eta > 0 $。则存在 $ \delta > 0 $ 使得

\[
|f(x - y) - f(x)| \le \eta/2m \quad \text{对 } x, y \in \mathbb{R}^{n} ,\quad |y| < \delta ,
\]

并且我们得到

\begin{equation}
\begin{aligned}
|\varphi_{\varepsilon} * f(x) - af(x)| &\le \int|f(x - y) - f(x)|\,|\varphi_{\varepsilon}(y)|\,dy \\ &\le \frac{\eta}{2m}\int_{[|y|<\delta]}|\varphi_{\varepsilon}(y)|\,dy + 2\|f\|_{\infty}\int_{[|y|\ge\delta]}|\varphi_{\varepsilon}(y)|\,dy \\ &\le \frac{\eta}{2} + 2\|f\|_{\infty}\int_{[|y|\ge\delta]}|\varphi_{\varepsilon}(y)|\,dy
\end{aligned}\tag{7.10}
\end{equation}

对 $ x \in \mathbb{R}^{n} $。换元法则于是给出

\[
\int_{[|y|\ge\delta]}|\varphi_{\varepsilon}(y)|\,dy = \varepsilon^{-n}\int_{[|y|\ge\delta]}|\varphi(y/\varepsilon)|\,dy = \int_{[|z|\ge\delta/\varepsilon]}|\varphi(z)|\,dz .
\]

由控制收敛定理，存在 $ \varepsilon_{0} > 0 $ 使得

\[
\int_{[|y|\ge\delta]}|\varphi_{\varepsilon}(y)|\,dy \le \frac{\eta}{4\|f\|_{\infty}} \quad \text{对 } \varepsilon \in (0,\varepsilon_{0}] .
\]

现在由 (7.10) 推出断言。\nobreak\hspace{0.35em}■

\end{subitems}

\end{theorem}

设 $ \varphi \in \mathcal{L}_{1} $ 满足 $ \int\varphi\,dx = 1 $，并置

\begin{equation}
\varphi_{\varepsilon}(x) := \varepsilon^{-n}\varphi(x/\varepsilon) \quad \text{对 } x \in \mathbb{R}^{n} \text{ 与 } \varepsilon > 0 .\tag{7.11}
\end{equation}

族 $ \{\varphi_{\varepsilon} ; \varepsilon > 0\} $ 称为\textbf{逼近核}或\textbf{恒等逼近}。若

\[
\varphi \in C^{\infty}(\mathbb{R}^{n},\mathbb{R}) , \quad \check{\varphi} = \varphi , \quad \varphi \ge 0 , \quad \operatorname{supp}(\varphi) \subset \overline{\mathbb{B}}^{n} , \quad \int\varphi\,dx = 1 ,
\]

则称 $ \{\varphi_{\varepsilon} ; \varepsilon > 0\} $ 为\textbf{磨光核}或\textbf{光滑核}。每个光滑核显然满足

\[
\operatorname{supp}(\varphi_{\varepsilon}) \subset \varepsilon\overline{\mathbb{B}}^{n} , \quad \|\varphi_{\varepsilon}\|_{1} = 1 \quad \text{对 } \varepsilon > 0 .
\]

7.12 例\footnote{在两个例子中，图形下方的面积始终为 1，因此 $ \varepsilon $ 越小，最大值就相应地越高。} (a) \textbf{高斯核}是由

\[
k(x) := (4\pi)^{-n/2}e^{-|x|^{2}/4} \quad \text{对 } x \in \mathbb{R}^{n}
\]

定义的族 $ \{k_{\varepsilon} ; \varepsilon > 0\} $。它是逼近核。

\bookfigure[0.40]{figures/6c7d70484ac69084.png}

\begin{proof}

由例 6.13(a)，我们知道

\[
\int g(x)\,dx = 1
\]

其中 $ g(x) := \pi^{-n/2}e^{-|x|^{2}} $。由于 $ k(x) = 2^{-n}g(x/2) $ 对 $ x \in \mathbb{R}^{n} $ 成立，由换元法则推出 $ \int k(x)\,dx = 1 $。\nobreak\hspace{0.35em}■

\begin{items}
\item[（b）]

设

\begin{equation*}
\varphi(x) := \begin{cases} c\,e^{1/(|x|^{2}-1)} & \text{若 } |x| < 1 , \\ 0 & \text{若 } |x| \ge 1 , \end{cases}
\end{equation*}

其中 $ c := \left(\int_{\mathbb{B}^{n}}e^{1/(|x|^{2}-1)}\,dx\right)^{-1} $ 的选取使 $ \varphi $ 的积分为 1。则族 $ \{\varphi_{\varepsilon} ; \varepsilon > 0\} $ 是光滑核。

\end{items}

\bookfigure[0.42]{figures/9bd60790e5c3e15a.png}

\proofhead
由于 $ x \mapsto |x|^{2} - 1 $ 在 $ \mathbb{R}^{n} $ 上光滑，例 IV.1.17 表明 $ \varphi $ 属于 $ C^{\infty}(\mathbb{R}^{n},\mathbb{R}) $（见习题 VII.5.16）。断言容易推出。\nobreak\hspace{0.35em}■
\end{proof}

\needspace{3\baselineskip}\subsection*{检验函数}

设 $ X $ 是度量空间，$ A $ 与 $ B $ 是 $ X $ 的子集。若 $ \overline{A} $ 紧且含于 $ B $ 的内部，则称 $ A $ \textbf{紧包含于} $ B $（记作 $ A \subset\subset B $）。

若 $ X $ 在 $ \mathbb{R}^{n} $ 中开且 $ E $ 是赋范向量空间，则称

\[
\mathcal{D}(X,E) := \{\varphi \in C^{\infty}(X,E) ; \operatorname{supp}(\varphi) \subset\subset X\}
\]

为 $ X $ 上的\textbf{（$ E $ 值）检验函数空间}。当 $ E = \mathbb{K} $ 时，按惯例记 $ \mathcal{D}(X) := \mathcal{D}(X,\mathbb{K}) $。显然 $ \mathcal{D}(X,E) $ 是 $ C^{\infty}(X,E) $ 与 $ C_{c}(X,E) $ 的向量子空间，且 $ \mathcal{D}(X,E) = C^{\infty}(X,E) \cap C_{c}(X,E) $。由于映射

\[
j : C_{c}(X,E) \to C_{c}(\mathbb{R}^{n},E) , \quad g \mapsto \tilde{g} ,
\]

是线性的且为单射，我们可以把 $ C_{c}(X,E) $ 与 $ C_{c}(\mathbb{R}^{n},E) $ 的一个向量子空间等同起来，并在需要时把前者的每个元素看作后者的元素。同样，我们把 $ \mathcal{D}(X,E) $ 与 $ \mathcal{D}(\mathbb{R}^{n},E) $ 的一个向量子空间等同起来。有了这些记号，对每个 $ p \in [1,\infty] $ 我们都有如下向量子空间的包含关系：

\[
\mathcal{D}(X,E) \subset \mathcal{D}(\mathbb{R}^{n},E) \subset C_{c}(\mathbb{R}^{n},E) \subset L_{p}(\mathbb{R}^{n},E) .
\]

\qyxffntextdef{1}{这里约定 $ \operatorname{dist}(\operatorname{supp}(g),\varnothing) := \infty $。}
\begin{theorem}

设 $ X $ 在 $ \mathbb{R}^{n} $ 中开，$ p \in [1,\infty) $。则 $ \mathcal{D}(X) $ 是 $ L_{p}(X) $ 与 $ C_{0}(X) $ 的稠密向量子空间。

证明

\begin{subitems}
\item

取 $ g \in C_{c}(X) $ 与 $ \eta > 0 $。又设 $ \{\varphi_{\varepsilon} ; \varepsilon > 0\} $ 是光滑核。由定理 7.8，$ \varphi_{\varepsilon} * g $ 属于 $ BUC^{k} $，从而对每个 $ k \in \mathbb{N} $ 属于 $ BUC^{\infty} $。由于 $ g $ 具有紧支集，存在 $ \varepsilon_{0} > 0 $ 使得\qyxffnmark{1} $ \operatorname{dist}(\operatorname{supp}(g),X^{c}) \ge \varepsilon_{0} $。由定理 7.10 推出

\[
\operatorname{supp}(\varphi_{\varepsilon} * g) \subset \operatorname{supp}(\varphi_{\varepsilon}) + \operatorname{supp}(g) \subset \operatorname{supp}(g) + \varepsilon\overline{\mathbb{B}}^{n} \quad \text{对 } \varepsilon > 0 .
\]

于是对 $ \varepsilon \in (0,\varepsilon_{0}) $，$ \varphi_{\varepsilon} * g $ 属于 $ \mathcal{D}(X) $。最后由定理 7.11，对每个 $ q \in [1,\infty] $ 都能找到某个 $ \varepsilon_{1} \in (0,\varepsilon_{0}) $ 使得 $ \|\varphi_{\varepsilon_{1}} * g - g\|_{q} < \eta/2 $。

\item

现在设 $ f \in L_{p}(X) $。由定理 5.1，可以找到 $ g \in C_{c}(X) $ 使得 $ \|f - g\|_{p} < \eta/2 $。由 (i)，存在 $ h \in \mathcal{D}(X) $ 使得 $ \|f - h\|_{p} < \eta $。

\item

对 $ f \in C_{0}(X) $，设 $ K $ 是 $ X $ 的紧子集，使得对 $ x \in X \setminus K $ 有 $ |f(x)| < \eta/2 $。由命题 4.13，可以选取 $ \varphi \in C_{c}(X) $ 使得 $ 0 \le \varphi \le 1 $ 且 $ \varphi|K = 1 $。置 $ g := \varphi f $。由于对 $ x \in K $ 有 $ f(x) = g(x) $，推出

\[
|f(x) - g(x)| = |f(x)|\,|1 - \varphi(x)| < \eta/2 \quad \text{对 } x \in X .
\]

因此 $ \|f - g\|_{\infty} \le \eta/2 $。断言于是由 (i) 推出。\nobreak\hspace{0.35em}■

\end{subitems}

\end{theorem}

\needspace{3\baselineskip}\subsection*{光滑单位分解}

在第 4 节中，我们证明了一般度量空间中连续乌雷松函数的存在性。这一结果在 $ \mathbb{R}^{n} $ 的特殊情形下可以明显改进，此时我们可以用磨光核真正构造出光滑截断函数。

\begin{proposition}

（光滑截断函数） 设 $ K \subset \mathbb{R}^{n} $ 紧，并置

\[
K_{\rho} := \{x \in \mathbb{R}^{n} ; \operatorname{dist}(x,K) < \rho\} \quad \text{对 } \rho > 0 .
\]

则对每个 $ \alpha \in \mathbb{N}^{n} $ 与每个 $ \rho > 0 $，存在正常数 $ c(\alpha) $ 与映射 $ \varphi \in \mathcal{D}(K_{\rho}) $ 使得 $ 0 \le \varphi \le 1 $、$ \varphi|K = 1 $，且 $ \|\partial^{\alpha}\varphi\|_{\infty} \le c(\alpha)\rho^{-|\alpha|} $。

\proofhead
设 $ \{\psi_{\varepsilon} ; \varepsilon > 0\} $ 是光滑核。置 $ \delta := \rho/3 $ 与 $ \varphi := \psi_{\delta} * \chi_{K_{\delta}} $。则 $ \varphi $ 属于 $ BUC^{\infty} $，且由定理 7.10 推出
\[
\operatorname{supp}(\varphi) \subset \operatorname{supp}(\psi_{\delta}) + \overline{K}_{\delta} \subset \delta\overline{\mathbb{B}}^{n} + \overline{K}_{\delta} \subset \overline{K}_{2\delta} \subset K_{3\delta} = K_{\rho} .
\]
\bookfigure[0.38]{figures/d1f82b83da527eb6.png}

因此 $ \varphi $ 属于 $ \mathcal{D}(K_{\rho}) $。此外

\[
\varphi(x) = \int\psi_{\delta}(x - y)\chi_{K_{\delta}}(y)\,dy \le \int\psi_{\delta}(x - y)\,dy = 1
\]

对 $ x \in \mathbb{R}^{n} $ 成立，从而 $ 0 \le \varphi \le 1 $。若 $ x $ 位于 $ K $ 中，则

\[
\varphi(x) = \int\psi_{\delta}(y)\chi_{K_{\delta}}(x - y)\,dy = \int\psi_{\delta}(y)\,dy = 1 ,
\]

因此 $ \varphi|K = 1 $。最后，由于 $ \partial^{\alpha}\psi_{\delta} = \delta^{-|\alpha|}(\partial^{\alpha}\psi_{1})_{\delta} $ 对 $ \alpha \in \mathbb{N}^{n} $ 成立，由定理 7.8(iv) 得到

\[
\partial^{\alpha}\varphi = \partial^{\alpha}(\psi_{\delta} * \chi_{K_{\delta}}) = \partial^{\alpha}\psi_{\delta} * \chi_{K_{\delta}} = \delta^{-|\alpha|}(\partial^{\alpha}\psi_{1})_{\delta} * \chi_{K_{\delta}} .
\]

现在 $ c(\alpha) := 3^{|\alpha|}\|\partial^{\alpha}\psi_{1}\|_{1} $ 与 $ \delta > 0 $ 无关，于是由杨不等式推出 $ \|\partial^{\alpha}\varphi\|_{\infty} \le c(\alpha)\rho^{-|\alpha|} $。\nobreak\hspace{0.35em}■

\end{proposition}

设 $ K \subset \mathbb{R}^{n} $ 紧，并用 $ \{X_{j} ; 0 \le j \le m\} $ 表示 $ K $ 的一个有限开覆盖。若对每个 $ j \in \{0,\ldots,m\} $ 都存在 $ \varphi_{j} \in C^{\infty}(\mathbb{R}^{n}) $ 使得

\begin{subitems}
\item

$ 0 \le \varphi_{j} \le 1 $，

\item

$ \operatorname{supp}(\varphi_{j}) \subset X_{j} $，且

\item

$ \sum_{j=0}^{m}\varphi_{j}(x) = 1 $ 对 $ x \in K $，
则称 $ \{\varphi_{j} ; 0 \le j \le m\} $ 为 $ K $ 上从属于覆盖 $ \{X_{j} ; 0 \le j \le m\} $ 的\textbf{光滑单位分解}。
若 $ X_{0} $ 在 $ \mathbb{R}^{n} $ 中开且 $ K \subset X_{0} $，则 $ \operatorname{dist}(K,X_{0}^{c}) > 0 $，而命题 7.14（取 $ \rho := \operatorname{dist}(K,X_{0}^{c}) $）保证了 $ K $ 上从属于单元素覆盖 $ \{X_{0}\} $ 的光滑单位分解的存在性。为处理有限覆盖的一般情形，我们需要一个技术性结果：

\end{subitems}

\begin{lemma}

（收缩引理） 设 $ \{X_{j} ; 0 \le j \le m\} $ 是 $ \mathbb{R}^{n} $ 的紧子集 $ K $ 的有限开覆盖。则存在 $ K $ 的开覆盖 $ \{U_{j} ; 0 \le j \le m\} $ 使得 $ U_{j} \subset\subset X_{j} $ 对 $ j \in \{0,\ldots,m\} $ 成立。

\proofhead
给定 $ x \in K $，选取 $ j \in \{0,\ldots,m\} $ 使得 $ x \in X_{j} $，并取 $ r_{x} > 0 $ 使得 $ V_{x} := \mathbb{B}^{n}(x,r_{x}) $ 紧且含于 $ X_{j} $。则 $ \{V_{x} ; x \in K\} $ 是 $ K $ 的开覆盖，于是存在 $ k \in \mathbb{N} $ 与 $ \{x_{0},\ldots,x_{k}\} \subset K $ 使得 $ K \subset \bigcup_{i=0}^{k}V_{x_{i}} $。取 $ U_{j} := \bigcup\{V_{x_{i}} ; V_{x_{i}} \subset X_{j}\} $（$ j \in \{0,\ldots,m\} $），我们便得到具有所需性质的族 $ \{U_{j} ; 0 \le j \le m\} $。\nobreak\hspace{0.35em}■
\end{lemma}

\begin{theorem}

（光滑单位分解） 若 $ K $ 是 $ \mathbb{R}^{n} $ 的紧子集，则 $ K $ 的每个有限开覆盖都有从属于它的光滑单位分解。

\proofhead
设 $ \{X_{j} ; 0 \le j \le m\} $ 是 $ K $ 的有限开覆盖。由引理 7.15，存在开覆盖 $ \{U_{j} ; 0 \le j \le m\} $ 使得 $ U_{j} \subset\subset X_{j} $ 对 $ j \in \{0,\ldots,m\} $ 成立。我们定义 $ K_{j} := \overline{U}_{j} $。则 $ K_{j} $ 紧，且 $ \operatorname{dist}(K_{j},X_{j}^{c}) $ 对每个 $ j \in \{0,\ldots,m\} $ 都为正。命题 7.14 于是表明存在 $ \psi_{j} \in \mathcal{D}(X_{j}) $ 使得 $ 0 \le \psi_{j} \le 1 $ 且 $ \psi_{j}|K_{j} = 1 $。置
\[
\varphi_{0} := \psi_{0} \quad \text{与} \quad \varphi_{k} := \psi_{k}\prod_{j=0}^{k-1}(1 - \psi_{j}) \quad \text{对 } 1 \le k \le m ,
\]

容易用归纳法验证 $ \sum_{j=0}^{m}\varphi_{j} = 1 - \prod_{j=0}^{m}(1 - \psi_{j}) $。断言现在由 $ K \subset \bigcup_{j=0}^{m}K_{j} $ 推出。\nobreak\hspace{0.35em}■

\end{theorem}

下面我们给出定理 7.16 的一些简单应用。更多更复杂的情形将在后续各章中描述。

\begin{application}

\begin{items}
\item

设 $ X $ 在 $ \mathbb{R}^{n} $ 中开。则对 $ f \in L_{0}(X) $，下列命题等价：

\begin{subitems}
\item

$ f \in L_{1,loc}(X) $；

\item

对每个 $ \varphi \in \mathcal{D}(X) $ 有 $ \varphi f \in L_{1}(X) $；

\item

对每个 $ K = \overline{K} \subset\subset X $ 有 $ \widetilde{f|K} \in L_{1}(X) $。

\proofhead
设 $ (U_{j})_{j\in\mathbb{N}} $ 是 $ X $ 的一列递增的相对紧开子集，且 $ X = \bigcup_{j}U_{j} $（见注 1.16(d) 与 (e)）。则（见习题 4.3）
\[
L_{1,loc}(X) = \{f \in L_{0}(X) ; \chi_{U_{j}}f \in L_{1}(X),\ j \in \mathbb{N}\} .
\]

“(i)$ \Rightarrow $(ii)” 设 $ \varphi \in \mathcal{D}(X) $。由于 $ K := \operatorname{supp}(\varphi) $ 紧且 $ (U_{j})_{j\in\mathbb{N}} $ 递增，存在 $ k \in \mathbb{N} $ 使得 $ K \subset U_{k} $。借助命题 7.14，我们找到 $ \psi \in \mathcal{D}(U_{k}) $ 使得 $ 0 \le \psi \le 1 $ 且 $ \psi|K = 1 $。则

\[
\int_{X}|\varphi f|\,dx = \int_{X}|\varphi\psi f|\,dx \le \|\varphi\|_{\infty}\int_{X}|\psi f|\,dx \le \|\varphi\|_{\infty}\|\chi_{U_{k}}f\|_{1} < \infty .
\]

因此 $ \varphi f $ 属于 $ L_{1}(X) $。
“(ii)$ \Rightarrow $(iii)” 取 $ K = \overline{K} \subset\subset X $ 与 $ \varphi \in \mathcal{D}(X) $ 使得 $ \varphi|K = 1 $。则

\[
\int_{K}|f| = \int_{K}|\varphi f| \le \|\varphi f\|_{1} < \infty ,
\]

因此 $ \widetilde{f|K} \in L_{1}(X) $。
“(iii)$ \Rightarrow $(i)” 这一蕴含是显然的，因为每个 $ \overline{U}_{j} $ 都紧。\nobreak\hspace{0.35em}■

\end{subitems}

\item

设 $ X $ 在 $ \mathbb{R}^{n} $ 中开。则 $ C(X) \subset L_{1,loc}(X) $。

\proofhead
取 $ f \in C(X) $ 与 $ \varphi \in \mathcal{D}(X) $。则 $ \varphi f $ 属于 $ C_{c}(X) $。由定理 5.1 有 $ \varphi f \in L_{1}(X) $，断言由 (a) 推出。\nobreak\hspace{0.35em}■
\item

群 $ (\mathbb{R}^{n},+) $ 在 $ BUC^{k} $ 中的线性表示是单射，从而是到其像 $ \mathcal{T}_{BUC^{k}} $ 上的群同构。

\proofhead
对 $ a \in \mathbb{R}^{n} $，设 $ \tau_{a} = \mathrm{id}_{BUC^{k}} $。取 $ r > |a| $ 以及关于 $ r\overline{\mathbb{B}}^{n} $ 的截断函数 $ \varphi \in \mathcal{D}(\mathbb{R}^{n}) $。则 $ f_{j} := \varphi\,\mathrm{pr}_{j} $ 属于 $ BUC^{k} $，并且我们得到
\[
-a_{j} = \tau_{a}f_{j}(0) = f_{j}(0) = 0 \quad \text{对 } j \in \{1,\ldots,n\} .
\]

因此 $ a = 0 $。这意味着表示 $ a \mapsto \tau_{a} $ 是单射。\nobreak\hspace{0.35em}■

\item

设 $ X $ 在 $ \mathbb{R}^{n} $ 中开且有界。又设 $ \{X_{j} ; 0 \le j \le m\} $ 是 $ X $ 的有限开覆盖，$ \{\varphi_{j} ; 0 \le j \le m\} $ 是从属于它的光滑单位分解。最后设 $ k \in \mathbb{N} $ 且

\[
|||u|||_{BC^{k}} := \sum_{j=0}^{m}\|\varphi_{j}u\|_{BC^{k}} \quad \text{对 } u \in BC^{k}(X) .
\]

则 $ |||\cdot|||_{BC^{k}} $ 是 $ BC^{k}(X) $ 上的等价范数。

\proofhead
取 $ u \in BC^{k}(X) $。显然
\[
\|u\|_{BC^{k}} = \left\|\sum_{j=0}^{m}\varphi_{j}u\right\|_{BC^{k}} \le \sum_{j=0}^{m}\|\varphi_{j}u\|_{BC^{k}} = |||u|||_{BC^{k}} .
\]

由莱布尼茨法则（见习题 VII.5.21），我们得到

\begin{equation*}
\begin{aligned}
|||u|||_{BC^{k}} &= \sum_{j=0}^{m}\max_{|\alpha|\le k}\|\partial^{\alpha}(\varphi_{j}u)\|_{\infty} = \sum_{j=0}^{m}\max_{|\alpha|\le k}\left\|\sum_{\beta\le\alpha}\binom{\alpha}{\beta}\partial^{\beta}\varphi_{j}\,\partial^{\alpha-\beta}u\right\|_{\infty} \\ &\le \sum_{j=0}^{m}c_{k}\|\varphi_{j}\|_{BC^{k}}\|u\|_{BC^{k}} \le C\|u\|_{BC^{k}} ,
\end{aligned}
\end{equation*}

其中我们置 $ c_{k} := \max_{|\alpha|\le k}\sum_{\beta\le\alpha}\binom{\alpha}{\beta} $ 与 $ C := c_{k}\sum_{j=0}^{m}\|\varphi_{j}\|_{BC^{k}} $。\nobreak\hspace{0.35em}■

\end{items}

\end{application}

\needspace{3\baselineskip}\subsection*{E 值函数的卷积}

回顾前面的证明可见，当两个函数之一取值于巴拿赫空间 $ F $ 而另一个为标量值时，卷积 $ f * g $ 同样可以定义。只要关于等距映射的换元法则对 $ F $ 值函数仍然成立，所有证明都可原样通过\footnote{当然，交换性公式 $ f * g = g * f $ 必须作正确的解释。}。情况确实如此，这一点我们将在下一节说明。特别地，定理 7.11 中的关键逼近结果对空间 $ L_{p}(\mathbb{R}^{n},F) $ 与 $ BUC^{k}(\mathbb{R}^{n},F) $ 仍然成立，其中 $ 1 \le p < \infty $、$ k \in \mathbb{N} $。由此得到的一个推论是定理 7.13 的类似结论：$ \mathcal{D}(X,F) $ 是 $ C_{0}(X,F) $ 与 $ L_{p}(X,F) $（$ 1 \le p < \infty $）的稠密向量子空间。

\needspace{3\baselineskip}\subsection*{分布}

设 $ X $ 是 $ \mathbb{R}^{n} $ 的非空开子集\footnote{本节其余部分意在提供应用与更深入理论的概览；初读时可以跳过。}。众所周知，$ X $ 上的标量函数是把实数或复数指派给 $ X $ 中每一点的规则。但这个定义只是抽象的说法，因为 $ X $ 的各个点在实践中无法分辨。例如，若要确定占据集合 $ X $ 的某种介质的温度分布，我们就必须依靠实验探针。然而这样的探针有一定大小，只能确定 $ f $ 在一个扩展区域上的值；它赋给 $ f(x_{0}) $ 的值并不代表 $ x_{0} $ 处的真实值（如果这种说法在物理上确实有意义），而是代表 $ x_{0} $ 附近某种平均值：在数学上就是积分 $ \int_{X}\varphi f\,dx $，其中 $ \varphi $ 是依赖于探针的“检验函数”。当然，检验函数 $ \varphi $ 越是集中在 $ x_{0} $ 附近，测量就越好地逼近精确值 $ f(x_{0}) $，也就是说，探针把数据抹得越少。

若要声称完全了解 $ f(x_{0}) $，人们可以想象动用一切可能的探针，换言之，确定 $ \int_{X}\varphi f\,dx $ 对所有可能的检验函数 $ \varphi $ 的取值。用数学语言来说，我们就是用定义在所有检验函数构成的空间上的一个泛函来代替逐点函数 $ f : X \to \mathbb{K} $，即映射

\begin{equation}
T_{f} : \mathcal{D}(X) \to \mathbb{K} , \quad \varphi \mapsto \int_{X}\varphi f\,dx .\tag{7.12}
\end{equation}

我们选取 $ \mathcal{D}(X) $ 作为检验函数空间在很大程度上是任意的。为概念上的简单起见，我们也许想用 $ C_{c}(X) $ 代替 $ \mathcal{D}(X) $。同时，我们还想避免进行比必要的更多的“测量”；这就要求选一个在某种意义下小的检验空间。但这个空间必须足够大，使得平均值 $ \int_{X}\varphi f\,dx $ 确实能确定 $ f $。也就是说，我们希望：对所有检验函数 $ \varphi $，$ \int_{X}\varphi f\,dx $ 与 $ \int_{X}\varphi g\,dx $ 相等，就能推出 $ f = g $。

下一个定理表明，若我们选取 $ \mathcal{D}(X) $ 作为检验空间并在 $ L_{1,loc}(X) $ 中工作，情况确实如此。（由应用 7.17(a)，$ L_{1,loc}(X) $ 是 $ L_{0}(X) $ 的最大的向量子空间 $ E $，使得 $ \int_{X}\varphi f\,dx $ 对所有 $ f \in E $ 与所有 $ \varphi \in \mathcal{D}(X) $ 都有定义。）

\begin{theorem}

设 $ f \in L_{1,loc}(X) $。若

\begin{equation}
\int_{X}\varphi f\,dx = 0 \quad \text{对 } \varphi \in \mathcal{D}(X) ,\tag{7.13}
\end{equation}

则 $ f = 0 $。

\proofhead
设 $ f \ne 0 $，并设 $ f^{\ast} \in \mathcal{L}_{1,loc}(X) $ 是 $ f $ 的一个代表。由正则性，存在 $ X $ 的具有正测度的紧子集 $ K $，使得对 $ x \in K $ 有 $ f^{\ast}(x) \ne 0 $。取 $ \eta \in \mathcal{D}(X) $ 使得 $ \eta|K = 1 $，并置 $ g := \eta f^{\ast} $。由应用 7.17(a)，$ g $ 属于 $ \mathcal{L}_{1} $。又对 $ x \in K $ 有 $ g(x) \ne 0 $。设 $ \{\varphi_{\varepsilon} ; \varepsilon > 0\} $ 是光滑核。则 $ \lim_{\varepsilon \to 0}\varphi_{\varepsilon} * g = g $ 在 $ \mathcal{L}_{1} $ 中成立。由推论 4.7，存在零序列 $ (\varepsilon_{j}) $ 与勒贝格测度为零的集合 $ N $，使得
\begin{equation}
\lim_{j \to \infty}\varphi_{\varepsilon_{j}} * g(x) = g(x) \quad \text{对 } x \in N^{c} .\tag{7.14}
\end{equation}

给定 $ x_{0} \in K \cap N^{c} $，置 $ \psi_{j} := \eta\tau_{x_{0}}\varphi_{\varepsilon_{j}} \in \mathcal{D}(X) $（$ j \in \mathbb{N} $）。由于由注 7.1(d) 有 $ \check{\varphi}_{\varepsilon} = \varphi_{\varepsilon} $，等式 (7.13) 给出

\begin{equation*}
\begin{aligned}
\varphi_{\varepsilon_{j}} * g(x_{0}) &= \int g(y)\varphi_{\varepsilon_{j}}(x_{0} - y)\,dy = \int_{X}(\eta f^{\ast})(y)\varphi_{\varepsilon_{j}}(x_{0} - y)\,dy \\ &= \int_{X}f^{\ast}(y)\psi_{j}(y)\,dy = 0 .
\end{aligned}
\end{equation*}

然而由 (7.14)，这与 $ g(x_{0}) \ne 0 $ 矛盾。断言成立，因为 $ f $ 的代表 $ f^{\ast} $ 是任意选取的。\nobreak\hspace{0.35em}■

\end{theorem}

显然映射 $ T_{f} $ 是 $ \mathcal{D}(X) $ 上的线性泛函。要使 $ T_{f}\varphi = \int_{X}\varphi f\,dx $ 作为一种测量值有意义，$ T_{f}\varphi $ 必须“连续地”依赖于测量装置；也就是说，探针的小扰动（从而检验函数 $ \varphi $ 的小扰动）只应引起测量值的微小变化。用数学语言来说，这意味着 $ T_{f} $ 必须是 $ \mathcal{D}(X) $ 上的连续线性泛函。因此我们必须在 $ \mathcal{D}(X) $ 上引入拓扑。

由于这里的处理是引论性的，我们只限于陈述序列在 $ \mathcal{D}(X) $ 中收敛是什么意思。这种收敛应当与 $ \mathcal{D}(X) $ 上的向量结构相容，因此只需考虑极限为 0 的情形。

我们说序列 $ (\varphi_{j}) $ 在 $ \mathcal{D}(X) $ 中\textbf{收敛到 0}（或\textbf{是零序列}），是指满足下列条件：

($ \mathcal{D}_{1} $) 存在 $ K \subset\subset X $ 使得 $ \operatorname{supp}(\varphi_{j}) \subset K $ 对 $ j \in \mathbb{N} $ 成立。

($ \mathcal{D}_{2} $) 对每个 $ k \in \mathbb{N} $ 有 $ \varphi_{j} \to 0 $ 于 $ BC^{k}(X) $。

显然 ($ \mathcal{D}_{2} $) 等价于：

\begin{equation}
\text{序列 } (\partial^{\alpha}\varphi_{j})_{j\in\mathbb{N}} \text{ 对每个 } \alpha \in \mathbb{N}^{n} \text{ 一致收敛到 } 0.\tag{7.15}
\end{equation}

因此，为使 $ \varphi_{j} $ 在 $ \mathcal{D}(X) $ 中收敛到 0，(7.15) 必须成立，而且函数 $ \varphi_{j} $ 的支集还必须都含于 $ X $ 的一个固定紧子集之中。

线性泛函 $ T : \mathcal{D}(X) \to \mathbb{K} $ 称为\textbf{连续的}，若对 $ \mathcal{D}(X) $ 中的每个零序列 $ (\varphi_{j}) $ 都有 $ T\varphi_{j} \to 0 $。$ \mathcal{D}(X) $ 上的连续线性泛函也称为 $ X $ 上的 \textbf{Schwartz 分布}，或简称\textbf{分布}。$ X $ 上分布的集合记作 $ \mathcal{D}'(X) $；它显然是 $ \operatorname{Hom}(\mathcal{D}(X),\mathbb{K}) $ 的向量子空间。

（在泛函分析中——更确切地说，在拓扑向量空间理论中——人们证明 $ \mathcal{D}(X) $ 上恰有一个豪斯多夫拓扑，它局部凸\footnote{这是说原点有一个由凸集组成的开邻域基。}、与向量结构相容，并且序列按上述意义收敛到 0 当且仅当它们按该拓扑收敛到 0。关于这个拓扑，$ \mathcal{D}'(X) $ 是 $ \mathcal{D}(X) $ 的对偶空间，即 $ \mathcal{D}(X) $ 上所有连续线性泛函构成的空间。例如可参看 [Sch66] 或 [Yos65]。）

\begin{example}

\begin{items}
\item

对每个 $ f \in L_{1,loc}(X) $，由 (7.12) 定义的线性泛函 $ T_{f} $ 是 $ X $ 上的分布。

\proofhead
设 $ (\varphi_{j}) $ 是 $ \mathcal{D}(X) $ 中的序列，使得 $ \varphi_{j} \to 0 $ 于 $ \mathcal{D}(X) $。则存在 $ X $ 的紧子集 $ K $ 使得 $ \operatorname{supp}(\varphi_{j}) \subset X $ 对 $ j \in \mathbb{N} $ 成立。推出
\[
|T_{f}\varphi_{j}| = \left|\int_{X}\varphi_{j}f\,dx\right| \le \int_{K}|\varphi_{j}|\,|f|\,dx \le \|f\|_{L_{1}(K)}\|\varphi_{j}\|_{\infty}
\]

对 $ j \in \mathbb{N} $ 成立。因为 $ \|f\|_{L_{1}(K)} < \infty $，我们得到 $ T_{f}\varphi_{j} \to 0 $ 于 $ \mathbb{K} $，因为 ($ \mathcal{D}_{2} $) 蕴含 $ \|\varphi_{j}\|_{\infty} \to 0 $。\nobreak\hspace{0.35em}■

\item

设 $ \mu $ 是 $ X $ 上的拉东测度。则

\[
\mathcal{D}(X) \to \mathbb{K} , \quad \varphi \mapsto \int_{X}\varphi\,d\mu
\]

定义了 $ X $ 上的一个分布。

\proofhead
设 $ (\varphi_{j}) $ 是 $ \mathcal{D}(X) $ 中的序列，使得 $ \varphi_{j} \to 0 $ 于 $ \mathcal{D}(X) $。又设 $ K = \overline{K} \subset\subset X $ 包含所有 $ j \in \mathbb{N} $ 的 $ \operatorname{supp}(\varphi_{j}) $。则
\[
\left|\int_{X}\varphi_{j}\,d\mu\right| \le \int_{K}|\varphi_{j}|\,d\mu \le \mu(K)\|\varphi_{j}\|_{\infty} \quad \text{对 } j \in \mathbb{N} .
\]

如 (a) 的证明中那样，这说明 $ \varphi $ 是 $ X $ 上的分布。\nobreak\hspace{0.35em}■

\item

设 $ \delta $ 是 $ \mathbb{R}^{n} $ 上支集在 0 处的狄拉克测度。则

\[
\varphi \mapsto \langle\delta,\varphi\rangle := \int_{X}\varphi\,d\delta = \varphi(0) \quad \text{对 } \varphi \in \mathcal{D}(\mathbb{R}^{n}) ,
\]

是 $ \mathbb{R}^{n} $ 上的一个分布，即\textbf{狄拉克分布}

\[
\delta : \mathcal{D}(\mathbb{R}^{n}) \to \mathbb{K} , \quad \varphi \mapsto \varphi(0) .
\]

不存在 $ u \in L_{1,loc}(\mathbb{R}^{n}) $ 使得 $ T_{u} = \delta $。

\proofhead
第一个断言是 (b) 的特例。
现在设 $ u \in L_{1,loc}(\mathbb{R}^{n}) $ 且 $ T_{u} = \delta $，即
\begin{equation}
\int_{\mathbb{R}^{n}}\varphi u\,dx = \varphi(0) \quad \text{对 } \varphi \in \mathcal{D}(\mathbb{R}^{n}) .\tag{7.16}
\end{equation}

只取使 $ \operatorname{supp}(\varphi) \subset\subset X := \mathbb{R}^{n} \setminus \{0\} $ 的 $ \varphi \in \mathcal{D}(\mathbb{R}^{n}) $，则 $ \varphi(0) = 0 $，于是由定理 7.18 推出 $ u|X = 0 $ 于 $ L_{1,loc}(X) $。但 $ X $ 与 $ \mathbb{R}^{n} $ 只在一个测度为零的集合上不同（单个点！），所以 $ u = 0 $ 于 $ L_{1,loc}(\mathbb{R}^{n}) $，这与 (7.16) 矛盾。\nobreak\hspace{0.35em}■

\item

设 $ \alpha \in \mathbb{N}^{n} $。则

\[
S_{\alpha} : \mathcal{D}(\mathbb{R}^{n}) \to \mathbb{K} , \quad \varphi \mapsto \partial^{\alpha}\varphi(0)
\]

定义一个分布。不存在 $ u \in L_{1,loc}(\mathbb{R}^{n}) $ 使得 $ T_{u} = S_{\alpha} $。

\proofhead
设 $ (\varphi_{j}) $ 是 $ \mathcal{D}(\mathbb{R}^{n}) $ 中的序列，使得 $ \varphi_{j} \to 0 $ 于 $ \mathcal{D}(\mathbb{R}^{n}) $，并设 $ K = \overline{K} \subset\subset \mathbb{R}^{n} $ 满足 $ \operatorname{supp}(\varphi_{j}) \subset K $ 对 $ j \in \mathbb{N} $ 成立。我们可以假定 0 位于 $ K $ 中。则我们有估计
\[
|\partial^{\alpha}\varphi_{j}(0)| \le \max_{x\in K}|\partial^{\alpha}\varphi_{j}(x)| \le \|\varphi_{j}\|_{BC^{|\alpha|}} \quad \text{对 } j \in \mathbb{N} .
\]

于是 ($ \mathcal{D}_{2} $) 蕴含 $ \partial^{\alpha}\varphi_{j}(0) \to 0 $ 于 $ \mathbb{K} $，这说明 $ S_{\alpha} $ 是一个分布。第二个断言的证明与 (c) 相同。\nobreak\hspace{0.35em}■

\end{items}

下述关键结果现在是定理 7.18 的一个简单推论。

\end{example}

\begin{theorem}

映射

\[
L_{1,loc}(X) \to \mathcal{D}'(X) , \quad f \mapsto T_{f}
\]

是线性的且为单射。

\proofhead
例 7.19(a) 表明该映射有定义。它是线性的，因为积分是线性的。由定理 7.18 它是单射。\nobreak\hspace{0.35em}■
\end{theorem}

由定理 7.20，我们可以把 $ L_{1,loc}(X) $ 与其在 $ \mathcal{D}'(X) $ 中的像等同起来。换言之，我们可以把 $ L_{1,loc}(X) $ 看作所有 Schwartz 分布构成的空间的向量子空间，办法是把函数 $ f \in L_{1,loc}(X) $ 与分布

\[
T_{f} = \left(\varphi \mapsto \int_{X}\varphi f\,dx\right) \in \mathcal{D}'(X)
\]

等同起来。

在这个意义上，每个 $ f \in L_{1,loc}(X) $ 都是一个分布。$ L_{1,loc}(X) $ 的元素称为\textbf{正则分布}。所有其他分布都是\textbf{奇异的}。例 7.19(c) 与 (d) 说明了奇异分布。

分布理论在高等分析中起着重要作用，特别是在偏微分方程的研究中，以及在理论物理中。我们无法在此详述，但例如可参看 [Sch65]、[RS72]。

\needspace{3\baselineskip}\subsection*{线性微分算子}

设 $ X $ 在 $ \mathbb{R}^{n} $ 中开。给定函数 $ a_{\alpha} \in C^{\infty}(X) $，对每个满足 $ |\alpha| \le m \in \mathbb{N} $ 的 $ \alpha \in \mathbb{N}^{n} $，我们置

\[
\mathcal{A}(\partial)u := \sum_{|\alpha|\le m}a_{\alpha}\partial^{\alpha}u \quad \text{对 } u \in \mathcal{D}(X) .
\]

显然 $ \mathcal{A}(\partial) $ 是 $ \mathcal{D}(X) $ 到自身的线性映射；我们称它是 $ X $ 上的\textbf{线性微分算子}，其\textbf{阶不超过 $ m $}（具有光滑系数）。若

\[
\sum_{|\alpha|=m}\|a_{\alpha}\|_{\infty} \ne 0 ,
\]

即 $ \mathcal{A}(\partial) $ 的\textbf{主部} $ \sum_{|\alpha|=m}a_{\alpha}\partial^{\alpha} $ 中至少有一个系数 $ a_{\alpha} $ 不恒为零，则说它具有\textbf{阶 $ m $}。我们用 $ \mathcal{D}\mathrm{iffop}(X) $ 表示 $ X $ 上所有线性微分算子的集合；阶不超过 $ m $ 的那些记作 $ \mathcal{D}\mathrm{iffop}_{m}(X) $。

线性映射 $ T : \mathcal{D}(X) \to \mathcal{D}(X) $ 称为\textbf{连续的}\footnote{泛函分析中证明，这些定义与我们先前关于连续性和 $ \mathcal{L}(E) $ 的定义是相容的。}，若对 $ \mathcal{D}(X) $ 中每个满足 $ \varphi_{j} \to 0 $ 于 $ \mathcal{D}(X) $ 的序列 $ (\varphi_{j}) $ 都有 $ T\varphi_{j} \to 0 $ 于 $ \mathcal{D}(X) $。$ \mathcal{D}(X) $ 的所有连续自同态构成的集合是 $ \operatorname{End}(\mathcal{D}(X)) $ 的一个向量子空间，我们把它记作 $ \mathcal{L}(\mathcal{D}(X)) $。

\begin{proposition}

$ \mathcal{D}\mathrm{iffop}(X) $ 是 $ \mathcal{L}(\mathcal{D}(X)) $ 的向量子空间，且 $ \mathcal{D}\mathrm{iffop}_{m}(X) $ 是 $ \mathcal{D}\mathrm{iffop}(X) $ 的向量子空间。

\proofhead
设 $ m \in \mathbb{N} $，取 $ \mathcal{A}(\partial) := \sum_{|\alpha|\le m}a_{\alpha}\partial^{\alpha} \in \mathcal{D}\mathrm{iffop}_{m}(X) $。设 $ (\varphi_{j}) $ 是 $ \mathcal{D}(X) $ 中的零序列，并设 $ K = \overline{K} \subset\subset X $ 包含所有 $ j \in \mathbb{N} $ 的 $ \operatorname{supp}(\varphi_{j}) $。则 $ \operatorname{supp}(\mathcal{A}(\partial)\varphi_{j}) \subset K $ 对 $ j \in \mathbb{N} $ 成立。对 $ \beta \in \mathbb{N}^{n} $，莱布尼茨法则给出
\[
\|\partial^{\beta}(a_{\alpha}\partial^{\alpha}\varphi_{j})\|_{C(K)} = \left\|\sum_{\gamma\le\beta}\binom{\beta}{\gamma}\partial^{\gamma}a_{\alpha}\,\partial^{\beta-\gamma+\alpha}\varphi_{j}\right\|_{C(K)} \le c(\alpha,\beta)\max_{\gamma\le\beta}\|\partial^{\gamma}a_{\alpha}\|_{C(K)}\|\partial^{\beta-\gamma+\alpha}\varphi_{j}\|_{\infty} .
\]

由此对 $ k \in \mathbb{N} $ 我们导出不等式

\[
\|\mathcal{A}(\partial)\varphi_{j}\|_{BC^{k}(X)} \le c(k)\sum_{|\alpha|\le m}\|a_{\alpha}\|_{BC^{k}(K)}\|\varphi_{j}\|_{BC^{k+m}(X)} \quad \text{对 } j \in \mathbb{N} ,
\]

其中常数 $ c(k) $ 与 $ j $ 无关。现在由 $ (D_2) $ 可得 $ \mathcal{A}(\partial)\varphi_j \to 0 $ 在 $ BC^k(X) $ 中成立。因为这对每个 $ k \in \mathbb{N} $ 都成立，我们看到 $ \mathcal{A}(\partial)\varphi_j \to 0 $ 在 $ \mathcal{D}(X) $ 中成立。这就证明了 $ \mathcal{D}\mathrm{iffop}_m(X) \subset \mathcal{L}(\mathcal{D}(X)) $。其余断言是显然的。 \nobreak\hspace{0.35em}■

\end{proposition}

设 $ (\cdot \mid \cdot) $ 表示 $ L_2(X) $ 中的内积，并设 $ \mathcal{A}(\partial) $ 属于 $ \mathcal{D}\mathrm{iffop}(X) $。若存在微分算子 $ \mathcal{A}^{\sharp}(\partial) \in \mathcal{D}\mathrm{iffop}(X) $ 使得

\[
(\mathcal{A}(\partial)u \mid v) = (u \mid \mathcal{A}^{\sharp}(\partial)v) \qquad \text{对 } u, v \in \mathcal{D}(X) ,
\]

我们就说 $ \mathcal{A}^{\sharp}(\partial) $ 是 $ \mathcal{A}(\partial) $ 的 \textbf{形式伴随}。因为

\[
(u \mid \mathcal{A}^{\sharp}(\partial)v) = \int_X u\,\overline{\mathcal{A}^{\sharp}(\partial)v}\, dx
\]

且对 $ v \in \mathcal{D}(X) $ 有 $ \overline{\mathcal{A}^{\sharp}(\partial)v} \in \mathcal{D}(X) \subset L_{1,\mathrm{loc}}(X) $，故由定理 7.18 容易推出 $ \mathcal{A}(\partial) $ 至多有一个形式伴随。若 $ \mathcal{A}(\partial) $ 有一个形式伴随 $ \mathcal{A}^{\sharp}(\partial) $ 与它自身相合，则称 $ \mathcal{A}(\partial) $ 是 \textbf{形式自伴的}。

下面我们证明每个 $ \mathcal{A}(\partial) \in \mathcal{D}\mathrm{iffop}(X) $ 都有与它形式伴随的微分算子，并导出 $ \mathcal{A}^{\sharp}(\partial) $ 的显式形式。为此先需要下述结果：

\begin{proposition}

（分部积分） 设 $ f \in C^1(X) $ 且 $ g \in C_c^1(X) $，则

\[
\int_X (\partial_j f)g\, dx = -\int_X f\partial_j g\, dx \qquad \text{对 } j \in \{1,\dots,n\} .
\]

\proofhead
我们只需考虑 $ j = 1 $ 的情形；由推论 6.10，一般情形可由坐标的置换得到。于是写 $ x = (x_1, x') \in \mathbb{R} \times \mathbb{R}^{n-1} $。由于 $ fg $ 具有紧支集，由分部积分可得
\[
\int_{-\infty}^{\infty} [\partial_1 f(x_1, x')]g(x_1, x')\, dx_1 = -\int_{-\infty}^{\infty} f(x_1, x')\partial_1 g(x_1, x')\, dx_1
\]

对每个 $ x' \in \mathbb{R}^{n-1} $ 成立。由富比尼定理，现在得到

\begin{equation*}
\begin{aligned}
\int_X (\partial_1 f)g\, dx &= \int_{\mathbb{R}^n} (\partial_1 f)g\, dx \\ &= \int_{\mathbb{R}^{n-1}} \left( \int_{-\infty}^{\infty} \partial_1 f(x_1, x')g(x_1, x')\, dx_1 \right) dx' \\ &= -\int_{\mathbb{R}^{n-1}} \left( \int_{-\infty}^{\infty} f(x_1, x')\partial_1 g(x_1, x')\, dx_1 \right) dx' \\ &= -\int_{\mathbb{R}^n} f\partial_1 g\, dx = -\int_X f\partial_1 g\, dx . \quad \blacksquare
\end{aligned}
\end{equation*}
\end{proposition}

\begin{corollary}

设 $ f \in C^k(X) $ 且 $ g \in C_c^k(X) $，则

\[
\int_X (\partial^{\alpha} f)g\, dx = (-1)^{|\alpha|} \int_X f\partial^{\alpha} g\, dx
\]

对 $ \alpha \in \mathbb{N}^n $ 且 $ |\alpha| \le k $ 成立。

分部积分也是下一个结果的证明的核心。

\end{corollary}

\begin{proposition}

每个微分算子

\[
\mathcal{A}(\partial) = \sum_{|\alpha| \le m} a_\alpha \partial^{\alpha} \in \mathcal{D}\mathrm{iffop}(X)
\]

都有唯一的形式伴随，它由下式显式给出：

\begin{equation}
\mathcal{A}^{\sharp}(\partial)v = \sum_{|\alpha| \le m} (-1)^{|\alpha|}\partial^{\alpha}(\overline{a_\alpha}v) \qquad \text{对 } v \in \mathcal{D}(X) .\tag{7.17}
\end{equation}

若 $ \mathcal{A}(\partial) $ 的阶为 $ m $，则 $ \mathcal{A}^{\sharp}(\partial) $ 也是 $ m $ 阶微分算子。

\proofhead
我们已经知道形式伴随至多有一个，故只需证明存在性以及 (7.17) 的正确性。

取 $ u, v \in \mathcal{D}(X) $。分部积分给出
\begin{equation*}
\begin{aligned}
(\mathcal{A}(\partial)u \mid v) &= \int_X (\mathcal{A}(\partial)u)\overline{v}\, dx = \sum_{|\alpha| \le m} \int_X (a_\alpha \partial^{\alpha} u)\overline{v}\, dx \\ &= \sum_{|\alpha| \le m} (-1)^{|\alpha|} \int_X u\partial^{\alpha}(a_\alpha \overline{v})\, dx = \int_X u\,\overline{\sum_{|\alpha| \le m} (-1)^{|\alpha|}\partial^{\alpha}(\overline{a_\alpha}v)}\, dx .
\end{aligned}
\end{equation*}

因此

\[
(\mathcal{A}(\partial)u \mid v) = (u \mid \mathcal{A}^{\sharp}(\partial)v) \qquad \text{对 } u, v \in \mathcal{D}(X)
\]

只要 $ \mathcal{A}^{\sharp}(\partial)v $ 如 (7.17) 所定义。由莱布尼茨法则，存在 $ b_\alpha \in C^\infty(X) $（$ \alpha \in \mathbb{N}^n $，$ |\alpha| \le m-1 $）使得

\[
\mathcal{A}^{\sharp}(\partial) = (-1)^m \sum_{|\alpha| = m} \overline{a_\alpha}\partial^{\alpha} + \sum_{|\alpha| \le m-1} b_\alpha \partial^{\alpha} .
\]

因此 $ \mathcal{A}^{\sharp}(\partial) $ 属于 $ \mathcal{D}\mathrm{iffop}(X) $。断言至此显然。 \nobreak\hspace{0.35em}■

\end{proposition}

对于描述系统时间演化的微分算子，通常把时间当作一个特殊的变量。例如在变量 $ (t, x) \in \mathbb{R} \times \mathbb{R}^n $ 中，我们回忆波动算子 $ \partial_t^2 - \Delta_x $ 与热算子 $ \partial_t - \Delta_x $（见习题 VII.5.10）。另一个例子是\textbf{薛定谔算子} $ (1/i)\partial_t - \Delta_x $。这三个算子都是二阶微分算子。

\begin{example}

\begin{items}
\item

波动算子与薛定谔算子是形式自伴的。\footnote{这些事实在数学物理中特别重要。}

\item

热算子的伴随是 $ (\partial_t - \Delta_x)^{\sharp} = -\partial_t - \Delta_x $。因此它不是形式自伴的。

\item

对 $ \mathcal{A}(\partial) := \partial_t - \sum_{j=1}^n \partial_j $，有 $ \mathcal{A}^{\sharp}(\partial) = -\mathcal{A}(\partial) $。

\item

设 $ a_{jk}, a_j, a_0 \in C^\infty(X, \mathbb{R}) $ 满足

\[
\sum_{j,k=1}^n \|a_{jk}\|_\infty \ne 0 , \qquad a_{jk} = a_{kj} \qquad \text{对 } j, k \in \{1,\dots,n\} .
\]

再定义 $ \mathcal{A}(\partial) \in \mathcal{D}\mathrm{iffop}_2(X) $ 为

\[
\mathcal{A}(\partial)u := \sum_{j,k=1}^n \partial_j(a_{jk}\partial_k u) + \sum_{j=1}^n a_j\partial_j u + a_0 u \qquad \text{对 } u \in \mathcal{D}(X) .
\]

这时我们说 $ \mathcal{A}(\partial) $ 是\textbf{散度型}算子。\footnote{这种说法的理由将在第 XI.6 节中说明。} 在此情形下，我们有

\[
\mathcal{A}^{\sharp}(\partial)v = \sum_{j,k=1}^n \partial_j(a_{jk}\partial_k v) - \sum_{j=1}^n a_j\partial_j v + \left( a_0 - \sum_{j=1}^n \partial_j a_j \right) v \qquad \text{对 } v \in \mathcal{D}(X) .
\]

因此形式伴随也是散度型的，并且 $ \mathcal{A}(\partial) $ 是形式自伴的当且仅当 $ a_j = 0 $（$ j = 1,\dots,n $）。

\proofhead
这由命题 7.22 容易推出。 \nobreak\hspace{0.35em}■
\item

拉普拉斯算子 $ \Delta $ 是形式自伴的二阶散度型微分算子。

\proofhead
在 (d) 中取 $ a_{jk} = \delta_{jk} $（克罗内克 δ）即得。 \nobreak\hspace{0.35em}■
\end{items}

\end{example}

\needspace{3\baselineskip}\subsection*{弱导数}

现在我们简要说明导数概念可以如何推广，使得在经典意义下不可微的函数也能被赋予某种广义导数。

设 $ X $ 是 $ \mathbb{R}^n $ 中的开集。我们说 $ u \in L_{1,\mathrm{loc}}(X) $ 是\textbf{弱可微的}，若存在 $ u_j \in L_{1,\mathrm{loc}}(X) $ 使得

\begin{equation}
\int_X (\partial_j \varphi)u\, dx = -\int_X \varphi u_j\, dx \qquad \text{对 } \varphi \in \mathcal{D}(X) \text{ 且 } 1 \le j \le n .\tag{7.18}
\end{equation}

更一般地，若 $ m \ge 2 $ 是整数，我们说 $ u \in L_{1,\mathrm{loc}}(X) $ 在 $ X $ 上是 \textbf{m 次弱可微的}，若存在 $ u_\alpha \in L_{1,\mathrm{loc}}(X) $ 使得

\begin{equation}
\int_X (\partial^{\alpha}\varphi)u\, dx = (-1)^{|\alpha|} \int_X \varphi u_\alpha\, dx \qquad \text{对 } \varphi \in \mathcal{D}(X) ,\tag{7.19}
\end{equation}

对所有满足 $ |\alpha| \le m $ 的 $ \alpha \in \mathbb{N}^n $ 成立。若确是如此，则由定理 7.18 立即推出 $ u_\alpha \in L_{1,\mathrm{loc}}(X) $ 由 $ u $（以及 $ \alpha $）唯一确定。我们称 $ u_\alpha $ 为 \textbf{α 阶弱偏导数}，并记 $ \partial^{\alpha} u := u_\alpha $。在 $ m = 1 $ 的情形，我们记 $ \partial_j u := u_j $。这些记号由下面第一条注即可说明其合理性。

\begin{remark}

\begin{items}
\item

设 $ m \in \mathbb{N}^{\times} $。则每个 $ u \in C^m(X) $ 都是 $ m $ 次弱可微的，且弱导数与经典（或通常）的偏导数一致。

\proofhead
这由推论 7.23 推出。 \nobreak\hspace{0.35em}■
\item

设 $ W_{1,\mathrm{loc}}^m(X) $ 为 $ X $ 上所有 $ m $ 次弱可微函数之集。则 $ W_{1,\mathrm{loc}}^m(X) $ 是 $ L_{1,\mathrm{loc}}(X) $ 的向量子空间，且对每个满足 $ |\alpha| \le m $ 的 $ \alpha \in \mathbb{N}^n $，映射

\[
W_{1,\mathrm{loc}}^m(X) \to W_{1,\mathrm{loc}}^{m-|\alpha|}(X) , \qquad u \mapsto \partial^{\alpha} u
\]

是定义良好的并且是线性的。

\proofhead
这个简单的证明留给读者作为练习。 \nobreak\hspace{0.35em}■
\item

对 $ u \in W_{1,\mathrm{loc}}^m(X) $ 与满足 $ |\alpha| + |\beta| \le m $ 的 $ \alpha, \beta \in \mathbb{N}^n $，有 $ \partial^{\alpha}\partial^{\beta} u = \partial^{\beta}\partial^{\alpha} u $。

\proofhead
这立即由定义方程 (7.19) 以及光滑函数的性质推出。 \nobreak\hspace{0.35em}■
\item

设 $ u \in L_{1,\mathrm{loc}}(\mathbb{R}) $ 由 $ u(x) := |x| $（$ x \in \mathbb{R} $）定义。则 $ u $ 是弱可微的，且 $ \partial u = \mathrm{sign} $。

\proofhead
首先，绝对值函数 $ |\cdot| $ 在 $ \mathbb{R}^{\times} $ 上光滑，其导数在 $ \mathbb{R}^{\times} $ 上为 $ \mathrm{sign} $。现在设 $ \varphi \in \mathcal{D}(\mathbb{R}) $。分部积分给出
\begin{equation*}
\begin{aligned}
\int_{\mathbb{R}} \varphi' u\, dx &= \int_0^{\infty} \varphi' u\, dx + \int_{-\infty}^0 \varphi' u\, dx \\ &= \varphi(x)x\Big|_0^{\infty} - \int_0^{\infty} \varphi(x)\, dx - \varphi(x)x\Big|_{-\infty}^0 + \int_{-\infty}^0 \varphi(x)\, dx \\ &= -\int_{\mathbb{R}} \varphi(x)\,\mathrm{sign}(x)\, dx .
\end{aligned}
\end{equation*}

断言成立，因为 $ \mathrm{sign} $ 属于 $ L_{1,\mathrm{loc}}(\mathbb{R}) $。 \nobreak\hspace{0.35em}■

\item

函数 $ \mathrm{sign} $ 属于 $ L_{1,\mathrm{loc}}(\mathbb{R}) $，且在 $ \mathbb{R}^{\times} $ 上光滑。然而它不是弱可微的。因此 (d) 中的绝对值函数不是二次弱可微的。

\proofhead
对 $ \varphi \in \mathcal{D}(\mathbb{R}) $，我们有
\begin{equation}
\int_{\mathbb{R}} \varphi'\,\mathrm{sign}\, dx = \int_0^{\infty} \varphi'(x)\, dx - \int_{-\infty}^0 \varphi'(x)\, dx = -2\varphi(0) .\tag{7.20}
\end{equation}

若 $ \mathrm{sign} $ 是弱可微的，则存在 $ v \in L_{1,\mathrm{loc}}(\mathbb{R}) $ 使得

\[
\int_{\mathbb{R}} \varphi v\, dx = 2\varphi(0) \qquad \text{对 } \varphi \in \mathcal{D}(\mathbb{R}) ,
\]

而这是不成立的：见例 7.19(c)。 \nobreak\hspace{0.35em}■

\end{items}

用狄拉克分布 $ \delta $ 来说，(7.20) 取下述形式：

\[
\int_{\mathbb{R}} \varphi'\,\mathrm{sign}\, dx = -2\delta(\varphi) \qquad \text{对 } \varphi \in \mathcal{D}(X) .
\]

按通常方式把对偶配对记为

\[
\langle \cdot , \cdot \rangle : \mathcal{D}'(X) \times \mathcal{D}(X) \to \mathbb{K} ,
\]

于是 $ \langle T, \varphi \rangle $ 是连续线性泛函 $ T $ 在元素 $ \varphi $ 处的值；我们有

\begin{equation}
\langle \mathrm{sign}, \varphi' \rangle = -2\langle \delta, \varphi \rangle \qquad \text{对 } \varphi \in \mathcal{D}(\mathbb{R}) ,\tag{7.21}
\end{equation}

这里我们把 $ \mathrm{sign} \in L_{1,\mathrm{loc}}(\mathbb{R}) $ 与正则分布 $ T_{\mathrm{sign}} \in \mathcal{D}'(X) $ 等同起来，正如定理 7.20 的证明之后所讨论的那样。比较 (7.19) 与 (7.21) 提示下述定义：设 $ S, T \in \mathcal{D}'(X) $，$ \alpha \in \mathbb{N}^n $。若

\[
\langle T, \partial^{\alpha}\varphi \rangle = (-1)^{|\alpha|}\langle S, \varphi \rangle \qquad \text{对 } \varphi \in \mathcal{D}(X) ,
\]

则称 $ S $ 是 $ T $ 的 \textbf{α 阶分布导数}。

在此情形下，$ S $ 显然由 $ T $（以及 $ \alpha $）确定，故可记 $ \partial^{\alpha} T := S $。我们容易看出，每个分布都有任意阶的分布导数，且对每个 $ \alpha \in \mathbb{N}^n $，分布导数

\[
\partial^{\alpha} : \mathcal{D}'(X) \to \mathcal{D}'(X) , \qquad T \mapsto \partial^{\alpha} T
\]

是线性映射。\footnote{见习题 13。} 特别地，(7.21) 表明，在分布的意义下，

\[
\partial(\mathrm{sign}) = 2\delta .
\]

关于分布理论我们在此不能再深入，但我们想简要介绍索伯列夫空间。设 $ m \in \mathbb{N} $ 且 $ 1 \le p \le \infty $。因为 $ L_p(X) \subset L_{1,\mathrm{loc}}(X) $，所以每个 $ u \in L_p(X) $ 都有任意阶的分布导数。我们令\footnote{若 $ X $ 是 $ \mathbb{R} $ 中的区间，则可以证明 $ W_1^1(X) $ 与习题 5.6 中引入的空间一致。}

\[
W_p^m(X) := \{ u \in L_p(X) ; \partial^{\alpha} u \in L_p(X), |\alpha| \le m \} ,
\]

其中 $ \partial^{\alpha} $ 表示 $ \alpha $ 阶分布导数。再令

\begin{equation}
\|u\|_{m,p} := \begin{cases} \left( \sum_{|\alpha| \le m} \|\partial^{\alpha} u\|_p^p \right)^{1/p} , & 1 \le p < \infty , \\ \max_{|\alpha| \le m} \|\partial^{\alpha} u\|_\infty , & p = \infty . \end{cases}\tag{7.22}
\end{equation}

我们容易验证

\[
W_p^m(X) := (W_p^m(X), \|\cdot\|_{m,p})
\]

是赋范向量空间，称为 $ m $ 阶\textbf{索伯列夫空间}。特别地，$ W_p^0(X) = L_p(X) $。

\end{remark}

\begin{theorem}

\begin{subitems}
\item

$ W_p^m(X) $ 连续嵌入 $ L_p(X) $ 中，且 $ u \in L_p(X) $ 属于 $ W_p^m(X) $ 当且仅当 $ u $ 是 $ m $ 次弱可微的并且所有阶数 $ \le m $ 的弱导数都属于 $ L_p(X) $。

\item

$ W_p^m(X) $ 是巴拿赫空间。

\proofhead
(i) 这是显然的。
\end{subitems}

(ii) 设 $ (u_j) $ 是 $ W_p^m(X) $ 中的柯西序列。由 (7.22) 立即推出，对每个满足 $ |\alpha| \le m $ 的 $ \alpha \in \mathbb{N}^n $，$ (\partial^{\alpha} u_j)_{j \in \mathbb{N}} $ 是 $ L_p(X) $ 中的柯西序列。由于 $ L_p(X) $ 是完备的，存在唯一的 $ u_\alpha \in L_p(X) $ 使得当 $ j \to \infty $ 时 $ \partial^{\alpha} u_j \to u_\alpha $ 在 $ L_p(X) $ 中成立，且 $ |\alpha| \le m $。我们令 $ u := u_0 $。则由 (7.19) 推出，对所有 $ j \in \mathbb{N} $，

\begin{equation}
\int_X (\partial^{\alpha}\varphi)u_j\, dx = (-1)^{|\alpha|} \int_X \varphi \partial^{\alpha} u_j\, dx \qquad \text{对 } \varphi \in \mathcal{D}(X) \text{ 且 } |\alpha| \le m .\tag{7.23}
\end{equation}

由赫尔德不等式，我们推出

\[
\left| \int_X (\partial^{\alpha}\varphi)u_j\, dx - \int_X (\partial^{\alpha}\varphi)u\, dx \right| = \left| \int_X \partial^{\alpha}\varphi(u_j - u)\, dx \right| \le \|\partial^{\alpha}\varphi\|_{p'} \|u_j - u\|_p ,
\]

这表明

\[
\int_X (\partial^{\alpha}\varphi)u_j\, dx \to \int_X (\partial^{\alpha}\varphi)u\, dx \qquad \text{对 } \varphi \in \mathcal{D}(X) .
\]

类似地，我们得到

\[
\int_X \varphi \partial^{\alpha} u_j\, dx \to \int_X \varphi u_\alpha\, dx \qquad \text{对 } \varphi \in \mathcal{D}(X) .
\]

于是由 (7.23) 可得

\[
\int_X (\partial^\alpha \varphi) u \, dx = (-1)^{|\alpha|} \int_X \varphi u_\alpha \, dx \quad \text{对 } \varphi \in \mathcal{D}(X) .
\]

这表明 $ u_\alpha $ 是 $ u $ 的 $ \alpha $ 阶弱导数，并且我们看到 $ u $ 是 $ m $ 次弱可微的。因为对 $ |\alpha| \le m $ 有 $ u_\alpha \in L_p(X) $，所以我们也有 $ u \in W_p^m(X) $，并且显然 $ u_j \to u $ 在 $ W_p^m(X) $ 中。因此 $ W_p^m(X) $ 是完备的。\nobreak\hspace{0.35em}■

\end{theorem}

\begin{corollary}

$ W_2^m(X) $ 是希尔伯特空间，其内积为

\[
(u|v)_m := \sum_{|\alpha|\le m} (\partial^\alpha u|\partial^\alpha v) \quad \text{对 } u, v \in W_2^m(X) .
\]

我们将通过证明 Sobolev 空间的所谓迹定理来结束本节。对 $ m \in \mathbb{N} $ 与 $ 1 \le p < \infty $，令

\[
\widehat{H}_p^m(X) := ( \{ u|X ; u \in C_c^m(\mathbb{R}^n) \}, \|\cdot\|_{m,p} ) .
\]

显然 $ \widehat{H}_p^m(X) $ 是 $ W_p^m(X) $ 的向量子空间。若 $ X $ 的边界 $ \partial X $ 足够好（例如，若 $ \overline{X} \subset \mathbb{R}^n $ 是带边的 $ n $ 维子流形（见 XI.1 节））就可以证明 $ \widehat{H}_p^m(X) $ 在 $ W_p^m(X) $ 中稠密。特别地，$ X := \mathbb{R}^n $ 或 $ X := \mathbb{H}^n $ 即属此情形。

\end{corollary}

\begin{theorem}

（迹定理） 设 $ 1 \le p < \infty $ 且 $ X = \mathbb{R}^n $ 或 $ X = \mathbb{H}^n $。则存在唯一的迹算子 $ \gamma \in \mathcal{L}(W_p^1(X), L_p(\mathbb{R}^{n-1})) $，使得对 $ u \in \mathcal{D}(\mathbb{R}^n) $ 有 $ \gamma u = u|\mathbb{R}^{n-1} $（更精确地说，对 $ u \in \widehat{H}_p^1(X) $）。这里 $ \mathbb{R}^{n-1} $ 与 $ \mathbb{R}^{n-1} \times \{0\} \subset \mathbb{R}^n $ 等同。

\proofhead
因为 $ \widehat{H}_p^1(X) $ 在 $ W_p^1(X) $ 中稠密，所以结论由命题 6.24、注 6.25 与定理 VI.2.6 推出。\nobreak\hspace{0.35em}■
\end{theorem}

本定理特别说明，每个元素 $ u \in W_p^1(\mathbb{H}^n) $ 都有边界值 $ \gamma u \in L_p(\partial\mathbb{H}^n) $。由于 $ u $ 在 $ \overline{\mathbb{H}}^n $ 上一般并不连续，$ \gamma u $ 不能简单地通过限制来确定。

迹的存在性是用泛函分析方法处理偏微分方程边值问题的基础。

\exerciseshead

\begin{exercises}
\item

对 $ a > 0 $，计算 $ \chi_{[-a,a]} * \chi_{[-a,a]} $ 与 $ \chi_{[-a,a]} * \chi_{[-a,a]} * \chi_{[-a,a]} $。

\item

设 $ p, p' \in (1, \infty) $ 满足 $ 1/p + 1/p' = 1 $。证明：

\begin{items}
\item

对 $ (f, g) \in \mathcal{L}_p \times \mathcal{L}_{p'} $，$ f * g $ 属于 $ C_0 $，且 $ \|f * g\|_\infty \le \|f\|_p \|g\|_{p'} $。

\item

卷积是从 $ L_p \times L_{p'} $ 到 $ C_0 $ 的良定义、双线性、连续映射。

\end{items}

\item

设 $ p, q, r \in [1, \infty] $ 满足 $ 1/p + 1/q = 1 + 1/r $。验证

\[
* : \mathcal{L}_p \times \mathcal{L}_q \to \mathcal{L}_r , \quad (f, g) \mapsto f * g
\]

是良定义、双线性、连续映射。并验证\textbf{广义杨不等式}

\[
\|f * g\|_r \le \|f\|_p \|g\|_q \quad \text{对 } (f, g) \in \mathcal{L}_p \times \mathcal{L}_q .
\]

（提示：$ r = 1 $ 与 $ r = \infty $ 的情形分别由定理 7.3 与习题 2 覆盖。对 $ r \in (1, \infty) $，考虑

\[
|f(x - y)g(y)| = |f(x - y)|^{1-p/r} \left( |f(x - y)|^p |g(y)|^q \right)^{1/r} |g(y)|^{1-q/r}
\]

然后应用赫尔德不等式。）

\item

证明对 $ (f, g) \in C_c^k \times L_{1,\mathrm{loc}} $，$ f * g $ 属于 $ C^k $。

\item

设 $ f \in \mathcal{L}_{1,\mathrm{loc}} $ 对给定的 $ \alpha \in \mathbb{N}^n $ 满足 $ \partial^\alpha f \in \mathcal{L}_{1,\mathrm{loc}} $。验证

\[
\partial^\alpha (f * \varphi) = (\partial^\alpha f) * \varphi = f * \partial^\alpha \varphi \quad \text{对 } \varphi \in BC^\infty .
\]
\item

给出 $ \mathrm{Funct} $ 的一个向量子空间，使得 $ (\mathbb{R}, +) $ 在其中不是线性可表示的。

\item

给定 $ p \in [1, \infty) $，设 $ K \subset L_p $ 是紧的。证明对每个 $ \varepsilon > 0 $ 都存在 $ \delta > 0 $，使得对一切 $ f \in K $ 与一切满足 $ |a| < \delta $ 的 $ a \in \mathbb{R}^n $ 都有 $ \|\tau_a f - f\|_p < \varepsilon $。（提示：回忆定理 III.3.10 与定理 5.1(iv)。）

\item

证明 $ (L_1, *) $ 的每个非平凡理想都在 $ L_1 $ 中稠密。

\item

设 $ p \in [1, \infty] $，并用 $ k $ 表示例 7.12(a) 中的高斯核。证明：

\begin{items}
\item

对 $ \alpha \in \mathbb{N}^n $ 有 $ \partial^\alpha k \in \mathcal{L}_p $。

\item

对 $ u \in L_p $ 有 $ k * u \in BUC^\infty $。

\end{items}

\item

设 $ f \in \mathcal{L}_1 $，并设对某个 $ \alpha \in \mathbb{N}^n $ 有 $ \partial^\alpha f \in \mathcal{L}_1 $。证明

\[
\int (\partial^\alpha f) \varphi \, dx = (-1)^{|\alpha|} \int f \partial^\alpha \varphi \, dx \quad \text{对 } \varphi \in BC^\infty .
\]
\item

设 $ V \in \{ \mathrm{Funct}, B, L_p ; 1 \le p \le \infty \} $。证明 $ (\mathbb{R}^n, +) $ 通过平移在 $ V $ 上的线性表示是群同构。

\item

设 $ f, g, h \in \mathcal{L}_0 $，假设 $ f $ 与 $ g $ 可卷积、$ g $ 与 $ h $ 可卷积。若 $ f * g $ 与 $ h $ 可卷积且 $ f $ 与 $ g * h $ 可卷积，证明 $ (f * g) * h = f * (g * h) $。于是卷积在 $ L_1 $ 上是结合的。

\item

证明分布导数

\[
\partial^\alpha : \mathcal{D}'(X) \to \mathcal{D}'(X) , \quad T \mapsto \partial^\alpha T
\]

对每个 $ \alpha \in \mathbb{N}^n $ 都是良定义的线性映射。

\item

证明当 $ 1 \le p \le \infty $、$ m \in \mathbb{N} $ 且 $ u \in W_p^m(X) $ 时，$ (f \mapsto f u) \in \mathcal{L}(BC^m(X), W_p^m(X)) $。

\item

设 $ (T_j) $ 是 $ \mathcal{D}'(X) $ 中的序列，且 $ T \in \mathcal{D}'(X) $。若

\[
\lim_{j} \langle T_j, \varphi \rangle = \langle T, \varphi \rangle \quad \text{对 } \varphi \in \mathcal{D}(X) ,
\]

则称 $ (T_j) $ \textbf{在 $ \mathcal{D}'(X) $ 中收敛}于 $ T $。设 $ \{ \varphi_\varepsilon ; \varepsilon > 0 \} $ 是一个恒等逼近，并设 $ (\varepsilon_j) $ 是零序列。证明 $ (\varphi_{\varepsilon_j}) $ 在 $ \mathcal{D}'(\mathbb{R}^n) $ 中收敛于 $ \delta $。

\end{exercises}

\section{换元法则}

在我们处理柯西–黎曼积分时，曾把定理 VI.5.1 的换元法则作为计算积分的基本工具。引入新变量，即选取适当的坐标，也是高维积分中的一种重要技巧。毫不意外，此情形下换元法则的证明更为困难。不过，我们已经奠定了基础，其形式就是线性映射的换元法则，它已在定理 IX.5.25 中推导过。

除了证明 $ n $ 维勒贝格积分的普遍换元法则之外，本节还将借助若干重要例子说明它的意义。同一个定理也是流形上积分理论的基石，而后者是我们最后一章的主题。

以下假定：

- $ X $ 与 $ Y $ 是 $ \mathbb{R}^{n} $ 的开子集；$ E $ 是巴拿赫空间。

\needspace{3\baselineskip}\subsection*{勒贝格测度的拉回}

设 $ (X, \mathcal{A}) $ 是可测空间，$ (Y, \mathcal{B}, \nu) $ 是测度空间。若 $ f : X \to Y $ 是双射且满足 $ f(\mathcal{A}) \subset \mathcal{B} $，即其逆映射是 $ \mathcal{B} $-$ \mathcal{A} $-可测的，则容易验证

\[
f^{*}\nu : \mathcal{A} \to [0, \infty] , \quad A \mapsto \nu(f(A))
\]

定义了 $ \mathcal{A} $ 上的一个测度，即测度 $ \nu $ 在 $ f $ 下的\textbf{拉回}（或\textbf{原像}）。在特殊情形 $ (X, \mathcal{A}) = (\mathbb{R}^{n}, \mathcal{L}(n)) $ 与 $ (Y, \mathcal{B}, \nu) = (\mathbb{R}^{n}, \mathcal{L}(n), \lambda_{n}) $ 下，定理 IX.5.25 所涵盖的换元法则的特殊情形描述了 $ \lambda_{n} $ 在 $ \mathbb{R}^{n} $ 的自同构下的拉回：

\[
\Phi^{*}\lambda_{n} = |\det \Phi| \lambda_{n} \quad \text{对 } \Phi \in \mathcal{L}\mathrm{aut}(\mathbb{R}^{n}) .
\]

利用这一结果，我们现在来确定勒贝格测度在任意 $ C^{1} $-微分同胚下的拉回。为此需要一个技术性结果：

\begin{lemma}

设 $ \Phi \in \mathrm{Diff}^{1}(X, Y) $。则

\[
\lambda_{n}(\Phi(J)) \le \int_{J} |\det \partial\Phi| \, dx
\]

对每个形如 $ [a, b) $ 的区间 $ J \subset\subset X $，其中 $ a, b \in \mathbb{Q}^{n} $。

证明

\begin{subitems}
\item

先考虑中心为 $ x_{0} \in X $、棱长为 $ r > 0 $ 的立方体 $ J = [x_{0} - (r/2)\mathbf{1}, x_{0} + (r/2)\mathbf{1}) $。再令 $ \mathbb{R}^{n}_{\infty} := (\mathbb{R}^{n}, |\cdot|_{\infty}) $ 以及

\[
K := \max_{x\in \overline{J}} \|\partial\Phi(x)\|_{\mathcal{L}(\mathbb{R}^{n}_{\infty})} .
\]

由中值定理可得

\[
|\Phi(x) - \Phi(x_{0})|_{\infty} \le K |x - x_{0}|_{\infty} \quad \text{对 } x \in J .
\]

因此 $ \Phi(J) $ 包含于 $ \overline{\mathbb{B}}^{n}_{\infty}(\Phi(x_{0}), Kr/2) $，并且我们得到

\begin{equation}
\lambda_{n}(\Phi(J)) \le (Kr)^{n} = K^{n}\lambda_{n}(J) .\tag{8.1}
\end{equation}
\item

设 $ J \subset\subset X $ 形如 $ [a, b) $，其中 $ a, b \in \mathbb{Q}^{n} $。取 $ \varepsilon > 0 $ 并令 $ M := \max_{x\in \overline{J}} \|[\partial\Phi(x)]^{-1}\|_{\mathcal{L}(\mathbb{R}^{n}_{\infty})} $。由于 $ \partial\Phi $ 在 $ \overline{J} $ 上一致连续，存在 $ \delta > 0 $ 使得

\begin{equation}
\|\partial\Phi(x) - \partial\Phi(y)\|_{\mathcal{L}(\mathbb{R}^{n}_{\infty})} \le \varepsilon/M\tag{8.2}
\end{equation}

对所有满足 $ |x - y| < \delta $ 的 $ x, y \in \overline{J} $ 成立。由于 $ a, b \in \mathbb{Q}^{n} $，我们可以把 $ J $（按棱剖分）分解为 $ N $ 个不交的立方体 $ J_{k} $，其形如 $ [\alpha, \beta)^{n} $，其中 $ 0 < \beta - \alpha < \delta $。现在选取 $ x_{k} \in \overline{J}_{k} $ 使得

\[
|\det \partial\Phi(x_{k})| = \min_{y\in \overline{J}_{k}} |\det \partial\Phi(y)|
\]

并令 $ T_{k} := \partial\Phi(x_{k}) $ 以及 $ \Phi_{k} := T_{k}^{-1} \circ \Phi $。因为

\[
\partial\Phi_{k}(y) = T_{k}^{-1}\partial\Phi(y) = \mathbf{1}_{n} + [\partial\Phi(x_{k})]^{-1}[\partial\Phi(y) - \partial\Phi(x_{k})]
\]

由 (8.2) 与 $ M $ 的定义可得

\begin{equation}
\max_{y\in \overline{J}_{k}} \|\partial\Phi_{k}(y)\|_{\mathcal{L}(\mathbb{R}^{n}_{\infty})} \le 1 + \varepsilon \quad \text{对 } k \in \{1, \dots, N\} .\tag{8.3}
\end{equation}

由定理 IX.5.25 中处理的换元法则的特殊情形，我们有

\[
\lambda_{n}(\Phi(J_{k})) = \lambda_{n}(T_{k} T_{k}^{-1}\Phi(J_{k})) = |\det T_{k}| \lambda_{n}(\Phi_{k}(J_{k})) .
\]

于是 (8.1) 与 (8.3) 蕴含

\[
\lambda_{n}(\Phi(J_{k})) \le (1 + \varepsilon)^{n} |\det T_{k}| \lambda_{n}(J_{k}) \quad \text{对 } k \in \{1, \dots, N\} .
\]

考虑到 $ \Phi $ 的双射性以及 $ x_{k} $ 的选取，我们得到

\begin{equation*}
\begin{aligned}
\lambda_{n}(\Phi(J)) &= \lambda_{n}\left(\bigcup_{k=1}^{N} \Phi(J_{k})\right) = \sum_{k=1}^{N} \lambda_{n}(\Phi(J_{k})) \\ &\le (1 + \varepsilon)^{n} \sum_{k=1}^{N} |\det T_{k}| \lambda_{n}(J_{k}) \le (1 + \varepsilon)^{n} \sum_{k=1}^{N} \int_{J_{k}} |\det \partial\Phi| \, dx \\ &= (1 + \varepsilon)^{n} \int_{J} |\det \partial\Phi| \, dx .
\end{aligned}
\end{equation*}

令 $ \varepsilon \to 0 $ 取极限即得结论。\nobreak\hspace{0.35em}■

\end{subitems}

\end{lemma}

\begin{proposition}

设 $ \Phi \in \mathrm{Diff}^{1}(X, Y) $。则

\[
\Phi^{*}\lambda_{n}(A) = \lambda_{n}(\Phi(A)) = \int_{A} |\det \partial\Phi| \, dx \quad \text{对 } A \in \mathcal{L}(n)|X .
\]

证明

\begin{subitems}
\item

由单调收敛定理容易推出

\[
\mu_{\Phi} : \mathcal{L}(n)|X \to [0, \infty] , \quad A \mapsto \int_{A} |\det \partial\Phi| \, dx
\]

是完备测度（见习题 2.11）。

\item

设 $ U $ 是开集且紧包含于 $ X $ 中。由命题 IX.5.6，存在一列不交的区间 $ (J_{k}) $，形如 $ [a, b) $，其中 $ a, b \in \mathbb{Q}^{n} $，使得 $ U = \bigcup_{k} J_{k} $。由 (i) 与引理 8.1 可得

\begin{equation*}
\begin{aligned}
\lambda_{n}(\Phi(U)) &= \lambda_{n}\left(\bigcup_{k} \Phi(J_{k})\right) = \sum_{k} \lambda_{n}(\Phi(J_{k})) \le \sum_{k} \int_{J_{k}} |\det \partial\Phi| \, dx \\ &= \sum_{k} \mu_{\Phi}(J_{k}) = \mu_{\Phi}\left(\bigcup_{k} J_{k}\right) = \mu_{\Phi}(U) = \int_{U} |\det \partial\Phi| \, dx .
\end{aligned}
\end{equation*}
\item

设 $ U $ 是 $ X $ 中的开集。由注 1.16(d) 与 (e)，存在 $ X $ 的一列开子集 $ (U_{k}) $，使得 $ U_{k} \subset\subset U_{k+1} $ 且 $ U = \bigcup_{k} U_{k} $。由 (ii) 以及测度 $ \lambda_{n} $ 与 $ \mu_{\Phi} $ 的下连续性可得

\[
\lambda_{n}(\Phi(U)) = \lim_{k} \lambda_{n}(\Phi(U_{k})) \le \lim_{k} \mu_{\Phi}(U_{k}) = \mu_{\Phi}(U) = \int_{U} |\det \partial\Phi| \, dx .
\]
\item

设 $ A \in \mathcal{L}(n)|X $ 有界。利用推论 IX.5.5，我们找到 $ X $ 的一列有界开子集 $ (U_{k}) $，使得 $ G := \bigcap_{k} U_{k} \supset A $ 且 $ \lambda_{n}(G) = \lambda_{n}(A) $。由 (iii) 以及测度 $ \lambda_{n} $ 与 $ \mu_{\Phi} $ 的上连续性，我们有

\begin{equation*}
\begin{aligned}
\lambda_{n}(\Phi(G)) &= \lim_{k} \lambda_{n}\left(\Phi\left(\bigcap_{j=0}^{k} U_{j}\right)\right) \le \lim_{k} \mu_{\Phi}\left(\bigcap_{j=0}^{k} U_{j}\right) \\ &= \mu_{\Phi}(G) = \int_{G} |\det \partial\Phi| \, dx .
\end{aligned}
\end{equation*}

注意到 $ A \subset G $ 且 $ \lambda_{n}(A) = \lambda_{n}(G) $，我们得到

\[
\lambda_{n}(\Phi(A)) \le \lambda_{n}(\Phi(G)) \le \int_{G} |\det \partial\Phi| \, dx = \int_{A} |\det \partial\Phi| \, dx .
\]
\item

取任意 $ A \in \mathcal{L}(n)|X $，并令 $ A_{k} := A \cap k\mathbb{B}^{n} $，$ k \in \mathbb{N} $。由 (iv) 以及测度的下连续性，我们得到

\[
\lambda_{n}(\Phi(A)) = \lim_{k} \lambda_{n}(\Phi(A_{k})) \le \lim_{k} \mu_{\Phi}(A_{k}) = \mu_{\Phi}(A) = \int_{A} |\det \partial\Phi| \, dx .
\]
\item

设 $ f \in \mathcal{S}(Y, \mathbb{R}^{+}) $ 具有标准形 $ f = \sum_{j=0}^{k} \alpha_{j}\chi_{A_{j}} $。由 (v)，

\begin{equation*}
\begin{aligned}
\int_{Y} f \, dy &= \sum_{j=0}^{k} \alpha_{j}\lambda_{n}(A_{j}) = \sum_{j=0}^{k} \alpha_{j}\lambda_{n}(\Phi(\Phi^{-1}(A_{j}))) \\ &\le \sum_{j=0}^{k} \alpha_{j} \int_{\Phi^{-1}(A_{j})} |\det \partial\Phi| \, dx = \int_{X} (f \circ \Phi)|\det \partial\Phi| \, dx .
\end{aligned}
\end{equation*}
\item

设 $ X $ 有界。给定 $ f \in \mathcal{L}_{0}(Y, \mathbb{R}^{+}) $，设 $ (f_{k}) $ 是 $ \mathcal{S}(Y, \mathbb{R}^{+}) $ 中的序列，使得 $ f_{k} \uparrow f $（见定理 1.12）。则 $ f_{k} \circ \Phi $ 属于 $ \mathcal{S}(X, \mathbb{R}^{+}) $。因为序列 $ (f_{k} \circ \Phi)_{k} $ 递增收敛到 $ f \circ \Phi $，我们知道 $ (f \circ \Phi)|\det \partial\Phi| $ 属于 $ \mathcal{L}_{0}(X, \mathbb{R}^{+}) $。现在 (vi) 与单调收敛定理蕴含

\[
\int_{Y} f \, dy = \lim_{k} \int_{Y} f_{k} \, dy \le \lim_{k} \int_{X} (f_{k} \circ \Phi)|\det \partial\Phi| \, dx = \int_{X} (f \circ \Phi)|\det \partial\Phi| \, dx .
\]
\item

设 $ X $ 任意，取 $ f \in \mathcal{L}_{0}(Y, \mathbb{R}^{+}) $。鉴于注 1.16(d) 与 (e)，我们可以找到 $ X $ 的一列递增的相对紧开子集 $ X_{k} $，使得 $ X = \bigcup_{k=0}^{\infty} X_{k} $。根据 (vii)，$ g_{k} := \chi_{X_{k}} f|\det \Phi| $ 属于 $ \mathcal{L}_{0}(X, \mathbb{R}^{+}) $，并且 $ g_{k} \uparrow g := f|\det \Phi| $。因此 $ g \in \mathcal{L}_{0}(X, \mathbb{R}^{+}) $。令 $ Y_{k} := \Phi(X_{k}) $，由 (vii) 得到

\[
\int_{Y_{k}} f \, dy \le \int_{X_{k}} (f \circ \Phi)|\det \partial\Phi| \, dx .
\]

现在 $ Y = \bigcup_{k=0}^{\infty} Y_{k} $，而单调收敛定理给出

\begin{equation}
\int_{Y} f \, dy \le \int_{X} (f \circ \Phi)|\det \partial\Phi| \, dx .\tag{8.4}
\end{equation}
\item

设 $ A \in \mathcal{L}(n)|X $。我们在 (viii) 中交换 $ X $ 与 $ Y $ 的角色，并把 (8.4) 应用于 $ C^{1} $-微分同胚 $ \Phi^{-1} : Y \to X $ 以及函数 $ (\chi_{\Phi(A)} \circ \Phi)|\det \partial\Phi| $，后者属于 $ \mathcal{L}_{0}(X, \mathbb{R}^{+}) $。则

\begin{equation*}
\begin{aligned}
\int_{X} (\chi_{\Phi(A)} \circ \Phi)|\det \partial\Phi| \, dx &\le \int_{Y} \left[ ((\chi_{\Phi(A)} \circ \Phi)|\det \partial\Phi|) \circ \Phi^{-1} \right] |\det \partial\Phi^{-1}| \, dy \\ &= \int_{Y} \chi_{\Phi(A)} |\det[(\partial\Phi \circ \Phi^{-1})\partial\Phi^{-1}]| \, dy .
\end{aligned}
\end{equation*}

进一步注意到

\begin{equation}
\mathbf{1}_{n} = \partial(\mathrm{id}_{Y}) = \partial(\Phi \circ \Phi^{-1}) = (\partial\Phi \circ \Phi^{-1})\partial\Phi^{-1}\tag{8.5}
\end{equation}

以及 $ \chi_{\Phi(A)} \circ \Phi = \chi_{A} $，我们得到

\[
\int_{A} |\det \partial\Phi| \, dx \le \int_{Y} \chi_{\Phi(A)} \, dy = \lambda_{n}(\Phi(A)) .
\]

由 (v)，即得结论。\nobreak\hspace{0.35em}■

\end{subitems}

\end{proposition}

\begin{example}

定义 $ X := \{ (r, \varphi) \in \mathbb{R} \times (0, 2\pi) ; 0 < r < \varphi/2\pi \} $ 以及

\[
\Phi : X \to \mathbb{R}^{2} , \quad (r, \varphi) \mapsto (r \cos\varphi, r \sin\varphi) .
\]

则 $ Y := \Phi(X) $ 是 $ \mathbb{R}^{2} $ 中的开集，且 $ \Phi \in \mathrm{Diff}^{\infty}(X, Y) $ 满足

\begin{equation*}
\left[\partial\Phi(r, \varphi)\right] = \begin{bmatrix} \cos\varphi & -r\sin\varphi \\ \sin\varphi & r\cos\varphi \end{bmatrix} .
\end{equation*}
\bookfigure[0.88]{figures/81a6c75297a89a6c.png}

因此 $ \det \partial\Phi(r, \varphi) = r $。又 $ \mathrm{pr}_{2}(X) = (0, 2\pi) $，且对 $ \varphi \in (0, 2\pi) $ 有 $ X^{[\varphi]} = (0, \varphi/2\pi) $。于是由命题 8.2 与托内利定理，

\[
\lambda_{2}(Y) = \int_{X} r \, d(r, \varphi) = \int_{0}^{2\pi} \int_{0}^{\varphi/2\pi} r \, dr \, d\varphi = \pi/3 .
\]
\end{example}

\needspace{3\baselineskip}\subsection*{换元法则：一般情形}

有了这些准备之后，证明微分同胚的换元法则就不再困难。我们先考虑标量情形，其证明甚至对跳过向量值函数的富比尼定理之证明的读者也是可读的。本节末尾再处理一般情形。

\begin{theorem}

（换元法则） 设 $ \Phi \in \mathrm{Diff}^{1}(X, Y) $。

\begin{subitems}
\item

对 $ f \in \mathcal{L}_{0}(Y, \mathbb{R}^{+}) $，

\begin{equation}
\int_{Y} f \, dy = \int_{X} (f \circ \Phi)|\det \partial\Phi| \, dx .\tag{8.6}
\end{equation}
\item

函数 $ f : Y \to \mathbb{K} $ 可积当且仅当 $ (f \circ \Phi)|\det \partial\Phi| $ 属于 $ \mathcal{L}_{1}(X) $。此时 (8.6) 成立。

\proofhead
(i) 定理 IX.5.12 蕴含 $ \Phi(\mathcal{L}_{X}) \subset \mathcal{L}_{Y} $。因此由推论 1.5，$ f \circ \Phi $ 可测。由于 $ |\det \partial\Phi| $ 连续，因而可测，注 1.2(d) 蕴含 $ g := (f \circ \Phi)|\det \partial\Phi| $ 也可测。由 (8.5) 我们得到 $ f = (g \circ \Phi^{-1})|\det \partial\Phi^{-1}| $。现在把 (8.4) 中的 $ (X, \Phi, f) $ 换成 $ (Y, \Phi^{-1}, g) $，给出
\[
\int_{X} (f \circ \Phi)|\det \partial\Phi| \, dx \le \int_{Y} f \, dy .
\]

由 (8.4)，这蕴含 (8.6)。现在 (ii) 由 (i)、推论 2.12 的 (ii) 与 (iii) 部分以及定理 3.14 推出。\nobreak\hspace{0.35em}■

\end{subitems}

\end{theorem}

就第 VIII.3 节定义的函数的拉回而言，换元法则 (8.6) 具有易于记忆的形式

\[
\int_{Y} f \, d\lambda_{n} = \int_{\Phi^{-1}(Y)} (\Phi^{*} f) \, d(\Phi^{*}\lambda_{n}) .
\]

这由命题 8.2 与习题 2.12 推出。

对于许多应用来说，$ \Phi $ 是微分同胚这一假设过于严格。我们在定理 8.4 的这个简单却重要的推广中把它稍加减弱：\footnote{关于进一步的推广，见习题 7。}

\begin{corollary}

设 $ M $ 是 $ X $ 的可测子集，使得 $ M \setminus \mathring{M} $ 的勒贝格测度为零。设 $ \Phi \in C^{1}(X, \mathbb{R}^{n}) $ 使得 $ \Phi|_{\mathring{M}} $ 是从 $ \mathring{M} $ 到 $ \Phi(\mathring{M}) $ 的微分同胚。

\begin{subitems}
\item

对每个 $ f \in \mathcal{L}_{0}(M, \mathbb{R}^{+}) $，

\begin{equation}
\int_{\Phi(M)} f \, dy = \int_{M} (f \circ \Phi)|\det \partial\Phi| \, dx .\tag{8.7}
\end{equation}
\item

函数 $ f : \Phi(M) \to \mathbb{K} $ 属于 $ \mathcal{L}_{1}(\Phi(M)) $ 当且仅当 $ (f \circ \Phi)|\det \partial\Phi| $ 属于 $ \mathcal{L}_{1}(M) $。此时 (8.7) 成立。

\proofhead
因为 $ \lambda_{n}(M \setminus \mathring{M}) = 0 $，由推论 IX.5.10，集合 $ \Phi(M) \setminus \Phi(\mathring{M}) \subset \Phi(M \setminus \mathring{M}) $ 的测度也为零。于是结论由引理 2.15 与定理 8.4 推出。\nobreak\hspace{0.35em}■
\end{subitems}

\end{corollary}

显然，这个推论给出了定理 VI.5.1 的换元法则的一个（部分）推广，尽管局限于微分同胚。与先前考虑的一维情形有一个明显的差别：现在导数项（即函数行列式）以绝对值的形式出现。原因是先前的结论使用了定向积分。

\needspace{3\baselineskip}\subsection*{平面极坐标}

应用中特别重要的一个特殊情形是由极坐标诱导的微分同胚，我们现在就来引入它。我们从二维情形开始。

设

\[
f_{2} : \mathbb{R}^{2} \to \mathbb{R}^{2} , \quad (r, \varphi) \mapsto (x, y) := (r \cos\varphi, r \sin\varphi)
\]

为（\textbf{平面}）极坐标映射\footnote{见结论 III.6.21(d)。}，并令 $ V_{2} := (0, \infty) \times (0, 2\pi) $。

\bookfigure[0.95]{figures/696f6e600f0740e9.png}

则 $ f_{2} $ 光滑，且 $ \det \partial f_{2}(r, \varphi) = r $，如例 8.3 所示。显然 $ \overline{V}_{2} \setminus V_{2} $ 的测度为零；此外

\begin{equation}
f_{2}(\overline{V}_{2}) = \mathbb{R}^{2} , \quad f_{2}(V_{2}) = \mathbb{R}^{2} \setminus (\mathbb{R}^{+} \times \{0\})\tag{8.8}
\end{equation}

以及

\begin{equation}
f_{2}|_{V_{2}} \in \mathrm{Diff}^{\infty}(V_{2}, f_{2}(V_{2})) .\tag{8.9}
\end{equation}

因此推论 8.5 适用于 $ M := \overline{V}_{2} $：

\begin{proposition}

（极坐标中的积分）

\begin{subitems}
\item

对 $ g \in \mathcal{L}_{0}(\mathbb{R}^{2}, \mathbb{R}^{+}) $，我们有

\begin{equation}
\begin{aligned}
\int_{\mathbb{R}^{2}} g(x, y) \, d(x, y) &= \int_{0}^{2\pi} \int_{0}^{\infty} g(r \cos\varphi, r \sin\varphi) r \, dr \, d\varphi \\ &= \int_{0}^{\infty} r \int_{0}^{2\pi} g(r \cos\varphi, r \sin\varphi) \, d\varphi \, dr .
\end{aligned}\tag{8.10}
\end{equation}
\item

函数 $ g : \mathbb{R}^{2} \to \mathbb{K} $ 可积当且仅当映射

\[
(0, \infty) \times (0, 2\pi) \to \mathbb{K} , \quad (r, \varphi) \mapsto g(r \cos\varphi, r \sin\varphi) r
\]

可积。此时 (8.10) 成立。

\proofhead
这由推论 8.5 连同 (8.8)、(8.9) 以及富比尼–托内利定理推出。\nobreak\hspace{0.35em}■
\end{subitems}

\end{proposition}

当 $ f $ 只依赖 $ |x| $，即只依赖 $ r $ 时，这些积分会简化。为说明这一点，我们给出高斯误差积分的一个优雅计算，它不需要关于 $ \Gamma $-函数的知识（比较应用 VI.9.7）。

\begin{example}

$ \displaystyle\int_{-\infty}^{\infty} e^{-x^{2}} dx = \sqrt{\pi} $。

\proofhead
托内利定理蕴含
\begin{equation*}
\begin{aligned}
\left(\int_{-\infty}^{\infty} e^{-x^{2}} dx\right)^{2} &= \int_{-\infty}^{\infty} e^{-x^{2}} dx \int_{-\infty}^{\infty} e^{-y^{2}} dy = \int_{\mathbb{R}} \left(\int_{\mathbb{R}} e^{-(x^{2}+y^{2})} dx\right) dy \\ &= \int_{\mathbb{R}^{2}} e^{-(x^{2}+y^{2})} \, d(x, y) .
\end{aligned}
\end{equation*}

因此命题 8.6(i) 表明

\[
\left(\int_{-\infty}^{\infty} e^{-x^{2}} dx\right)^{2} = \int_{0}^{2\pi} \int_{0}^{\infty} r e^{-r^{2}} \, dr \, d\varphi = 2\pi \int_{0}^{\infty} \frac{d}{dr}\left[-e^{-r^{2}}/2\right] dr = \pi ,
\]

于是结论得证。\nobreak\hspace{0.35em}■

\end{example}

\needspace{3\baselineskip}\subsection*{高维极坐标}

对 $ n \ge 1 $，我们通过以下方式递归地定义 $ h_{n} : \mathbb{R}^{n} \to \mathbb{R}^{n+1} $：

\begin{equation}
h_{1}(z) := (\cos z, \sin z) \quad \text{对 } z \in \mathbb{R}\tag{8.11}
\end{equation}

以及

\begin{equation}
h_{n+1}(z) := (h_{n}(z') \sin z_{n+1}, \cos z_{n+1}) \quad \text{对 } z = (z', z_{n+1}) \in \mathbb{R}^{n} \times \mathbb{R} .\tag{8.12}
\end{equation}

显然 $ h_{n} $ 光滑，并且通过归纳法可以验证

\begin{equation}
|h_{n}(z)| = 1 \quad \text{对 } z \in \mathbb{R}^{n} .\tag{8.13}
\end{equation}

现在我们对 $ n \ge 2 $ 定义 $ f_{n} : \mathbb{R}^{n} \to \mathbb{R}^{n} $：

\begin{equation}
f_{n}(y) := y_{1} h_{n-1}(z) \quad \text{对 } y = (y_{1}, z) \in \mathbb{R} \times \mathbb{R}^{n-1} .\tag{8.14}
\end{equation}

则 $ f_{n} $ 也光滑，并且我们有

\begin{equation}
h_{n-1}(z) = f_{n}(1, z) , \quad |f_{n}(y)| = |y_{1}| .\tag{8.15}
\end{equation}

我们通常按惯例把 $ y $ 坐标重新命名如下：

\[
(r, \varphi, \vartheta_{1}, \dots, \vartheta_{n-2}) := (y_{1}, y_{2}, y_{3}, \dots, y_{n}) .
\]

通过归纳法容易验证

\begin{equation}
f_{n} : \mathbb{R}^{n} \to \mathbb{R}^{n} , \quad (r, \varphi, \vartheta_{1}, \dots, \vartheta_{n-2}) \mapsto (x_{1}, x_{2}, x_{3}, \dots, x_{n})\tag{8.16}
\end{equation}

由下式给出

\begin{equation}
\left. \begin{array}{r@{\;}l} x_{1} & = r\cos\varphi\sin\vartheta_{1}\sin\vartheta_{2}\cdots\sin\vartheta_{n-2}, \\ x_{2} & = r\sin\varphi\sin\vartheta_{1}\sin\vartheta_{2}\cdots\sin\vartheta_{n-2}, \\ x_{3} & = r\cos\vartheta_{1}\sin\vartheta_{2}\cdots\sin\vartheta_{n-2}, \\ & \vdots \\ x_{n-1} & = r\cos\vartheta_{n-3}\sin\vartheta_{n-2}, \\ x_{n} & = r\cos\vartheta_{n-2}. \end{array} \right\}\tag{8.17}
\end{equation}

因此 $ f_{2} $ 与平面极坐标映射一致，而 $ f_{3} $ 是例 VII.9.11(a) 的\textbf{球坐标映射}。在一般情形下，$ f_{n} $ 是 \textbf{$ n $ 维极坐标映射}。由 (8.12) 与 (8.14)，递推关系

\begin{equation}
f_{n}(y) = \left( f_{n-1}(y')\sin y_{n}, y_{1}\cos y_{n} \right) \qquad \text{对 } y = (y', y_{n}) \in \mathbb{R}^{n-1}\times\mathbb{R}\tag{8.18}
\end{equation}

对 $ n \ge 3 $ 成立。对 $ n \ge 2 $，我们令

\begin{equation}
W_{n-1} := (0, 2\pi)\times(0, \pi)^{n-2}, \quad V_{n} := (0, \infty)\times W_{n-1},\tag{8.19}
\end{equation}

以及

\begin{equation}
V_{n}(r) := (0, r)\times W_{n-1} \qquad \text{对 } r > 0.\tag{8.20}
\end{equation}

若把闭的 $ (n-1) $ 维半空间记为

\begin{equation}
H_{n-1} := \mathbb{R}^{+}\times\{0\}\times\mathbb{R}^{n-2} \subset \mathbb{R}^{n},\tag{8.21}
\end{equation}

则得

\begin{equation}
h_{n-1}(W_{n-1}) = S^{n-1}\setminus H_{n-1}, \quad f_{n}(V_{n}(r)) = r\mathbb{B}^{n}\setminus H_{n-1}\tag{8.22}
\end{equation}

以及

\begin{equation}
h_{n-1}(\overline{W}_{n-1}) = S^{n-1}, \quad f_{n}(\overline{V_{n}(r)}) = r\overline{\mathbb{B}^{n}}.\tag{8.23}
\end{equation}

又有

\begin{equation}
f_{n}(V_{n}) = \mathbb{R}^{n}\setminus H_{n-1}, \quad f_{n}(\overline{V}_{n}) = \mathbb{R}^{n}.\tag{8.24}
\end{equation}

此外，映射 $ h_{n-1}|W_{n-1} $ 与 $ f_{n}|V_{n} $ 都是到其像上的双射。

\bookfigure[0.95]{figures/2f2d6e6af6e39ffb.png}

这些结论容易由归纳法得到。

\begin{lemma}

设 $ n \ge 3 $ 且 $ r > 0 $。则映射 $ f_{n} $ 是从 $ V_{n}(r) $ 到 $ r\mathbb{B}^{n}\setminus H_{n-1} $ 上、以及从 $ V_{n} $ 到 $ \mathbb{R}^{n}\setminus H_{n-1} $ 上的 $ C^{\infty} $ 微分同胚。此外

\[
\det \partial f_{n}(r, \varphi, \vartheta_{1}, \ldots, \vartheta_{n-2}) = (-1)^{n}r^{n-1}\sin\vartheta_{1}\sin^{2}\vartheta_{2}\cdots\sin^{n-2}\vartheta_{n-2}
\]

对 $ (r, \varphi, \vartheta_{1}, \ldots, \vartheta_{n-2}) \in \overline{V}_{n} $。

\proofhead
鉴于前面所述，我们只需计算函数行列式 $ \det \partial f_{n}(y) $ 的值。我们递推地进行。由 (8.12) 与 (8.14) 得
\begin{equation*}
[\partial f_{n+1}(y)] = \begin{bmatrix} h_{n-1}(z')\sin z_{n} & \left[ r\partial_{z}\left( h_{n-1}(z')\sin z_{n} \right) \right] \\ \cdots\cdots\cdots\cdots\cdots\cdots\cdots\cdots & \cdots\cdots\cdots\cdots\cdots\cdots\cdots\cdots\cdots \\ \cos z_{n} & 0 \quad \cdots \quad 0 \quad -r\sin z_{n} \end{bmatrix}
\end{equation*}
\begin{equation*}
= \begin{bmatrix} \left[ \partial f_{n}(y')\sin z_{n} \right] & \begin{matrix} * \\ \vdots \\ * \end{matrix} \\ \cdots\cdots\cdots\cdots\cdots & \cdots\cdots\cdots\cdots \\ * \quad \cdots \quad * & -r\sin z_{n} \end{bmatrix} ,
\end{equation*}

其中 $ y = (r, z) = (y', z_{n}) $ 且 $ z = (z', z_{n}) \in \mathbb{R}^{n} $。按最后一行展开，得

\begin{equation}
\det \partial f_{n+1}(y) = (-1)^{n}\cos z_{n}\det S - r\sin^{n+1}z_{n}\det \partial f_{n}(y'),\tag{8.25}
\end{equation}

其中 $ S := [r\partial_{z}(h_{n-1}(z')\sin z_{n})] $。我们可以假定 $ \sin z_{n} \ne 0 $；否则结论是平凡的。在 $ S $ 的最后一列中有 $ r h_{n-1}(z')\cos z_{n} $。这个向量与矩阵 $ T := [\partial f_{n}(y')\sin z_{n}] $ 的第一列向量 $ h_{n-1}(z')\sin z_{n} $ 只相差因子 $ r\cot z_{n} $。$ S $ 的前 $ n-1 $ 列也与 $ T $ 的后 $ n-1 $ 列按同样的顺序一致。因此

\[
\det S = (-1)^{n-1}r\cot z_{n}\det T = (-1)^{n-1}r\cos z_{n}\sin^{n-1}z_{n}\det \partial_{n}f(y').
\]

于是由 (8.25) 推出

\[
\det \partial f_{n+1}(y) = -r\sin^{n-1}z_{n}\det \partial f_{n}(y').
\]

现在结论成立，因为 $ \det \partial f_{2}(r, \varphi) = r $。\nobreak\hspace{0.35em}■

\end{lemma}

为简短起见，令

\[
w_{n}(\vartheta) := \sin\vartheta_{1}\sin^{2}\vartheta_{2}\cdots\sin^{n-2}\vartheta_{n-2}, \quad \vartheta := (\vartheta_{1}, \ldots, \vartheta_{n-2}) \in [0, \pi]^{n-2}.
\]

\begin{proposition}

（极坐标下的积分） 设 $ n \ge 3 $。

\begin{subitems}
\item

对 $ g \in \mathcal{L}_{0}(\mathbb{R}^{n}, \mathbb{R}^{+}) $，有

\begin{equation}
\int_{\mathbb{R}^{n}} g\,dx = \int_{V_{n}} (g\circ f_{n})(r, \varphi, \vartheta)r^{n-1}w_{n}(\vartheta)\, d(r, \varphi, \vartheta).\tag{8.26}
\end{equation}
\item

映射 $ g: \mathbb{R}^{n} \to \mathbb{K} $ 可积当且仅当

\[
V_{n} \to \mathbb{K}, \quad (r, \varphi, \vartheta) \mapsto (g\circ f_{n})(r, \varphi, \vartheta)r^{n-1}w_{n}(\vartheta)
\]

可积。此时 (8.26) 成立。

\proofhead
因为 $ \lambda_{n}(\overline{V}_{n}\setminus V_{n}) = 0 $，结论由 (8.24)、推论 8.5 与引理 8.8 推出。\nobreak\hspace{0.35em}■
\end{subitems}

\end{proposition}

\begin{example}

\begin{items}
\item

对 $ g \in \mathcal{L}_{0}(\mathbb{R}^{3}, \mathbb{R}^{+}) $，有

\begin{equation}
\int_{\mathbb{R}^{3}} g(x, y, z)\, d(x, y, z) = \int_{0}^{\infty}\int_{0}^{2\pi}\int_{0}^{\pi} g(r\cos\varphi\sin\vartheta, r\sin\varphi\sin\vartheta, r\cos\vartheta)r^{2}\sin\vartheta\, d\vartheta\, d\varphi\, dr.\tag{8.27}
\end{equation}

右端的积分可以按任意顺序进行。

\proofhead
这由命题 8.9(i) 与托内利定理推出。\nobreak\hspace{0.35em}■
\item

映射 $ g: \mathbb{R}^{3} \to \mathbb{K} $ 可积当且仅当

\[
V_{3} \to \mathbb{K}, \quad (r, \varphi, \vartheta) \mapsto g(r\cos\varphi\sin\vartheta, r\sin\varphi\sin\vartheta, r\cos\vartheta)r^{2}\sin\vartheta
\]

可积。这样的映射满足 (8.27)，且其中的积分可以按任意顺序进行。

\proofhead
这是命题 8.9(ii) 与富比尼–托内利定理的推论。\nobreak\hspace{0.35em}■
\item

对 $ n \ge 3 $，有

\[
2\pi\int_{[0, \pi]^{n-2}} w_{n}(\vartheta)\, d\vartheta = n\omega_{n},
\]

其中 $ \omega_{n} = \pi^{n/2}/\Gamma(1 + n/2) $ 是 $ \mathbb{B}^{n} $ 的体积。

\proofhead
由 (8.22)、(8.23)、命题 8.9 与托内利定理推出
\[
\omega_{n} = \int_{\mathbb{B}^{n}} dx = \int_{\mathbb{B}^{n}} 1\,dx = \int_{V_{n}(1)} (1\circ f_{n})(r, \varphi, \vartheta)r^{n-1}w_{n}(\vartheta)\, d(r, \varphi, \vartheta)
\]
\[
= \int_{0}^{1} r^{n-1}dr\int_{0}^{2\pi} d\varphi\int_{[0, \pi]^{n-2}} w_{n}(\vartheta)\, d\vartheta = \frac{2\pi}{n}\int_{[0, \pi]^{n-2}} w_{n}(\vartheta)\, d\vartheta,
\]

于是证毕。\nobreak\hspace{0.35em}■

\end{items}

\end{example}

\needspace{3\baselineskip}\subsection*{旋转对称函数的积分}

设 $ 0 \le r_{0} < r_{1} \le \infty $，并令 $ R(r_{0}, r_{1}) := \{ x \in \mathbb{R}^{n} ;\ r_{0} < |x| < r_{1} \} $。称函数 $ g: R(r_{0}, r_{1}) \to E $ 是\textbf{旋转对称的}，如果存在映射 $ \dot{g}: (r_{0}, r_{1}) \to E $ 使得

\[
g(x) = \dot{g}(|x|) \qquad \text{对 } x \in R(r_{0}, r_{1}).
\]

这种情形出现当且仅当 $ g $ 在每个球面 $ rS^{n-1} $（$ r_{0} < r < r_{1} $）上是常数。对这样的函数，$ \dot{g} $ 由 $ g $ 唯一确定（反之亦然）。

正如我们在例 8.7 中所见，对旋转对称函数而言，积分问题会大大简化。

\begin{theorem}

设 $ 0 \le r_{0} < r_{1} \le \infty $。

\begin{subitems}
\item

若 $ g \in \mathcal{L}_{0}(R(r_{0}, r_{1}), \mathbb{R}^{+}) $ 是旋转对称的，则

\begin{equation}
\int_{R(r_{0}, r_{1})} g\,dx = n\omega_{n}\int_{r_{0}}^{r_{1}} \dot{g}(r)r^{n-1}dr,\tag{8.28}
\end{equation}

其中 $ \omega_{n} := \lambda_{n}(\mathbb{B}^{n}) = \pi^{n/2}/\Gamma(1 + n/2) $。

\item

旋转对称函数 $ g: R(r_{0}, r_{1}) \to \mathbb{K} $ 可积当且仅当

\[
(r_{0}, r_{1}) \to \mathbb{K}, \quad r \mapsto \dot{g}(r)r^{n-1}
\]

可积。此时 (8.28) 成立。

\proofhead
$ n = 1 $ 的情形是清楚的（见习题 5.12）。对 $ n \ge 2 $，由 (8.15) 以及 $ g $ 的旋转对称性得
\[
g\circ f_{n}(r, \varphi, \vartheta) = \dot{g}(r) \qquad \text{对 } r_{0} < r < r_{1} \text{ 与 } (\varphi, \vartheta) \in W_{n-1}.
\]

现在，结论由命题 8.6 与 8.9（应用于 $ g $ 的平凡延拓）以及例 8.10(c) 推出。\nobreak\hspace{0.35em}■

\end{subitems}

\end{theorem}

\begin{example}

\begin{items}
\item

设 $ f: \mathbb{R}^{n} \to \mathbb{K} $ 可测，且存在 $ c \ge 0 $、$ \rho > 0 $ 与 $ \varepsilon > 0 $ 使得

\begin{equation*}
|f(x)| \le \begin{cases} c\,|x|^{-n+\varepsilon} & \text{若 } 0 < |x| \le \rho, \\ c\,|x|^{-n-\varepsilon} & \text{若 } |x| \ge \rho. \end{cases}
\end{equation*}

则 $ f $ 可积。

\proofhead
令
\[
g(x) := c\left( |x|^{-n+\varepsilon}\chi_{\rho\mathbb{B}^{n}}(x) + |x|^{-n-\varepsilon}\chi_{(\rho\mathbb{B}^{n})^{c}}(x) \right) \qquad \text{对 } x \in \mathbb{R}^{n}\setminus\{0\} = R(0, \infty).
\]

则 $ g $ 是旋转对称的，且 $ |f(x)| \le g(x) $ 对 $ x \in R(0, \infty) $。由例 VI.8.4(a) 与 (b)，$ r \mapsto \dot{g}(r)r^{n-1} $ 属于 $ \mathcal{L}_{1}(\mathbb{R}^{+}) $。因此定理 8.11 推出 $ g $ 也属于 $ \mathcal{L}_{1}(R(0, \infty)) = \mathcal{L}_{1}(\mathbb{R}^{n}) $。现在由定理 3.14 得结论。\nobreak\hspace{0.35em}■

\item

设 $ \mu \in \mathcal{L}_{\infty}(\mathbb{R}^{n}) $ 具有紧支集。再定义

\[
\frac{1}{r}: \mathbb{R}^{n}\setminus\{0\} \to \mathbb{R}^{+}, \quad x \mapsto \frac{1}{|x|}.
\]

则对 $ \alpha < n $，$ (1/r)^{\alpha}*\mu $ 存在，且

\[
\left( \frac{1}{r} \right)^{\alpha}*\mu(x) = \int_{\mathbb{R}^{n}} \frac{\mu(y)}{|x - y|^{\alpha}}\, dy \qquad \text{对 } x \in \mathbb{R}^{n}.
\]

\proofhead
取 $ x \in \mathbb{R}^{n} $，令 $ K := \operatorname{supp}(\mu) $ 且 $ g_{x}(y) := \|\mu\|_{\infty}|y|^{-\alpha}\chi_{x-K}(y) $ 对 $ y \ne 0 $。则 $ \widetilde{g}_{x} $ 属于 $ \mathcal{L}_{0}(\mathbb{R}^{n}) $，且
\[
|\mu(x - y)|\,|y|^{-\alpha} \le g_{x}(y) \qquad \text{对 } y \ne 0.
\]

因为 $ \alpha < n $，由 (a) 可知 $ \widetilde{g}_{x} $ 可积。现在结论由定理 7.8(ii) 推出。\nobreak\hspace{0.35em}■

\end{items}

对 $ n \ge 3 $，按 (b) 的记号，函数 $ u_{n} := (1/r)^{n-2}*\mu $ 称为与密度 $ \mu $ 相联系的\textbf{牛顿势}或\textbf{库仑势}。由习题 3.6 我们知道 $ u_{n} $ 在 $ K^{c} $ 中是光滑且调和的，而 (b) 表明 $ u_{n} $ 在整个 $ \mathbb{R}^{n} $ 上有定义。

\end{example}

\needspace{3\baselineskip}\subsection*{向量值函数的换元法则}

我们现在对向量值函数证明定理 8.4 的换元公式。

\begin{lemma}

设 $ f \in \mathcal{S}_{c}(Y, E) $ 且 $ \Phi \in \operatorname{Diff}^{1}(X, Y) $。则 $ (f\circ\Phi)|\det \partial\Phi| $ 属于 $ \mathcal{L}_{1}(X, E) $，且

\[
\int_{Y} f\,dy = \int_{X} (f\circ\Phi)|\det \partial\Phi|\,dx.
\]

\proofhead
因为 $ \operatorname{supp}(f\circ\Phi) = \Phi^{-1}(\operatorname{supp}(f)) $，$ f\circ\Phi $ 的支集是紧的。特别地，$ f\circ\Phi $ 属于 $ \mathcal{S}_{c}(X, E) $。由此容易推出 $ (f\circ\Phi)|\det \partial\Phi| $ 可积。此外，定理 2.11(iii) 表明，对 $ e \in E $ 与 $ g \in \mathcal{L}_{1}(X, \mathbb{K}) $，函数 $ eg $ 属于 $ \mathcal{L}_{1}(X, E) $ 且 $ e\int_{X} g\,dx = \int_{X} eg\,dx $。令 $ \sum_{j=0}^{m} e_{j}\chi_{A_{j}} $ 为 $ f $ 的标准形，则由命题 8.2 可见
\[
\int_{Y} f\,dy = \sum_{j=0}^{m} e_{j}\lambda_{n}(A_{j}) = \sum_{j=0}^{m} e_{j}\int_{\Phi^{-1}(A_{j})} |\det \partial\Phi|\,dx = \sum_{j=0}^{m}\int_{\Phi^{-1}(A_{j})} e_{j}|\det \partial\Phi|\,dx = \int_{X} (f\circ\Phi)|\det \partial\Phi|\,dx.\quad\blacksquare
\]
\end{lemma}

\begin{theorem}

（换元法则） 设 $ \Phi\in\operatorname{Diff}^1(X,Y) $ 且 $ f\in E^Y $。则 $ f $ 属于 $ \mathcal{L}_1(Y,E) $ 当且仅当 $ (f\circ\Phi)|\det\partial\Phi| $ 属于 $ \mathcal{L}_1(X,E) $。此时有

\[
\int_Y f\,dy=\int_X (f\circ\Phi)|\det\partial\Phi|\,dx .
\]

证明

\begin{subitems}
\item

设 $ f\in\mathcal{L}_1(Y,E) $，并取序列 $ (f_j)\subset\mathcal{S}_c(Y,E) $，它在 $ \mathcal{L}_1(Y,E) $ 中几乎处处收敛到 $ f $，且满足 $ \lim_{j}\int_Y f_j=\int_Y f $（见引理 6.18、注 6.19(a) 与 (c)、定理 2.18）。令 $ g_j:=(f_j\circ\Phi)|\det\partial\Phi| $，$ j\in\mathbb{N} $。由引理 8.13 知 $ (g_j) $ 是 $ \mathcal{L}_1(X,E) $ 中的柯西序列，且 $ \int_Y f_j\,dy=\int_X g_j\,dx $。由于 $ \mathcal{L}_1(X,E) $ 完备，存在 $ g\in\mathcal{L}_1(X,E) $ 使 $ g_j\to g $ 于 $ \mathcal{L}_1(X,E) $。又由定理 2.18 知 $ \int_X g_j\,dx $ 收敛到 $ \int_X g\,dx $，且 $ (g_j) $ 的某个子序列 $ (g_{j_k})_{k\in\mathbb{N}} $ 在 $ X $ 中几乎处处收敛到 $ g $。于是 $ g $ 与 $ (f\circ\Phi)|\det\partial\Phi| $ 在 $ X $ 中几乎处处重合。由引理 2.15，$ (f\circ\Phi)|\det\partial\Phi| $ 属于 $ \mathcal{L}_1(X,E) $，且 $ \int_X g=\int_X (f\circ\Phi)|\det\partial\Phi| $。由此可得

\[
\int_Y f\,dy=\lim_{j}\int_Y f_j\,dy=\lim_{j}\int_X g_j\,dx=\int_X g\,dx=\int_X (f\circ\Phi)|\det\partial\Phi|\,dx .
\]
\item

反之，设 $ (f\circ\Phi)|\det\partial\Phi| $ 属于 $ \mathcal{L}_1(X,E) $。由 (8.5) 有

\[
f=\big((f\circ\Phi)|\det\partial\Phi|\big)\circ\Phi^{-1}|\det\partial(\Phi^{-1})| ,
\]

故 (i) 表明 $ f $ 属于 $ \mathcal{L}_1(Y,E) $。 \nobreak\hspace{0.35em}■

\end{subitems}

\end{theorem}

显然，推论 8.5 对 $ E $ 值映射也成立。由此可知，命题 8.6(ii)、8.9(ii) 与定理 8.11(ii) 对 $ E $ 值函数同样成立。

\exerciseshead

\begin{exercises}
\item

设 $ G\in\mathbb{R}^{n\times n} $ 对称且正定。证明

\[
\int_{\mathbb{R}^n} e^{-(Gx\mid x)}\,dx=\pi^{n/2}/\sqrt{\det G} .
\]

（提示：主轴变换。）

\item

证明：对 $ p\in\mathbb{C} $ 且 $ \operatorname{Re}p>n/2 $，有

\[
\int_{\mathbb{R}^n} (1+|x|^2)^{-p}\,dx=\pi^{n/2}\Gamma(p-n/2)/\Gamma(p) .
\]

（提示：参见例 6.13(b)。）

\item

设 $ D:=\{(x,y)\in\mathbb{R}^2 ; x,y\ge 0, x+y\le 1\} $ 且 $ p,q\in(0,\infty) $。证明：对 $ f:(0,1)\to\mathbb{R} $，函数

\[
D\to\mathbb{R}, \quad (x,y)\mapsto x^{p-1}y^{q-1}f(x+y)
\]

可积当且仅当 $ s\mapsto s^{p+q-1}f(s) $ 属于 $ \mathcal{L}_1((0,1)) $。此时有

\[
\int_D x^{p-1}y^{q-1}f(x+y)\,d(x,y)=\mathrm{B}(p,q)\int_0^1 s^{p+q-1}f(s)\,ds .
\]

（提示：考虑 $ (s,t)\mapsto(s(1-t),st) $。）

\item

设 $ 0\le\alpha<\beta\le 2\pi $，并设 $ f:[\alpha,\beta]\to(0,\infty) $ 可测。证明

\[
S(\alpha,\beta,f):=\{z\in\mathbb{C} ; \operatorname{arg}_N(z)\in[\alpha,\beta],\ |z|\le f(\operatorname{arg}_N(z))\}
\]

是勒贝格可测的，且

\[
\lambda_2(S(\alpha,\beta,f))=\frac{1}{2}\int_\alpha^\beta [f(\varphi)]^2\,d\varphi .
\]
\item

设 $ g\in\mathcal{L}^2_{\mathrm{sym}}(\mathbb{R}^n) $ 正定。计算由曲面 $ g^{-1}(1) $ 围成的立体椭球 $ g^{-1}([0,1]) $ 的体积（见注 VII.10.18）。

6（萨德引理）设 $ \Phi\in C^1(X,\mathbb{R}^n) $，并令 $ C:=\{x\in X ; \partial\Phi(x)\notin\mathcal{L}\mathrm{aut}(\mathbb{R}^n)\} $ 为 $ \Phi $ 的临界点集。证明 $ \Phi(C) $ 的测度为零。（提示：由于 $ C $ 是 $ \sigma $-紧的，只需验证对每个紧致的 $ n $ 维立方体 $ J $，$ \Phi(C\cap J) $ 的测度均为零。取 $ x_0\in C $ 与 $ r>0 $，使 $ J_0:=[x_0-(r/2)\mathbf{1},x_0+(r/2)\mathbf{1}] $ 紧包含在 $ X $ 中，并令

\[
\rho(r):=\max_{x\in J_0}\int_0^1\|\partial\Phi(x_0+t(x-x_0))\|\,dt .
\]

证明存在 $ c_n>0 $ 使 $ \lambda_n(\Phi(J_0))\le c_n r^n\rho(r) $。由于 $ \lim_{r\to 0}\rho(r)=0 $，将 $ J_0 $ 的各边细分即得结论。）

\item

设 $ \Phi\in C^1(X,\mathbb{R}^n) $ 且 $ C:=\{x\in X ; \partial\Phi(x)\notin\mathcal{L}\mathrm{aut}(\mathbb{R}^n)\} $。又设 $ \Phi|(X\setminus C) $ 是单射。证明：

\begin{subitems}
\item

对 $ f\in\mathcal{L}_0(X,\mathbb{R}^+) $，

\begin{equation}
\int_{\Phi(X)} f\,dy=\int_X (f\circ\Phi)|\det\partial\Phi|\,dx .\tag{8.29}
\end{equation}
\item

函数 $ f:\Phi(X)\to E $ 属于 $ \mathcal{L}_1(\Phi(X),E) $ 当且仅当 $ (f\circ\Phi)|\det\partial\Phi| $ 属于 $ \mathcal{L}_1(X,E) $。此时 (8.29) 成立。

\end{subitems}

\end{exercises}

\section{傅里叶变换}

为结束本章，我们引入最重要的积分变换，即所谓的傅里叶变换\footnote{本节的内容在本书的其余部分不会被用到}。对其基本性质的研究可以说就是对勒贝格积分理论的一次回顾：我们处处都会遇到勒贝格空间的完备性、控制收敛定理以及富比尼–托内利定理这样的基石。

尤其引人注目的是傅里叶变换与卷积、以及与 $ L_2 $ 的希尔伯特空间结构之间的相互作用。前者我们借助傅里叶乘子算子来说明，后者则借助普朗歇尔定理以及量子力学中位置算子与动量算子的应用来说明。

本节中我们只考虑定义在整个 $ \mathbb{R}^n $ 上的复值函数空间。因此，像第 7 节那样，我们在记号中省略 $ (\mathbb{R}^n,\mathbb{C}) $，例如把 $ \mathcal{L}_1(\mathbb{R}^n,\mathbb{C}) $ 写作 $ \mathcal{L}_1 $。此外，$ \int f\,dx $ 总是指 $ \int_{\mathbb{R}^n} f\,dx $，并且我们把 $ \mathbb{R}^n $ 与其对偶空间典范地等同，使得 $ \langle\cdot,\cdot\rangle $ 与欧几里得内积在形式上一致。

\needspace{3\baselineskip}\subsection*{定义与基本性质}

设 $ f \in \mathcal{L}_1 $。对每个 $ \xi \in \mathbb{R}^n $，映射 $ \mathbb{R}^n \to \mathbb{C} $，$ x \mapsto e^{-i\langle x,\xi\rangle}f(x) $ 都属于 $ \mathcal{L}_1 $。由

\begin{equation}
\hat{f}(\xi) := (2\pi)^{-n/2}\int_{\mathbb{R}^n} e^{-i\langle x,\xi\rangle} f(x)\,dx \quad \text{for } \xi \in \mathbb{R}^n\tag{9.1}
\end{equation}

定义的映射 $ \hat{f}: \mathbb{R}^n \to \mathbb{C} $ 称为 $ f $ 的\textbf{傅里叶变换}。映射 $ \mathcal{F} := (f \mapsto \hat{f}) $ 也称为傅里叶变换（或者在需要避免混淆时称为\textbf{傅里叶变换映射}）。

刚刚给出的定义中有不同的约定介入；人们常常看到不用 (9.1)，而把傅里叶变换定义为\footnote{见第 VIII.6 节}，

\[
\xi \mapsto \int e^{-i\langle x,\xi\rangle} f(x)\,dx \quad \text{或} \quad \xi \mapsto \int e^{-2\pi i\langle x,\xi\rangle} f(x)\,dx .
\]

显然，这些规范化上的差异对基本理论无关紧要；不过它们确实会使 $ 2\pi $ 的幂出现在后面的一些表达式中。阅读文献时应当留意这一点。这里选用的规范化有这样的优点：此类因子只在少数几处出现，而且普朗歇尔定理取特别简单的形式。

\begin{remark}

\begin{items}
\item

对 $ f \in \mathcal{L}_1 $，令 $ \mathcal{F}f := \hat{f} := \mathcal{F}f^{*} $，其中 $ f^{*} $ 是 $ f $ 的任意代表元。则 $ \mathcal{F}f $ 是良定义的，且 $ \mathcal{F} \in \mathcal{L}(\mathcal{L}_1, BC) $。

\proofhead
第一个断言是显然的。因为
\[
|\hat{f}(\xi)| \le (2\pi)^{-n/2}\|f\|_1 \quad \text{for } \xi \in \mathbb{R}^n ,
\]

第二个断言由定理 3.17（关于含参积分的连续性）与定理 VI.2.5 即易推出。\nobreak\hspace{0.35em}■

\item

对 $ f \in L_1 $，有 $ \widehat{\check{f}} = \check{\hat{f}} $。由

\[
\mathbb{R}^n \to \mathbb{C} , \quad \xi \mapsto \check{f}(\xi) = (2\pi)^{-n/2}\int e^{i\langle x,\xi\rangle} f(x)\,dx
\]

定义的函数称为 $ f $ 的\textbf{逆傅里叶变换}，其原因很快就会清楚；而映射

\[
\overline{\mathcal{F}} := (f \mapsto \check{f})
\]

是\textbf{逆傅里叶变换（逆傅里叶变换映射）}。由于反演 $ (f \mapsto \check{f}) $ 是 $ \mathcal{L}_1 $、$ L_1 $ 与 $ BC $ 上的连续自同构，逆傅里叶变换与傅里叶变换具有相同的连续性。

\proofhead
这由换元法则立即推出。\nobreak\hspace{0.35em}■
\item

对 $ \lambda > 0 $，我们用 $ \sigma_\lambda: \mathbb{R}^n \to \mathbb{R}^n $，$ x \mapsto \lambda x $ 表示按因子 $ \lambda $ 的\textbf{伸缩}。我们在 $ \operatorname{Funct} := \operatorname{Funct}(\mathbb{R}^n,\mathbb{C}) $ 上定义群 $ ((0,\infty),\cdot) $ 的一个作用

\begin{equation}
((0,\infty),\cdot) \times \operatorname{Funct} \to \operatorname{Funct} , \quad (\lambda,f) \mapsto \sigma_\lambda f ,\tag{9.2}
\end{equation}

其方式为令

\[
\sigma_\lambda f := f \circ \sigma_{1/\lambda} = (\sigma_{1/\lambda})^{*}f .
\]

若 $ V $ 是 $ \operatorname{Funct} $ 的向量子空间且在该作用下不变（即对 $ \lambda > 0 $ 有 $ \sigma_\lambda(V) \subset V $），则映射

\[
\sigma_\lambda: V \to V , \quad v \mapsto \sigma_\lambda v
\]

是线性的，并且对 $ \lambda,\mu > 0 $ 满足 $ \sigma_\lambda\sigma_\mu = \sigma_{\lambda\mu} $ 与 $ \sigma_1 = \operatorname{id}_V $。因此 $ \sigma_\lambda $ 是 $ V $ 的向量空间自同构，且对 $ \lambda > 0 $ 有 $ (\sigma_\lambda)^{-1} = \sigma_{1/\lambda} $。这表明

\[
((0,\infty),\cdot) \to \operatorname{Aut}(V) , \quad \lambda \mapsto \sigma_\lambda
\]

是乘法群 $ ((0,\infty),\cdot) $ 在 $ V $ 上的一个线性表示。特别地，$ \{\sigma_\lambda ; \lambda > 0\} $ 是 $ \operatorname{Aut}(V) $ 的子群，即 $ V $ 上的\textbf{伸缩群}。相应地，$ \sigma_\lambda v $ 就是 $ v $ 按因子 $ \lambda $ 的\textbf{伸缩}。与平移群的情形一样，若 $ V $ 在 (9.2) 下不变，我们就说 $ ((0,\infty),\cdot) $ 在 $ V $ 上\textbf{线性可表示}。
设 $ 1 \le p \le \infty $。则 $ ((0,\infty),\cdot) $ 在 $ L_p $ 上线性可表示，且

\[
\|\sigma_\lambda f\|_p = \lambda^{n/p}\|f\|_p .
\]

\proofhead
这由换元法则推出。\nobreak\hspace{0.35em}■
\item

对 $ \lambda > 0 $，$ \mathcal{F}\sigma_\lambda = \lambda^n\sigma_{1/\lambda}\mathcal{F} $。

\proofhead
设 $ f \in \mathcal{L}_1 $ 且 $ \lambda > 0 $。则
\[
\mathcal{F}\sigma_\lambda f(\xi) = (2\pi)^{-n/2}\int e^{-i\langle x,\xi\rangle} f(x/\lambda)\,dx = \lambda^n(2\pi)^{-n/2}\int e^{-\langle x/\lambda,\lambda\xi\rangle} f(x/\lambda)\lambda^{-n}\,dx
\]

对 $ \xi \in \mathbb{R}^n $。但换元法则表明最后一个表达式等于 $ \lambda^n\hat{f}(\lambda\xi) $。\nobreak\hspace{0.35em}■

\item

设 $ a \in \mathbb{R}^n $。则 $ (e^{i\langle a,\cdot\rangle}f)^{\wedge} = \tau_a\hat{f} $，$ f \in \mathcal{L}_1 $。\nobreak\hspace{0.35em}■

\end{items}

\end{remark}

\needspace{3\baselineskip}\subsection*{速降函数空间}

现在我们引入 $ \mathcal{L}_1 $ 的一个向量子空间，在它上面傅里叶变换特别容易处理。随后利用稠密性论证，我们就能把结果推广到更大的函数空间。

若对每个 $ (k,m) \in \mathbb{N}^2 $ 都存在 $ c_{k,m} > 0 $ 使得

\[
(1+|x|^2)^k|\partial^{\alpha}f(x)| \le c_{k,m} \quad \text{for } x \in \mathbb{R}^n, \alpha \in \mathbb{N}^n, \text{ and } |\alpha| \le m ,
\]

则称 $ f \in C^{\infty} $ 是\textbf{速降的}。

换言之，$ f \in C^{\infty} $ 速降是指：当 $ |x| \to \infty $ 时，每个导数 $ \partial^{\alpha}f $ 都趋于零的速度快于 $ 1/|x| $ 的任何幂。

现在令

\[
q_{k,m}(f) := \max_{|\alpha| \le m}\sup_{x \in \mathbb{R}^n}(1+|x|^2)^{k/2}|\partial^{\alpha}f(x)| \quad \text{for } f \in C^{\infty} \text{ and } k,m \in \mathbb{N} .
\]

空间

\[
\mathcal{S} := \{ f \in C^{\infty} ; q_{k,m}(f) < \infty \text{ for } k,m \in \mathbb{N} \}
\]

称为\textbf{施瓦茨空间}或\textbf{速降函数空间}。

\begin{remark}

\begin{items}
\item

$ \mathcal{S} $ 是 $ BUC^{\infty} $ 的向量子空间。每个 $ q_{k,m} $ 都是 $ \mathcal{S} $ 上的范数。

\proofhead
设 $ m \in \mathbb{N} $。由于 $ q_{0,m} $ 是 $ BC^m $ 上的范数，$ \mathcal{S} $ 是 $ BC^m $ 的向量子空间。设 $ \alpha \in \mathbb{N}^n $ 且 $ |\alpha| \le m $。则由中值定理易得 $ \partial^{\alpha}f $ 一致连续。这证明了第一个断言。第二个断言是显然的。\nobreak\hspace{0.35em}■
\item

对 $ (f,g) \in \mathcal{S}\times\mathcal{S} $，令

\[
d(f,g) := \sum_{k,m=0}^{\infty} 2^{-(k+m)}\frac{q_{k,m}(f-g)}{1+q_{k,m}(f-g)} .
\]

则 $ (\mathcal{S},d) $ 是度量空间。

\begin{subitems}
\item

\proofhead
(i) 显然，对每个 $ f \in \mathcal{S} $，二重级数 $ \sum 2^{-(k+m)}q_{k,m}(f)/(1+q_{k,m}(f)) $ 收敛。因此 $ d: \mathcal{S}\times\mathcal{S} \to \mathbb{R}^+ $ 是良定义的。此外 $ d $ 是对称的，且在 $ \mathcal{S}\times\mathcal{S} $ 的对角线上恒为零。
\item

因为 $ t \mapsto t/(1+t) $ 在 $ \mathbb{R}^+ $ 上递增，故对 $ r,s,t \in \mathbb{R}^+ $ 且 $ r \le s+t $ 有

\[
\frac{r}{1+r} \le \frac{s+t}{1+s+t} = \frac{s}{1+s+t} + \frac{t}{1+s+t} \le \frac{s}{1+s} + \frac{t}{1+t} .
\]

由此易得 $ d $ 满足三角不等式。\nobreak\hspace{0.35em}■

\end{subitems}

\item

对 $ f \in \mathcal{S} $ 及 $ \mathcal{S} $ 中的序列 $ (f_j) $，下列条件等价：

\begin{subitems}
\item

在 $ (\mathcal{S},d) $ 中 $ \lim f_j = f $；

\item

在 $ (\mathcal{S},d) $ 中 $ \lim(f-f_j) = 0 $；

\item

对 $ k,m \in \mathbb{N} $，$ \lim_j q_{k,m}(f-f_j) = 0 $。
于是，序列 $ (f_j) $ 在 $ \mathcal{S} $ 中收敛于 $ f $ 当且仅当 $ (f_j - f) $ 关于每个半范数 $ q_{k,m} $ 收敛于零。

\proofhead
“(i)⇒(ii)” 这一蕴含是显然的。
“(ii)⇒(iii)” 取 $ \varepsilon \in (0,1] $ 及 $ k,m \in \mathbb{N} $。存在 $ N \in \mathbb{N} $ 使得不等式 $ d(f,f_j) < \varepsilon/2^{k+m+1} $ 对 $ j \ge N $ 成立。于是
\[
\frac{2^{-(k+m)}q_{k,m}(f-f_j)}{1+q_{k,m}(f-f_j)} < \frac{\varepsilon}{2^{k+m+1}} ,
\]

故当 $ j \ge N $ 时 $ q_{k,m}(f-f_j) < \varepsilon $。
“(iii)⇒(i)” 取 $ \varepsilon > 0 $。存在 $ N \in \mathbb{N} $ 使得

\[
\sum_{k+m=N+1}^{\infty}\frac{2^{-(k+m)}q_{k,m}(f-f_j)}{1+q_{k,m}(f-f_j)} \le \sum_{\ell=N+1}^{\infty} 2^{-\ell} < \frac{\varepsilon}{2} .
\]

依假设，存在 $ M \in \mathbb{N} $ 使得对 $ j \ge M $ 及 $ k+m \le N $ 有 $ q_{k,m}(f-f_j) \le \varepsilon/4 $。因此

\[
d(f,f_j) \le \sum_{k,m=0}^{N}\frac{2^{-(k+m)}q_{k,m}(f-f_j)}{1+q_{k,m}(f-f_j)} + \frac{\varepsilon}{2} \le \varepsilon \quad \text{for } j \ge M .
\]
\end{subitems}

\item

$ \mathcal{D} $ 是 $ \mathcal{S} $ 的稠密向量子空间。函数 $ \mathbb{R}^n \to \mathbb{R} $，$ x \mapsto e^{-|x|^2} $ 属于 $ \mathcal{S} $ 但不属于 $ \mathcal{D} $。

\proofhead
显然 $ \mathcal{D} $ 是 $ \mathcal{S} $ 的向量子空间。设 $ f \in \mathcal{S} $。选取 $ \varphi \in \mathcal{D} $ 使 $ \varphi | \overline{\mathbb{B}}^n = 1 $，并令
\[
f_j(x) := f(x)\varphi(x/j) \quad \text{for } x \in \mathbb{R}^n , \quad j \in \mathbb{N}^{\times} .
\]

则 $ f_j $ 属于 $ \mathcal{D} $，且

\[
f(x) - f_j(x) = f(x)(1-\varphi(x/j)) \quad \text{for } x \in \mathbb{R}^n .
\]

因此对 $ x \in j\overline{\mathbb{B}}^n $ 及 $ \alpha \in \mathbb{N}^n $ 有 $ \partial^{\alpha}(f-f_j)(x) = 0 $。现在莱布尼茨法则表明存在 $ c = c(\varphi,m) > 0 $ 使得

\begin{equation*}
\begin{aligned}
|\partial^{\alpha}(f-f_j)(x)| &= \left|\sum_{\beta \le \alpha}\binom{\alpha}{\beta}\partial^{\beta}f(x)\partial^{\alpha-\beta}(1-\varphi)(x/j)j^{-|\alpha-\beta|}\right| \le c\max_{\beta \le \alpha}|\partial^{\beta}f(x)| \\ &\le c\,q_{k+1,m}(f)(1+|x|^2)^{-(k+1)/2}
\end{aligned}
\end{equation*}

对 $ x \in \mathbb{R}^n $，$ j \in \mathbb{N}^{\times} $，$ k \in \mathbb{N} $ 及 $ |\alpha| \le m $。令 $ C := c\,q_{k+1,m}(f) $，得

\begin{equation*}
\begin{aligned}
q_{k,m}(f-f_j) &= \max_{|\alpha| \le m}\sup_{|x| \ge j}(1+|x|^2)^{k/2}|\partial^{\alpha}(f-f_j)(x)| \\ &\le c\,q_{k+1,m}(f)\sup_{|x| \ge j}(1+|x|^2)^{-1/2} \le C/j ,
\end{aligned}
\end{equation*}

且当 $ j \to \infty $ 时，第一个断言由 (c) 推出。第二个断言是显然的。\nobreak\hspace{0.35em}■

\item

对 $ m \in \mathbb{N} $，有 $ \mathcal{S} \hookrightarrow BUC^m $。

\proofhead
这由 (a) 与 (c) 推出。\nobreak\hspace{0.35em}■
\item

$ \mathcal{S} $ 是 $ C_0 $ 的稠密向量子空间。

\proofhead
设 $ f \in \mathcal{S} $。由 (a) 以及 $ |f(x)| \le q_{1,0}(f)(1+|x|^2)^{-1/2} $（$ x \in \mathbb{R}^n $）可知 $ f $ 属于 $ C_0 $。因此 $ \mathcal{S} $ 是 $ C_0 $ 的向量子空间。定理 7.13 表明 $ \mathcal{D} $ 是 $ C_0 $ 的稠密向量子空间，于是由包含关系 $ \mathcal{D} \subset \mathcal{S} \subset C_0 $ 即得结论。\nobreak\hspace{0.35em}■
\item

对 $ k,m \in \mathbb{N} $，存在正常数 $ c $ 与 $ C $ 使得

\[
c\max_{|\alpha| \le m,\,|\beta| \le k}\sup_{x \in \mathbb{R}^n}|\partial^{\alpha}(x^{\beta}f(x))| \le q_{k,m}(f) \le C\max_{|\alpha| \le m,\,|\beta| \le k}\sup_{x \in \mathbb{R}^n}|x^{\beta}\partial^{\alpha}f(x)| \quad \text{for } f \in \mathcal{S} .
\]

\proofhead
这易由莱布尼茨法则推出。\nobreak\hspace{0.35em}■
\item

设 $ f \in \mathcal{S} $ 且 $ \alpha,\beta \in \mathbb{N}^n $。则 $ x \mapsto x^{\alpha}\partial^{\beta}f(x) $ 属于 $ \mathcal{S} $。

\proofhead
这是 (g) 的推论。\nobreak\hspace{0.35em}■
\item

反演 $ f \mapsto \check{f} $ 是 $ \mathcal{S} $ 的连续自同构。

\proofhead
这是显然的。\nobreak\hspace{0.35em}■
\end{items}

\end{remark}

\begin{theorem}

设 $ p \in [1,\infty) $。则 $ \mathcal{S} $ 是 $ L_p $ 的稠密向量子空间，且存在 $ c = c(n,p) > 0 $ 使得

\begin{equation}
\|f\|_p \le c\,q_{n+1,0}(f) \quad \text{for } f \in \mathcal{S} .\tag{9.3}
\end{equation}

\proofhead
对 $ f \in \mathcal{S} $，有
\begin{equation}
\begin{aligned}
\int |f|^p\,dx &= \int |f(x)|^p(1+|x|^2)^{(n+1)p/2}(1+|x|^2)^{-(n+1)p/2}\,dx \\ &\le (q_{n+1,0}(f))^p\int (1+|x|^2)^{-(n+1)p/2}\,dx .
\end{aligned}\tag{9.4}
\end{equation}

进一步，由定理 8.11(i) 及 $ (n+1)p > n $，有

\[
\int_{[|x| \ge 1]}|x|^{-(n+1)p}\,dx = n\omega_n\int_1^{\infty} r^{-((n+1)p-n+1)}\,dr < \infty .
\]

因此 $ \int(1+|x|^2)^{-(n+1)p/2}\,dx $ 也有限，且 (9.3) 由 (9.4) 推出。特别地，$ f $ 属于 $ L_p $，从而我们看到 $ \mathcal{S} $ 是 $ L_p $ 的向量子空间。由定理 7.13，$ \mathcal{D} $ 是 $ L_p $ 的稠密向量子空间，而由注 9.2(d)，它包含于 $ \mathcal{S} $。结论成立。\nobreak\hspace{0.35em}■

\end{theorem}

\needspace{3\baselineskip}\subsection*{卷积代数 $ \mathcal{S} $}

由注 9.2(a) 与定理 9.3，$ \mathcal{S}\times\mathcal{S} $ 包含于 $ BUC^{\infty}\times L_1 $。因此卷积在 $ \mathcal{S}\times\mathcal{S} $ 上有定义，且由推论 7.9 有

\begin{equation}
*: \mathcal{S}\times\mathcal{S} \to BUC^{\infty} .\tag{9.5}
\end{equation}

下一个结果表明，对 $ (f,g) \in \mathcal{S}\times\mathcal{S} $，$ f*g $ 实际上是速降的。

\begin{proposition}

卷积 $ \mathcal{S}\times\mathcal{S} $ 是映入 $ \mathcal{S} $ 的连续双线性映射。

证明

\begin{subitems}
\item

下面验证卷积 $ \mathcal{S}\times\mathcal{S} $ 映入 $ \mathcal{S} $。设 $ (f,g) \in \mathcal{S}\times\mathcal{S} $ 且 $ k,m \in \mathbb{N} $。由 (9.5)，只需检验 $ q_{k,m}(f*g) $ 有限。因为

\[
|x|^k \le (|x-y|+|y|)^k = \sum_{j=0}^{k}\binom{k}{j}|x-y|^j|y|^{k-j} \quad \text{for } x,y \in \mathbb{R}^n ,
\]

存在 $ c_k > 0 $ 使得

\begin{equation*}
\begin{aligned}
|x|^k|f*g(x)| &\le \int\sum_{j=0}^{k}\binom{k}{j}|x-y|^j|f(x-y)||y|^{k-j}|g(y)|\,dy \\ &\le c_k q_{k,0}(f)\int(1+|y|^2)^{k/2}|g(y)|\,dy .
\end{aligned}
\end{equation*}

注意到 $ \widetilde{c}_n := \int(1+|y|^2)^{-(n+1)/2}\,dy $ 有限，我们得到

\[
|x|^k|f*g(x)| \le c_k\widetilde{c}_n q_{k,0}(f)q_{k+n+1,0}(g) .
\]

于是由注 9.2(g)，存在 $ c = c(k,n) \ge 1 $ 使得

\begin{equation}
q_{k,0}(f*g) \le c\,q_{k,0}(f)q_{k+n+1,0}(g) .\tag{9.6}
\end{equation}

最后，由定理 7.8(iv)，有

\[
q_{k,m}(f*g) = \max_{|\alpha| \le m}q_{k,0}(\partial^{\alpha}(f*g)) = \max_{|\alpha| \le m}q_{k,0}((\partial^{\alpha}f)*g) ,
\]

而 (9.6) 蕴含

\begin{equation}
q_{k,m}(f*g) \le c\max_{|\alpha| \le m}q_{k,0}(\partial^{\alpha}f)q_{k+n+1,0}(g) = c\,q_{k,m}(f)q_{k+n+1,0}(g) .\tag{9.7}
\end{equation}
\item

显然卷积是双线性的。设 $ (f,g) \in \mathcal{S}\times\mathcal{S} $，$ ((f_j,g_j))_{j \in \mathbb{N}} $ 是 $ \mathcal{S}\times\mathcal{S} $ 中的序列，使得当 $ j \to \infty $ 时 $ (f_j,g_j) \to (f,g) $ 于 $ \mathcal{S}\times\mathcal{S} $。
又令

\[
\alpha := c(q_{k,m}(f) + q_{k+n+1,0}(g) + 1)
\]

其中 $ c $ 是 (9.7) 中的常数。设 $ \varepsilon \in (0,1] $。由注 9.2(c)，存在 $ N \in \mathbb{N} $ 使得

\[
q_{k,m}(f-f_j) < \varepsilon/\alpha , \quad q_{k+n+1,0}(g-g_j) < \varepsilon/\alpha \quad \text{for } j \ge N .
\]

因为

\[
f*g - f_j*g_j = (f-f_j)*g + (f_j-f)*(g-g_j) + f*(g-g_j) ,
\]

由 (9.7) 推得

\begin{equation*}
\begin{aligned}
q_{k,m}(f*g-f_j*g_j) &\le c\big(q_{k,m}(f-f_j)q_{k+n+1,0}(g) + q_{k,m}(f-f_j)q_{k+n+1,0}(g-g_j) \\ &\qquad + q_{k,m}(f)q_{k+n+1,0}(g-g_j)\big) < \varepsilon
\end{aligned}
\end{equation*}

对 $ j \ge N $。于是证毕。\nobreak\hspace{0.35em}■

\end{subitems}

\end{proposition}

\begin{corollary}

$ (\mathcal{S},+,*) $ 是可交换代数 $ (L_1,+,*) $ 的子代数。

\proofhead
这由命题 9.4 与定理 9.3 推出。\nobreak\hspace{0.35em}■
\end{corollary}

\needspace{3\baselineskip}\subsection*{傅里叶变换的计算}

现在导出关于导数的傅里叶变换以及傅里叶变换的导数的一些法则。为使这些公式的表述简化，令 $ \Lambda(x) := (1+|x|^2)^{1/2} $（$ x \in \mathbb{R}^n $），并令

\[
D_j := -i\partial_j , \quad j \in \{1,\ldots,n\} \quad \text{for } D^{\alpha} := D_1^{\alpha_1}\cdots D_n^{\alpha_n} , \quad \alpha \in \mathbb{N}^n ,
\]

其中 $ i $ 是虚数单位。照例，由多项式 $ p \in \mathbb{C}[X_1,\ldots,X_n] $ 诱导的多项式函数也用 $ p $ 表示。

\begin{proposition}

设 $ f \in \mathcal{L}_1 $。

\begin{subitems}
\item

对 $ \alpha \in \mathbb{N}^n $，设 $ D^{\alpha}f $ 存在且属于 $ \mathcal{L}_1 $。则 $ X^{\alpha}\hat{f} = \widehat{D^{\alpha}f} $。

\item

对 $ m \in \mathbb{N} $，设 $ \Lambda^m f $ 属于 $ \mathcal{L}_1 $。则 $ \hat{f} $ 属于 $ BC^m $，且

\[
D^{\alpha}\hat{f} = (-1)^{|\alpha|}\widehat{X^{\alpha}f} \quad \text{for } \alpha \in \mathbb{N}^n , \quad |\alpha| \le m .
\]

\proofhead
(i) 设 $ \{\varphi_{\varepsilon} ; \varepsilon > 0\} $ 是光滑核。由分部积分（见习题 7.10）推得
\begin{equation}
\begin{aligned}
\int \xi^{\alpha}e^{-i\langle x,\xi\rangle}(f*\varphi_{\varepsilon})(x)\,dx &= (-1)^{|\alpha|}\int D_x^{\alpha}(e^{-i\langle x,\xi\rangle})(f*\varphi_{\varepsilon})(x)\,dx \\ &= \int e^{-i\langle x,\xi\rangle}((D^{\alpha}f)*\varphi_{\varepsilon})(x)\,dx .
\end{aligned}\tag{9.8}
\end{equation}

定理 7.11 与定理 2.18(ii) 表明

\[
\lim_{\varepsilon\to0}(2\pi)^{-n/2}\int\xi^\alpha e^{-i\langle x,\xi\rangle}(f*\varphi_\varepsilon)(x)\,dx = \xi^\alpha\widehat{f}(\xi)
\]

以及

\[
\lim_{\varepsilon\to0}(2\pi)^{-n/2}\int e^{-i\langle x,\xi\rangle}\left((D^\alpha f)*\varphi_\varepsilon\right)(x)\,dx = \widehat{D^\alpha f}(\xi)
\]

对 $ \xi \in \mathbb{R}^n $ 成立。利用 (9.8)，这就证明了断言。

\end{subitems}

(ii) 我们令 $ h(x,\xi) := e^{-i\langle x,\xi\rangle}f(x) $，$ (x,\xi) \in \mathbb{R}^n \times \mathbb{R}^n $。则对每个 $ \xi \in \mathbb{R}^n $，$ h(\cdot,\xi) $ 属于 $ \mathcal{L}_1 $；对每个 $ x \in \mathbb{R}^n $，$ h(x,\cdot) $ 属于 $ C^\infty $。此外还有

\[
D_\xi^\alpha h(x,\xi) = (-1)^{|\alpha|}x^\alpha h(x,\xi) \qquad \text{对 } (x,\xi) \in \mathbb{R}^{2n}, \quad \alpha \in \mathbb{N}^n,
\]

从而

\begin{equation}
|D_\xi^\alpha h(x,\xi)| \le (1+|x|^2)^{|\alpha|/2}|h(x,\xi)| = \Lambda^{|\alpha|}(x)|f(x)| .\tag{9.9}
\end{equation}

于是由含参积分求导定理可知，$ \widehat{f} $ 属于 $ C^m $，并且

\begin{equation*}
\begin{aligned}
D^\alpha\widehat{f}(\xi) &= (2\pi)^{-n/2}\int D_\xi^\alpha h(x,\xi)\,dx = (2\pi)^{-n/2}(-1)^{|\alpha|}\int x^\alpha h(x,\xi)\,dx \\ &= (-1)^{|\alpha|}\widehat{X^\alpha f}(\xi)
\end{aligned}
\end{equation*}

对 $ \xi \in \mathbb{R}^n $ 及满足 $ |\alpha| \le m $ 的 $ \alpha \in \mathbb{N}^n $ 成立。最后，(9.9) 表明

\[
|D^\alpha\widehat{f}(\xi)| \le (2\pi)^{-n/2}\int|D_\xi^\alpha h(x,\xi)|\,dx \le (2\pi)^{-n/2}\|\Lambda^{|m|}f\|_1 < \infty \qquad \text{对 } \xi \in \mathbb{R}^n .
\]

因此 $ \widehat{f} $ 属于 $ BC^m $。\nobreak\hspace{0.35em}■

\end{proposition}

\begin{proposition}

傅里叶变换连续且线性地把 $ \mathcal{S} $ 映入自身。

证明

\begin{subitems}
\item

设 $ f \in \mathcal{S} $ 且 $ m \in \mathbb{N} $。则

\begin{equation*}
\begin{aligned}
\int\Lambda^m(x)|f(x)|\,dx &= \int(1+|x|^2)^{(m+n+1)/2}|f(x)|(1+|x|^2)^{-(n+1)/2}\,dx \\ &\le q_{m+n+1,0}(f)\int(1+|x|^2)^{-(n+1)/2}\,dx < \infty .
\end{aligned}
\end{equation*}

由命题 9.6(ii) 可知 $ \widehat{f} $ 属于 $ BC^m $，从而属于 $ BC^\infty $。

\item

设 $ k,m \in \mathbb{N} $ 且 $ \alpha,\beta \in \mathbb{N}^n $，满足 $ |\alpha| \le m $、$ |\beta| \le k $。又设 $ f \in \mathcal{S} $。则由注 9.2(h) 与定理 9.3 可知，$ \Lambda^m f $ 与 $ D^\beta(X^\alpha f) $ 属于 $ \mathcal{L}_1 $。因此命题 9.6 给出

\begin{equation}
\xi^\beta D^\alpha\widehat{f}(\xi) = (-1)^{|\alpha|}\xi^\beta\widehat{X^\alpha f}(\xi) = (-1)^{|\alpha|}\widehat{D^\beta(X^\alpha f)}(\xi) \qquad \text{对 } \xi \in \mathbb{R}^n .\tag{9.10}
\end{equation}

注 9.2(g) 表明存在 $ c > 0 $ 使得

\begin{equation*}
\begin{aligned}
|\xi^\beta D^\alpha\widehat{f}(\xi)| &\le (2\pi)^{-n/2}\int|D^\beta(X^\alpha f)(x)|(1+|x|^2)^{(n+1)/2}(1+|x|^2)^{-(n+1)/2}\,dx \\ &\le c\, q_{m+n+1,k}(f)
\end{aligned}
\end{equation*}

对 $ |\alpha| \le m $ 与 $ |\beta| \le k $ 成立。因此存在 $ C > 0 $ 使得

\begin{equation}
q_{k,m}(\widehat{f}) \le Cq_{m+n+1,k}(f) .\tag{9.11}
\end{equation}

所以 $ \widehat{f} $ 属于 $ \mathcal{S} $。傅里叶变换的连续性现在容易由 (9.11) 与注 9.2(c) 推出。这样就完成了。\nobreak\hspace{0.35em}■

\end{subitems}

\end{proposition}

\begin{corollary}

对 $ f \in \mathcal{S} $ 与 $ \alpha \in \mathbb{N}^n $，有

\[
\widehat{D^\alpha f} = X^\alpha\widehat{f} \qquad \text{与} \qquad \widehat{X^\alpha f} = (-1)^{|\alpha|}D^\alpha\widehat{f} .
\]

\proofhead
这些都是 (9.10) 的特殊情形。\nobreak\hspace{0.35em}■
\end{corollary}

命题 9.6 与推论 9.8 表明，傅里叶变换把微分变成乘以函数，反过来也一样。这一事实是其巨大效用的重要来源。

现在容易把注 9.1(a) 的陈述改进为：$ L_1 $ 在 $ \mathcal{F} $ 下的像已经落在 $ C_0 $ 中。

9.9 命题\footnote{也称为黎曼–勒贝格引理。}（黎曼–勒贝格）$ \mathcal{F} \in \mathcal{L}(L_1,C_0) $。

\begin{proof}

命题 9.7 与 $ \mathcal{S} \subset C_0 $ 蕴含 $ \mathcal{F}(\mathcal{S}) \subset C_0 $。由定理 9.3 知 $ \mathcal{S} $ 是 $ L_1 $ 的稠密向量子空间，而注 9.1(a) 保证 $ \mathcal{F} $ 连续地把空间 $ L_1 $ 映入 $ BC $。断言随之成立，因为 $ C_0 $ 是 $ BC $ 的闭向量子空间。\nobreak\hspace{0.35em}■

\end{proof}

\begin{example}

\begin{items}
\item

对 $ g := g_n\colon \mathbb{R}^n \to \mathbb{R} $，$ x \mapsto e^{-|x|^2/2} $，有 $ \widehat{g} = g $。

\begin{subitems}
\item

\proofhead
(i) 指数函数的一个性质蕴含
\[
g_n(x) = g_1(x_1)\cdot\ldots\cdot g_1(x_n) \qquad \text{对 } x = (x_1,\ldots,x_n) \in \mathbb{R}^n .
\]

为清楚起见，我们用 $ \mathcal{F}_n $ 表示 $ \mathbb{R}^n $ 上的傅里叶变换。则由富比尼–托内利定理可得

\begin{equation*}
\begin{aligned}
\mathcal{F}_n(g_n)(\xi) &= (2\pi)^{-n/2}\int_{\mathbb{R}^n}e^{-i\langle x,\xi\rangle}e^{-|x|^2/2}\,dx = (2\pi)^{-n/2}\int_{\mathbb{R}^n}\prod_{j=1}^n e^{-ix_j\xi_j}e^{-x_j^2/2}\,dx \\ &= \prod_{j=1}^n(2\pi)^{-1/2}\int_{\mathbb{R}}e^{-ix_j\xi_j}e^{-x_j^2/2}\,dx_j = \prod_{j=1}^n\mathcal{F}_1(g_1)(\xi_j) .
\end{aligned}
\end{equation*}

这表明只需处理一维情形。

\item

因此设 $ n = 1 $。对 $ f := \widehat{g} $，由例 8.7 得

\[
f(0) = \widehat{g}(0) = \frac{1}{\sqrt{2\pi}}\int_{-\infty}^{\infty}e^{-x^2/2}\,dx = 1 .
\]

因为 $ xe^{-x^2/2} = -\partial(e^{-x^2/2}) $，即因为 $ Xg = -\partial g = -iDg $，推论 9.8 给出

\[
\partial f = \partial\widehat{g} = iD\widehat{g} = -i\widehat{Xg} = -\widehat{Dg} = -X\widehat{g} = -Xf .
\]

因此 $ f $ 在 $ \mathbb{R} $ 上解线性初值问题 $ y'(t) = -ty(t) $、$ y(0) = 1 $；其唯一解是 $ g $。\nobreak\hspace{0.35em}■

\end{subitems}

\item

采用 (a) 与 (7.11) 的记号，有

\[
\widehat{g(\varepsilon\,\cdot)}(\xi) = g_\varepsilon(\xi) \qquad \text{对 } \xi \in \mathbb{R}^n, \quad \varepsilon > 0 .
\]

\proofhead
因为 $ g(\varepsilon\,\cdot) = \sigma_{1/\varepsilon}g $，这由 (a) 与注 9.1(d) 推出。\nobreak\hspace{0.35em}■
\item

设

\[
\varphi(x) := (2\pi)^{-n/2}e^{-|x|^2} \qquad \text{对 } x \in \mathbb{R}^n ,
\]

并令 $ \varepsilon > 0 $。则 $ \widehat{\varphi(\varepsilon\,\cdot)} = k_\varepsilon $，其中 $ k_1 = k $ 是\textbf{高斯核}。

\proofhead
由 $ \varphi = (2\pi)^{-n/2}\sigma_{1/\sqrt{2}}g $ 与注 9.1(c) 可得
\[
\varphi(\varepsilon\,\cdot) = \sigma_{1/\varepsilon}\varphi = (2\pi)^{-n/2}\sigma_{1/\sqrt{2}\,\varepsilon}g = (2\pi)^{-n/2}g(\sqrt{2}\,\varepsilon\,\cdot) .
\]

于是由 (b) 得到

\[
\widehat{\varphi(\varepsilon\,\cdot)}(x) = (2\pi)^{-n/2}g_{\sqrt{2}\,\varepsilon}(x) = \varepsilon^{-n}(4\pi)^{-n/2}e^{-|x|^2/4\varepsilon^2} = k_\varepsilon(x)
\]

对 $ x \in \mathbb{R}^n $ 成立。\nobreak\hspace{0.35em}■

\end{items}

\end{example}

\needspace{3\baselineskip}\subsection*{傅里叶积分定理}

为了对 $ L_1 $ 上的傅里叶变换作更深入的研究，我们先给出以下结果。

\begin{proposition}

设 $ f,g \in L_1 $。则 $ \widehat{f}g $ 与 $ f\widehat{g} $ 属于 $ L_1 $，且

\[
\int\widehat{f}g\,dx = \int f\widehat{g}\,dx .
\]

\proofhead
由命题 9.9 容易推出 $ \widehat{f}g $ 与 $ f\widehat{g} $ 属于 $ L_1 $。设 $ \overset{*}{f} $ 与 $ \overset{*}{g} $ 分别是 $ f $ 与 $ g $ 的代表元。则由引理 7.2 可知
\begin{equation}
h\colon \mathbb{R}^{2n} \to \mathbb{C}, \quad (x,y) \mapsto e^{-i\langle x,y\rangle}\overset{*}{f}(x)\overset{*}{g}(y)\tag{9.12}
\end{equation}

是可测的。因为

\begin{equation}
\iint|h(x,y)|\,dx\,dy = \|f\|_1\|g\|_1 ,\tag{9.13}
\end{equation}

我们可以对 $ h $ 应用富比尼–托内利定理，并得到

\begin{equation*}
\begin{aligned}
\int\widehat{\overset{*}{f}}(y)\overset{*}{g}(y)\,dy &= \int(2\pi)^{-n/2}\int e^{-i\langle x,y\rangle}\overset{*}{f}(x)\,dx\,\overset{*}{g}(y)\,dy \\ &= \int(2\pi)^{-n/2}\int e^{-i\langle x,y\rangle}\overset{*}{g}(y)\,dy\,\overset{*}{f}(x)\,dx = \int\widehat{\overset{*}{g}}(x)\overset{*}{f}(x)\,dx .
\end{aligned}
\end{equation*}

于是，注意到 $ \widehat{f} = \widehat{\overset{*}{f}} $ 与 $ \widehat{g} = \widehat{\overset{*}{g}} $，断言成立。\nobreak\hspace{0.35em}■

\end{proposition}

我们现在对函数及其变换作各种假设，来证明关于傅里叶变换之逆的定理。

\begin{theorem}

对 $ f \in L_1 $，以下陈述成立：

\begin{subitems}
\item

\[ \lim_{\varepsilon\to0}(2\pi)^{-n/2}\int e^{i\langle\cdot,\xi\rangle}\widehat{f}(\xi)e^{-\varepsilon^2|\xi|^2}\,d\xi = f \qquad \text{在 } L_1 \text{ 中} . \]

\item

（$ L_1 $ 的傅里叶积分定理）若 $ \widehat{f} $ 属于 $ L_1 $，则 $ f = \overline{\mathcal{F}}(\widehat{f}) $，其中 $ \overline{\mathcal{F}} $ 是傅里叶余变换。

\proofhead
(i) 我们采用例 9.10 的记号，并令
\[
\varphi^\varepsilon(\xi,y) := e^{i\langle\xi,y\rangle}\varphi(\varepsilon\xi) = (2\pi)^{-n/2}e^{i\langle y,\xi\rangle}e^{-\varepsilon^2|\xi|^2}
\]

对 $ \xi,y \in \mathbb{R}^n $ 与 $ \varepsilon > 0 $ 成立。对 $ y \in \mathbb{R}^n $，我们用 $ \widehat{\varphi^\varepsilon}(\cdot,y) $ 表示 $ \xi \mapsto \varphi^\varepsilon(\xi,y) $ 的傅里叶变换。由例 9.10(c) 与注 9.1(e) 可得

\[
\widehat{\varphi^\varepsilon}(x,y) = k_\varepsilon(y-x) \qquad \text{对 } x,y \in \mathbb{R}^n .
\]

因此命题 9.11 蕴含

\begin{equation*}
\begin{aligned}
(2\pi)^{-n/2}\int\widehat{f}(\xi)e^{i\langle y,\xi\rangle}e^{-\varepsilon^2|\xi|^2}\,d\xi &= \int\widehat{f}(\xi)\varphi^\varepsilon(\xi,y)\,d\xi \\ &= \int f(x)\widehat{\varphi^\varepsilon}(x,y)\,dx = k_\varepsilon*f(y)
\end{aligned}
\end{equation*}

对 $ y \in \mathbb{R}^n $ 成立。断言现在由定理 7.11 与例 7.12(a) 推出。

\end{subitems}

(ii) 若 $ \widehat{f} $ 属于 $ L_1 $，则由控制收敛定理可知

\[
\lim_{\varepsilon\to0}\int e^{i\langle y,\xi\rangle}\widehat{f}(\xi)e^{-\varepsilon^2|\xi|^2}\,d\xi = \int e^{i\langle y,\xi\rangle}\widehat{f}(\xi)\,d\xi = (2\pi)^{n/2}\mathcal{F}(\widehat{f})^{\sim}(y)
\]

对 $ y \in \mathbb{R}^n $ 成立。于是 (i)、注 9.1(b) 与定理 2.18(i) 完成了证明。\nobreak\hspace{0.35em}■

\end{theorem}

\qyxffnnumdef{1}
\begin{corollary}

\begin{subitems}
\item

（$ \mathcal{S} $ 的傅里叶积分定理）傅里叶变换是 $ \mathcal{S} $ 的连续自同构。它的逆是傅里叶余变换。

\item

傅里叶变换连续且单射地把 $ L_1 $ 映入 $ C_0 $，并且像集稠密。

\item

对 $ f \in L_1\cap BUC $，等式\qyxffnmark{1}

\[
f(x) = \lim_{\varepsilon\to0}(2\pi)^{-n/2}\int e^{i\langle x,\xi\rangle}\widehat{f}(\xi)e^{-\varepsilon^2|\xi|^2}\,d\xi
\]

关于 $ x \in \mathbb{R}^n $ 一致成立。

\item

对 $ f \in L_1\cap BUC $，设 $ \widehat{f} $ 属于 $ L_1 $。则

\[
f(x) = (2\pi)^{-n/2}\int e^{i\langle x,\xi\rangle}\widehat{f}(\xi)\,d\xi \qquad \text{对 } x \in \mathbb{R}^n .
\]

\proofhead
(i) 如同赋范向量空间的情形，我们用 $ \mathcal{L}(\mathcal{S}) $ 表示 $ \mathcal{S} $ 的全部连续自同态构成的向量空间；类似地，用 $ \mathcal{L}\mathrm{aut}(\mathcal{S}) $ 表示 $ \mathcal{S} $ 的自同构。则由注 9.2(i) 与命题 9.7 可知，$ \mathcal{F} $ 与 $ \overline{\mathcal{F}} $ 属于 $ \mathcal{L}(\mathcal{S}) $。由于 $ \mathcal{S} \subset L_1 $，定理 9.12(ii) 于是表明 $ \overline{\mathcal{F}} $ 是 $ \mathcal{F} $ 在 $ \mathcal{L}(\mathcal{S}) $ 中的左逆。又由 $ \widetilde{\widehat{u}} = \widehat{\widetilde{u}} $ 可得，对 $ f \in \mathcal{S} $ 有
\[
\mathcal{F}\overline{\mathcal{F}}f = \mathcal{F}(\mathcal{F}f)^{\sim} = \overline{\mathcal{F}}\mathcal{F}f = f .
\]

因此 $ \overline{\mathcal{F}} $ 也是 $ \mathcal{F} $ 在 $ \mathcal{L}(\mathcal{S}) $ 中的右逆，从而 $ \mathcal{F} \in \mathcal{L}\mathrm{aut}(\mathcal{S}) $。

\end{subitems}

\begin{subitems}
\item[（ii）]

若对 $ f \in L_1 $ 有 $ \widehat{f} = 0 $，则由定理 9.12(ii) 推出 $ f = 0 $。因此 $ \mathcal{F} $ 在 $ L_1 $ 上是单射；又由黎曼–勒贝格引理知 $ \mathcal{F} $ 属于 $ \mathcal{L}(L_1,C_0) $。由于 (i) 与 $ \mathcal{S} \subset L_1 $，我们有 $ \mathcal{S} = \mathcal{F}(\mathcal{S}) \subset \mathcal{F}(L_1) $。于是由注 9.2(f) 可知 $ \mathcal{F}(L_1) $ 在 $ C_0 $ 中稠密。

\item[（iii）]

由定理 9.12(i) 的证明与定理 7.11 推出。

\item[（iv）]

现在显然。\nobreak\hspace{0.35em}■

\end{subitems}

\end{corollary}
\qyxffntextof{1}{可以证明，对 $ f \in L_1\cap C $，(iii) 与 (iv) 仍然成立。}

\begin{remark}

\begin{items}
\item

对 $ f \in \mathcal{S} $，有 $ \widehat{\widehat{f}} = \widetilde{f} $。

\item

可以证明，在傅里叶变换下 $ L_1 $ 在 $ C_0 $ 中没有闭像（见 [Rud83]）。因此 $ \mathcal{F} \in \mathcal{L}(L_1,C_0) $ 不是满射。\nobreak\hspace{0.35em}■

\end{items}

\end{remark}

\needspace{3\baselineskip}\subsection*{卷积与傅里叶变换}

我们现在研究卷积在傅里叶变换下会变成什么。为此先引入另一个光滑函数空间；它在下一小节中将显得特别重要。

设 $ \varphi \in C^\infty $。若对每个 $ \alpha \in \mathbb{N}^n $，都存在常数 $ c_\alpha > 0 $ 与 $ k_\alpha \in \mathbb{N} $ 使得

\[
|\partial^\alpha\varphi(x)| \le c_\alpha(1+|x|^2)^{k_\alpha} \qquad \text{对 } x \in \mathbb{R}^n ,
\]

则称 $ \varphi $ 是\textbf{缓增的}。我们用 $ \mathcal{O}_M $ 表示具有该性质的全部函数之集，即\textbf{缓增函数空间}。

\begin{remark}

\begin{items}
\item

按向量子空间的意义，有包含关系 $ \mathcal{S} \subset \mathcal{O}_M \subset C^\infty $ 与 $ \mathbb{C}[X_1,\ldots,X_n] \subset \mathcal{O}_M $。

\item

$ (\mathcal{O}_M,+,\cdot) $ 是带单位元的交换代数。

\item

设 $ (\varphi,f) \in \mathcal{O}_M \times \mathcal{S} $。则 $ \varphi f $ 属于 $ \mathcal{S} $，且对每个 $ m \in \mathbb{N} $，存在 $ c = c(\varphi,m) > 0 $ 与 $ k' = k'(\varphi,m) \in \mathbb{N} $ 使得对 $ k \in \mathbb{N} $ 有 $ q_{k,m}(\varphi f) \le c\, q_{k+k',m}(f) $。

\proofhead
设 $ m \in \mathbb{N} $。则存在 $ c = c(\varphi,m) > 0 $ 与 $ k' = k'(\varphi,m) \in \mathbb{N} $ 使得
\[
|\partial^\alpha\varphi(x)| \le c(1+|x|^2)^{k'/2} \qquad \text{对 } x \in \mathbb{R}^n, \quad \alpha \in \mathbb{N}^n, \quad |\alpha| \le m .
\]

由莱布尼茨法则可得

\begin{equation*}
\begin{aligned}
q_{k,m}(\varphi f) &= \max_{|\alpha|\le m}\sup_{x\in\mathbb{R}^n}(1+|x|^2)^{k/2}\left|\sum_{\beta\le\alpha}\binom{\alpha}{\beta}\partial^\beta\varphi(x)\partial^{\alpha-\beta}f(x)\right| \\ &\le c\max_{|\alpha|\le m}\sup_{x\in\mathbb{R}^n}(1+|x|^2)^{(k+k')/2}|\partial^\alpha f(x)| = c\, q_{k+k',m}(f)
\end{aligned}
\end{equation*}

对 $ f \in \mathcal{S} $ 与 $ k \in \mathbb{N} $ 成立。\nobreak\hspace{0.35em}■

\item

设 $ \varphi \in \mathcal{O}_M $。则 $ f \mapsto \varphi f $ 是 $ \mathcal{S} $ 到自身的线性连续映射。

\proofhead
这由 (c) 与注 9.2(c) 推出。\nobreak\hspace{0.35em}■
\item

对每个 $ s \in \mathbb{R} $，$ \Lambda^s $ 属于 $ \mathcal{O}_M $。\nobreak\hspace{0.35em}■

\end{items}

我们现在可以证明傅里叶变换的另一条重要性质。

\end{remark}

\begin{theorem}

（卷积定理）

\begin{subitems}
\item

\[ (f*g)^{\wedge} = (2\pi)^{n/2}\widehat{f}\,\widehat{g} \qquad \text{对 } (f,g) \in \mathcal{L}_1\times\mathcal{L}_1 . \]

\item

\[ \widehat{\varphi}*\widehat{f} = (2\pi)^{n/2}\widehat{\varphi f} \qquad \text{对 } (\varphi,f) \in \mathcal{S}\times\mathcal{L}_1 . \]

\proofhead
(i) 由 (9.12) 与 (9.13) 可知，富比尼–托内利定理可以应用。于是由推论 7.9 得
\begin{equation*}
\begin{aligned}
(f*g)^{\wedge}(\xi) &= (2\pi)^{-n/2}\int e^{-i\langle x,\xi\rangle}\int f(x-y)g(y)\,dy\,dx \\ &= (2\pi)^{-n/2}\int g(y)\int e^{-i\langle x,\xi\rangle}f(x-y)\,dx\,dy .
\end{aligned}
\end{equation*}

因为

\[
\int e^{-i\langle x,\xi\rangle}f(x-y)\,dx = e^{-i\langle y,\xi\rangle}\int e^{-i\langle z,\xi\rangle}f(z)\,dz = e^{-i\langle y,\xi\rangle}(2\pi)^{n/2}\widehat{f}(\xi) ,
\]

我们便得到

\[
(f*g)^{\wedge}(\xi) = (2\pi)^{-n/2}\int(2\pi)^{n/2}\widehat{f}(\xi)e^{-i\langle y,\xi\rangle}g(y)\,dy = (2\pi)^{n/2}\widehat{f}(\xi)\widehat{g}(\xi) .
\]
\end{subitems}

(ii) 设 $ (\varphi,f) \in \mathcal{S}\times\mathcal{L}_1 $。由定理 9.3，存在 $ \mathcal{S} $ 中的序列 $ (f_j) $，使得在 $ \mathcal{L}_1 $ 中 $ f_j \to f $。命题 9.4 与 9.7 蕴含 $ \widehat{\varphi}*\widehat{f_j} $ 属于 $ \mathcal{S} $。由注 9.15(c)，$ \varphi f_j $ 也属于 $ \mathcal{S} $，故由 (i) 与注 9.14(a) 得

\[
(\widehat{\varphi}*\widehat{f_j})^{\wedge} = (2\pi)^{n/2}\widehat{\widehat{\varphi}}\,\widehat{\widehat{f_j}} = (2\pi)^{n/2}(\varphi f_j)^{\sim} \qquad \text{对 } j \in \mathbb{N} .
\]

由定理 9.12(ii)，我们便得到

\begin{equation}
\widehat{\varphi}*\widehat{f_j} = (2\pi)^{n/2}\widehat{\varphi f_j} \qquad \text{对 } j \in \mathbb{N} .\tag{9.14}
\end{equation}

因为 $ f_j \to f $ 在 $ L_1 $ 中成立，由注 9.1(a) 知 $ \widehat{f_j} \to \widehat{f} $ 在 $ BC $ 中成立。因此由推论 7.9，并利用 $ \widehat{\varphi} \in \mathcal{S} \subset L_1 $，可知序列 $ (\widehat{\varphi}*\widehat{f_j}) $ 在 $ BC $ 中收敛到 $ \widehat{\varphi}*\widehat{f} $。又因 $ \varphi f_j \to \varphi f $ 显然在 $ L_1 $ 中成立，由命题 9.9 推出序列 $ (\widehat{\varphi f_j}) $ 在 $ BC $ 中收敛到 $ \widehat{\varphi f} $。于是由注 9.1(a) 得到断言。\nobreak\hspace{0.35em}■

\end{theorem}

作为卷积定理的一个应用，我们证明一条引理，它构成傅里叶变换的 $ L_2 $ 理论的基础。

\begin{lemma}

对 $ f \in \mathcal{L}_1\cap\mathcal{L}_2 $，$ \widehat{f} $ 属于 $ C_0\cap\mathcal{L}_2 $，且 $ \|f\|_2 = \|\widehat{f}\|_2 $。

\proofhead
设 $ f \in \mathcal{L}_1\cap\mathcal{L}_2 $。由黎曼–勒贝格引理，$ \widehat{f} $ 属于 $ C_0 $，故只需验证 $ \|f\|_2 = \|\widehat{f}\|_2 $。于是令 $ g := f*\overline{\widetilde{f}} $。由定理 7.3(ii) 与习题 7.2 知 $ g $ 属于 $ \mathcal{L}_1\cap C_0 $，且
\[
g(0) = \int f(y)\overline{\widetilde{f}}(0-y)\,dy = \int f\overline{\widetilde{f}} = \|f\|_2^2 .
\]

由推论 9.13(iii) 得

\begin{equation}
\|f\|_2^2 = g(0) = \lim_{\varepsilon\to0}(2\pi)^{-n/2}\int\widehat{g}(\xi)e^{-\varepsilon^2|\xi|^2}\,d\xi .\tag{9.15}
\end{equation}

现在我们注意

\[
\widehat{\overline{\widetilde{f}}} = (2\pi)^{-n/2}\int e^{-i\langle x,\xi\rangle}\overline{f}(-x)\,dx = (2\pi)^{-n/2}\overline{\int e^{-i\langle -x,\xi\rangle}f(-x)\,dx} = \overline{\widehat{f}} ,
\]

这由积分的欧氏不变性推出。然后定理 9.16(i) 表明

\[
\widehat{g} = (f*\overline{\widetilde{f}})^{\wedge} = (2\pi)^{n/2}\widehat{f}\,\widehat{\overline{\widetilde{f}}} = (2\pi)^{n/2}|\widehat{f}|^2 .
\]

特别地，$ \widehat{g} $ 非负。因此 (9.15) 与单调收敛定理蕴含 $ \|f\|_2 = \|\widehat{f}\|_2 $。\nobreak\hspace{0.35em}■

\end{lemma}

\needspace{3\baselineskip}\subsection*{傅里叶乘子算子}

为了说明傅里叶变换的映射性质之意义，我们现在考虑常系数线性微分算子，并证明它们在“傅里叶域”中由乘子算子表示。

对 $ m \in \mathbb{N} $，我们用 $ \mathbb{C}_m[X_1,\ldots,X_n] $ 表示由次数 $ \le m $ 的全部多项式构成的向量子空间 $ \mathbb{C}[X_1,\ldots,X_n] $。对

\[
p = \sum_{|\alpha|\le m}a_\alpha X^\alpha \in \mathbb{C}_m[X_1,\ldots,X_n] ,
\]

我们令

\[
p(D) := \sum_{|\alpha|\le m}a_\alpha D^\alpha ,
\]

它是阶数 $ \le m $ 的常系数线性微分算子。这里 $ p $ 称为 $ p(D) $ 的\textbf{符号}。下面我们令

\[
\mathcal{D}\mathrm{iffop}^0 := \left\{ p(D) \,;\, p \in \mathbb{C}[X_1,\ldots,X_n] \right\} ,
\]

而 $ \mathcal{D}\mathrm{iffop}^0_m $ 是阶数不高于 $ m $ 的全部常系数线性微分算子构成的子集。

\begin{remark}

\begin{items}
\item

$ p(D) \in \mathcal{D}\mathrm{iffop}^0 $ 是 $ \mathcal{S} $ 到自身的线性连续映射，即 $ p(D) \in \mathcal{L}(\mathcal{S}) $。

\proofhead
这由注 9.2(c) 与 (h) 推出。\nobreak\hspace{0.35em}■
\end{items}

\begin{items}
\item[（b）]

映射

\begin{equation}
\mathbb{C}[X_1,\ldots,X_n] \to \mathcal{L}(\mathcal{S}), \quad p \mapsto p(D)\tag{9.16}
\end{equation}

是线性单射。

\proofhead
线性性是显然的。设 $ p = \sum_{|\alpha| \le m} a_\alpha X^\alpha \in \mathbb{C}[X_1,\ldots,X_n] $，且对一切 $ f \in \mathcal{S} $ 都有 $ p(D)f = 0 $。我们选取 $ \varphi \in \mathcal{D} $ 使得 $ \varphi | \mathbb{B}^n = 1 $。对 $ \beta \in \mathbb{N}^n $，由莱布尼茨法则可得
\[
D^\alpha(\varphi X^\beta) = \varphi D^\alpha X^\beta + \sum_{\gamma < \alpha} \binom{\alpha}{\gamma} D^{\alpha-\gamma}\varphi\, D^\gamma X^\beta .
\]

由于当 $ |x| < 1 $ 时 $ \varphi(x) = 1 $，我们由此得到

\begin{equation*}
D^\alpha(\varphi X^\beta)(0) = D^\alpha X^\beta(0) = \begin{cases} \beta! & \text{若 } \alpha = \beta , \\ 0 & \text{其他情形} \end{cases}
\end{equation*}

对 $ \alpha \in \mathbb{N}^n $。因为 $ \varphi X^\beta \in \mathcal{D} \subset \mathcal{S} $，我们于是得到 $ 0 = p(D)(\varphi X^\beta) = \beta!\, a_\beta $，对 $ \beta \in \mathbb{N}^n $ 且 $ |\beta| \le m $；因此 $ p = 0 $。这就证明了所述的单射性。 \nobreak\hspace{0.35em}■

\item[（c）]

$ p(D) $ 是形式自伴的当且仅当 $ p $ 有实系数。

\proofhead
令
\[
\mathcal{A}(\partial) := p(D) = \sum_{|\alpha| \le m} a_\alpha (-i)^{|\alpha|} \partial^\alpha
\]

则由命题 7.24 可得

\[
\mathcal{A}^{\sharp}(\partial) = \sum_{|\alpha| \le m} (-1)^{|\alpha|} \overline{a_\alpha (-i)^{|\alpha|}}\, \partial^\alpha = \sum_{|\alpha| \le m} (-i)^{|\alpha|} \overline{a_\alpha}\, \partial^\alpha ,
\]

这就完成了证明。 \nobreak\hspace{0.35em}■

\end{items}

由注 9.18(b)，我们可以把 $ \mathcal{D}\mathrm{iffop}^0 $ [或 $ \mathcal{D}\mathrm{iffop}^0_m $] 与 $ \mathbb{C}[X_1,\ldots,X_n] $ [或 $ \mathbb{C}_m[X_1,\ldots,X_n] $] 在映射 (9.16) 下的像等同起来。换言之，作为向量子空间，我们有

\[
\mathcal{D}\mathrm{iffop}^0_m \subset \mathcal{D}\mathrm{iffop}^0 \subset \mathcal{L}(\mathcal{S}) \quad \text{对 } m \in \mathbb{N} .
\]

对 $ a \in \mathcal{O}_M $ 与 $ f \in \mathcal{S} $，由推论 9.13(i) 与注 9.15(d) 可得 $ (f \mapsto a\hat{f}) \in \mathcal{L}(\mathcal{S}) $。再由推论 9.13(i) 可得

\[
a(D) := \mathcal{F}^{-1}a\mathcal{F}: \mathcal{S} \to \mathcal{S} , \quad f \mapsto \mathcal{F}^{-1}(a\hat{f})
\]

是 $ \mathcal{L}(\mathcal{S}) $ 的一个定义良好的元素，即一个以 $ a $ 为\textbf{符号}的\textbf{傅里叶乘子算子}。我们置

\[
\mathrm{Op} := \{ a(D) \in \mathcal{L}(\mathcal{S}) ;\ a \in \mathcal{O}_M \} .
\]
\end{remark}

\begin{proposition}

$ \mathrm{Op} $ 是 $ \mathcal{L}(\mathcal{S}) $ 的一个带单位的交换代数，且映射

\[
ev: (\mathcal{O}_M, +, \cdot) \to \mathrm{Op} , \quad a \mapsto a(D)
\]

是代数同构。

\proofhead
显然 $ \mathcal{O}_M := (\mathcal{O}_M, +, \cdot) $ 是代数 $ \mathbb{C}^{(\mathbb{R}^n)} $ 的带单位的交换子代数。也容易验证 $ ev $ 把向量空间 $ \mathcal{O}_M $ 线性地映入 $ \mathcal{L}(\mathcal{S}) $。

对 $ a, b \in \mathcal{O}_M $ 与 $ f \in \mathcal{S} $，我们有
\[
(ab)(D)f = \mathcal{F}^{-1}(ab\hat{f}) = \mathcal{F}^{-1}(a\mathcal{F}\mathcal{F}^{-1}(b\hat{f})) = \mathcal{F}^{-1}(a\widehat{b(D)f}) = a(D) \circ b(D)f .
\]

因此 $ ev $ 是满的代数同态。

最后设 $ a, b \in \mathcal{O}_M $ 且 $ a(D) = b(D) $。取 $ \xi \in \mathbb{R}^n $，并记 $ \varphi \in \mathcal{D} $ 为 $ \mathbb{B}^n(\xi,1) $ 的截断函数。则 $ f := \mathcal{F}^{-1}\varphi $ 属于 $ \mathcal{S} $，且 $ \hat{f}(\xi) = 1 $，并由推论 9.13(i) 可得

\[
a(\xi) = (a\hat{f})(\xi) = \mathcal{F}(a(D)f)(\xi) = \mathcal{F}(b(D)f)(\xi) = (b\hat{f})(\xi) = b(\xi) .
\]

由于这对每个 $ \xi \in \mathbb{R}^n $ 都成立，我们有 $ a = b $。因此 $ ev $ 是单射。 \nobreak\hspace{0.35em}■

\end{proposition}

\begin{corollary}

\begin{subitems}
\item

对 $ a, b \in \mathcal{O}_M $，我们有 $ ab(D) = a(D)b(D) = b(D)a(D) $。

\item

$ \mathbf{1}(D) = 1_{\mathcal{L}(\mathcal{S})} $。

\item

$ \mathcal{D}\mathrm{iffop}^0 $ 是 $ \mathbb{C}[X_1,\ldots,X_n] $ 在 $ ev $ 下的像。特别地，$ \mathcal{D}\mathrm{iffop}^0 $ 是 $ \mathrm{Op} $ 的带单位的交换子代数。

\proofhead
(i) 与 (ii) 是命题 9.19 的特殊情形。
\end{subitems}

(iii) 对 $ p \in \mathbb{C}[X_1,\ldots,X_n] \subset \mathcal{O}_M $ 且 $ p = \sum_{|\alpha| \le m} a_\alpha X^\alpha $，由命题 9.6(i) 得到

\begin{equation*}
\begin{aligned}
ev(p)f &= \mathcal{F}^{-1}p\mathcal{F}f = \mathcal{F}^{-1}(p\hat{f}) = \sum_{|\alpha| \le m} a_\alpha \mathcal{F}^{-1}(X^\alpha \hat{f}) \\ &= \sum_{|\alpha| \le m} a_\alpha \mathcal{F}^{-1}(\widehat{D^\alpha f}) = \sum_{|\alpha| \le m} a_\alpha D^\alpha f
\end{aligned}
\end{equation*}

对 $ f \in \mathcal{S} $。因此 $ ev(p) = \sum_{|\alpha| \le m} a_\alpha D^\alpha $，由此即得结论。 \nobreak\hspace{0.35em}■

\end{corollary}

这个推论表明，傅里叶变换可以用来求解常系数线性微分方程，办法是把它们化为简单的代数方程。这一事实是傅里叶变换的根本意义的一部分。下面的例子初步展示这些方法。

\begin{example}

\begin{items}
\item

设多项式 $ p \in \mathbb{C}[X_1,\ldots,X_n] $ 没有实零点。则 $ p(D) \in \mathcal{L}(\mathcal{S}) $ 是 $ \mathcal{S} $ 的自同构，且 $ [p(D)]^{-1} = (1/p)(D) $。

\proofhead
容易看出 $ 1/p $ 属于 $ \mathcal{O}_M $。现在由推论 9.20 推出
\[
1_{\mathcal{L}(\mathcal{S})} = \mathbf{1}(D) = (p \cdot 1/p)(D) = p(D)(1/p)(D) = (1/p)(D)p(D) .
\]

因为 $ a \in \mathcal{O}_M $ 时 $ a(D) \in \mathcal{L}(\mathcal{S}) $，这就证明了结论。 \nobreak\hspace{0.35em}■

\item

$ 1 - \Delta \in \mathcal{L}\mathrm{aut}(\mathcal{S}) $，且 $ (1 - \Delta)^{-1} = \Lambda^{-2}(D) $。

\proofhead
因为 $ 1 - \Delta = \Lambda^2(D) $，这由 (a) 推出。 \nobreak\hspace{0.35em}■
\end{items}

例 9.21(b) 说明，偏微分方程

\begin{equation}
-\Delta u + u = f\tag{9.17}
\end{equation}

对每个 $ f \in \mathcal{S} $ 都有唯一解 $ u \in \mathcal{S} $，且 $ u $ 在 $ \mathcal{S} $ 的拓扑下连续依赖于 $ f $。此外，我们也可以先对该方程作“傅里叶变换”来得到 (9.17) 的解 $ u \in \mathcal{S} $。按命题 9.6，这给出方程 $ (|\xi|^2 + 1)\hat{u}(\xi) = \Lambda^2(\xi)\hat{u}(\xi) = \hat{f}(\xi) $，$ \xi \in \mathbb{R}^n $。然后可以解出 $ \hat{u} $，得 $ \hat{u} = \Lambda^{-2}\hat{f} $，再作“傅里叶逆变换”，得 $ u = \mathcal{F}^{-1}(\Lambda^{-2}\hat{f}) = \Lambda^{-2}(D)f $。这种“傅里叶变换方法”在偏微分方程理论中起着突出作用。注意 $ \Lambda^{-2}(D) $，或更一般地 $ (1/p)(D) $，不是微分算子。

\end{example}

\needspace{3\baselineskip}\subsection*{普朗歇尔定理}

为结束本章，我们说明傅里叶变换也可以定义在 $ L_2 $ 上，并解释这一事实的几个推论。

设 $ H $ 是希尔伯特空间。若 $ T $ 是等距同构，则称 $ T: H \to H $ 是\textbf{酉的}。

\begin{remark}

设 $ H $ 是（实或复）希尔伯特空间，$ T: H \to H $ 是线性的。

\begin{items}
\item

若 $ T $ 是酉的，则 $ T $ 属于 $ \mathcal{L}\mathrm{aut}(H) $，且

\[
(Tx\,|\,Ty) = (x\,|\,y) \quad \text{对 } x, y \in H .
\]

\proofhead
第一个断言是显然的。因为 $ T $ 是等距，我们有
\[
4\operatorname{Re}(Tx\,|\,Ty) = \|T(x+y)\|^2 - \|T(x-y)\|^2 = \|x+y\|^2 - \|x-y\|^2 = 4\operatorname{Re}(x\,|\,y) ,
\]

因此对 $ x, y \in H $ 有 $ \operatorname{Re}(Tx\,|\,Ty) = \operatorname{Re}(x\,|\,y) $。把此恒等式中的 $ y $ 换成 $ iy $，我们得到

\[
\operatorname{Im}(Tx\,|\,Ty) = \operatorname{Re}(Tx\,|\,Tiy) = \operatorname{Re}(x\,|\,iy) = \operatorname{Im}(x\,|\,y) ,
\]

于是对 $ x, y \in H $ 有 $ (Tx\,|\,Ty) = (x\,|\,y) $。 \nobreak\hspace{0.35em}■

\item

若 $ H $ 是有限维的，则下列断言等价：

\end{items}

\begin{subitems}
\item

$ T $ 是酉的。

\item

对 $ x, y \in H $ 有 $ (Tx\,|\,Ty) = (x\,|\,y) $。

\item

$ T^*T = \operatorname{id}_H $。

\proofhead
“(i) $ \Rightarrow $ (ii)” 是 (a) 的推论。
“(ii) $ \Rightarrow $ (iii)” 设 $ \{b_1,\ldots,b_m\} $ 是 $ H $ 的规范正交基。则每个 $ y \in H $ 都可展开为 $ y = \sum_{j=1}^{m} (y\,|\,b_j)b_j $（见习题 II.3.12 与定理 VI.7.14）。由习题 VII.1.5 与 (ii) 可得
\[
T^*Tx = \sum_{j=1}^{m} (T^*Tx\,|\,b_j)b_j = \sum_{j=1}^{m} (Tx\,|\,Tb_j)b_j = \sum_{j=1}^{m} (x\,|\,b_j)b_j = x
\]

对每个 $ x \in H $。
“(iii) $ \Rightarrow $ (i)” 因为 $ T^*T = \operatorname{id}_H $，我们知道 $ T $ 是单射，因而由线性代数中的秩公式也是满射。对 $ x \in H $，我们还有

\[
\|Tx\|^2 = (Tx\,|\,Tx) = (T^*Tx\,|\,x) = (x\,|\,x) = \|x\|^2 .
\]

因此 $ T $ 是等距。 \nobreak\hspace{0.35em}■

\end{subitems}

\end{remark}

\begin{theorem}

（普朗歇尔） 傅里叶变换从 $ L_1 \cap L_2 $ 延拓为 $ L_2 $ 上的酉算子，这样的延拓是唯一的。

\proofhead
记 $ X_2 $ 为希尔伯特空间 $ L_2 $ 的向量子空间 $ L_1 \cap L_2 $。则由引理 9.17 可得 $ \mathcal{F} $ 属于 $ \mathcal{L}(X_2, L_2) $ 且是等距。因为 $ X_2 $ 包含空间 $ \mathcal{S} $，定理 9.3 与 VI.2.6 蕴涵存在唯一的等距延拓 $ \mathfrak{F} \in \mathcal{L}(L_2) $。作为等距，$ \mathfrak{F} $ 的像集是闭的，而由推论 9.13(i) 它包含空间 $ \mathcal{S} $。因此命题 V.4.4 蕴涵 $ \mathfrak{F} $ 是满射，从而是酉的。 \nobreak\hspace{0.35em}■
\end{theorem}

按照惯例，我们仍用符号 $ \mathcal{F} $ 表示 $ \mathfrak{F} $ 的唯一连续延拓，并同样称它为\textbf{傅里叶变换}。\footnote{在 $ L_2 $ 上，傅里叶变换有时也称为傅里叶–普朗歇尔变换，或普朗歇尔变换。}

下一个命题描述任意 $ f \in L_2 $ 的傅里叶变换 $ \mathcal{F}f $。

\begin{proposition}

对 $ f \in L_2 $，我们有

\[
\mathcal{F}f = \lim_{R \to \infty} \mathcal{F}(\chi_{R\mathbb{B}^n}f) = \lim_{R \to \infty} (2\pi)^{-n/2} \int_{[|x| \le R]} e^{-i\langle x,\cdot\rangle} f(x)\,dx \quad \text{于 } L_2 .
\]

\proofhead
对 $ R > 0 $，元素 $ f_R := \chi_{R\mathbb{B}^n}f $ 属于 $ L_1 \cap L_2 $，且控制收敛定理蕴涵
\[
\int |f - f_R|^2 dx = \int |f|^2 (1 - \chi_{R\mathbb{B}^n})^2 dx \to 0 \quad (R \to \infty) .
\]

因此在 $ L_2 $ 中 $ \lim_{R \to \infty} f_R = f $。于是由普朗歇尔定理，$ \mathcal{F}f_R $ 在 $ L_2 $ 中收敛到 $ \mathcal{F}f $。因为

\[
\mathcal{F}(f_R)(\xi) = (2\pi)^{-n/2} \int_{[|x| \le R]} e^{-i\langle x,\xi\rangle} f(x)\,dx \quad \text{对 } \xi \in \mathbb{R}^n ,
\]

结论得证。 \nobreak\hspace{0.35em}■

\end{proposition}

\begin{example}

设 $ n = 1 $ 且 $ a > 0 $。又设 $ f := \chi_{[-a,a]} \in \mathcal{L}_1(\mathbb{R}) $。则

\[
\hat{f}(\xi) = \frac{1}{\sqrt{2\pi}} \int_{-a}^{a} e^{-ix\xi} dx = \frac{-1}{\sqrt{2\pi}\, i\xi}\left(e^{-i\xi a} - e^{i\xi a}\right) = \sqrt{\frac{2}{\pi}}\, a\, \frac{\sin(a\xi)}{a\xi}
\]

对 $ \xi \in \mathbb{R} $。由于 $ \int |f|^2 dx = 2a $，普朗歇尔定理给出

\[
\int_{-\infty}^{\infty} \left[\frac{\sin(ax)}{ax}\right]^2 dx = \frac{\pi}{a} \quad \text{对 } a > 0 .
\]

注意 $ x \mapsto \sin(x)/x $ 不属于 $ \mathcal{L}_1(\mathbb{R}) $。 \nobreak\hspace{0.35em}■

\end{example}

\needspace{3\baselineskip}\subsection*{对称算子}

设 $ E $ 是 $ \mathbb{K} $ 上的巴拿赫空间。所谓 $ E $ 中的\textbf{线性算子} $ A $，是指映射 $ A: \operatorname{dom}(A) \subset E \to E $，其中 $ \operatorname{dom}(A) $ 是 $ E $ 的向量子空间且 $ A $ 是线性的。对线性算子 $ A_j: \operatorname{dom}(A_j) \subset E \to E $ 与 $ \lambda \in \mathbb{K}^{\times} $，我们通过

\[
\operatorname{dom}(A_0 + \lambda A_1) := \operatorname{dom}(A_0) \cap \operatorname{dom}(A_1) , \quad (A_0 + \lambda A_1)x := A_0 x + \lambda A_1 x .
\]

来定义 $ A_0 + \lambda A_1 $。

\textbf{乘积} $ A_0A_1 $ 由

\[
\operatorname{dom}(A_0A_1) := \{ x \in \operatorname{dom}(A_1) ;\ A_1x \in \operatorname{dom}(A_0) \} , \quad (A_0A_1)x := A_0(A_1x) .
\]

定义。

最后，由

\[
\operatorname{dom}([A_0,A_1]) := \operatorname{dom}(A_0A_1 - A_1A_0) , \quad [A_0,A_1]x := (A_0A_1 - A_1A_0)x
\]

定义的算子称为 $ A_0 $ 与 $ A_1 $ 的\textbf{交换子}。显然 $ A_0 + \lambda A_1 $、$ A_0A_1 $ 与 $ [A_0,A_1] $ 都是 $ E $ 中的线性算子，并且

\[
A_0 + \lambda A_1 = \lambda A_1 + A_0 , \quad \lambda A_0 = A_0(\lambda \operatorname{id}_E) , \quad [A_0,A_1] = -[A_1,A_0] .
\]

设 $ H $ 是希尔伯特空间，$ A: \operatorname{dom}(A) \subset H \to H $ 是 $ H $ 上的线性算子。若

\[
(Au\,|\,v) = (u\,|\,Av) \quad \text{对 } u, v \in \operatorname{dom}(A) ,
\]

则称 $ A $ 是\textbf{对称的}。

\begin{remark}

\begin{items}
\item

设 $ H $ 是复希尔伯特空间，$ A $ 是 $ H $ 上的线性算子。则下列断言等价：

\begin{subitems}
\item

$ A $ 是对称的。

\item

对 $ u \in \operatorname{dom}(A) $ 有 $ (Au\,|\,u) \in \mathbb{R} $。

\proofhead
“(i) $ \Rightarrow $ (ii)” 因为 $ A $ 是对称的，可得
\[
(Au\,|\,u) = (u\,|\,Au) = \overline{(Au\,|\,u)} \quad \text{对 } u \in \operatorname{dom}(A) ,
\]

因此 $ \operatorname{Im}(Au\,|\,u) = 0 $。
“(ii) $ \Rightarrow $ (i)” 对 $ u, v \in \operatorname{dom}(A) $，我们有

\begin{equation}
(A(u+v)\,|\,u+v) = (Au\,|\,u) + (Av\,|\,u) + (Au\,|\,v) + (Av\,|\,v) .\tag{9.18}
\end{equation}

由 (ii) 可得 $ \operatorname{Im}(Au\,|\,v) = -\operatorname{Im}(Av\,|\,u) $，因此

\[
\operatorname{Im}(Au\,|\,v) = -\operatorname{Im}(Av\,|\,u) = -\operatorname{Im}\overline{(u\,|\,Av)} = \operatorname{Im}(u\,|\,Av) .
\]

在 (9.18) 中把 $ u $ 换成 $ iu $，我们得到

\[
\operatorname{Re}(Au\,|\,v) = \operatorname{Im}(A(iu)\,|\,v) = \operatorname{Im}(iu\,|\,Av) = \operatorname{Re}(u\,|\,Av) .
\]

因此 $ (Au\,|\,v) = (v\,|\,Au) $。 \nobreak\hspace{0.35em}■

\end{subitems}

\item

设 $ p \in \mathbb{C}[X_1,\ldots,X_n] $，$ P $ 是 $ L_2 $ 上的线性算子，使得 $ \operatorname{dom}(P) = \mathcal{S} $ 且对 $ u \in \mathcal{S} $ 有 $ Pu := p(D)u $。则下列断言等价：

\begin{subitems}
\item

$ P $ 是对称的。

\item

$ p(D) $ 是形式自伴的。

\item

$ p $ 有实系数。

\proofhead
“(i) $ \Rightarrow $ (ii)” $ P $ 的对称性蕴涵
\[
(p(D)u\,|\,v) = (Pu\,|\,v) = (u\,|\,Pv) = (u\,|\,p(D)v) \quad \text{对 } u, v \in \mathcal{D} ,
\]

而这又由形式伴随算子的唯一性蕴涵 (ii)。
“(ii) $ \Rightarrow $ (iii)” 注 9.18(c)。
“(iii) $ \Rightarrow $ (i)” 设 $ p = \sum_{|\alpha| \le m} a_\alpha X^\alpha $。则推论 9.20(iii) 与普朗歇尔定理蕴涵

\[
(Pu\,|\,u) = (p(D)u\,|\,u) = (p\hat{u}\,|\,\hat{u}) = \sum_{|\alpha| \le m} a_\alpha \int \xi^\alpha |\hat{u}(\xi)|^2 d\xi \quad \text{对 } u \in \mathcal{S} .
\]

因此 $ (Pu\,|\,u) $ 是实的，而结论由 (a) 推出。 \nobreak\hspace{0.35em}■

\end{subitems}

\item

以 $ \mathcal{S} $ 为定义域时，拉普拉斯算子、波算子与薛定谔算子都是 $ L_2 $ 中的对称算子。

\proofhead
这由 (b) 与例 7.25(a) 和 (e) 推出。 \nobreak\hspace{0.35em}■
\end{items}

\end{remark}

\needspace{3\baselineskip}\subsection*{海森堡不确定性关系}

作为普朗歇尔定理的又一应用，我们以讨论量子力学中位置算子与动量算子的若干重要性质来结束本节。于是固定 $ j \in \{1,\ldots,n\} $ 并置

\[
\operatorname{dom}(A_j) := \{ u \in L_2 ;\ X_j\hat{u} \in L_2 \} , \quad \operatorname{dom}(B_j) := \{ u \in L_2 ;\ X_j u \in L_2 \} .
\]

然后，我们定义 $ L_2 $ 中的线性算子，即\textbf{动量算子} $ A_j $ 与\textbf{位置算子} $ B_j $（对第 $ j $ 个坐标），分别为

\[
A_j u := \mathcal{F}^{-1}(X_j\hat{u}) \quad \text{与} \quad B_j v := X_j v \quad \text{对 } u \in \operatorname{dom}(A_j) , \ v \in \operatorname{dom}(B_j) .
\]

\begin{remark}

\begin{items}
\item

我们有 $ \mathcal{S} \subset \operatorname{dom}(A_j) $，且

\[
A_j u = X_j(D)u = D_j u = -i\partial_j u \quad \text{对 } u \in \mathcal{S} .
\]

\proofhead
这由命题 9.7 与推论 9.8 推出。 \nobreak\hspace{0.35em}■
\item

我们有 $ \mathcal{F}(\operatorname{dom}(A_j)) = \operatorname{dom}(B_j) $ 以及一个交换图：

\begin{equation*}
\begin{array}{ccc} \operatorname{dom}(A_j) & \xrightarrow{\ A_j\ } & L_2 \\ \downarrow \mathcal{F} & & \downarrow \mathcal{F} \\ \operatorname{dom}(B_j) & \xrightarrow{\ B_j\ } & L_2 \end{array}
\end{equation*}

特别地，

\[
A_j u = \mathcal{F}^{-1}B_j\mathcal{F}u , \quad u \in \operatorname{dom}(A_j) \quad \text{与} \quad B_j u = \mathcal{F}A_j\mathcal{F}^{-1}u , \quad u \in \operatorname{dom}(B_j) .
\]

\proofhead
这些都是普朗歇尔定理的推论。 \nobreak\hspace{0.35em}■
\item

量子力学中的位置算子与动量算子是对称的。

\proofhead
设 $ u \in \operatorname{dom}(A_j) $。则 (b) 与普朗歇尔定理蕴涵
\[
(A_j u\,|\,u) = (\mathcal{F}^{-1}B_j\mathcal{F}u\,|\,u) = (B_j\hat{u}\,|\,\hat{u}) = \int \xi_j |\hat{u}(\xi)|^2 d\xi ,
\]

于是结论由注 9.26(a) 推出。 \nobreak\hspace{0.35em}■

\item

对 $ u \in \operatorname{dom}([A_j,B_j]) $，我们有 $ ([A_j,B_j]u\,|\,u) = 2i\operatorname{Im}(A_jB_j u\,|\,u) $。

\proofhead
由 (b)、(c) 与普朗歇尔定理，对 $ u \in \operatorname{dom}([A_j,B_j]) $ 得到
\begin{equation*}
\begin{aligned}
([A_j,B_j]u\,|\,u) &= (A_jB_j u - B_jA_j u\,|\,u) \\ &= (\mathcal{F}^{-1}B_j\mathcal{F}B_j u - B_j\mathcal{F}^{-1}B_j\mathcal{F}u\,|\,u) \\ &= (\mathcal{F}B_j u\,|\,B_j\mathcal{F}u) - (B_j\mathcal{F}u\,|\,\mathcal{F}B_j u) \\ &= 2i\operatorname{Im}(\mathcal{F}B_j u\,|\,B_j\mathcal{F}u) = 2i\operatorname{Im}(\mathcal{F}^{-1}B_j\mathcal{F}B_j u\,|\,u) \\ &= 2i\operatorname{Im}(A_jB_j u\,|\,u) . \qquad \blacksquare
\end{aligned}
\end{equation*}
\item

算子 $ i[A_j,B_j] $ 在 $ L_2 $ 中是对称的。

\proofhead
这由 (d) 推出。 \nobreak\hspace{0.35em}■
\item

我们有 $ \mathcal{S} \subset \operatorname{dom}([A_j,B_j]) $，且对 $ u \in \mathcal{S} $ 有 $ [A_j,B_j]u = -iu $。

\proofhead
第一个断言容易由命题 9.7 与注 9.2(h) 推出。此外 (a) 表明
\[
[A_j,B_j]u = D_j(X_j u) - X_j D_j u = (D_j X_j)u = -iu
\]

对 $ u \in \mathcal{S} $。 \nobreak\hspace{0.35em}■

\item

（$ \mathcal{S} $ 上的海森堡不确定性关系）对 $ j \in \{1,\ldots,n\} $，我们有

\[
\|u\|_2^2 \le 2\|\partial_j u\|_2 \|X_j u\|_2 \quad \text{对 } u \in \mathcal{S} .
\]

\proofhead
设 $ u \in \mathcal{S} $。由 (d) 与 (f) 有
\[
-i\|u\|_2^2 = -i(u\,|\,u) = ([A_j,B_j]u\,|\,u) = 2i\operatorname{Im}(A_jB_j u\,|\,u) .
\]

柯西–施瓦茨不等式于是给出

\[
\|u\|_2^2 = 2\left|\operatorname{Im}(A_jB_j u\,|\,u)\right| \le 2\left|(A_jB_j u\,|\,u)\right| = 2\left|(B_j u\,|\,A_j u)\right| \le 2\|A_j u\|_2 \|B_j u\|_2 ,
\]

于是由 (a) 即得结论。 \nobreak\hspace{0.35em}■

\end{items}

我们以把 $ \mathcal{S} $ 上的海森堡不确定性关系的适用范围推广到 $ \operatorname{dom}(A_j) \cap \operatorname{dom}(B_j) $ 来结束本节。为此先需要一个引理。

\end{remark}

\begin{lemma}

对每个 $ u \in \operatorname{dom}(A_j) \cap \operatorname{dom}(B_j) $，存在 $ \mathcal{S} $ 中的序列 $ (u_m) $ 使得

\[
\lim_{m \to \infty} (u_m, A_j u_m, B_j u_m) = (u, A_j u, B_j u) \quad \text{于 } L_2^3 .
\]

证明

\begin{subitems}
\item

设 $ u \in \operatorname{dom}(A_j) \cap \operatorname{dom}(B_j) $，并设 $ \{k_\varepsilon ;\ \varepsilon > 0\} $ 是高斯逼近核。我们置 $ u^\varepsilon := k_\varepsilon * u $。由习题 8(iv)，$ u^\varepsilon $ 属于 $ \mathcal{S} $，且定理 7.11 表明在 $ L_2 $ 中 $ \lim_{\varepsilon \to 0} u^\varepsilon = u $。

\item

因为 $ \check{k} = k $，由例 9.10(c) 可得

\[
\widehat{k_\varepsilon}(\xi) := \widehat{\check{k}_\varepsilon}(\xi) = \mathcal{F}^{-1}k_\varepsilon(\xi) = \varphi(\varepsilon\xi) = (2\pi)^{-n/2} e^{-\varepsilon^2|\xi|^2} \quad \text{对 } \xi \in \mathbb{R}^n .
\]

按定理 9.3，我们可以找到 $ \mathcal{S} $ 中的序列 $ (v_m) $ 使得在 $ L_2 $ 中 $ \lim_m v_m = u $。卷积定理于是表明

\[
\widehat{(k_\varepsilon * v_m)}(\xi) = (2\pi)^{n/2}\widehat{k_\varepsilon}(\xi)\widehat{v_m}(\xi) = e^{-\varepsilon^2|\xi|^2}\widehat{v_m}(\xi) \quad \text{对 } \xi \in \mathbb{R}^n .
\]

令 $ m \to \infty $ 则给出 $ \widehat{u^\varepsilon} = e^{-\varepsilon^2|\cdot|^2}\widehat{u} $（见推论 7.9 与定理 9.23）。因为

\[
\|A_j u - A_j u^\varepsilon\|_2^2 = \|X_j\hat{u} - X_j\widehat{u^\varepsilon}\|_2^2 = \int |\xi_j\hat{u}(\xi)|^2\left(1 - e^{-\varepsilon^2|\xi|^2}\right)^2 d\xi ,
\]

由控制收敛定理可得在 $ L_2 $ 中 $ \lim_{\varepsilon \to 0} A_j u^\varepsilon = A_j u $。

\item

设 $ \ddot{u} $ 是 $ u $ 的一个代表。我们置

\[
d_\varepsilon(x,z) := x_j\left(\ddot{u}(x) - \ddot{u}(x - \varepsilon z)\right) , \quad g_\varepsilon(x,z) := d_\varepsilon(x,z)k(z)
\]

对 $ \varepsilon > 0 $ 与 $ (x,z) \in \mathbb{R}^n \times \mathbb{R}^n $。则如 (7.7) 中那样（或由积分的闵可夫斯基不等式）可得

\begin{equation}
\|X_j u - X_j u^\varepsilon\|_2 \le \left(\int\left[\int |g_\varepsilon(x,z)|\,dz\right]^2 dx\right)^{1/2} \le \int \|d_\varepsilon(\cdot,z)\|_2 k(z)\,dz ,\tag{9.19}
\end{equation}

其中对最后一个不等式我们用了 $ g_\varepsilon = (d_\varepsilon\sqrt{k})\sqrt{k} $ 与 $ \int k\,dx = 1 $ 以及柯西–施瓦茨不等式。进一步注意到

\[
d_\varepsilon(\cdot,z) = X_j\ddot{u} - \tau_{\varepsilon z}(X_j\ddot{u}) - \varepsilon z_j \tau_{\varepsilon z}\ddot{u} ,
\]

由平移群在 $ L_2 $ 上的强连续性与积分的平移不变性可得

\[
\lim_{\varepsilon \to 0} \|d_\varepsilon(\cdot,z)\|_2 k(z) = 0 \quad \text{对 } z \in \mathbb{R}^n ,
\]

以及

\[
\|d_\varepsilon(\cdot,z)\|_2 k(z) \le 2\max\{\|X_j u\|_2, \|u\|_2\}(1 + |z_j|)k(z) \quad \text{对 } \varepsilon \in (0,2] , \ z \in \mathbb{R}^n .
\]

因为 $ z \mapsto (1 + |z_j|)k(z) $ 属于 $ \mathcal{L}_1 $，结论由 (9.19) 与控制收敛定理推出。 \nobreak\hspace{0.35em}■

\end{subitems}

\end{lemma}

\begin{corollary}

（海森堡不确定性关系） 对 $ 1 \le j \le n $，我们有

\[
\|u\|_2^2 \le 2\|A_j u\|_2 \|B_j u\|_2 \quad \text{对 } u \in \operatorname{dom}(A_j) \cap \operatorname{dom}(B_j) .
\]

\proofhead
这由注 9.27(a) 和 (g) 以及引理 9.28 推出。 \nobreak\hspace{0.35em}■
\end{corollary}

由注 9.27(a) 与引理 9.28 容易推出（如同定理 7.27 的证明中那样），对 $ u \in \operatorname{dom}(A_j) $，分布意义下的导数 $ \partial_j u $ 属于 $ L_2 $，因此是一个弱 $ L_2 $-导数。又 $ A_j u = -i\partial_j u $。从而，对 $ u \in \operatorname{dom}(A_j) \cap \operatorname{dom}(B_j) $，我们也可以把海森堡不确定性关系写成

\[
\left(\frac{1}{2}\int |u|^2 dx\right)^2 \le \int |\partial_j u|^2 dx \int |X_j u|^2 dx
\]

的形式，只要把 $ \partial_j u $ 理解为弱意义下的导数。对算子 $ A_j $ 与 $ B_j $ 的这种拓宽了的解释的重要性，在泛函分析所发展的希尔伯特空间上无界自伴算子理论中得到了阐明。由位置算子与动量算子构成的自伴算子，特别是薛定谔算子，被用于量子

力学（例如 [RS72]）。关于海森堡不确定性关系的解释，我们请读者参阅物理学文献。

\exerciseshead

\begin{exercises}
\item

设 $ a > 0 $。求下列函数的傅里叶变换：

\begin{subitems}
\item

$ \sin(ax)/x $，(ii) $ 1/(a^2 + x^2) $，(iii) $ e^{-a|x|} $，(iv) $ (1 - |x|/a)\chi_{[-a,a]}(x) $，(v) $ (\sin(ax)/x)^2 $。
（提示：见 §VIII.6。）

\end{subitems}

\item

设 $ f(x) := (\sin(x)/x)^2 $，$ g(x) := e^{2ix}f(x) $，$ x \in \mathbb{R}^\times $。试证 $ f * g = 0 $。

（提示：应用习题 1 与定理 9.16。）

\item

证明：若 $ f \in \mathcal{L}_1 $ 满足 $ f * f = f $ 或 $ f * f = 0 $，则 $ f = 0 $。

\item

设 $ \{\varphi_\varepsilon ; \varepsilon > 0\} $ 是恒等逼近，$ (\varepsilon_j) $ 是零序列。证明 $ (\mathcal{F}(\varphi_{\varepsilon_j})) $ 在 $ \mathcal{D}'(\mathbb{R}^n) $ 中收敛到 $ (2\pi)^{-n/2}\mathbf{1} $。

\item

设 $ a, f \in \mathcal{S} $，证明 $ a(D)f = \widehat{a} * f $。

\item

设 $ s \ge 0 $，定义 $ H^s := \{u \in L_2 ; \Lambda^s \widehat{u} \in L_2\} $，并对 $ u, v \in H^s $ 定义 $ (u|v)_{H^s} := (\Lambda^s \widehat{u} | \widehat{v})_{L_2} $。证明

\begin{subitems}
\item

$ H^s := (H^s ; (\cdot|\cdot)_{H^s}) $ 是满足 $ H^0 = L_2 $ 的 Hilbert 空间，且

\[
\mathcal{S} \overset{d}{\hookrightarrow} H^s \overset{d}{\hookrightarrow} H^t \overset{d}{\hookrightarrow} L_2 \qquad \text{当 } s > t > 0 \text{ 时} ;
\]
\item

对 $ m \in \mathbb{N} $，$ H^m = W_2^m $。

\end{subitems}

\item

设 $ s > n/2 $，证明

\begin{subitems}
\item

$ \mathcal{F}(H^s) \subset L_1 $，以及

\item

$ H^2 \overset{d}{\hookrightarrow} C_0 $（索伯列夫嵌入定理）。
（提示：(i) 对 $ \Lambda^s |\widehat{u}| \Lambda^{-s} $ 应用柯西–施瓦茨不等式。(ii) 黎曼–勒贝格定理。）

\end{subitems}

\item

设 $ \sigma \ge 0 $，$ \{k_\varepsilon ; \varepsilon > 0\} $ 是高斯逼近核。证明：

\begin{subitems}
\item

$ T(t) := [f \mapsto k_{\sqrt{t}} * f] $ 对每个 $ t > 0 $ 属于 $ \mathcal{L}(H^\sigma) $。

\item

$ T(t + s) = T(t)T(s) $，$ s, t > 0 $。

\item

对 $ f \in H^\sigma $，$ \lim_{t \to 0} T(t)f = f $。

\item

$ T(t)(L_2) \subset \mathcal{S} $，$ t > 0 $。

\item

设 $ f \in L_2 \cap C $，$ u(t, x) := T(t)f(x) $，$ (t, x) \in [0, \infty) \times \mathbb{R}^n $。证明 $ u $ 求解 $ \mathbb{R}^n $ 中热方程的初值问题，即

\begin{equation}
\partial_t u - \Delta u = 0 \quad \text{在 } (0, \infty) \times \mathbb{R}^n \text{ 中} \quad \text{且} \quad u(0, \cdot) = f \quad \text{在 } \mathbb{R}^n \text{ 上},\tag{9.20}
\end{equation}

其中 $ u \in C^\infty((0, \infty) \times \mathbb{R}^n) \cap C(\mathbb{R}^+ \times \mathbb{R}^n) $，且 $ u $ 逐点满足 (9.20)。
注 令 $ T(0) := \operatorname{id}_{H^\sigma} $。则称 $ \{T(t) ; t \ge 0\} $ 为（$ H^\sigma $ 的）\textbf{高斯–魏尔斯特拉斯半群}。
（提示：(v) 为得到常微分方程的初值问题，对 (9.20) 关于 $ x \in \mathbb{R}^n $ 作傅里叶变换。）

\end{subitems}

\item

设 $ n = 1 $，$ p_y(x) := \sqrt{2/\pi}\, y/(x^2 + y^2) $，$ (x, y) \in \mathbb{H}^2 $。又设 $ \sigma \ge 0 $。证明下列命题：

\begin{subitems}
\item

$ P(y) := [f \mapsto p_y * f] $ 对每个 $ t > 0 $ 属于 $ \mathcal{L}(H^\sigma) $。

\item

对 $ y, z > 0 $，$ P(y + z) = P(y)P(z) $。

\item

对 $ f \in H^\sigma $，$ \lim_{y \to 0} P(y)f = f $。

\item

$ P(y)(L_2) \subset \mathcal{S} $。

\item

设 $ f \in L_2 \cap C $，令

\[
u(x, y) := (P(y)f)(x) \qquad \text{对 } (x, y) \in \mathbb{H}^2 .
\]

则 $ u $ 属于 $ C^2(\mathbb{H}^2) \cap C(\overline{\mathbb{H}^2}) $，且求解半平面的 Dirichlet 边值问题：

\[
\Delta u = 0 \quad \text{在 } \mathbb{H}^2 \text{ 中} \quad \text{且} \quad u(\cdot, 0) = f \quad \text{在 } \mathbb{R} \text{ 上} .
\]

注 令 $ P(0) := \operatorname{id}_{H^\sigma} $，则称 $ \{P(y) ; y \ge 0\} $ 为（$ H^\sigma $ 的）\textbf{泊松半群}。
（提示：(ii) 习题 1。(v) 例 9.21(b)。）

\end{subitems}

\item

设 $ X $ 是 $ \mathbb{R}^n $ 中的开集，$ (X_k) $ 是 $ X $ 的递增相对紧开子集序列，且 $ X = \bigcup_k X_k $（见注 1.16(d) 与 (e)）。又设

\[
q_k(f) := \max_{|\alpha| \le k} \|\partial^\alpha f\|_{\infty, \overline{X}_k} \qquad \text{对 } f \in C^\infty(X) \text{ 与 } k \in \mathbb{N} ,
\]

以及

\[
d(f, g) := \sum_{k=0}^{\infty} 2^{-k} \frac{q_k(f - g)}{1 + q_k(f - g)} \qquad \text{对 } f, g \in C^\infty(X) .
\]

证明 $ (C^\infty(X), d) $ 是完备度量空间。（提示：为证明完备性，应用对角线序列原理（注 III.3.11(a)）。）

\item

证明 $ \mathcal{D} \overset{d}{\hookrightarrow} C^\infty $ 与 $ \mathcal{S} \overset{d}{\hookrightarrow} C^\infty $。

（提示：考虑 $ \varphi(\varepsilon \cdot) $，其中 $ \varphi $ 是 $ \overline{\mathbb{B}}^n $ 的截断函数。）

\item

对 $ f \in \mathcal{D} $，令

\[
F(z) := \int e^{-i(z|x)_{\mathbb{C}^n}} f(x) \, dx \qquad \text{对 } z \in \mathbb{C} .
\]

证明 $ F $ 属于 $ C^\omega(\mathbb{C}, \mathbb{C}) $。

（提示：参照注 V.3.4(c)，应用推论 3.19。）

\item

证明：对 $ f \in \mathcal{D} \setminus \{0\} $，$ \widehat{f} $ 不属于 $ \mathcal{D} $。（提示：回忆习题 12 以及解析函数的唯一性定理（定理 V.3.13）。）

\end{exercises}
